% Options for packages loaded elsewhere
\PassOptionsToPackage{unicode}{hyperref}
\PassOptionsToPackage{hyphens}{url}
\documentclass[
  ukrainian,
]{article}
\usepackage{xcolor}
\usepackage[margin=2.5cm]{geometry}
\usepackage{amsmath,amssymb}
\setcounter{secnumdepth}{-\maxdimen} % remove section numbering
\usepackage{iftex}
\ifPDFTeX
  \usepackage[T1]{fontenc}
  \usepackage[utf8]{inputenc}
  \usepackage{textcomp} % provide euro and other symbols
\else % if luatex or xetex
  \usepackage{unicode-math} % this also loads fontspec
  \defaultfontfeatures{Scale=MatchLowercase}
  \defaultfontfeatures[\rmfamily]{Ligatures=TeX,Scale=1}
\fi
\usepackage{lmodern}
\ifPDFTeX\else
  % xetex/luatex font selection
\fi
% Use upquote if available, for straight quotes in verbatim environments
\IfFileExists{upquote.sty}{\usepackage{upquote}}{}
\IfFileExists{microtype.sty}{% use microtype if available
  \usepackage[]{microtype}
  \UseMicrotypeSet[protrusion]{basicmath} % disable protrusion for tt fonts
}{}
\makeatletter
\@ifundefined{KOMAClassName}{% if non-KOMA class
  \IfFileExists{parskip.sty}{%
    \usepackage{parskip}
  }{% else
    \setlength{\parindent}{0pt}
    \setlength{\parskip}{6pt plus 2pt minus 1pt}}
}{% if KOMA class
  \KOMAoptions{parskip=half}}
\makeatother
\usepackage{longtable,booktabs,array}
\newcounter{none} % for unnumbered tables
\usepackage{calc} % for calculating minipage widths
% Correct order of tables after \paragraph or \subparagraph
\usepackage{etoolbox}
\makeatletter
\patchcmd\longtable{\par}{\if@noskipsec\mbox{}\fi\par}{}{}
\makeatother
% Allow footnotes in longtable head/foot
\IfFileExists{footnotehyper.sty}{\usepackage{footnotehyper}}{\usepackage{footnote}}
\makesavenoteenv{longtable}
\usepackage{graphicx}
\graphicspath{{/Users/kalinicenkovitalijmikolajovic/.codex/worktrees/5c15/Script_for_translate_book/Output_Final/}}
\makeatletter
\newsavebox\pandoc@box
\newcommand*\pandocbounded[1]{% scales image to fit in text height/width
  \sbox\pandoc@box{#1}%
  \Gscale@div\@tempa{\textheight}{\dimexpr\ht\pandoc@box+\dp\pandoc@box\relax}%
  \Gscale@div\@tempb{\linewidth}{\wd\pandoc@box}%
  \ifdim\@tempb\p@<\@tempa\p@\let\@tempa\@tempb\fi% select the smaller of both
  \ifdim\@tempa\p@<\p@\scalebox{\@tempa}{\usebox\pandoc@box}%
  \else\usebox{\pandoc@box}%
  \fi%
}
% Set default figure placement to htbp
\def\fps@figure{htbp}
\makeatother
\ifLuaTeX
\usepackage[bidi=basic,shorthands=off,provide=*]{babel}
\else
\usepackage[bidi=default,shorthands=off,provide=*]{babel}
\fi
\ifLuaTeX
  \usepackage{selnolig} % disable illegal ligatures
\fi
\setlength{\emergencystretch}{3em} % prevent overfull lines
\providecommand{\tightlist}{%
  \setlength{\itemsep}{0pt}\setlength{\parskip}{0pt}}
\usepackage{bookmark}
\IfFileExists{xurl.sty}{\usepackage{xurl}}{} % add URL line breaks if available
\urlstyle{same}
\hypersetup{
  pdflang={uk},
  hidelinks,
  pdfcreator={LaTeX via pandoc}}

\author{}
\date{}

\begin{document}

ЧАСТКОВО ДИФЕРЕНЦІАЛ РІВНЯННЯ для Вчені та Інженери

Стенлі Дж. Фарлоу, професор математики, Університет Мену

\section{С. Фарлоу}\label{ux441.-ux444ux430ux440ux43bux43eux443}

\section{РІВНЯННЯ З ПОХІДНИМИ ПОХІДНИМИ ДЛЯ НАУКОВЦІВ І
ІНЖЕНЕРІВ}\label{ux440ux456ux432ux43dux44fux43dux43dux44f-ux437-ux43fux43eux445ux456ux434ux43dux438ux43cux438-ux43fux43eux445ux456ux434ux43dux438ux43cux438-ux434ux43bux44f-ux43dux430ux443ux43aux43eux432ux446ux456ux432-ux456-ux456ux43dux436ux435ux43dux435ux440ux456ux432}

Перекладено з англійської \textbf{A.I.} Плісом

Редактор

С. І. Похожаєва

5\textbar{} 66K 22.161.6 φ24 υ ΥΠΚ 517.2

\section{Фарлоу С.}\label{ux444ux430ux440ux43bux43eux443-ux441.}

F 24 Рівняння з частими похідними для науковців та інженерів. Москва,
видавництво Мир, 1985, 384 с.

King's American Mathematician, який є підручником з теорії
диференціальних рівнянь у похідних похідних. Він характеризується
компактністю, ясністю та ясністю подачі, а також неформальним підходом
до подачі матеріалу. Вона містить багато ілюстрацій, графіків і діаграм,
але замість суворих доказів часто даються роздуми, засновані на інтуїції
або аналогіях.

Для інженерів і нематематиків --- біологів, диміків, а також студентів
університету.
\(\Phi = \frac{1702050000-380}{041(01)-85} 26-85, \text{ q. 1}\) BBK
22.161.6

Редакційна колегія журналу Literature on Mathematical Sciences

© переклад російською, Мир, 1985

© John Wiley \& Sons, Inc.~1982. Всі права захищені. Авторизований
переклад з англомовного видання, опублікований John Wiley \& Sons, Inc.

\section{З РЕДАКТОРА
ПЕРЕКЛАДУ}\label{ux437-ux440ux435ux434ux430ux43aux442ux43eux440ux430-ux43fux435ux440ux435ux43aux43bux430ux434ux443}

Книга, запропонована для читачів, слід розглядати як підручник для
курсів: «Диференціальні рівняння в похідних похідних» та «Рівняння
математичної фізики» для студентів технічних університетів. Вона також
буде корисною для інженерів різних спеціальностей. У книзі представлено
традиційні методи розв'язання диференціальних рівнянь у похідних
похідних, включаючи деякі чисельні методи. Водночас особлива увага
приділяється методу розділення змінних і методу інтегральних
перетворень, і цю частину книги можна використовувати на практичних
заняттях в університетах

Орієнтація книги на широку «нематематичну» аудиторію визначила стиль
викладання розглянутих питань: інтуїтивний підхід і велика увага до
фізичного значення не лише самих рівнянь, а й меж і початкових умов для
різних задач. Вважається, що читачі знайомі лише з основами
диференціального та інтегрального числення та елементами теорії
звичайних диференціальних рівнянь.

Фахівці можуть не задовольнитися інтуїтивним описом певних понять і
методів автором, але більш суворий виклад неминуче призведе до
ускладнення книги і звузить коло її читачів, що, очевидно, не було
наміром автора.

Більшим запереченням може бути відсутність загальних результатів щодо
задач для диференціальних рівнянь у похідних з змінними коефіцієнтами,
включно з рівняннями вищого порядку. Однак читач має чітко розуміти, що
ця книга є вступним курсом з теорії диференціальних рівнянь у часткових
похідних. Питання лінійної теорії диференціальних рівнянь, які автор не
включив, можна знайти в літературі, наведеній наприкінці книги.

Матеріал книги подається автором з великою методологічною майстерністю.
Однак не все було однаково успішним, і винний читач знайде це,
наприклад, у вивченні принципу Дюамеля, сферичних заповітів та деяких
інших уривків у книзі.

Без сумніву, запропонована книга С. Фарлоу знайде читача не лише серед
студентів і інженерів, а й серед викладачів, які можуть багато чого
навчитися з конкретного способу презентації та розробки спеціальних
математичних курсів.

\section{ПЕРЕДМОВА}\label{ux43fux435ux440ux435ux434ux43cux43eux432ux430}

За останні кілька років спостерігається значне зростання кількості
учнів, які починають вивчати диференціальні рівняння з похідними у своїх
одинадцятих класах. Крім того, багато з них спеціалізуються на сферах,
де інтуїція важливіша за математичну строгість. Тому, пишучи цю книгу, я
прагнув створити текст, який стимулює інтуїтивне мислення читачів, не
втрачаючи при цьому значної математичної строгості. Можливо, на одному
полюсі, подати предмет на високому рівні «епсилон-дельта», але тоді, як
це зазвичай буває, студенти не знатимуть, що робити з отриманою
інформацією. Інша крайність --- повне ігнорування математичних тонкощів,
але тоді ні учні, ні вчителі не знатимуть, як діяти. Я намагався знайти
розумний компроміс між цими крайнощами.

Ця книга виникла з курсу лекцій, які я читав протягом останніх п'яти
років. Розташування матеріалу не є загальноприйнятим --- вони майже
незалежні один від одного

Друг лекцій.

Основна увага приділяється двом найважливішим аналітичним методам:
розділенню змінних та інтегральним перетворенням. Також представлені
деякі нестандартні розділи, такі як метод Монте-Карло, варіаційне
числення, теорія керування, теорія потенціалів, інтегральні рівняння. Це
пов'язано з тим, що більшості студентів доведеться стикатися з цими
питаннями під час навчання, і якщо вони не ознайомляться з ними зараз,
навряд чи коли-небудь систематично їх вивчатимуть.

Цю книгу можна використовувати як навчальний посібник для одно- або
двосеместрового курсу як для початківців, так і для старших студентів.
Читачі повинні знати лише диференціальне та інтегральне числення та
звичайні диференціальні рівняння. Матеріал більшості лекцій подається в
одному-двох уроках. Типовий односеместровий курс може включати лекції
1-13, 15-17, 19-20, 22-23, 25-27, 30-32, 37-39. Усі 47 лекцій легко
читати протягом двох семестрів. Водночас часу буде достатньо для
вирішення проблем.

Автор вдячний Вайлі за пропозицію написати книгу, а також рецензентам
--- професорам Крісу Рорресу та М. Куршиду Алі, які зробили низку
корисних коментарів до тексту книги. Будь-які пропозиції щодо подальшого
покращення цього підручника будуть вдячні як від учнів, так і від
викладачів. Я також вдячний Дороті, Сюзан, Александер і Дезі Фарлоу.

Стенлі Дж. Фарлоу

\section{ВСТУП ДО ТЕОРІЇ ДИФЕРЕНЦІАЛЬНИХ РІВНЯНЬ У ПОХІДНИХ
ПОХІДНИХ}\label{ux432ux441ux442ux443ux43f-ux434ux43e-ux442ux435ux43eux440ux456ux457-ux434ux438ux444ux435ux440ux435ux43dux446ux456ux430ux43bux44cux43dux438ux445-ux440ux456ux432ux43dux44fux43dux44c-ux443-ux43fux43eux445ux456ux434ux43dux438ux445-ux43fux43eux445ux456ux434ux43dux438ux445}

МЕТА ЛЕКЦІЇ: Показати, що таке диференціальні рівняння в похідних
похідних, де і як вони виникають і як їх розв'язують. Коротко розгляньте
методи класифікації рівнянь і надайте список базових понять, які будуть
детально вивчені на наступних лекціях.

Більшість фізичних явищ у таких сферах, як гідродинаміка, електрика та
магнетизм, механіка, оптика та теплопередача, загалом можна описати за
допомогою рівнянь у похідних похідних (PNP). Більшість рівнянь у
математичній фізиці є диференціальними рівняннями з похідними похідними.
Правда, з деякими спрощеними припущеннями ці рівняння зводяться до
звичайних диференціальних рівнянь, але повний опис таких систем неминуче
призводить до використання диференціальних рівнянь у похідних похідних.

\section{Що таке часткові диференціальні
рівняння?}\label{ux449ux43e-ux442ux430ux43aux435-ux447ux430ux441ux442ux43aux43eux432ux456-ux434ux438ux444ux435ux440ux435ux43dux446ux456ux430ux43bux44cux43dux456-ux440ux456ux432ux43dux44fux43dux43dux44f}

Рівняння з частковими похідними --- це рівняння, яке містить
диференціали у похідних похідних. На відміну від звичайних
диференціальних рівнянь (TAC), де невідома функція залежить лише від
однієї змінної, у рівняннях часткових похідних вона залежить від кількох
змінних (наприклад, температура \(u(x, t)\) залежить від координати x і
часу t).

Тепер наведемо деякі з найважливіших диференціальних рівнянь у похідних
похідних. Для спрощення запису ми використаємо наступні нотації

\[
u_t \equiv \frac{\partial u}{\partial t}
\]

, \(u_x \equiv \frac{\partial u}{\partial x}\) ,
\(u_{xx} \equiv \frac{\partial^2 u}{\partial x^2}\) , \ldots{}

\section{Деякі диференціальні рівняння в похідних
похідних}\label{ux434ux435ux44fux43aux456-ux434ux438ux444ux435ux440ux435ux43dux446ux456ux430ux43bux44cux43dux456-ux440ux456ux432ux43dux44fux43dux43dux44f-ux432-ux43fux43eux445ux456ux434ux43dux438ux445-ux43fux43eux445ux456ux434ux43dux438ux445}

\[
u_t = u_{xx}
\]

(одновимірне рівняння теплопровідності), \(u_t = u_{xx} + u_{yy}\)
(двовимірне рівняння теплопровідності),
\(u_{rr} + \frac{1}{r}u_r + \frac{1}{r^2}u_{00} = 0\) (рівняння Лапласа
за полярними координатами), \(u_{tt} = u_{xx} + u_{yy} + u_{zz}\)
(тривимірне хвильове рівняння), \(u_{tt} = u_{xx} + zu_t + \beta u\)
(телеграфне рівняння).

\section{Примітки щодо наведених
прикладів}\label{ux43fux440ux438ux43cux456ux442ux43aux438-ux449ux43eux434ux43e-ux43dux430ux432ux435ux434ux435ux43dux438ux445-ux43fux440ux438ux43aux43bux430ux434ux456ux432}

У всіх наведених вище прикладах функція невідома і залежить від більше
ніж однієї змінної. Змінна і (яку ми диференціюємо) називається залежною
змінною. Змінні, за якими відбувається диференціювання, називаються
незалежними змінними. Наприклад, у рівнянні.

\[
u_t = u_{xx}
\]

Залежна змінна \(u(x, t)\) є функцією двох незалежних змінних x і t; та
в рівнянні

\[
u_t = u_{rr} + \frac{1}{r} u_r + \frac{1}{r^2} u_{00}
\]

змінна \(u(r, \theta, t)\) залежить від \(r, \theta, t\) .

\section{Чому необхідно вивчати диференціальні рівняння у похідних
похідних?}\label{ux447ux43eux43cux443-ux43dux435ux43eux431ux445ux456ux434ux43dux43e-ux432ux438ux432ux447ux430ux442ux438-ux434ux438ux444ux435ux440ux435ux43dux446ux456ux430ux43bux44cux43dux456-ux440ux456ux432ux43dux44fux43dux43dux44f-ux443-ux43fux43eux445ux456ux434ux43dux438ux445-ux43fux43eux445ux456ux434ux43dux438ux445}

Більшість фізичних законів природи можна сформулювати через рівняння з
похідними похідними. Прикладами є рівняння Максвелла, закон
теплопередачі Ньютона, рівняння Нав'є-Стокса, рівняння руху Ньютона та
рівняння Шредінгера в кватичній механіці. У всіх цих рівняннях фізичні
явища описуються мовою просторових і часових похідних. Похідні
з'являються в рівняннях, оскільки вони описують найважливіші фізичні
величини (такі як швидкість, прискорення, сила, тертя, потік, струм
тощо). Таким чином, виникають диференціальні рівняння в частинних
похідних, які містять невідому функцію, яку потрібно визначити.

Мета цієї книги --- показати читачеві:

\begin{enumerate}
\def\labelenumi{\arabic{enumi}.}
\item
  Як сформулювати фізичну задачу у формі диференціального рівняння у
  похідних похідних (побудова математичної моделі).
\item
  Як розв'язати диференціальне рівняння в похідних (з урахуванням
  початкових і граничних умов).
\end{enumerate}

Перш ніж перейти до побудови математичних моделей, давайте коротко
розглянемо методи розв'язання рівнянь у похідних похідних.

\section{Як розв'язати диференціальні рівняння в похідних
похідних?}\label{ux44fux43a-ux440ux43eux437ux432ux44fux437ux430ux442ux438-ux434ux438ux444ux435ux440ux435ux43dux446ux456ux430ux43bux44cux43dux456-ux440ux456ux432ux43dux44fux43dux43dux44f-ux432-ux43fux43eux445ux456ux434ux43dux438ux445-ux43fux43eux445ux456ux434ux43dux438ux445}

Гарне питання. Виявляється, існує цілий арсенал методів, придатних для
практичного використання. Найважливішими є ті, у яких диференціальні
рівняння з похідними похідними зводяться до звичайних диференціальних
рівнянь. Перелічимо десять методів розв'язання диференціальних рівнянь у
похідних похідних:

\begin{enumerate}
\def\labelenumi{\arabic{enumi}.}
\tightlist
\item
  Метод розділення змінних. Рівняння з факторами, отриманими з n
  незалежними змінними, зводиться до n звичайних-
\end{enumerate}

диференціальнірівняння 1).

\begin{enumerate}
\def\labelenumi{\arabic{enumi}.}
\setcounter{enumi}{1}
\item
  Метод інтегральних перетворень. Рівняння в похідних з n незалежними
  змінними зводиться до рівняння з похідними у (n-1) незалежних змінних;
  Відповідно, рівняння з частковими похідними з двома незалежними
  змінними можна звести до звичайного диференціального рівняння \(^2\)
  ).
\item
  Метод перетворення координат. Початкове диференціальне рівняння в
  похідних похідних зводиться до звичайного рівняння з частковими
  похідними або до іншого, простішого рівняння у похідних похідних
  шляхом відповідного перетворення координат (наприклад, обертання осей
  координат тощо).
\item
  Перетворення залежної змінної. Початкове диференціальне рівняння в
  похідних похідних перетворюється у таке рівняння з похідними похідними
  для іншої невідомої функції, яка
\end{enumerate}

легше розв'язати, ніж початкові 3).

\begin{enumerate}
\def\labelenumi{\arabic{enumi}.}
\setcounter{enumi}{4}
\tightlist
\item
  Чисельні методи. Початкове диференціальне рівняння в похідних похідних
  зводиться до системи різницей, яку розв'язують методом ітерацій на
  комп'ютері. У багатьох випадках це єдиний спосіб розв'язати
  диференціальне рівняння з похідними похідними. Окрім різницьких
  методів для розв'язання рівнянь з частковими про-
\end{enumerate}

\begin{enumerate}
\def\labelenumi{\alph{enumi})}
\setcounter{enumi}{4}
\tightlist
\item
  Ефективніше використовувати загальні перетворення незалежних і
  залежних змінних, включаючи похідні. Ед.
\end{enumerate}

4) Це стосується повного розділення змінних. При частковому розділенні
змінних одне диференціальне рівняння в похідних похідних зводиться до
кількох рівнянь у часткових похідних з меншою кількістю незалежних
змінних. --- Ед. Ед.

2) Це відбувається у випадку одномірного інтегрального перетворення. У
випадку k-вимірного інтегрального перетворення LCP з n незалежними
змінними зводиться до рівняння часткових з n---k незалежними змінними.
--- Ед. Ед.

Існують також інші чисельні методи, зокрема ті, що базуються на
наближенні розв'язків за поліноміальними поверхнями (наближення за
сплайнами).

\begin{enumerate}
\def\labelenumi{\arabic{enumi}.}
\setcounter{enumi}{5}
\tightlist
\item
  Метод теорії збурень. Початкова нелінійна задача зводиться до
  послідовності лінійних задач, апрокенн-
\end{enumerate}

нелінійна задача 1).

\begin{itemize}
\tightlist
\item
  \begin{enumerate}
  \def\labelenumi{\arabic{enumi}.}
  \setcounter{enumi}{6}
  \tightlist
  \item
    Метод функцій Гріна. Початкові та граничні умови замінюються
    системою найпростіших джерел, і задача розв'язується для кожного
    найпростішого джерела. Повне розв'язання початкової задачі
    отримується шляхом підсумовування розв'язків елементарних джерел.
  \end{enumerate}
\item
  \begin{enumerate}
  \def\labelenumi{\arabic{enumi}.}
  \setcounter{enumi}{7}
  \tightlist
  \item
    Метод інтегральних рівнянь. Рівняння з частковими похідними
    зводяться до інтегрального рівняння (рівняння, в якому невідома
    функція знаходиться під знаком інтеграла). Існує багато різних
    методів розв'язання інтегральних рівнянь.
  \end{enumerate}
\item
  \begin{enumerate}
  \def\labelenumi{\arabic{enumi}.}
  \setcounter{enumi}{8}
  \tightlist
  \item
    Варіаційні методи. Замість рівняння з частковими похідними
    диференціальними розв'язують задачу мінімізації. Виявляється, що
    функція, яка дає мінімум певному виразу (наприклад, загальній
    енергії системи), одночасно є розв'язком початкового рівняння
    похідниху похідних 2.
  \end{enumerate}
\item
  \begin{enumerate}
  \def\labelenumi{\arabic{enumi}.}
  \setcounter{enumi}{9}
  \tightlist
  \item
    Метод розкладу за власними функціями. Розв'язок рівняння часткових
    похідних шукається як послідовність власних функцій. Ці власні
    функції знаходяться як розв'язки так званої задачі власних значень,
    що відповідає початковій задачі для рівняння часткових похідних.
  \end{enumerate}
\end{itemize}

\section{Типи рівнянь у похідних
похідних}\label{ux442ux438ux43fux438-ux440ux456ux432ux43dux44fux43dux44c-ux443-ux43fux43eux445ux456ux434ux43dux438ux445-ux43fux43eux445ux456ux434ux43dux438ux445}

Рівняння в похідних похідних можна класифікувати за багатьма
характеристиками. Класифікація рівнянь важлива, оскільки, як виявилося,
кожен клас має свою загальну теорію та методи розв'язання рівнянь.

Тут ми наведемо шість основних методів класифікації

рівняння.

\begin{enumerate}
\def\labelenumi{\arabic{enumi}.}
\tightlist
\item
  Порядок рівняння. Порядок рівняння --- це найвищий порядок часткових
  похідних, включених до рівняння. Наприклад,
\end{enumerate}

Примітка. Ед.

1) Точніше, це метод лінеаризації нелінійних задач. Загальна теорія
збурень містить лінійні та нелінійні теорії збурень, кожна з яких, у
свою чергу, є регулярною та особливою. --- Ед. Ед.

Ці методи застосовні до спеціальних задач, що зумовлено існуванням
відповідних виразів --- функціоналів для цих задач.

Застосовувався раніше

\[
u_t = u_{xx}
\]

(рівняння другого порядку), \(u_t = u_x\) (рівняння першого порядку),
\(u_1 = uu_{xxx} + \sin x\) (рівняння трьох порядків).

\begin{enumerate}
\def\labelenumi{\arabic{enumi}.}
\setcounter{enumi}{1}
\tightlist
\item
  Кількість змінних. Кількість змінних --- це кількість незалежних
  змінних. Наприклад,
\end{enumerate}

\[
u_t = u_{xx}
\]

(рівняння з двома змінними \(x\) і \(t\) ),
\(u_t = u_{rr} + \frac{1}{r} u_r + \frac{1}{r^2} u_{\theta\theta}\)
(рівняння з трьома змінними \(r\) ,

\begin{enumerate}
\def\labelenumi{\arabic{enumi}.}
\setcounter{enumi}{2}
\tightlist
\item
  Лінійність. Рівняння з частковими похідними можуть бути лінійними та
  нелінійними. У лінійних рівняннях залежна змінна та всі її часткові
  похідні включаються лінійно, зокрема, вони не множаться одна на одну,
  не множаться в квадраті тощо). Точніше, лінійне рівняння другого
  порядку з двома незалежними змінними є рівнянням вигляду
\end{enumerate}

\[
(1.1) Au_{xx} + Bu_{xy} + Cu_{yy} + Du_x + Eu_y + Fu = G
\]

де A, B, C, D, E, F і G --- це константи або задані функції
пояснювальних змінних x і y.

\[
u_{tt} = e^{-t}u_{xx} + \sin t
\]

(лінійне рівняння), \(uu_{xx} + u_t = 0\) (нелінійне рівняння),
\(u_{xx} + yu_{yy} = 0\) (лінійне рівняння), \(xu_x + yu_y + u^2 = 0\)
(нелінійне рівняння).

\begin{itemize}
\tightlist
\item
  \begin{enumerate}
  \def\labelenumi{\arabic{enumi}.}
  \setcounter{enumi}{3}
  \tightlist
  \item
    Однорідність. Рівняння (1.1) називають однорідним, якщо права
    частина \(G(x, y)\) ідентична нулю для всіх x і y. Якщо \(G(x, y)\)
    не ідентичне нулю, то рівняння називають неоднорідним.
  \end{enumerate}
\item
  \begin{enumerate}
  \def\labelenumi{\arabic{enumi}.}
  \setcounter{enumi}{4}
  \tightlist
  \item
    Типи шансів. Якщо коефіцієнти A, B, C, D, E і F рівняння (1.1) є
    сталими, то це рівняння називається рівнянням зі сталими
    коефіцієнтами (інакше рівнянням зі змінними коефіцієнтами).
  \end{enumerate}
\item
  \begin{enumerate}
  \def\labelenumi{\arabic{enumi}.}
  \setcounter{enumi}{5}
  \tightlist
  \item
    Три основні типи лінійних рівнянь. Усі лінійні рівняння з частковими
    похідними другого порядку виду (1.1) належать до одного з трьох
    типів: а) параболічних, б) гіперболічних, в) еліптичних.
  \end{enumerate}
\end{itemize}

Параболічний тип. Рівняння параболічного типу описують процеси
теплопровідності та дифузії і визначаються умовою \(B^2 - 4AC = 0\) .

Гіперболічний тип. Рівняння гіперболічного типу описів

осцилюючі системи та хвильові рухи визначаються умовою \(B^2-4AC>0\) .

Еліпічний тип. Рівняння еліптичного типу описують усталені процеси і
визначаються умовним \(B^2-4AC<0\) . Приклади

\begin{enumerate}
\def\labelenumi{\alph{enumi})}
\tightlist
\item
\end{enumerate}

\[
u_i = u_{xx}
\]

, \(B^2 - 4AC = 0\) (параболічна), (гіперболічна), (гіперболічна),
(гіперболічна \(y > 0\) ), (гіперболічна), (еліптична), (еліптична),
параболічна на \(y = 0\) , гіперболічна на \(y < 0\) .

(У випадку змінних коефіцієнтів тип рівняння може змінюватися від точки
до точки.)

\section{ЗАУВАЖЕННЯ}\label{ux437ux430ux443ux432ux430ux436ux435ux43dux43dux44f}

\begin{itemize}
\tightlist
\item
  \begin{enumerate}
  \def\labelenumi{\arabic{enumi}.}
  \tightlist
  \item
    Загалом, значення \(B^2-4AC\) є функцією незалежних змінних.
    Відповідно, базовий тип рівняння може змінюватися в межах визначення
    рівняння, хоча це не є обов'язковим.
  \end{enumerate}
\item
  \begin{enumerate}
  \def\labelenumi{\arabic{enumi}.}
  \setcounter{enumi}{1}
  \tightlist
  \item
    У рівнянні (1.1) незалежними змінними є x і y. У багатьох задачах
    однією з двох незалежних змінних є час, і рівняння (1.1) можна
    записати через x і t.
  \end{enumerate}
\end{itemize}

\pandocbounded{\includegraphics[keepaspectratio]{images/_page_10_Figure_9.jpeg}}

PHC 1.1. Діаграма класифікації диференціальних рівнянь у похідних
похідних.

\begin{enumerate}
\def\labelenumi{\arabic{enumi}.}
\setcounter{enumi}{2}
\tightlist
\item
  Рис. 1.1. Представлено діаграму класифікації для диференціальних
  рівнянь у похідних похідних.
\end{enumerate}

\section{ЗАВДАННЯ}\label{ux437ux430ux432ux434ux430ux43dux43dux44f}

\begin{itemize}
\tightlist
\item
  \begin{enumerate}
  \def\labelenumi{\arabic{enumi}.}
  \tightlist
  \item
    Класифікуємо наступні рівняння відповідно до всіх характеристик,
    зазначених у днаграмі на рис. 1.1:
  \end{enumerate}
\item
  \begin{enumerate}
  \def\labelenumi{\alph{enumi})}
  \tightlist
  \item
    \(u_t = u_{xx} + 2u_x + u\) ,
  \end{enumerate}
\end{itemize}

\begin{enumerate}
\def\labelenumi{\arabic{enumi})}
\setcounter{enumi}{5}
\tightlist
\item
  \(u_t = u_{xx} + e^{-t}\) , B) \(u_{xx} + 3u_{xy} + u_{yy} = \sin x\)
  , \(u_{tt} = uu_{xxx} + e^{-t}\) .
\end{enumerate}

\begin{itemize}
\tightlist
\item
  \begin{enumerate}
  \def\labelenumi{\arabic{enumi}.}
  \setcounter{enumi}{1}
  \tightlist
  \item
    Скільки розв'язків \(u_1 = u_{xx}\) рівняння існує? Спробуйте знайти
    рішення, як \(u(x, t) = e^{a\lambda + bt}\) .
  \end{enumerate}
\item
  \begin{enumerate}
  \def\labelenumi{\arabic{enumi}.}
  \setcounter{enumi}{2}
  \tightlist
  \item
    Якщо функції \(u_1(x, t)\) і \(u_2(x, t)\) задовольняють рівняння
    (1.1), чи задовольняє сума цих функцій його? Доведи це.
  \end{enumerate}
\item
  \begin{enumerate}
  \def\labelenumi{\arabic{enumi}.}
  \setcounter{enumi}{3}
  \tightlist
  \item
    Ймовірно, найпростіше диференціальне рівняння в похідних ---
  \end{enumerate}
\end{itemize}

\[
\frac{\partial u(x, y)}{\partial x} = 0.
\]

Чи можете ви розв'язати це рівняння? (Знайдіть усі функції, що
задовольняють це рівняння.)

\begin{enumerate}
\def\labelenumi{\arabic{enumi}.}
\setcounter{enumi}{4}
\tightlist
\item
  Що можна сказати про рівняння
\end{enumerate}

\[
\frac{\partial^2 u(x, y)}{\partial x \partial y} = 0
\]

?

Чи можете ви знайти всі розв'язки цього рівняння? Скільки їх? Порівняйте
з кількістю розв'язків звичайного диференціального рівняння

\[
\frac{d^2y}{dx^2} = 0.
\]

\section{ЗАДАЧІ ТИПУ ДИФУЗІЇ (ПАРАБОЛІЧНІ
РІВНЯННЯ)}\label{ux437ux430ux434ux430ux447ux456-ux442ux438ux43fux443-ux434ux438ux444ux443ux437ux456ux457-ux43fux430ux440ux430ux431ux43eux43bux456ux447ux43dux456-ux440ux456ux432ux43dux44fux43dux43dux44f}

МЕТА ЛЕКЦІЇ: Показати, як параболічні рівняння використовуються для
розв'язання задач теплопровідності та дифузії. Наведено кілька прикладів
параболічних рівнянь, а також пояснюється фізичне значення різних членів
\((u_1, u_2, u_3, u_4, u_6)\) , включених у ці рівняння.

На конкретному прикладі розглядається основна ідея зміщеної задачі для
параболічного рівняння. Одна з головних цілей цієї лекції --- надати
читачеві інтуїтивне уявлення

Розуміння параболічних проблем.

Почнемо з розгляньмо просту фізичну задачу і те, як її можна описати
математичною моделлю з рівнянням у похідних похідних.

Потім ми ускладнимо задачу і покажемо, як нові фізичні умови призводять
до появи нових рівнянь у похідних похідних. У цій лекції диференціальні
рівняння не виводяться і не розв'язуються. Про це ми поговоримо пізніше.

\section{Простий експеримент з
теплопровідністю}\label{ux43fux440ux43eux441ux442ux438ux439-ux435ux43aux441ux43fux435ux440ux438ux43cux435ux43dux442-ux437-ux442ux435ux43fux43bux43eux43fux440ux43eux432ux456ux434ux43dux456ux441ux442ux44e}

Припустимо, що ми провели простий експеримент, який розбили на такі
кроки:

КРОК 1. Візьміть достатньо довгий (скажімо, \(L=2\,\mathrm{m}\) )
стрижень (наприклад, мідний) діаметром \(2\,\mathrm{cm}\) , бічна
поверхня (але не кінці) якої термічно ізольована. Замість стрижня можна
використовувати мідну трубку, покриту теплоізоляцією зовні та зсередини.

Іншими словами, тепло може проходити лише через кінці стрижня і не може
проходити через його бічну поверхню.

КРОК 2. Тепер помістимо цей стрижень у пристрій із фіксованою
температурою \(T_{\mathfrak{o}}\) (градуси Цельсія) на достатньо довгий
час, щоб температура всередині стрижня стала

Так само, як і температура пристрою. Для простоти припустимо, що
температура пристрою \(T_0 = 10^{\circ}\) C.

КРОК 3. Давайте витягнемо стрижень з нашого пристрою у певний момент
часу, що зручно взяти як нуль t = 0, і приєднаємо до нього дві термопари
з кінців. Завдання цих елементів --- підтримувати фіксовані температури
\(T_1\) і \(T_2\) на кінцях (наприклад, \(T_1=0^{\circ}\) C і
\(T_2=50^{\circ}\) C). Іншими словами, два термостати постійно
контролюють температуру на кінцях стрижня, і якщо температура
відхиляється від встановлених значень \(T_1\) і \(T_2\) , активуються
потужні нагрівальні або охолоджувальні елементи (пристрої) для
відповідного регулювання температури. Наш експеримент ілюструє RPS, 2 1,

\pandocbounded{\includegraphics[keepaspectratio]{images/_page_13_Figure_4.jpeg}}

Рис. 2.1. Схема експерименту: a --- це термопара, яка підтримує
встановлену температуру \(T_1\) на лівому кінці стрижня; \(\delta\) ---
це термопара, яка підтримує встановлену температуру \(T_2\) на правому
кінці стрижня.

КРОК 4. Тепер ми будемо слідкувати за температурним профілем у стрижні
за допомогою пристрою візуалізації. Чому ми хочемо провести такий
експеримент? Це інше питання, про яке ми поговоримо пізніше.

Це кінець опису нашого експерименту.

Головна мета цієї лекції --- показати, як цю фізичну проблему та її
модифікації можна змоделювати мовою параболічних рівнянь.

\section{Математична модель
теплопровідності}\label{ux43cux430ux442ux435ux43cux430ux442ux438ux447ux43dux430-ux43cux43eux434ux435ux43bux44c-ux442ux435ux43fux43bux43eux43fux440ux43eux432ux456ux434ux43dux43eux441ux442ux456}

Опис нашої фізичної проблеми вимагає трьох типів відносин.

\begin{itemize}
\tightlist
\item
  \begin{enumerate}
  \def\labelenumi{\arabic{enumi}.}
  \tightlist
  \item
    Рівняння з частковими похідними, які описують фізичне явище
    теплопровідності.
  \end{enumerate}
\item
  \begin{enumerate}
  \def\labelenumi{\arabic{enumi}.}
  \setcounter{enumi}{1}
  \tightlist
  \item
    Граничні умови, що описують процеси теплопередачі на межах.
  \end{enumerate}
\item
  \begin{enumerate}
  \def\labelenumi{\arabic{enumi}.}
  \setcounter{enumi}{2}
  \tightlist
  \item
    Початкові твердження, які описують стан системи на початку процесу.
  \end{enumerate}
\end{itemize}

\section{Рівняння
теплопровідності}\label{ux440ux456ux432ux43dux44fux43dux43dux44f-ux442ux435ux43fux43bux43eux43fux440ux43eux432ux456ux434ux43dux43eux441ux442ux456}

Базове одновимірне рівняння теплопровідності записується у вилці

(2.1)

\[
u_t = \alpha^2 u_{xx}, \quad 0 < x < L, \quad 0 < t < \infty.
\]

Це рівняння пов'язує значення \(u_t\) --- швидкість зміни температури з
часом (виміряна в градусах/с)

та величину \(u_{xx}\) --- це увігнутий профілю температури \(u(x, t)\)
(який є мірою різниці між температурою в заданій точці та температурою в
сусідніх точках).

\pandocbounded{\includegraphics[keepaspectratio]{images/_page_14_Figure_14.jpeg}}

Рис. 2.2. Стрілки показують зміну температури відповідно до рівняння
\(u_t = \alpha^2 u_{xx}\) .

Це рівняння буде виведено з базового рівняння збереження тепла в
наступних лекціях, але наразі ми розглянемо його заданим. Це рівняння
каже, що температура \(u\left(x,\ t\right)\) (у певній точці t у певній
точці x стрижня) зростає \((u_t>0)\) або зменшується \((u_t<0)\) залежно
від того, чи є друга похідна \(u_{xx}\) позитивною чи негативною. Фіг.
Рисунок 2.2 ілюструє зміну температури в різних точках на елементі.

Давайте подивимось, як значення \(u_{xx}\) можна інтерпретувати мовою
теплопровідності. Припустимо, ми наблизимо значення \(u_{xx}\) фінальних
постів

\[
u_{xx}(x, t) \cong \frac{1}{\Delta x^2} [u(x + \Delta x, t) - 2u(x, t) + u(x - \Delta x, t)].
\]

Це співвідношення можна переписати у вигляді

\[
u_{\lambda x}(x, t) \simeq \frac{2}{\Delta x^2} \left[ \frac{u(x + \Delta x, t) + u(x - \Delta x, t)}{2} - u(x, t) \right].
\]

Тепер ми можемо дати наступну інтерпретацію значення \(u_{xx}\) .

\begin{enumerate}
\def\labelenumi{\arabic{enumi}.}
\item
  Якщо температура \(u(x, t)\) менша за середнє значення температури у
  двох сусідніх точках, то \(u_{xx} > 0\) (тут загальний тепловий потік
  по осі x додатний).
\item
  Якщо температура \(u(x, t)\) дорівнює середньому значенню температур у
  двох сусідніх точках, то \(u_{xx} = 0\) (тут загальний тепловий потік
  по осі x дорівнює нулю).
\item
  Якщо температура \(u\left(x,\ t\right)\) вища за середню температуру у
  двох сусідніх точках, то \(u_{xx}<0\) (тут --- повний потік
\end{enumerate}

тепла вздовж осі x є від'ємною).

Усе це ілюструється на рис. 2.2. Іншими словами, якщо температура в
точці x більша за середню температуру у двох сусідніх точках
\(x - \Delta x\) і \(x + \Delta x\) , то температура в точці x
зменшиться. Відповідно, точна швидкість зниження температури, а саме
величина \(u_t\) , пропорційна цій різниці. Коефіцієнт пропорційності
\(\alpha^2\) визначається властивостями матеріалу. Ми розглянемо цю
константу детальніше у майбутніх лекціях.

\section{Граничні умови
(ГІ)}\label{ux433ux440ux430ux43dux438ux447ux43dux456-ux443ux43cux43eux432ux438-ux433ux456}

Усі фізичні проблеми характеризуються наявністю певних меж. Отже, щоб
рухатися вперед, ми повинні включити межі в математичну модель, щоб
адекватно описати фізичну проблему. У нашому експерименті граничні умови
(ГІ) дуже прості. Оскільки температура на кінцях x=0 і x=L завжди
залишається сталою і дорівнює \(T_1\) та \(T_2\) , можливо записати

(2.2)

\begin{verbatim}
\left\{ \begin{array}{l} u\left(0,\ t\right) = T_{t}, \\ u\left(L,\ t\right) = T_{2}, \end{array} \right. \quad 0 < t < \infty,
\end{verbatim}

\section{Початкові умови
(NU)}\label{ux43fux43eux447ux430ux442ux43aux43eux432ux456-ux443ux43cux43eux432ux438-nu}

Усі фізичні процеси мають починатися у певний момент часу (зазвичай
приймається як нуль t=0), тому ми повинні встановити фізичні умови саме
в цю точку. Відтоді, як ми почали контролювати температуру з моменту,
коли паличка мала постійну температуру \(T_{\mathrm{ }o}\) , ми

(2.3)

\[
u(x, 0) = T_0, \quad 0 \le x \le L
\]

Тепер ми створили математичну модель експерименту. Система співвідношень
(2.1), (2.2) та (2.3) називається задачею зміщення для рівняння
теплопровідності і зазвичай записується як

(2.4)

\[
\begin{aligned} & (\text{YUII}) \quad u_t = \alpha^2 u_{xx}, & 0 < x < L, \ 0 < t < \infty, \\ & \left( \Gamma \text{Y} \right) \quad \begin{cases} u \ (0, \ t) = T_1, \\ u \ (L, \ t) = T_2, \end{cases} & 0 < t < \infty, \\ & (\text{HY}) \quad u \ (x, \ 0) = T_0, \quad 0 \leqslant x \leqslant L. \end{aligned}
\]

Цікаво, що це зовсім не очевидно: чи існує одна функція \(u(x, t)\), яка
задовольняє умови (2.4), і чи описує ця функція температуру в стрижні?
Отже, наше безпосереднє завдання --- знайти цю функцію \(u(x, t)\) з
(2.4).

Перш ніж завершити цю лекцію, давайте розглянемо деякі варіації основної
проблеми. Почнемо з невеликих змін рівняння теплопровідності.

\section{Деякі рівняння дифузійного
типу}\label{ux434ux435ux44fux43aux456-ux440ux456ux432ux43dux44fux43dux43dux44f-ux434ux438ux444ux443ux437ux456ux439ux43dux43eux433ux43e-ux442ux438ux43fux443}

Теплопередача через бічну поверхню пропорційна різниці температур

Рівняння

\[
u_t = \alpha^2 u_{xx} - \beta (u - u_0)
\]

Описує теплопровідність елемента, враховуючи не лише дифузію теплових
\(\alpha^2 u_{xx}\) вздовж елемента, а й теплопередачу через бічну
поверхню елемента. Відтік тепла ( \(\beta > 0\) ) або його приплив (
\(\beta < 0\) ) пропорційний різниці між температурою u стрижня та
температурою навколишнього середовища \(u_0\) ( \(\beta\) --- це сталий
коефіцієнт пропорційності). Якщо коефіцієнт \(\beta\) великий порівняно
з \(\alpha^2\) , то теплопотік уздовж стрижня буде малим порівняно з

Протікання через бічну поверхню, і тому тепло буде проходити через неї
відповідно до закону \(u_t = -\beta (u - u_a)\) .

У хімії, де величина відіграє роль концентрації, рівняння

\[
u_1 = \alpha^2 u_{xx} - \beta (u - u_0)
\]

означає, що швидкість зміни \(u_i\) кількості речовини залежить як від
дифузії \(\alpha^2 u_{xx}\) (у напрямку осі x), так і від появи (
\(\beta < 0\) ) або розпаду ( \(\beta > 0\) ) речовини в хімічній
реакції, і пропорційна різниці між двома концентраціями u і
\(u_{\theta}\) .

\section{Внутрішнє джерело
тепла}\label{ux432ux43dux443ux442ux440ux456ux448ux43dux454-ux434ux436ux435ux440ux435ux43bux43e-ux442ux435ux43fux43bux430}

Неоднородне рівняння

\[
u_t = \alpha^2 u_{xx} + f(x, t)
\]

відповідає випадку, коли всередині елемента є джерело тепла (розташоване
вздовж усього елемента і діє в будь-який момент t). Наприклад, може бути
так, що дріт з електричним струмом проходить через стрижень і опір є
постійним джерелом тепла \(f(x,t)\) = K.

\section{Рівняння конвективної
дифузії}\label{ux440ux456ux432ux43dux44fux43dux43dux44f-ux43aux43eux43dux432ux435ux43aux442ux438ux432ux43dux43eux457-ux434ux438ux444ux443ux437ux456ux457}

Припустимо, що домішка поширюється вздовж потоку, що рухається зі
швидкістю v. Очевидно, що концентрація \(u\left(x,\,t\right)\) домішок
змінюється залежно від x (позитивний напрямок осі x обирається вздовж
потоку) і часу t. Швидкість зміни концентрації \(u_t\) визначається
рівнянням конвективної дифузії

\[
u_t = \alpha^2 u_{xx} - v u_x.
\]

Термін \(\alpha^2 u_{xx}\) описує внесок дифузії, а \(vu_x\) ---
конвективний компонент. Що важливіше, дифузія або конвекція залежить від
співвідношення \(\alpha^2\) до v. Димові частинки піднімаються
конвективно через рух гарячого повітря і одночасно розсіюють через
вихорові рухи повітря. Окрім цих модифікацій рівняння теплопровідності,
можливо змінювати граничні умови на кінцях стрижня відповідно до
фізичної ситуації. Деякі з цих змін ми обговоримо у третій лекції.

\section{ЗАУВАЖЕННЯ}\label{ux437ux430ux443ux432ux430ux436ux435ux43dux43dux44f-1}

Рівняння теплопровідності \(u_t = \alpha^2(x) u_{xx}\) з змінним
коефіцієнтом \(\alpha(x)\) відповідає задачі, в якій дифузія залежить
від x (тобто матеріал є гетерогенним). Наприклад, мідь і сталь---

\pandocbounded{\includegraphics[keepaspectratio]{images/_page_18_Picture_4.jpeg}}

PHC. 2.3.

Пластини контактують, як показано на рис. 2.3. Припустимо, що
температура \(u(0,t)=0^{\circ}\) C підтримується на лівому боці міді, а
температура \(u(L,t)=20^{\circ}\) C --- з правого боку сталі, тоді
рівняння, що описує теплопровідність у цій системі, виглядає так:

\[
u_t = \alpha^2(x) \, u_{_{\! A \, X}}, \quad 0 < x < L,
\]

де

\[
\alpha(x) = \begin{cases} \alpha_1 & \text{(коэффициент диффузии в меди)}, \quad 0 < x < L/2, \\ \alpha_2 & \text{(коэффициент диффузии в стали)}, \quad L/2 < x < L. \end{cases}
\]

\# ЗАВДАННЯ

\begin{enumerate}
\def\labelenumi{\arabic{enumi}.}
\tightlist
\item
  Якщо початкова температура в стрижні
\end{enumerate}

\[
u(x, 0) = \sin \pi x, \quad 0 \leqslant x \leqslant 1,
\]

а граничні умови такі:

\[
u(0, t) = 0,
\]

\(u(1, t) = 0,\) Тоді як поводитиметься температура в стрижні
\(U(x, t)\) при T \textgreater{} 0? 2. Припустимо, що в стрижні є
постійне внутрішнє джерело тепла, так що рівняння теплопровідності в
стрижні записується як

\[
u_t = \alpha^2 u_{xx} + 1, \quad 0 < x < 1.
\]

Припустимо, що ми зафіксували температуру на межі з значеннями
\(u(0, t)\) = 0 і \(u(1, t)\) = 1. Якою буде стаціонарна температура у
стрижні? Іншими словами, який c-незалежний розподіл \(U(x, t)\)
перешкоджує \(u(x, t)\)?

\begin{enumerate}
\def\labelenumi{\arabic{enumi}.}
\setcounter{enumi}{2}
\tightlist
\item
  Припустимо, що металевий стрижень втрачає тепло через бічну поверхню
  згідно з рівнянням
\end{enumerate}

\[
u_1 = \alpha^2 u_{xx} - \beta u
\]

, \(0 < x < 1\) .

І припустимо, що ми підтримуємо кінці стрижня при температурах
\(U(0, T)\) = 1, \(U(1, T)\) = 1. Знайдіть стаціонарний розподіл
температури по стрижню. Побудуйте графік цього розподілу. У яких
моментах відбувається тепловий потік?

\begin{enumerate}
\def\labelenumi{\arabic{enumi}.}
\setcounter{enumi}{3}
\tightlist
\item
  Припустимо, що термоізольований металевий стрижень довжини L = 1 має
  початкову температуру \(\sin(3\pi x)\) , а фіксовану температуру --- 0
  і \(10\,^{\circ}\) C на лівому та правому кінцях.
\end{enumerate}

\section{Лекція 3}\label{ux43bux435ux43aux446ux456ux44f-3}

\section{ГРАНИЧНІ УМОВИ В ЗАДАЧАХ ДИФУЗІЙНОГО
ТИПУ}\label{ux433ux440ux430ux43dux438ux447ux43dux456-ux443ux43cux43eux432ux438-ux432-ux437ux430ux434ux430ux447ux430ux445-ux434ux438ux444ux443ux437ux456ux439ux43dux43eux433ux43e-ux442ux438ux43fux443}

МЕТА ЛЕКЦІЇ: Показати, як виникають різні граничні умови у задачах
теплопровідності та дифузії, а також ввести важливе поняття теплового
потоку.

Розглядаються три найважливіші типи граничних умов:

\begin{itemize}
\tightlist
\item
  \begin{enumerate}
  \def\labelenumi{\arabic{enumi}.}
  \tightlist
  \item
    u = \(g(t)\) (температура задається на межі),
  \end{enumerate}
\item
  \texttt{2.\ \textbackslash{}\ \textbackslash{}frac\{\textbackslash{}partial\ u\}\{\textbackslash{}partial\ n\}\ +\ \textbackslash{}lambda\ u\ =\ g\textbackslash{}left(t\textbackslash{}right)}
  (встановлення температури навколишнього середовища);
\end{itemize}

n --- вектор зовнішньої нормалі до межі),
\texttt{3.\ \textbackslash{}frac\{\textbackslash{}partial\ u\}\{\textbackslash{}partial\ n\}\ =\ g(t)}
(визначений тепловий потік через межу).

У задачах теплопровідності зазвичай використовують три основні типи
граничних умов. Ця лекція наводить приклади того, як різні фізичні умови
експерименту призводять до різних типів граничних умов.

\section{Граничні умови першого виду (температура задається на
межі)}\label{ux433ux440ux430ux43dux438ux447ux43dux456-ux443ux43cux43eux432ux438-ux43fux435ux440ux448ux43eux433ux43e-ux432ux438ux434ux443-ux442ux435ux43cux43fux435ux440ux430ux442ux443ux440ux430-ux437ux430ux434ux430ux454ux442ux44cux441ux44f-ux43dux430-ux43cux435ux436ux456}

Розглянемо тепловий потік в одномірному стрижні, показаний на рис. 3 1,
і припустимо, що температура кінців змінюється відповідно до законів
\(g_1(t)\) і \(g_2(t)\) .

\pandocbounded{\includegraphics[keepaspectratio]{images/_page_20_Picture_4.jpeg}}

Рис. 31. Температура встановлюється на межі.

Як згадувалося в попередній лекції, кожен пристрій, що підтримує
встановлену температуру на кінцях, складається з термостата та
нагрівального елемента для підтримки відповідних теплових умов. Проблеми
з граничними умовами першого виду дуже поширені. У деяких випадках
складність полягає у пошуку контролерів граничних умов, тобто температур
\(g_1(t)\) та \(g_2(t)\) меж, які призводять до певної зміни температури
всередині стрижня. Наприклад, у металургії часто необхідно обирати
граничні умови контролю, щоб температура металу всередині печі
змінювалася з часом, а температурний градієнт був невеликим.

\pandocbounded{\includegraphics[keepaspectratio]{images/_page_20_Picture_7.jpeg}}

Рис. 3.2. Коливання температури на межі.

Граничні умови першого виду також виникають у багатовимірних задачах
теплопровідності. Як приклад можна навести цікаву задачу визначення
температурного поля всередині кругового диска з радіусом \emph{R} (див.
рис. 3.2), коли гранична температура задається у полярних координатах
формулою \(u(R, \theta, t) = \cos t \sin \theta\) .

Звісно, щоб розв'язати задачу, потрібно знати початкову температуру, але
в цьому випадку початкові умови перестають впливати на розв'язок через
дуже короткий час, і отримана температура всередині диска визначається
лише значеннями температури на межі.

\section{Граничні умови другого виду (задана температура навколишнього
середовища)}\label{ux433ux440ux430ux43dux438ux447ux43dux456-ux443ux43cux43eux432ux438-ux434ux440ux443ux433ux43eux433ux43e-ux432ux438ux434ux443-ux437ux430ux434ux430ux43dux430-ux442ux435ux43cux43fux435ux440ux430ux442ux443ux440ux430-ux43dux430ux432ux43aux43eux43bux438ux448ux43dux44cux43eux433ux43e-ux441ux435ux440ux435ux434ux43eux432ux438ux449ux430}

Припустимо, що ми знову розглядаємо термічно ізольований мідний
стержень, але тепер відмовляємося від заданого температурного режиму на
кінцях стрижня. Натомість давайте наблизимо кінці стрижня до контакту з
двома середовищами. Нехай температура одного з них \(g_1(t)\) , а іншого
\(g_2(t)\) . Іншими словами, припустимо, що лівий кінець стрижня
поміщений у контейнер із рідиною, температура якої змінюється відповідно
до закону \(g_1(t)\) . А правий кінець поміщається в іншу рідину з
температурою \(g_2(t)\) (рис. 3, 3)

\pandocbounded{\includegraphics[keepaspectratio]{images/_page_21_Picture_6.jpeg}}

Рис. 3.3. Конвективна теплопередача через межі, a --- це рідина при
температурі \(g_1(t)\) , \(\delta\) --- рідина при температурі
\(g_2(t)\) .

Задаючи граничні умови такого типу, ми не можемо вважати, що граничні
температури стрижня збігаються з температурами рідини. Але ми знаємо
закон Ньютона: якщо температура одного з кінців стрижня менша за
температуру відповідної рідини, тепло потече в стрижень зі швидкістю,
пропорційною різницям температур. Іншими словами, для одномірного
стрижня з межами x=0 і x=1 закон теплопередачі Ньютона формулюється так:

(3.1)

\begin{verbatim}
\{
\end{verbatim}

вихідний тепловий потік на \(x=0\) дорівнює
\(h\left[u\left(0,\,t\right)-g_{1}\left(t\right)\right]\) , вихідний
тепловий потік при \(x=L\) дорівнює
\(h\left[u\left(L,\,t\right)-g_{2}\left(t\right)\right]\) ,

де h --- коефіцієнт теплопередачі, який показує, скільки калорій
проходить через межу за одну секунду при різниці температур у один
градус. Вихідний потік дорівнює кількості кало-

Пройшов крізь палочний спис за одну секунду. Зверніть увагу, що вихідний
тепловий потік буде позитивним на кінці стрижня, де температура стрижня
вища за температуру навколишнього середовища. Рівняння (3.1) разом із
відомим законом теплопровідності Фур'є тепер можна використовувати для
отримання Р и с

\pandocbounded{\includegraphics[keepaspectratio]{images/_page_22_Figure_3.jpeg}}

. 3.4. Ілюстрація закону Фур'є: \(\frac{\partial u}{\partial n} > 0\)
означає, що тепло надходить у цю область;
\(\frac{\partial u}{\partial n} < 0\) означає, що тепло виходить із цієї
зони; n --- напрямок зовнішньої нормалі;
\(\frac{\partial u}{\partial n} -\) Зміна температури в напрямку N;
\(\frac{\partial u}{\partial n} = -\frac{\partial u}{\partial (-n)}\) є
похідною у напрямку внутрішньої нормальної --- це похідна, взята
\(\mathbf{e}\) протилежному знаку протилежного знака пормалі.

граничні умови. Закон Фур'є дає друге представлення (перше міститься у
(3.1)) для вихідного теплового потоку. Прирівнюючи ці два вирази до
потоку, ми отримуємо граничні умови, які шукаємо. Сформулюємо закон
Фур'є (встановлений експериментально):

(3.2) Тепловий потік, що проходить через межу площі, пропорційний
нормальній похідній температури у напрямку внутрішньої норми.

Цей закон стверджує, що якщо температура швидко підвищується у напрямку
зовнішньої температури до межі D (рисунок 3.4), то потік буде текти з
навколишнього середовища до області D.

У нашій одновимірній задачі закон Фур'є має вигляд:

(3.3)

\[
\begin{cases} \text{вытекающий поток тепла при } x=0 \text{ равен } k \frac{\partial u}{\partial x}, \\ \text{вытекающий поток тепла при } x=' \text{ равен } -k \frac{\partial u}{\partial x}, \end{cases}
\]

де k --- теплопровідність матеріалу, яка є мірою того, наскільки добре
матеріал проводить тепло. Погано провідні матеріали мають коефіцієнт
близький до нуля в одиницях GHS, тоді як мідь та алюміній мають значення
значно вище нуля.

\pandocbounded{\includegraphics[keepaspectratio]{images/_page_23_Figure_5.jpeg}}

Рис. 3.5. Ще одна ілюстрація закону Фур'є.

Закон Фур'є (3.3) фактично описує теплопровідність всюди всередині
стрижня, а не лише на його межі; Наприклад, (див. рис. 3.5).

(3.4) Потік, що проходить через точку

x

зліва направо \(=-k\frac{\partial n}{\partial x}\) Закон Фур'є (3.4)
стверджує: якщо \(u_x < 0\) , то в точці \(x_0\) тепло тече зліва
направо, якщо \(u_x > 0\) , то в точці \(x_0\) тепло тече справа наліво.
(Тепло завжди переходить від вищих температур до нижчих.)

Нарешті, якщо використовувати вирази (3.1) і (3.3) для теплового потоку,
то отримуємо бажані граничні умови для задачі, описаної на рис. 3.3, у
чисто математичній формі

\[
(\Gamma Y) \quad \begin{cases} \frac{\partial u\left(0,\,t\right)}{\partial x} = \frac{h}{k}\left[u\left(0,\,t\right) - g_{1}\left(t\right)\right], \\ \frac{\partial u\left(L,\,t\right)}{\partial x} = -\frac{h}{k}\left[u\left(L,\,t\right) - g_{2}\left(t\right)\right]. \end{cases} \quad 0 < t < \infty,
\]

Дуже часто константа h/k позначається як \(\lambda\) , а граничні умови
теплового потоку на межі записуються у вигляді

(3.5)

\[
u_{x}(0, t) = \lambda [u(0, t) - g_{1}(t)], u_{x}(L, t) = -\lambda [u(L, t) - g_{2}(t)].
\]

Для випадку вищих розмірностей отримують подібні граничні умови.
Наприклад, якщо межа кругового диска осяяна рухомою рідиною з
температурою \(g(\theta,t)\) , то граничні умови записуються як

\[
\frac{\partial u}{\partial r}(R, \theta, t) = -\frac{h}{k} [u(R, \theta, t) - g(\theta, t)].
\]

Тут \(\frac{\partial u}{\partial r}(R, \theta, t)\) позначає зовнішню
нормальну похідну (у позитивному напрямку r-осі) значення u, обчислену в
точці \((R, \theta)\) межі. Ми назвемо такі граничні умови лінійними
(оскільки вони лінійні в U і \(u_r\) ), але неоднорідними, оскільки
функція \(g(\theta, t)\) знаходиться праворуч.

\section{Граничні умови третього виду (враховуючи потік; зокрема,
враховується випадок термічно ізольованих
меж)}\label{ux433ux440ux430ux43dux438ux447ux43dux456-ux443ux43cux43eux432ux438-ux442ux440ux435ux442ux44cux43eux433ux43e-ux432ux438ux434ux443-ux432ux440ux430ux445ux43eux432ux443ux44eux447ux438-ux43fux43eux442ux456ux43a-ux437ux43eux43aux440ux435ux43cux430-ux432ux440ux430ux445ux43eux432ux443ux454ux442ux44cux441ux44f-ux432ux438ux43fux430ux434ux43eux43a-ux442ux435ux440ux43cux456ux447ux43dux43e-ux456ux437ux43eux43bux44cux43eux432ux430ux43dux438ux445-ux43cux435ux436}

Жоден потік не проходить через термоізольовані межі, тому нормальна
похідна (зовнішня чи внутрішня) має зрівнятися з нулем на межі, оскільки
потік пропорційний нормальній похідній. У випадку одномірного елемента з
ізольованими кінцями x=0 і x=L, граничні умови мають вигляд

\[
u_{\lambda}(0, t) = 0,
\]

\(u_{\lambda}(L, t) = 0.\) \(0 < t < \infty,\) У двовимірній області
ізоляція межі означає, що нормальна похідна температури на межі
перетворюється на нуль. Припустимо, наприклад, що круглий диск термічно
ізольований уздовж межі. Тоді граничну умову слід записати так:
\(u_r(R, \theta, t) = 0\) для всіх \(0 \le \theta < 2\pi\) і всіх
\(0 < t < \infty\) .

З іншого боку, якщо вказати кількість тепла, що проходить через межу
нашого диска, то граничні умови набувають вигляду

\[
u_r(R, \theta, t) = f(\theta, t),
\]

де \(f(0, t)\) --- це кількість тепла, що переходить у круглий диск із
зовнішнього джерела.

Тепер проілюструємо різні типи граничних умов.

\section{Типові граничні умови для задачі одномірної
теплопровідності}\label{ux442ux438ux43fux43eux432ux456-ux433ux440ux430ux43dux438ux447ux43dux456-ux443ux43cux43eux432ux438-ux434ux43bux44f-ux437ux430ux434ux430ux447ux456-ux43eux434ux43dux43eux43cux456ux440ux43dux43eux457-ux442ux435ux43fux43bux43eux43fux440ux43eux432ux456ux434ux43dux43eux441ux442ux456}

Припустимо, що у нас є мідний стрижень довжиною 200 см, бікова поверхня
якого термічно ізольована, а початкова температура становить 0 °C.
Припустимо також, що верхній кінець стрижня x=0 термічно ізольований, а
нижній кінець x=200 змивається рухомою водою, яка має постійну
температуру \(g_2(t)=20\) °C (див. рис. 3.6).

\pandocbounded{\includegraphics[keepaspectratio]{images/_page_25_Picture_6.jpeg}}

Рис. 3.6. Змішане завдання.

Математична модель цієї задачі описується наступними чотирма
відношеннями:

\[
(3.6) \qquad \begin{array}{llll} (\mathrm{Y}\mathrm{H}\Pi) & u_t = \alpha^2 u_{xx}, & 0 < x < 200, & 0 < t < \infty, \\ u_x (0, \, t) = 0, & \\ u_x (200, \, t) = -\frac{h}{k} \left[ u \, (200, \, t) - 20 \right], & 0 < t < \infty, \\ (\mathrm{HY}) & u \, (x, \, 0) = 0 \, ^{\mathrm{o}}\mathrm{C}, & 0 \leqslant x \leqslant 200, & \end{array}
\]

де \(\alpha^2 = 1,16\) см 2/с --- теплова дифузія міді;

k = 0,93 кал/см·с·°C --- теплопровідність міді;

h --- коефіцієнт теплопередачі. Визначення цього співвідношення ---
велика проблема. Він вимірюється швидкістю виходу тепла від другого
кінця стрижня до води. Це залежить від швидкості потоку води, стану
поверхні, яку слід змивати тощо.

Читач може спробувати провести експеримент, щоб визначити

ділення значення цього коефіцієнта.

\section{ЗАУВАЖЕННЯ}\label{ux437ux430ux443ux432ux430ux436ux435ux43dux43dux44f-2}

\begin{enumerate}
\def\labelenumi{\arabic{enumi}.}
\tightlist
\item
  Типова задача внутрішньої теплопровідності для квадрата показана на
  рис. 3.7.
\end{enumerate}

У цій задачі початкова температура \(u\left(x,\,y,\,t\right)\) при t = 0
задається всередині квадрата, і рівняння з похідними з частковими має
вигляд Р и с

\pandocbounded{\includegraphics[keepaspectratio]{images/_page_26_Figure_9.jpeg}}

. 3.7. Типові граничні умови для задачі дифузії в квадраті: a ---
кількість тепла, \(f_1(t)\) проходить через межу; b --- термоізольована
межа; a --- це навколишня температура Равії \(10^{\circ}\) C; z ---
температура підтримується рівною \(10^{\circ}\) C.

а граничні умови будуть виконані на \(0 < t < \infty\) і визначатимуть
температуру в будь-який момент. Якою б не була ця температура, вона
повинна відповідати граничним умовам, зазначеним на рис. 3.7.

\begin{enumerate}
\def\labelenumi{\arabic{enumi}.}
\setcounter{enumi}{1}
\tightlist
\item
  Зверніть увагу, що гранична умова
\end{enumerate}

\[
u_r(R, \theta, t) = -\frac{h}{k} [u(R, \theta, t) - g(\theta, t)]
\]

На колі не означає, що гранична температура дорівнює \(g(\theta,t)\) ,
але якщо коефіцієнт теплопередачі h великий, то ці умови фактично
еквівалентні вимогі \(g(\theta,t)\) температури на межі.

\section{ЗАВДАННЯ}\label{ux437ux430ux432ux434ux430ux43dux43dux44f-1}

\begin{enumerate}
\def\labelenumi{\arabic{enumi}.}
\item
  Намалюйте контури графіків для розв'язання змішаної задачі (3.6) для
  різних моментів часу. Чи відповідають ваші графи визначеним граничним
  умовам? Яка стаціонарна температура стрижня? Очевидно, ваша відповідь
  базується на інтуїції?
\item
  Як би ви інтерпретували наступне змішане завдання?
\end{enumerate}

\[
\begin{array}{lll} (\text{YHII}) & u_t = \alpha^2 u_{xx}, & 0 < x < 1, & 0 < t < \infty, \\ (\text{FY}) & \begin{cases} u\left(0,\,t\right) = 0, \\ u_x\left(1,\,t\right) = 1, \\ u\left(x,\,0\right) = \sin\left(\pi x\right), & 0 \leqslant x \leqslant 1. \\ \end{cases} \end{array}
\]

Чи можете ви графічно показати рішення цієї проблеми в різні моменти
часу? Чи буде це рішення, як правило, стаціонарним? Це очевидно?

\begin{enumerate}
\def\labelenumi{\arabic{enumi}.}
\setcounter{enumi}{2}
\tightlist
\item
  Яку фізичну інтерпретацію ви можете дати проблемі
\end{enumerate}

\[
\begin{array}{llll} (\text{УЧП}) & u_t = \alpha^2 u_{xx}, & 0 < x < 1, & 0 < t < \infty, \\ (\text{ГУ}) & \begin{cases} u_x \, (0, \, t) = 0, & 0 < t < \infty, \\ u_x \, (1, \, t) = 0, & 0 < t < \infty, \end{cases} \\ (\text{НУ}) & u \, (x, \, 0) = \sin (\pi x), & 0 \leqslant x \leqslant 1? \\ \end{array}
\]

Чи можете ви графічно показати рішення цієї проблеми в різні моменти
часу?

Що можна сказати про стабільну температуру?

\begin{enumerate}
\def\labelenumi{\arabic{enumi}.}
\setcounter{enumi}{3}
\tightlist
\item
  Припустимо, що бічна поверхня металевого стрижня не ізольована і має
  початкову температуру 20°C, але температура одного кінця стає 50°C.
  Потім стрижень промивають рідиною з температурою 30°C. Як виглядає
  змішана задача, що відповідає цьому випадку?
\end{enumerate}

\section{ВИВЕДЕННЯ РІВНЯННЯ
ТЕПЛОПРОВІДНОСТІ}\label{ux432ux438ux432ux435ux434ux435ux43dux43dux44f-ux440ux456ux432ux43dux44fux43dux43dux44f-ux442ux435ux43fux43bux43eux43fux440ux43eux432ux456ux434ux43dux43eux441ux442ux456}

МЕТА ЛЕКЦІЇ: Показати, як можна отримати одномірне рівняння
теплопровідності

\[
u_t = \alpha^2 u_{xx} + f(x, t),
\]

спираючись лише на фундаментальний закон збереження кількості тепла.
Розглядається залежність швидкості теплопередачі від основних
термофізичних параметрів: теплопровідності, теплоємності та густини
матеріалу. Розглядаються різні модифікації базового рівняння
теплопровідності.

У кожній галузі знань існує набір базових принципів, правильність яких
очевидна. Усі інші твердження мають бути виведені з цих базових
принципів. Звісно, те, що очевидно одній людині, може здатися сумнівним
для іншої. Історія будь-якої науки полягає у постійному пошуку більш
фундаментальних принципів, що лежать в основі цієї науки. Основні
принципи мають бути настільки універсальними, щоб ніхто не сумнівався.

? --- Припущення A --- Припущення B --- Припущення C --- Припущення D
--- ? Рис. 4.1. Аксіоматичний метод.

Наприклад, одна людина може вважати, що всі факти певної науки можна
вивести з одного базового припущення, яке ми позначимо як B. На основі
припущення B вона або вона може довести теорему C, що призведе до
теореми D, а яка, у свою чергу, призведе до доведення багатьох інших
теорем (див. рисунок 4.1). Ось як виглядає прогрес у науці --- дослідник
може отримувати все більше нових результатів. Найвидатніші фізики,
хіміки та біологи або саме на цьому шляху.

З іншого боку, замість доведення нових теорем можна спробувати, якщо
можливо, знайти новий принцип (назвемо його A), який є більш
фундаментальним за B. Припущення A природно вважається більш
фундаментальним, ніж B.

Якщо B можна вивести з A. Отже, якщо можна знайти новий принцип A, тоді
межа наших знань зміщується до сфери більш фундаментальних ідей про
певну науку. У теорії теплопровідності основним принципом є закон
збереження енергії (теплова енергія). Усі інші твердження випливають із
цього базового принципу (рис. 4.2).

Закон збереження енергії

\[
u_t = \alpha^2 u_{xx} + f(x,t)
\]

Інші властивості теплопровідності Рис. 4.2. Енергозбереження: наріжний
камінь теплопровідності.

Звісно, ми могли б забути цю лекцію і використати рівняння
теплопровідності як фундаментальне поняття (дехто може так подумати,
якщо це очевидно), але тоді ми обдуримо більш серйозних студентів, адже
закон збереження вважається найфундаментальнішим принципом у науці.
Науковці, моделюючи свої задачі, часто записують відношення, що
виражають закони збереження, а потім переписують їх у вигляді рівнянь у
похідних похідних.

Повернемося до головної мети нашої лекції --- вивести рівняння
теплопровідності з рівняння збереження енергії.

\section{Виведення рівняння
теплопровідності}\label{ux432ux438ux432ux435ux434ux435ux43dux43dux44f-ux440ux456ux432ux43dux44fux43dux43dux44f-ux442ux435ux43fux43bux43eux43fux440ux43eux432ux456ux434ux43dux43eux441ux442ux456-1}

Розглянемо однорідний стрижень довжини L, щодо якого зробимо такі
припущення:

\begin{enumerate}
\def\labelenumi{\arabic{enumi}.}
\tightlist
\item
  Стрижень виготовлений з одного однорідного провідного матеріалу.
\end{enumerate}

\pandocbounded{\includegraphics[keepaspectratio]{images/_page_29_Figure_10.jpeg}}

Рис. 4.3. Тонкий теплопровідний стержень.

\begin{enumerate}
\def\labelenumi{\arabic{enumi}.}
\setcounter{enumi}{1}
\item
  Бікова поверхня стрижня термічно ізольована (тепло може поширюватися
  лише вздовж осі x).
\item
  Стрижень тонка, тобто температура всіх точок у кожному поперечному
  перерізі стрижня є сталою. Якщо розглядати частину вудки на сегменті
  \([x, x+\Delta x]\) і використовувати
\end{enumerate}

закон збереження кількості тепла, тоді можна написати:

(4.1) Загальна зміна кількості тепла на сегменті \([x, x + \Delta x] =\)
дорівнює загальній кількості тепла, що пройшло через межі + Нульова
кількість тепла, утвореного всередині сегмента \([x, x + \Delta x]\) .

Оскільки загальна кількість тепла (у калоріях) всередині сегмента
\(\{x, x + \Delta x\}\) в будь-який момент часу t обчислюється формулою:

Загальна кількість тепла всередині сегмента
\([x, x+\Delta x] = \int\limits_{-\infty}^{x+\Delta x} coAu\) (s, t) ds,

де c --- питома теплоємність матеріалу (показує здатність матеріалу
накопичувати тепло),

Щільність матеріалу,

A --- це площа поперечного перерізу стрижня.

\textbf{va}con збереження енергії (4.1) можна надати у такій
математичній формі:

(4.2)

\[
\frac{d}{dt} \int_{x}^{x+\Delta x} c\rho Au(s, t) ds = c\rho A \int_{x}^{x+\Delta x} u_{t}(s, t) ds = eA \left[ u_{x}(x+\Delta x, t) - u_{x}(x, t) \right] + A \int_{x+\Delta x}^{x+\Delta x} f(s, t) ds.
\]

\textbf{де} k --- теплопровідність матеріалу (вказує на здатність
матеріалу проводити тепло), \(f(x, t)\) --- об'ємна потужність
зовнішнього джерела тепла (калорії/см·с).

Задача полягає в тому, щоб записати рівняння (4.2) у формі, яка не
містить інтегралів. Щоб розв'язати цю задачу, нагадаємо читачеві теорему
про середнє значення з ходу інтегральної нечисельності.

\section{Теорема про
значення}\label{ux442ux435ux43eux440ux435ux43cux430-ux43fux440ux43e-ux437ux43dux430ux447ux435ux43dux43dux44f}

Якщо функція \(f(x)\) є неперервною на сегменті {[}a, b{]}, то існує
принаймні одна точка \(\xi \in [a, b]\) така, що

\[
\int_{a}^{b} f(x) dx = f(\xi) (b-a).
\]

Застосовуючи цей результат до рівняння (4.2), ми отримуємо наступний
результат

\[
coAu_{t}(\xi, t) \Delta x = kA\{u_{x}(x + \Delta x, t) - u_{x}(x, t)\} + Af(\xi, t) \Delta x,
x < \xi < x + \Delta x,
\]

або

\[
u_t(\xi, t) = \frac{k}{c\rho} \left\{ \frac{u_x(x + \Delta x, t) - u_x(x, t)}{\Delta x} \right\} + \frac{1}{c\rho} f(\xi, t).
\]

Спрямовуючи \(\Delta\) x до 0, отримуємо потрібне рівняння

(4.3)

\[
u_t(x, t) = \alpha^2 u_{xx}(x, t) + F(x, t),
\]

де

\[
\alpha^{2}=\frac{k}{c\rho}
\]

--- коефіцієнт теплової дифузії,
F\left(x,~t\right)\(=\frac{1}{c\rho}f\left(x,\ t\right)\)--- густина
джерел тепла.

Ми могли б зупинитися на цьому. Однак ми також розглянемо випадок, коли
бічна поверхня стрижня не термічно ізольована. Припустимо, що кількість
теплового потоку через бічну поверхню стрижня в цьому випадку
пропорційна різниці між температурою стрижня \(u(x,t)\) і температурою
навколишнього середовища, яка залишається сталою і дорівнює нулю. У
цьому випадку закон збереження кількості тепла веде до рівняння

\[
(4.4) u_t = \alpha^2 u_{xx} - \beta u + F(x, t),
\]

де \(\beta\) --- коефіцієнт пропорційності потоку через бічну поверхню.

\section{ЗАУВАЖЕННЯ}\label{ux437ux430ux443ux432ux430ux436ux435ux43dux43dux44f-3}

\begin{enumerate}
\def\labelenumi{\arabic{enumi}.}
\tightlist
\item
  Як зазначалося вище, константа k називається теплопровідністю речовини
  і чисельно дорівнює кількості тепла (у калоріях), що проходить за
  секунду через пластину товщиною 1 см і площею поперечного перерізу 1
  см³ при різниці температур на протилежних гранях пластини рівною 1°C.
  Числові значення теплопровідності для різних матеріалів можна знайти в
  довідниках з фізики або хімії. Типові значення теплопровідності
  варіюються від 1 (для міді) до значень, близьких до нуля для хороших
  теплоізоляторів. Якщо стрижень зроблений із однорідного матеріалу, то
  значення k не залежить від x. Для деяких матеріалів значення k
  залежить від температури, і тому рівняння теплопровідності в цьому
  випадку стає нелінійним:
\end{enumerate}

\[
u_{t} = \frac{1}{c\rho} \frac{\partial}{\partial x} \{ k (u) | u_{x} \}.
\]

Однак найчастіше теплопровідність k дуже слабо впливає на температуру u,
і цю нелінійність можна ігнорувати.

\begin{enumerate}
\def\labelenumi{\arabic{enumi}.}
\setcounter{enumi}{1}
\item
  Значення c зазвичай називають питомою теплоємністю речовини. Він
  характеризує здатність речовини зберігати теплову енергію. Наприклад,
  запечена картопля має високу теплоємність, оскільки може зберігати
  велику кількість тепла на одиницю маси (саме тому вона зберігає високу
  температуру так довго). Чисельна теплоємність дорівнює кількості
  калорій тепла, які потрібно передати 1 граму речовини, щоб підвищити
  її температуру на 1°C. Для більшості розглянутих задач значення c
  можна вважати сталим, не записуючи ні x, ні i. Значення питомої
  теплоємності для конкретних матеріалів можна знайти у фізичних
  довідниках.
\item
  Одиниці вимірювання основних величин, включених у рівняння
\end{enumerate}

теплопровідність (у системі GHS) наступна:

i-температура (градуси Цельсія), \(u_t\) --- швидкість зміни температури
(°C/s), \(u_x\) --- нахил температурної кривої (°C/см), \(u_{xx}\) ---
опуклість температурної кривої (°C/см²),

\begin{verbatim}
c = щільність тепла (калор/г \cdot °C)

k — теплопровідність (кал/см \cdot при \cdot °C), \rho — щільність (г/<sup>см3</sup>), $\alpha$^2 — це теплова дифузія (см<sup>²</sup>/с).

4. Слід зазначити, що теплова дифузія матеріалу $\alpha$^2 = \frac{k}{\epsilon \rho} прямо пропорційна теплопровідності матеріалу і обернено пропорційна густині \rho та питому теплоємності \epsilon ; Ми сподіваємося, що цей зв'язок відповідає інтуїтивним ідеям читача.
\end{verbatim}

\section{ЗАВДАННЯ}\label{ux437ux430ux432ux434ux430ux43dux43dux44f-2}

\begin{verbatim}
1. Підставити розмірність кожної з u,\ u_t,\ \dots величин у рівняння
\end{verbatim}

\[
u_t = \alpha^2 u_{xx} - \beta u
\]

і переконайтеся, що всі терміни мають однаковий розмір
\{\}\^{}\{\circ\}\mathrm{C}/c .

\begin{enumerate}
\def\labelenumi{\arabic{enumi}.}
\setcounter{enumi}{1}
\tightlist
\item
  Підставити розмірність кожної з величин у рівняння
\end{enumerate}

\[
u_t = \alpha^2 u_{xx} - v u_x,
\]

де \upsilon вимірюється в одиницях швидкості, і переконайтеся, що всі
терми мають однакову розмірність.

3 Отримаємо рівняння теплопровідності

\[
u_{t} = \frac{1}{c\rho} \frac{\partial}{\partial x} [h(x) u_{x}] + f(x, t)
\]

Для випадку, коли теплопровідність \(k\left( x\right)\) залежить від
координати X.

\begin{enumerate}
\def\labelenumi{\arabic{enumi}.}
\setcounter{enumi}{3}
\tightlist
\item
  Припустимо, що \(u(x, t)\) --- це значення концентрації речовини в
  потоку, що рухається зі швидкістю v. Припустимо, що концентрація
  змінюється внаслідок процесів дифузії та конвекції. Виведемо рівняння
\end{enumerate}

\[
u_1 = \alpha^2 u_{xx} - v u_x,
\]

Виходячи з того, що в будь-який момент часу в певному {[}x, x +
\(\Delta\) x{]} сегменті матерії не з'являється і не зникає.

ПРИМІТКА. Закон збереження матерії виглядає так:

\begin{verbatim}
Зміни маси речовини на сегменті [x, x + $\Delta$ x] = = Зміна внаслідок дифузії через межі + Зміна внаслідок переносу через межі.
\end{verbatim}

Лекція 5

\section{ЗМІННЕ
РОЗДІЛЕННЯ}\label{ux437ux43cux456ux43dux43dux435-ux440ux43eux437ux434ux456ux43bux435ux43dux43dux44f}

МЕТА ЛЕКЦІЇ: Познайомити читача з потужним методом розділення змінних і
показати, як цей метод можна використати для розв'язання добре відомої
проблеми дифузії. Оскільки цей метод створює труднощі для деяких учнів
(головним чином через громіздкі алгебранічні перетворення), ми
супроводжуємо презентацію інтуїтивними міркуваннями.

Основна ідея методу полягає в тому, щоб розкласти початкову умову на
найпростіші компоненти, знайти відгук системи на кожен найпростіший
компонент і підсумувати всі відповіді. Таким чином можна знайти
відповідь на довільну початкову умову.

Існуючий метод поділу змінних певною мірою зменшує можливість такої
інтерпретації методу, але все ж не виключає його.

Метод розділення змінних є одним із найвідоміших методів розв'язання
зсувних задач і використовується, якщо:

\begin{itemize}
\tightlist
\item
  \begin{enumerate}
  \def\labelenumi{\arabic{enumi}.}
  \tightlist
  \item
    Рівняння є лінійним і однорідним (не обов'язково з постійними
    коефіцієнтами).
  \end{enumerate}
\item
  \begin{enumerate}
  \def\labelenumi{\arabic{enumi}.}
  \setcounter{enumi}{1}
  \tightlist
  \item
    Граничні умови визначаються як
  \end{enumerate}
\end{itemize}

\[
\alpha u_x(0, t) + \beta u(0, t) = 0,
\]

\gamma \(u_x\)(1, t) + \delta \(u(1, t)\) = 0, де \alpha , \beta ,
\gamma і \delta --- константи (граничні умови, задані в такій формі,
називаються лінійними однорідними граничними умовами).

Метод був створений за часів Фур'є (зазвичай відомий як метод Фур'є) і,
здається, є найпопулярнішим (де це застосовно).

\pandocbounded{\includegraphics[keepaspectratio]{images/_page_34_Figure_8.jpeg}}

Рис. 5.1. Схематичне представлення задачі дифузії.

Замість того, щоб вивчати метод у загальному випадку, давайте спочатку
розглянемо конкретну задачу (загальний випадок обговоримо пізніше).
Розглянемо змішану задачу типу дифузії: щоб знайти розв'язок

(PPP)

\[
u_t = \alpha^2 u_{xx}
, 0 < x < 1 , 0 < t < \infty ,
\]

що задовольняє граничні умови

(GI)

\[
\begin{cases} u(0, t) = 0, & 0 < t < \infty, \\ u(1, t) = 0, & 0 < t < \infty, \end{cases}
\]

та початкова умова

(ГУ)

\[
u(x, 0) = \varphi(x), 0 \le x \le 1.
\]

Перш ніж почати розділяти змінні, давайте дамо фізичну інтерпретацію
нашої проблеми. Отже, існує стрижень скінченної довжини, кінці якого
підтримуються при сталій температурі, рівній нулю (насправді кінці можна
підтримувати при значно вищій температурі, значення якої вважається
початковою точкою). Додаткові дані про завдання подаються як початкова
умова. Наша мета --- знайти розподіл температури \(u\left(x,t\right)\) в
наступний час.

\section{Загальні принципи методу розділення
змінних}\label{ux437ux430ux433ux430ux43bux44cux43dux456-ux43fux440ux438ux43dux446ux438ux43fux438-ux43cux435ux442ux43eux434ux443-ux440ux43eux437ux434ux456ux43bux435ux43dux43dux44f-ux437ux43cux456ux43dux43dux438ux445}

Для найпростішого диференціального рівняння в похідних похідних
розділення змінних є пошуком розв'язків у вигляді

\[
u(x, t) = X(x)T(t),
\]

де \(X(x)\) --- це функція, що залежить лише від змінної x, а \(T(t)\)
--- функція, що залежить лише від t. Це розв'язок у певному сенсі є
найпростішим, оскільки температура \(u(x, t)\), представлена у такій
формі, зберігатиме «форму» профілю у різний час часу t (див. рисунок
5.2).

\pandocbounded{\includegraphics[keepaspectratio]{images/_page_35_Figure_6.jpeg}}

РИСУНОК 5.2 Графік функції X(x)\$\$T(t) у різних точках часу t.

\begin{verbatim}
Загальна ідея полягає в тому, щоб знайти нескінченну кількість таких часткових диференціальних розв'язків (які задовольняють граничні умови). Ці найпростіші функції u_n\left(x,\ t\right) = X_n\left(x\right)T_n\left(t\right) (які називаються фундаментальними розв'язками) є, так би мовити, елементарними будівельними блоками, з яких будується розв'язок нашої проблеми. Розв'язок нашої задачі u\left(x,\ t\right) у вигляді такої лінійної комбінації фундаментальних розв'язків X_n\left(x\right)T_n\left(t\right) , що отримана сума
\end{verbatim}

\[
\sum_{n=1}^{\infty} A_n X_n(x) T_n(t)
\]

задовольняє початкові умови. Оскільки ця сума задовольняє рівняння та
граничні умови, вона є розв'язком початкової задачі. Ми маємо робити всі
ці розрахунки достатньо детально.

\section{Розділення
змінних}\label{ux440ux43eux437ux434ux456ux43bux435ux43dux43dux44f-ux437ux43cux456ux43dux43dux438ux445}

КРОК 1. (Пошук елементарних розв'язків рівняння з частковими похідними.)

Ми хочемо знайти функцію \(u(x, t)\), яка є розв'язком

Наступного завдання:

\[
\begin{array}{ll} (\mathrm{Y}^{\mathbf{U}}\Pi) & u_t = \alpha^2 u_{xx}, & 0 < x < 1, & 0 < t < \infty, \\ (\Gamma \mathrm{Y}) & \begin{cases} u\left(0,\,t\right) = 0, & 0 < t < \infty, \\ u\left(1,\,t\right) = 0, & 0 < t < \infty, \end{cases} \\ (\mathrm{H}\mathrm{Y}) & u\left(x,\,0\right) = \varphi\left(x\right), & 0 < x \leq 1. \end{array}
\]

Пошукаємо розв'язки, які представлені у вигляді \(u(x, t)\) = \(X(x)\)
\(T(t)\). Для цього підставляємо вираз \(X(x)\) \(T(t)\) у рівняння.
Внаслідок підстановки отримуємо

\[
X(x) T'(t) = q^2 X''(x) T(t).
\]

\textbf{Тепер} виконаємо операцію, властиву цьому методу:
\emph{поділимо} обидві сторони останнього рівняння на \alpha\^{}2
\(X(x)\) \(T(t)\) , у результаті чого отримуємо

\[
\frac{T'(t)}{\alpha^2 T(t)} = \frac{X''(x)}{X(x)}.
\]

Цей вираз вважається таким, що розділив змінні, тобто ліва частина
рівняння залежить лише від t, а права --- лише від x. Оскільки x і t не
є розділенням комунікації один від одного, кожна частина рівняння має
бути сталою. Позначимо цю константу k, тоді

\[
\frac{T'}{\alpha^2 T} = \frac{X''}{X} = k
\]

або

\[
T'-k\alpha^2T=0,
X''-kX=0.
\]

Тепер можливо розв'язати кожне з цих звичайних диференціальних рівнянь.
Добуток відповідних розв'язків задовольнятиме початкове рівняння з
похідними похідними. (Зверніть увагу, що ми спростили оригінальне
диференціальне рівняння другого порядку до двох звичайних
диференціальних рівнянь.)

Тепер звернемо увагу на таку важливу обставину: константа відстань k має
бути від'ємною \(^1\) ) (іншими словами, функції \(T(t)\) мають прагнути
до нуля в точці \(t \to \infty\) ). З огляду на це, давайте розглянемо
позначення \(k = -\lambda^2\) , де \(\lambda\) не дорівнює нулю (у цьому
випадку вираз \(-\lambda^2\) завжди буде від'ємним). Враховуючи нову
нотацію для константи відстань двох звичайних диференціальних рівнянь,

1) Інакше рівняння X'\,' - kX = 0 з граничними умовами \(X(0)\) =
\(X(1)\) = 0 має лише гривну розв'язку \(X(x)\) = 0. Ед.

записуються у формі

\[
T' + \lambda^2 \alpha^2 T = 0,
X'' + \lambda^2 X = 0.
\]

Розв'яжемо ці рівняння. Це рівняння стандартного типу. Їхні загальні
рішення фіксуються у формі

\[
T(t) = Ae^{-\gamma *_{O} *_{I}}
\]

(A --- довільна константа),

\[
X(x) = A \sin(\lambda x) + B \cos(\lambda x)
\]

(A, B --- довільні константи).

Відповідно, функції виду

\[
u(x, t) = e^{-\lambda^2 \alpha^2 t} [A \sin(\lambda x) + B \cos(\lambda x)]
\]

(де A, B і \lambda --- довільні константи) задовольняють NPP
\(u_t = \alpha^2 u_{xx} (\)пропонується перевірити це у задачі 1,
розміщеній наприкінці цієї лекції). Отже, маємо нескінченну множину
функцій, які задовольняють початкове диференціальне рівняння в похідних
похідних.

КРОК 2. (Пошук розв'язків, що задовольняють граничним умовам.)

Зараз ситуація така: ми маємо нескінченну кількість розв'язків
початкового рівняння, але не всі вони задовольняють граничні або
початкові умови. Наступний крок --- вибрати таку підмножину розв'язків
вигляду

(5.1)

\[
e^{-\lambda^2 \alpha^{2\ell}} [A \sin(\lambda x) + B \cos(\lambda x)],
\]

які задовольняють граничні умови

\[
u(0, t) = 0,
\]

\(u(1, t)\) = 0. Для цього підставляємо розв'язки (5.1) у ці граничні
умови. У результаті ми отримуємо

\[
u(0, t) = Be^{-\lambda^2 \alpha^{2t}} = 0 \Rightarrow B = 0,
u(1, t) = Ae^{-\lambda^2 \alpha^{2t}} \sin \lambda = 0 \Rightarrow \sin \lambda = 0.
\]

Друга гранична умова накладає обмеження на можливі значення константи
\(\lambda\) розділення: вона має бути коренем рівняння \(\sin\lambda=0\)
. Іншими словами, щоб виконати умову \(u(1, t)\)=0, необхідно вимагати
виконання

\[
\lambda = \pm \pi
\]

, \(\pm 2\pi\) , \ldots{} або

\[
\lambda_n = \pm n\pi
\]

, \(n = 1, 2, \ldots\) Зверніть увагу, що друга гранична умова може бути
виконана, якщо поставити A = 0, але в цьому випадку розв'язок (5.1) буде
ідентичним нулю.

\pandocbounded{\includegraphics[keepaspectratio]{images/_page_38_Figure_3.jpeg}}

Рис. 5.3. Фундаментальні рішення
\(n_n(x, t) = A_n e^{-(\pi n \alpha)^{n/t}} \sin{(\pi n x)}\) .

Тепер, коли ми зробили другий крок, маємо нескінченний набір ознак

(5.2)

\[
u_n(x, t) = A_n e^{-(n\pi\alpha)^2 t} \sin(n\pi x), \quad n = 1, 2 \dots
\]

кожна з яких задовольняє рівняння з частковими похідними та граничними
умовами). Отримані функції --- це будівельні блоки, з яких ми будемо
створювати рішення нашої проблеми. Це рішення буде певною сумою

2) Зверніть увагу, що функції \(u_n\) і \(u_{-n}\) відрізняються лише
знаком.

цих найпростіших функцій. Конкретний тип суми залежить від початкового
стану. Фіг. Рисунок 5.3 показує графіки кількох таких фундаментальних
функцій

КРОК 3. (Пошук розв'язку, який задовольняє рівняння, межу та початкові
умови.)

Останній (і, ймовірно, найцікавіший з математичної точки зору) крок ---
знайти таку суму фундаментальних розв'язків

\[
u(x, t) = \sum_{n=1}^{\infty} A_n e^{-(n\pi\alpha)^n t} \sin(n\pi x),
\]

тобто при виборі таких коефіцієнтів \(A_n\) , що функція задовольнятиме
початкову умову

\[
u\left( x,\ 0\right) =\varphi\left( x\right) .
\]

Підставляючи суму в початковому стані, отримуємо

(5.3)

\[
q(x) = \sum_{n=1}^{\infty} A_n \sin(n\pi x).
\]

Це рівняння підводить нас до цікавого питання: чи можна розкласти
початкову температуру \(\phi(x)\) на ряд елементарних функцій вигляду

\[
A_1 \sin(\pi x) + A_2 \sin(2\pi x) + A_3 \sin(3\pi x) + \dots
\]

?

Позитивну відповідь на це питання дав французький математик Жозеф Фур'є.
Виявилося, що для достатньо хороших функцій така декомпозиція можлива
\(^1\) ). Тоді виникає нове питання: знайти коефіцієнти розкладу \(A_n\)
?

Насправді це дуже просто. Для цього використаємо властивість системи
функцій

\[
\{\sin(n\pi x); n=1, 2, \ldots\},
\]

відомий як ортогональність.

У нашому випадку (див. завдання 2) ці функції задовольняють умови

\[
\int_{0}^{1} \sin(m\pi x) \sin(n\pi x) dx = \begin{cases} 0, & m \neq n, \\ 1/2, & m = n. \end{cases}
\]

Умови розкладу функцій у рядах Фур'є наведені в лекції 11. Для глибшого
ознайомлення з цим питанням можна звернутися до численних освітніх
літератур (див., наприклад: Г. П. Толстов, Серія Фур'є. --- Москва:
Наука, 1980). --- Ед. Ед.

Ця властивість чітко показана на рис. 5.4, яка показує графіки деяких
функцій із цієї системи.

Отже, ми хочемо знайти коефіцієнти у розкладі

\[
\varphi(x) = A_1 \sin(\pi x) + A_2 \sin(2\pi x) + A_3 \sin(3\pi x) + \dots
\]

\pandocbounded{\includegraphics[keepaspectratio]{images/_page_40_Figure_5.jpeg}}

РИСУНОК 5.4. Послідовність ортогональних функцій

Помножте обидві сторони цього співвідношення на \(\sin{(m\pi x)}\) (m
--- довільне ціле число) і інтегруйте від нуля до одиниці. У результаті
ми отримуємо

\[
\int_{0}^{1} \varphi(x) \sin(m\pi x) \, dx = A_{m} \int_{0}^{1} \sin^{2}(m\pi x) \, dx = \frac{A_{m}}{2}
\]

(усі інші члени зрівнялися з нулем через ортогональність). Розв'язуючи
рівняння відносно \(A_m\) , отримуємо

\[
A_m = 2 \int_0^1 \varphi(x) \sin(m\pi x) dx.
\]

Отже, отримано, що розв'язок записується у вигляді

(5.4)

\[
u(x, t) = \sum_{n=1}^{\infty} A_n e^{-(n\pi\alpha)^2 t} \sin(\pi n x),
\]

де коефіцієнти \(A_n\) визначаються формулами

\[
A_n = 2\int_0^1 \varphi(x) \sin(\pi nx) dx
\]

Ви можете переконатися, що отримане рішення задовольняє всі умови
початкової задачі. На цьому завершується третій крок.

Багато студентів розчаровані, коли бачать, наскільки складним виглядає
рішення, і майже не намагаються дивитися на нього знову (виглядає дуже
погано). Однак це рішення не здаватиметься таким складним, якщо ви
витратите час на його аналіз. Насправді, складніша форма розв'язку
робить його більш інформативним. Нижче наведено кілька нотаток, які
допоможуть вам інтерпретувати це рішення.

\section{ЗАУВАЖЕННЯ}\label{ux437ux430ux443ux432ux430ux436ux435ux43dux43dux44f-4}

Зверніть увагу, що єдина різниця між синусоїдальним розкладом функції
φ(x) і розв'язком (5.4) --- це наявність часового фактора e-(nπα)31 у
кожному члені ряду. Отже, якщо початкова умова має просту декомпозицію,
наприклад, вигляду

\[
\varphi(x) = \sin(\pi x) + \frac{1}{2}\sin(3\pi x),
\]

тоді рішення можна негайно записати

\[
u(x, t) = e^{-(\pi\alpha)^2 t} \sin(\pi x) + \frac{1}{2} e^{-(3\pi\alpha)^2 t} \sin(3\pi x).
\]

У цьому випадку очевидно, що якщо ми розкладемо \(\phi(x)\) на ряд Фур'є
вздовж синусів, то отримуємо

\[
A_1 = 1,
\]

\(A_2 = 0,\) \(A_3 = 1/2,\) \(A_4 = A_5 = \dots = 0.\) - 2. Розв'язок
(5.4) можна інтерпретувати так: ми представляємо початкову температуру
\(\varphi(x)\) як суму найпростіших функцій виду \(A_n \sin(\pi nx)\) ;
кожна така функція породжує відповідь
\(A_n e^{-(\pi n\alpha)^2 t} \sin(\pi nx)\) . Сумуючи всі такі
відповіді, отримуємо розв'язок, який відповідає початковій умові
\(u(x, 0) = \varphi(x)\) . - 3. Кожен член у декомпозиції

\[
u(x, t) = A_1 e^{-(\pi\alpha)^2 t} \sin(\pi x) + A_2 e^{-(2\pi\alpha)^2 t} \sin(2\pi x) + \cdots
\]

є функцією від x і t. Зверніть увагу, що внесок членів з великими
числами в t\textgreater0 дуже малий через множник
\(e^{-(n\pi\alpha)^2 t}\) . Відповідно, після закінчення досить великого
Р и с

\pandocbounded{\includegraphics[keepaspectratio]{images/_page_42_Figure_2.jpeg}}

. 5.5. У задачах дифузії вищі порядки терміни швидко затухають;
\(\phi(x)\) --- початкова температура. Слід зазначити, що по-перше, у
розчині зникають високі частоти --- по-перше, зникли найбільші коливання
міді.

Повний розв'язок приблизно збігається з першим членом, який є півхвилі
синусоїди, що з часом згасає (рис. 1). 5.5.)

\section{ЗАВДАННЯ}\label{ux437ux430ux432ux434ux430ux43dux43dux44f-3}

\begin{enumerate}
\def\labelenumi{\arabic{enumi}.}
\tightlist
\item
  Показати, що функції
\end{enumerate}

\[
u(x, t) = e^{-\lambda^2 \alpha' t} \left[ A \sin(\lambda x) + B \cos(\lambda x) \right]
\]

задовольняють рівняння \(u_l = \alpha^2 u_{xx}\) для довільних значень
A, B і \(\lambda\) .

\begin{enumerate}
\def\labelenumi{\arabic{enumi}.}
\setcounter{enumi}{1}
\tightlist
\item
  Покажи це
\end{enumerate}

\[
\int_0^1 \sin(m\pi x) \sin(n\pi x) dx = \begin{cases} 0, & m \neq n, \\ 1/2, & m = n. \end{cases}
\]

ПРИМІТКА. Використання ідентичності

\[
\sin(mx)\sin(nx) = \frac{1}{2}[\cos(m-n)x - \cos(m+n)x].
\]

\begin{enumerate}
\def\labelenumi{\arabic{enumi}.}
\setcounter{enumi}{2}
\tightlist
\item
  Знайти розклад у рядах Фур'є синусами функції \(\varphi(x) = 1\) на
  сегменті {[}0, 1{]}. Побудуйте перші три терміни розкладу.
\end{enumerate}

Використовуючи результати завдання 3, знайдіть розв'язок наступної
змішаної задачі:

(PPP)

\[
u_t = u_{xx}
\]

\(0 < x < 1\) , (GI)

\[
\begin{cases} u(0, t) = 0, \\ u(1, t) = 0, \end{cases}
\]

\(0 < t < \infty\) , (НУ) \(u(x, 0) = 1, \quad 0 \le x \le 1\) .

(Зверніть увагу, що ця проблема фізично не має значення, оскільки
припускається, що температура на кінцях стрижня миттєво знижується з
одного до нуля \(\phi(x)\) .

\begin{enumerate}
\def\labelenumi{\arabic{enumi}.}
\setcounter{enumi}{4}
\tightlist
\item
  Знайти розв'язок проблеми 4, якщо початкова умова встановлена у
  вигляді
\end{enumerate}

\[
u(x, 0) = \sin(2\pi x) + \frac{1}{3}\sin(4\pi x) + \frac{1}{5}\sin(6\pi x).
\]

\begin{enumerate}
\def\labelenumi{\arabic{enumi}.}
\setcounter{enumi}{5}
\tightlist
\item
  Розв'яжіть задачу 4, якщо початкова умова
  \(u\left( x,\ 0\right) =x-x^{2},\quad 0< x<1.\) Лекція 6
\end{enumerate}

\section{ПЕРЕТВОРЕННЯ НЕОДНОРІДНИХ ГРАНИЧНИХ УМОВ У ОДНОРІДНІ ГРАНИЧНІ
УМОВИ}\label{ux43fux435ux440ux435ux442ux432ux43eux440ux435ux43dux43dux44f-ux43dux435ux43eux434ux43dux43eux440ux456ux434ux43dux438ux445-ux433ux440ux430ux43dux438ux447ux43dux438ux445-ux443ux43cux43eux432-ux443-ux43eux434ux43dux43eux440ux456ux434ux43dux456-ux433ux440ux430ux43dux438ux447ux43dux456-ux443ux43cux43eux432ux438}

МЕТА ЛЕКЦІЇ: Показати, як виникає змішана задача

\[
\begin{array}{lll} (\text{УЧП}) & u_t - a^2 u_{xx} = f\left(x, \ t\right) & \frac{\beta_t}{97} \\ (\text{ГУ}) & \begin{cases} \alpha_1 u_x\left(0, \ t\right) + \beta_1 u\left(0, \ t\right) = g_1\left(t\right), \\ a_2 u_x\left(L, \ t\right) + \beta_2 u\left(L, \ t\right) = g_2\left(t\right), \end{cases} \\ (\text{HУ}) & u\left(x, \ 0\right) = \varphi\left(x\right) \end{array}
\]

можна звести до задачі вигляду

\[
U_{t} - \alpha^{2}U_{xx} = F(x, t),
\alpha_{1}U_{x}(0, t) + \beta_{1}U(0, t) = 0,
\alpha_{2}U_{x}(L, t) + \beta_{2}U(L, t) = 0,
U(x, 0) = \varphi(x)
\]

з нульовими граничними умовами. Цей новий виклик можна подолати

\begin{itemize}
\item
  \begin{enumerate}
  \def\labelenumi{\arabic{enumi}.}
  \tightlist
  \item
    Шляхом ділення змінних (якщо нова задача виявиться однорідною, тобто
    \(F(x, t)\) = 0).
  \end{enumerate}

  \begin{itemize}
  \tightlist
  \item
    \begin{enumerate}
    \def\labelenumi{\arabic{enumi}.}
    \setcounter{enumi}{1}
    \tightlist
    \item
      Методом інтегральних перетворень і розкладу власними функціями,
      якщо F(x, t) \neq 0 .
    \end{enumerate}
  \end{itemize}
\end{itemize}

Метод розділення змінних, про який йшлося в попередній лекції, є досить
потужним, а розв'язки, отримані за його допомогою, подаються у зручній
формі. Однак читач повинен розуміти, що цей метод не застосовний до всіх
задач. Щоб метод розділення змінних був застосовним, граничні умови
мають бути лінійними, однорідними, тобто представленими як

(6.1)

\[
\begin{aligned} \alpha_1 u_x(0, t) + \beta_1 u(0, t) &= 0, \\ \alpha_2 u_x(L, t) + \beta_2 u(L, t) &= 0. \end{aligned}
\]

Мета цієї лекції --- показати, як виникає задача з неоднорідними
граничними умовами форми

(6.2)

\[
\begin{aligned} & (\mathsf{У}\mathsf{Ч}\Pi) & u_t = \alpha^2 u_{xx}, \\ & (\mathsf{\Gamma}\mathsf{Y}) & \begin{cases} \alpha_1 u_x \, (0, \, t) + \beta_1 u \, (0, \, t) = g_1 \, (t) \\ \alpha_2 u_x \, (L, \, t) + \beta_2 u \, (L, \, t) = g_2 \, (t) \end{cases} \text{ (неоднородные } \mathsf{\Gamma}\mathsf{Y}), \\ & (\mathsf{H}\mathsf{Y}) & u \, (x, \, 0) = \varphi \, (x) \end{aligned}
\]

можна розв'язати, зведивши її до задачі з однорідними граничними
умовами.

Нове завдання можна розв'язати іншими методами

(наприклад, шляхом розкладу за власними функціями).

Почнемо з надзвичайно простого прикладу перетворення гетерогенних
граничних умов на однорідні.

\section{Перетворення неоднорідних граничних умов у однорідні граничні
умови}\label{ux43fux435ux440ux435ux442ux432ux43eux440ux435ux43dux43dux44f-ux43dux435ux43eux434ux43dux43eux440ux456ux434ux43dux438ux445-ux433ux440ux430ux43dux438ux447ux43dux438ux445-ux443ux43cux43eux432-ux443-ux43eux434ux43dux43eux440ux456ux434ux43dux456-ux433ux440ux430ux43dux438ux447ux43dux456-ux443ux43cux43eux432ux438-1}

Розглянемо проблему поширення тепла в термоізольованому стрижні, кінці
якого підтримуються при сталих температурах \(k_1\) і \(k_2\) , тобто

\[
\begin{array}{lll} (\mathrm{Y}\mathrm{H}\Pi) & u_t = \alpha^2 u_{xx}, & 0 < x < L, & 0 < t < \infty, \\ (6.3) & (\Gamma\mathrm{Y}) & \begin{cases} u \ (0, \ t) = k_t, \\ u \ (L, \ t) = k_2, \end{cases} & 0 < t < \infty, \\ (\mathrm{H}\mathrm{Y}) & u \ (x, \ 0) = \varphi \ (x), & 0 \leqslant x \leqslant L. \end{array}
\]

Складність цієї задачі полягає в тому, що, оскільки граничні умови в ній
є гетерогенними, ми не можемо розв'язати її шляхом розділення змінних.
Однак очевидно, що в t \to \infty розв'язок нашої задачі схиляється до
стаціонарного розв'язку, який лінійно (вздовж x) змінюється від
температури \(k_1\) до температури \(k_2\) (рис. 6.1). Іншими словами,
розумно припустити, що температуру в нашій задачі можна представити як
суму

\section{з двох
термів}\label{ux437-ux434ux432ux43eux445-ux442ux435ux440ux43cux456ux432}

\[
u\left(x,\,t\right)=
\]

стаціонарна + тимчасова = \uparrow \uparrow Частина розв'язку, що
залежить від HY і має тенденцію до нуля з помноженням на
=\left[$k_1$+\frac{x}{L}$$\left($k_2$-$k_1$\right)\right]+U\left(x,,t\right).
У цьому випадку наше завдання --- знайти температуру переходу U\left(
x,,t\right) . Заміна

\[
u(x, t) = \left[k_1 + \frac{x}{L}(k_2 - k_1)\right] + U(x, t)
\]

У початковій задачі (6.3) ми отримуємо нову задачу відносно невідомої
функції U\left(x,~t\right) . Ви можете вирішити проблему Р и с

\pandocbounded{\includegraphics[keepaspectratio]{images/_page_45_Figure_7.jpeg}}

\begin{verbatim}
. 6.1. Розв'язання задачі (6.3) у різний час часу a- початкова температура u\left(x,\ 0\right)=$\phi$\left(x\right); b - стаціонарна температура u\left(x,\ \infty\right)=k_{1}+\frac{x}{L}\left(k_{2}-k_{1}\right); b - перехідна температура.

Відносно нова невідома функція U\left(x,\,t\right) і додає її до стаціонарного розв'язку, що дає бажану функцію u\left(x,\,t\right) . Виконавши ці прості перетворення з (6.3),
\end{verbatim}

Ми отримаємо

(6.4)

\[
U_{t} = \alpha^{2}U_{xx} \quad 0 < x < L,
(\Gamma Y) \quad \begin{cases} U(0, t) = 0, \\ U(L, t) = 0, \end{cases} \quad 0 < t < \infty,
(HY) \quad U(x, 0) = \varphi(x) - \left[k_{1} + \frac{x}{L}(k_{2} - k_{1})\right] = \overline{\varphi}(x),
\]

де \$\(\overline{\)\phi\(}\)(x)\$ --- нова, але відома початкова умова.

Ця задача стосується не лише однорідного рівняння, а, на щастя,
однорідних граничних умов, що дозволяє розв'язувати її методом
розділення змінних. Насправді, читач, ймовірно, пам'ятає рішення:

(6.5)

\[
U(x, t) = \sum_{n=1}^{\infty} a_n e^{-(n\pi\alpha)^4 t} \sin(n\pi x/L),
\]

рей

\[
a_n = \frac{2}{L} \int_{-\infty}^{L} \overline{\phi}(\xi) \sin(\pi n\xi/L) d\xi.
\]

Це все для сердечників із фіксованою температурою на кінцях. А як щодо
більш реалістичних GU з правою стороною, залежною від часу? Основні ідеї
тут ті ж, що й у попередній задачі, але вони набагато складніші.

\section{Перетворення граничних умов, залежних від часу, у
нульові}\label{ux43fux435ux440ux435ux442ux432ux43eux440ux435ux43dux43dux44f-ux433ux440ux430ux43dux438ux447ux43dux438ux445-ux443ux43cux43eux432-ux437ux430ux43bux435ux436ux43dux438ux445-ux432ux456ux434-ux447ux430ux441ux443-ux443-ux43dux443ux43bux44cux43eux432ux456}

Розглянемо типову задачу

(PPP)

\[
u_t = \alpha^2 u_{xx}
\]

, 0 \textless{} x \textless{} L , 0 \textless{} t \textless{} \infty , 0
\textless{} t \textless{} \infty , 0 \textless{} t \textless{} \infty ,
0 \textless{} t \textless{} \infty , 0 \textless{} t \textless{}
\infty , 0 \textless{} t \textless{} \infty , 0 \textless{} t
\textless{} \infty , 0 \textless{} t \textless{} \infty , 0 \textless{}
t \textless{} \infty , 0 \textless{} t \textless{} \infty , 0
\textless{} t \textless{} \infty , 0 \textless{} t \textless{} \infty ,
0 \textless{} t \textless{} \infty , 0 \textless{} t \textless{}
\infty , 0 \textless{} t \textless{} \infty , 0 \textless{} t
\textless{} \infty , 0 \textless{} t \textless{} \infty , 0 \textless{}
t \textless{} \infty , 0 \textless{} t \textless{} \infty , 0
\textless{} t \textless{} \infty , 0 \textless{} t \textless{} \infty ,
0 \textless{} t \textless{} \infty , 0 \textless{} t \textless{}
\infty , 0 \textless{} t \textless{} \infty , 0 \textless{} t
\textless{} \infty , 0 \textless{} t \textless{} \infty , , 0
\textless{} t \textless{} \infty , 0 \textless{} t \textless{} \infty ,
0 \textless{} t \textless{} \infty , 0 \textless{} t \textless{}
\infty , 0 \textless{} t \textless{} \infty , 0 \textless{} t
\textless{} \infty , 0 \textless{} t \textless{} \infty , 0 \textless{}
t \textless{} \infty , 0 \textless{} t \textless{} \infty , 0
\textless{} t \textless{} \infty , 0 \textless{} t \textless{} \infty ,
0 \textless{} t \textless{} \infty , 0 \textless{} t \textless{}
\infty , 0 \textless{} t \textless{} \infty , 0 \textless{} t
\textless{} \infty , 0 \textless{} t \textless{} \infty , 0 \textless{}
t \textless{} \infty , 0 \textless{} t \textless{} \infty , 0
\textless{} t \textless{} \infty , 0 \textless{} t \textless{} \infty ,
0 \textless{} t \textless{} \infty , 0 \textless{} t \textless{}
\infty , 0 \textless{} t \textless{} \infty , 0 \textless{} t
\textless{} \infty , 0 \textless{} t \textless{} \infty , 0 \textless{}
t \textless{} \infty , 0 \textless{} t \textless{} \infty , 0
\textless{} t \textless{} \infty , 0 \textless{} t \textless{} \infty ,
0 \textless{} t \textless{} \infty , 0 \textless{} t \textless{}
\infty , 0 \textless{} t \textless{} \infty , 0 \textless{} t
\textless{} \infty , 0 \textless{} t \textless{} \infty , 0 \textless{}
t \textless{} \infty , 0 \textless{} t \textless{} \infty , 0
\textless{} t \textless{} \infty , 0 \textless{} t \textless{} \infty ,
0 \textless{} t \textless{} \infty , 0 \textless{} t \textless{}
\infty , 0 \textless{} t \textless{} \infty , 0 \textless{} t
\textless{} \infty , 0 \textless{} t \textless{} \infty , 0 \textless{}
t \textless{} \infty , 0 \textless{} t \textless{} \infty , 0
\textless{} t \textless{} \infty , 0 \textless{} t \textless{} \infty ,
0 \textless{} t \textless{} \infty , 0 \textless{} t \textless{}
\infty , 0 \textless{} t \textless{} \infty , 0 \textless{} t
\textless{} \infty , 0 \textless{} t \textless{} \infty , 0 \textless{}
t \textless{} \infty , 0 \textless{} t \textless{} \infty , 0
\textless{} t \textless{} \infty , 0 \textless{} t \textless{} \infty ,
0 \textless{} t \textless{} \infty , 0 \textless{} t \textless{}
\infty , 0 \textless{} t \textless{} \infty , 0 \textless{} t
\textless{} \infty , 0 \textless{} t \textless{} \infty , 0 \textless{}
t \textless{} \infty , 0 \textless{} t \textless{} \infty , 0
\textless{} t \textless{} \infty , 0 \textless{} t \textless{} \infty ,
0 \textless{} t \textless{} \infty , 0 \textless{} t \textless{}
\infty , 0 \textless{} t \textless{} \infty , 0 \textless{} t
\textless{} \infty , 0 \textless{} t \textless{} \infty , 0 \textless{}
t \textless{} \infty , 0 \textless{} t \textless{} \infty , 0
\textless{} t \textless{} \infty , 0 \textless{} t \textless{} \infty ,
0 \textless{} t \textless{} \infty , 0 \textless{} t \textless{}
\infty , 0 \textless{} t \textless{} \infty , 0 \textless{} t
\textless{} \infty , 0 \textless{} t \textless{} \infty , 0 \textless{}
t \textless{} \infty , 0 \textless{} t \textless{} \infty , 0
\textless{} t \textless{} \infty , 0 \textless{} t \textless{} \infty ,
0 \textless{} t \textless{} \infty , 0 \textless{} t \textless{}
\infty , 0 \textless{} t \textless{} \infty , 0 \textless{} t
\textless{} \infty , 0 \textless{} t \textless{} \infty , 0 \textless{}
t \textless{} \infty , 0 \textless{} t \textless{} \infty , 0
\textless{} t \textless{} \infty , 0 \textless{} t \textless{} \infty ,
0 \textless{} t \textless{} \infty , 0 \textless{} t \textless{}
\infty , 0 \textless{} t \textless{} \infty , 0 \textless{} t
\textless{} \infty , 0 \textless{} t \textless{} \infty , 0 \textless{}
t \textless{} \infty , 0 \textless{} t \textless{} \infty , 0
\textless{} t \textless{} \infty , 0 \textless{} t \textless{} \infty ,
0 \textless{} t \textless{} \infty , 0 \textless{} t \textless{}
\infty , 0 \textless{} t \textless{} \infty , 0 \textless{} t
\textless{} \infty , 0 \textless{} t \textless{} \infty , 0 \textless{}
t \textless{} \infty , 0 \textless{} t \textless{} \infty , 0
\textless{} t \textless{} \infty , 0 \textless{} t \textless{} \infty ,
, 0 \textless{} t \textless{} \infty , 0 \textless{} t \textless{}
\infty , 0 \textless{} t \textless{} \infty , 0 \textless{} t
\textless{} \infty , 0 \textless{} t \textless{} \infty , 0 \textless{}
t \textless{} \infty , 0 \textless{} t \textless{} \infty , 0
\textless{} t \textless{} \infty , 0 \textless{} t \textless{} \infty ,
0 \textless{} t \textless{} \infty , 0 \textless{} t \textless{}
\infty , 0 \textless{} t \textless{} \infty , 0 \textless{} t
\textless{} \infty , 0 \textless{} t \textless{} \infty , 0 \textless{}
t \textless{} \infty , 0 \textless{} t \textless{} \infty , 0
\textless{} t \textless{} \infty , 0 \textless{} t \textless{} \infty ,
0 \textless{} Щоб перетворити ці граничні умови на нульові умови, ми
(після деяких спроб і помилок) зупинилися на такій формі розв'язку:

(6.7)

\[
u(x, t) = A(t)[1-x/L] + B(t)[x/L] + U(x, t),
\]

де функції \(A\left(t\right)\) і \(B\left(t\right)\) обираються так, що
«квазі-стаціонарна» частина розв'язку (6.7)

(6.8)

\[
S(x, t) = A(t)[1 - x/L] + B(t)[x/L]
\]

задовольняло граничні умови початкової задачі. У цьому випадку функція
U\left(x,~t\right) задовольнятиме однорідні граничні умови.

Підставляючи функцію \(S(x, t)\) у граничні умови

\[
S(0, t) = g_1(t),
\]

\(S_x\)(L, t) + hS(L, t) = \(g_2\)(t) призводить до двох рівнянь, з яких
можна визначити \(A(t)\) і \(B(t)\). У результаті ми отримуємо

(6.9)

\[
A(t) = g_1(t), B(t) = \frac{g_1(t) + Lg_2(t)}{1 + Lh}.
\]

Отже.

\[
u(x, t) = g_1(t) [1 - x/L] + \frac{g_1(t) + Lg_2(t)}{1 + Lh} [x/L] + U(x, t).
\]

Якщо ми підставимо ці вирази для you(x, t) у початкову задачу (це
просять читача), отримаємо нову задачу для невідомої функції
\(U(x, t)\):

\begin{verbatim}
\begin{array}{llllllllllllllllllllllllllllllllllll
\end{verbatim}

Тепер ми стикаємося з новою проблемою однорідних граничних умов (на
жаль, рівняння стало гетерогенним). Ми не можемо розв'язати цю задачу
методом розділення змінних, але якщо читач погоджується чекати, то за
кілька лекцій він ознайомиться з розв'язком цієї задачі методом цілісних
перетворень і розкладання за власними функціями.

\section{ЗАУВАЖЕННЯ
.}\label{ux437ux430ux443ux432ux430ux436ux435ux43dux43dux44f-.}

\begin{enumerate}
\def\labelenumi{\arabic{enumi}.}
\item
  Як зазначалося раніше, мета цієї лекції --- навчити читача, як
  перетворювати задачі з неунірольними граничними умовами на задачі з
  нульовими граничними умовами. Після цього може статися, що нове
  диференціальне рівняння в похідних похідних стане однорідним. Якщо це
  станеться, ми можемо вважати себе щасливцями, оскільки в цьому випадку
  проблему можна розв'язати методом розділення змінних (див. перший
  приклад з цієї лекції). Якщо перетворене рівняння виявиться
  гетерогенним, тоді нову задачу потрібно розв'язати іншим методом.
\item
  Лінійні неоднорідні граничні умови найзагальної форми
\end{enumerate}

\[
\alpha_1 u_x(0, t) + \beta_1 u(0, t) = g_1(t),
\]

\alpha\_2 \(u_x\)(L, t) + \beta\_2 \(u(L, t)\) = \(g_2\)(t) Ви також
можете перетворити умови на нульові граничні умови, використовуючи
метод, описаний у другому прикладі. Звісно, нове рівняння, ймовірно,
буде гетерогенним.

\begin{itemize}
\tightlist
\item
  \begin{enumerate}
  \def\labelenumi{\arabic{enumi}.}
  \setcounter{enumi}{2}
  \tightlist
  \item
    Деякі методи розв'язку не містять вимог щодо унікальності граничних
    умов і, отже, не потребують попередніх перетворень. Пізніше, коли ми
    ознайомимося з перетворенням Лапласа, побачимо, що нульові граничні
    умови є необов'язковими (у випадку нульових граничних умов розв'язок
    дещо простіший, ніж в інших випадках).
  \end{enumerate}
\item
  \begin{enumerate}
  \def\labelenumi{\arabic{enumi}.}
  \setcounter{enumi}{3}
  \tightlist
  \item
    Для граничних умов типу
  \end{enumerate}
\end{itemize}

\[
u(0, t) = g_1(t),
\]

\(u(L, t)\) = \(g_2\)(t) Метод, описаний у другому прикладі, призводить
до наступної форми розв'язання задач:

\[
u(x, t) = \left\{ g_1(t) + \frac{x}{L} \left[ g_2(t) - g_1(t) \right] \right\} + U(x, t).
\]

\section{ЗАВДАННЯ}\label{ux437ux430ux432ux434ux430ux43dux43dux44f-4}

\begin{enumerate}
\def\labelenumi{\arabic{enumi}.}
\tightlist
\item
  Знайти розв'язок наступної змішаної задачі:
\end{enumerate}

\[
\begin{array}{lll} (\text{УЧП}) & u_t = \alpha^2 u_{xx} + x \cos t, & 0 < x < 1, \\ (\Gamma \text{У}) & \begin{cases} u \ (0, \ t) = 1, \\ u_x \ (1, \ t) + hu \ (1, \ t) = 1, \end{cases} & 0 < t < \infty, \\ (\text{HY}) & u \ (x, \ 0) = \sin (\pi x), & 0 \leqslant x \leqslant 1, \end{array}
\]

попередньо перетворюючи його на задачу з нульовими граничними умовами.
Чи відповідає рішення, яке ви отримуєте, вашим задумам щодо цієї
проблеми?

Завдання

(PPP)

\[
u_t = u_{xx}, \quad 0 < x < 1,
\]

(GI)

\[
\begin{cases} u(0, t) = 0, & 0 < t < \infty, \\ u(1, t) = 1, & 0 < t < \infty, \end{cases}
\]

(НУ) \(u(x, 0)\) = x\^{}2, \quad 0 \le x \le 1, Перетворіть у задачу з
нульовими граничними умовами і розв'яжіть її. Як виглядатиме рішення у
різних

Миті часу? Чи відповідає це вашим інтуїтивним ідеям? Що таке стаціонарне
рішення? Як виглядає перехідне рішення?

\begin{enumerate}
\def\labelenumi{\arabic{enumi}.}
\setcounter{enumi}{2}
\tightlist
\item
  Завдання
\end{enumerate}

\[
\begin{array}{ll} (\text{УЧП}) & u_t = u_{xx}, & 0 < x < 1, \\ (\Gamma \text{У}) & \begin{cases} u_x(0, t) = 0, \\ u_x(1, t) + hu(1, t) = 1, \\ (\text{HY}) & u(x, 0) = \sin(\pi x), \end{cases} & 0 < x < 1, \\ 0 < t < \infty, \\ 0 \le x \le 1, \end{array}
\]

Перетворимо на задачу з нульовими граничними умовами. Чи буде нове
рівняння однорідним?

Лекція 7

\section{РОЗВ'ЯЗАННЯ СКЛАДНІШИХ ЗАДАЧ ШЛЯХОМ РОЗДІЛЕННЯ
ЗМІННИХ}\label{ux440ux43eux437ux432ux44fux437ux430ux43dux43dux44f-ux441ux43aux43bux430ux434ux43dux456ux448ux438ux445-ux437ux430ux434ux430ux447-ux448ux43bux44fux445ux43eux43c-ux440ux43eux437ux434ux456ux43bux435ux43dux43dux44f-ux437ux43cux456ux43dux43dux438ux445}

МЕТА ЛЕКЦІЇ: Показати, як більш складні проблеми теплопровідності можна
розв'язати методом розділення змінних. У цій лекції ми головним чином
розглядаємо такі задачі, розв'язання яких дозволить читачеві краще
ознайомитися з методом ділення змінних. Сподіваємося, що читач зможе
екстраполювати викладені тут ідеї для власних завдань.

У тій же лекції читач ознайомиться з задачею за її власними функціями,
або, як її ще називають, проблемою Штурма-Ліувіля. Обговорюються деякі
властивості проблеми Штурма-Ліувіля.

Як згадувалося вище, мета цієї лекції --- розв'язати змішану задачу
шляхом розділення змінних. Можливо, саме таким методом читач вирішить
свою проблему. Водночас ми сподіваємося, що він зможе перенести метод,
розглянутий у цій лекції, для власних завдань.

Почнемо з одномірної задачі теплопровідності, в якій одна з граничних
умов містить похідну.

\section{Задача теплопровідності з похідною в граничній
умові}\label{ux437ux430ux434ux430ux447ux430-ux442ux435ux43fux43bux43eux43fux440ux43eux432ux456ux434ux43dux43eux441ux442ux456-ux437-ux43fux43eux445ux456ux434ux43dux43eux44e-ux432-ux433ux440ux430ux43dux438ux447ux43dux456ux439-ux443ux43cux43eux432ux456}

Розглянемо пристрій, показаний на рис. 7.1. Припустимо, що верхній
кінець стрижня підтримується при нульовій температурі, тобто \(u(0, t)\)
= 0, а нижній кінець занурений у воду, температура якої також стала і
дорівнює нулю (у цьому випадку --- нуль Р и с

\pandocbounded{\includegraphics[keepaspectratio]{images/_page_50_Picture_4.jpeg}}

. 7.1. Схематичне представлення зсувної задачі.

\pandocbounded{\includegraphics[keepaspectratio]{images/_page_50_Figure_6.jpeg}}

Рис. 7.2. Природа кривих, які задовольняють

\[
(\Gamma Y) = \begin{cases} u(0, t) = 0, \\ u_x(1, t) = -hu(1, t). \end{cases}
\]

відповідає певній еталонній температурі). Теплообмін між стрижнем і
водою відбувається згідно із законом Ньютона, який записується у x=1 як
наступна гранична умова:

\[
u_x(1, t) = -hu(1, t).
\]

Тепер припустимо, що початкова температура стрижня дорівнює \(u(x, 0)\)
= x, але одразу при t \textgreater{} 0 включаємо граничні умови. Щоб
визначити температуру в стрижні, необхідно розв'язати проблему зміщення

\begin{verbatim}
\begin{array}{lll} (\mathrm{УЧ}\Pi) & u_t = $\alpha$^2 u_{xx}, & 0 < x < 1, & 0 < t < \infty, \\ (7.1) & (\Gamma\mathrm{У}) & \begin{cases} u_t(0, t) = 0 \\ u_x(1, t) + hu_t(1, t) = 0 \end{cases} & \text{(однородные } \Gamma\mathrm{У}), \\ (\mathrm{H}\mathrm{У}) & u_t(x, 0) = x, & 0 \leq x \leq 1. \end{array}
\end{verbatim}

Щоб використати метод розділення змінних, розбиємо процес розв'язання на
кілька етапів.

\textbf{КРОК} 1. (Перетворення диференціального рівняння в похідних у
два звичайних диференціальних рівняння)

\begin{verbatim}
Підставляючи u(x, t) = X(x)T(t) у рівняння (7.1), отримуємо XT' = $\alpha$^2 X''T Поділивши інші частини цього рівняння на $\alpha$^2 XT отримуємо
\end{verbatim}

\[
\frac{T'}{\alpha^2 T} = \frac{X''}{X} .
\]

Оскільки ліва частина останнього рівняння залежить лише від часу, а
права --- лише від x (і оскільки x і t незалежні), обидві сторони
рівняння мають бути константами. Прирівнюючи обидві частини до однієї
сталої \mu , отримуємо

(7.2)

\[
T' - \mu \alpha^{2} T = 0, X'' - \mu X = 0.
\]

Процес розщеплення змінних завершено.

КРОК 2. (Визначення константи відокремлення)

\begin{verbatim}
По-перше, константа і не повинна бути додатною, інакше ^{1} ) функція T\left( t\right) буде зростати експоненційно, що призведе до нескінченного підвищення температури u=XT. Очевидно, що таке рішення не має фізичного значення.

По-друге, припустимо, що \mu = 0 . Тоді
\end{verbatim}

\[
X'' = 0
\]

і, отже.

\[
X(x) = A + \hat{B}x.
\]

Однак, оскільки граничні умови вимагають, що

\[
u(0, t) = X(0) T(t) = 0,
\]

\(u(1, t)\) = (hX(1) + X'(1)) \(T(t)\) = 0, Ми одразу дійшли висновку,
що

\[
X(0) = 0
\Rightarrow A = 0 X'(1) + hX(1) = 0 \Rightarrow B = 0 .
\]

\begin{verbatim}
Це означає, що u(x, t) \equiv 0 . Іншими словами, лише нульове розв'язання відповідає нульовій константі розділення. Ми відкинемо його, оскільки нас цікавлять ненульові розв'язки.

Якщо, нарешті, \mu < 0 , можна поставити \mu = -\lambda^2 і переписати рівняння (7.2) у вигляді
\end{verbatim}

\[
T' + \lambda^2 \alpha^2 T = 0,
X' + \lambda^2 X = 0.
\]

1) Розв'язок \(X(x)\) відповідної задачі з граничним значенням буде лише
тривіальним, \(X(x)\) = 0. --- Ед. Ед.

Загальні розв'язки цих рівнянь легко знайти:

\[
T(t) = Ae^{-(\lambda \alpha)^2 t},
X(x) = B\sin(\lambda x) + C\cos(\lambda x).
\]

Отже, будь-яка функція провідника

(7.3)

\[
u(x, t) = e^{-(\lambda \alpha)^2 t} [A \sin(\lambda x) + B \cos(\lambda x)]
\]

при будь-яких значеннях \lambda A і B задовольняє рівняння (читач може
перевірити це самостійно простою заміною).

Тепер нам потрібно з'ясувати, які з цих функцій задовольняють граничним
умовам

(7.4)

\[
u(0, t) = 0, u_x(1, t) + hu(1, t) = 0.
\]

Підстановка розв'язків виду (7.3) у граничні умови (7.4) дозволяє знайти
відношення, яким мають задовольняти значення \lambda , \Lambda та B:

\[
Be^{-(\lambda\alpha)^2 t} = 0 \Rightarrow B = 0,
\]

A\lambda \(e^{-(\lambda\alpha)^2 t}\) \cos \lambda +
hAe\textsuperscript{\{-(\lambda\alpha)}2 t\} \sin \lambda = 0. Після
простих алгебранхічних перетворень останнє рівняння зводиться до вигляду

\[
tg \lambda = -\lambda/\hbar
\]

.

\begin{verbatim}
Іншими словами, щоб визначити \lambda , потрібно знайти точки перетину графів функцій tg \lambda і — \lambda/h (рисунок 7.3).

Числові значення значень \lambda_1 , \lambda_2 , ... можна знайти за допомогою комп'ютера, якщо задано значення h Р и с
\end{verbatim}

\pandocbounded{\includegraphics[keepaspectratio]{images/_page_52_Figure_15.jpeg}}

\begin{verbatim}
. 7.3. Перетин графів функцій tg(\lambda) і -\lambda/h .
\end{verbatim}

викликати власні значення задачі граничного значення

(7.5)

\[
X'' + \lambda^{2}X = 0, \begin{cases} X(0) = 0 \\ X'(1) + hX(1) = 0. \end{cases}
\]

Іншими словами, з цими значеннями параметра \lambda існують нетривіальні
розв'язки задачі (7.5). Власні значення задачі (7.5) визначаються як
корені рівняння \log \lambda = -\lambda/h . Їхні числові значення
знаходилися на комп'ютері для випадку h = 1. Таблиця 7.1 показує
результати розрахунків перших кількох власних значень.

\begin{verbatim}
| n | \lambda_B | rt . | \lambda_n |
\end{verbatim}

\textbar-------------\textbar----------------------\textbar--------\textbar----------------\textbar{}

123 \textbar{} 2,024,917,98 \textbar{} 45 \textbar{} 11,0814,20
\textbar{}

\begin{verbatim}
Розв'язки задачі (7.5), що відповідають власним значенням \lambda_n , називаються власними функціями X_n(x) . Для розглянутого завдання власні функції виглядають так:
\end{verbatim}

\[
X_n(x) = \sin(\lambda_n x)
\]

і зображені на рис. 7.4.

\pandocbounded{\includegraphics[keepaspectratio]{images/_page_53_Figure_10.jpeg}}

Рис. 7.4. Власні функції \(X_n\)(x) завдання (7.5) при h=1; Усі власні
функції задовольняють граничні умови.

КРОК 3, (Пошук фундаментальних рішень)

Тепер у нас є нескінченний набір функцій (фундаментальних розв'язків) у
нашому розпорядженні

\[
u_n(x, t) = X_n(x) T_n(t) = e^{-(\lambda_n x)^2 t} \sin(\lambda_n x),
\]

кожна з яких задовольняє рівняння та граничні умови. Залишився ще один
крок --- загалом

\[
u(x, t) = \sum_{n=1}^{\infty} a_n X_n(x) T_n(t) =
= \sum_{n=1}^{\infty} a_n e^{-(\lambda_n c t^2 t)} \sin(\lambda_n x)
\]

Вибрати коефіцієнти \(a_n\) так, щоб функція \(u(x, t)\) задовольняла
початкову умову \(u(x, 0)\) = x (у цьому випадку функція \(u(x, t)\)
задовольнятиме рівняння та граничні умови, оскільки, по-перше, кожен
член суми їх задовольняє, а по-друге, і рівняння, і граничні умови є
лінійними та однорідними).

Отже, ймовірність \(a_n\) слід визначати з цієї умови

(7 6)

\[
u(x, 0) = x = \sum_{n=1}^{\infty} a_n \sin(\lambda_n x).
\]

Ми маємо зробити останній крок.

КРОК 4. (Розклад початкової умови на ряд за власними функціями)

Щоб знайти коефіцієнти \(a_n\) у розкладі (7,6) на власні функції,
помножте обидві частини цього співвідношення на \(\sin(\lambda_m x)\) і
інтегруйте від 0 до 1. У результаті ми отримуємо

\[
\int_{0}^{2} \xi \sin(\lambda_{m} \xi) d\xi = \sum_{n=1}^{\infty} a_{n} \int_{0}^{1} \sin(\lambda_{n} \xi) \sin(\lambda_{m} \xi) d\xi =
= a_{m} \int_{0}^{1} \sin^{2}(\lambda_{m} \xi) d\xi =
= a_{m} \left( \frac{\lambda_{m} - \sin \lambda_{m} \cos \lambda_{m}}{2\lambda_{m}} \right).
\]

Розв'язавши останнє відношення відносно \(a_m\) і замінивши m на n,
отримуємо бажаний результат:

(7.7)

\[
a_n = \frac{2\lambda_n}{(\lambda_n - \sin \lambda_n \cos \lambda_n)} \int_0^{\infty} \xi \sin (\lambda_n \xi) d\xi.
\]

Таблиця 7.2.

\section{\texorpdfstring{Шанси \(a_n\) у формулі
(7.8)}{Шанси a\_n у формулі (7.8)}}\label{ux448ux430ux43dux441ux438-a_n-ux443-ux444ux43eux440ux43cux443ux43bux456-7.8}

{\def\LTcaptype{none} % do not increment counter
\begin{longtable}[]{@{}llll@{}}
\toprule\noalign{}
n & an & n & \(a_n\) \\
\midrule\noalign{}
\endhead
\bottomrule\noalign{}
\endlastfoot
123 & 0,240,22-0,03 & 45 & -0.11-0.09  \\
\end{longtable}
}

Отже, розв'язок задачі (7.1) записується у вигляді

(7.8)

\[
u(x, t) = \sum_{n=1}^{\infty} a_n e^{-i\lambda_n \alpha_n^{2} t} \sin(\lambda_n x),
\]

де коефіцієнти \(a_n\) задаються формулами (7.7). Перші п'ять
коефіцієнтів \(a_n\) нашої задачі були знайдені чисельно. Їхні значення
наведені в Таблиці 1. 7.2.

\pandocbounded{\includegraphics[keepaspectratio]{images/_page_55_Figure_8.jpeg}}

Рис. 7.5. Графічне представлення функції (7.8); межа умовно
задовольняється у всі моменти часу.

Відповідно, перші три члени розв'язання змішаної задачі

(7.9)

\[
\begin{aligned} u_t &= u_{xx}, & 0 < x < 1, & 0 < t < \infty, \\ u_t &= 0, & 0 < t < \infty, \\ u_x &= 0, & 0 < t < \infty, \\ u_x &= 0, & 0 < t < \infty, \\ u_x &= 0, & 0 < t < \infty, \\ u_x &= 0, & 0 < t < \infty, \\ u_x &= 0, & 0 < t < \infty, \\ u_x &= 0, & 0 < t < \infty, \\ u_x &= 0, & 0 < t < \infty, \\ u_x &= 0, & 0 < t < \infty, \\ u_x &= 0, & 0 < t < \infty, \\ u_x &= 0, & 0 < t < \infty, \\ u_x &= 0, & 0 < t < \infty, \\ u_x &= 0, & 0 < t < \infty, \\ u_x &= 0, & 0 < t < \infty, \\ u_x &= 0, & 0 < t < \infty, \\ u_x &= 0, & 0 < t < \infty, \\ u_x &= 0, & 0 < t < \infty, \\ u_x &= 0, & 0 < t < \infty, \\ u_x &= 0, & 0 < t < \infty, \\ u_x &= 0, & 0 < t < \infty, \\ u_x &= 0, & 0 < t < \infty, \\ u_x &= 0, & 0 < t < \infty, \\ u_x &= 0, & 0 < t < \infty, \\ u_x &= 0, & 0 < t < \infty, \\ u_x &= 0, & 0 < t < \infty, \\ u_x &= 0, & 0 < t < \infty, \\ u_x &= 0, & 0 < t < \infty, \\ u_x &= 0, & 0 < t < \infty, \\ u_x &= 0, & 0 < t < \infty, \\ u_x &= 0, & 0 < t < \infty, \\ u_x &= 0, & 0 < t < \infty, \\ u_x &= 0, & 0 < t < \infty, \\ u_x &= 0, & 0 < t < \infty, \\ u_x &= 0, & 0 < t < \infty, \\ u_x &= 0, & 0 < t < \infty, \\ u_x &= 0, & 0 < t < \infty, \\ u_x &= 0, & 0 < t < \infty, \\ u_x &= 0, & 0 < t < \infty, \\ u_x &= 0, & 0 < t < \infty, \\ u_x &= 0, & 0 < t < \infty, \\ u_x &= 0, & 0 < t < \infty, \\ u_x &= 0, & 0 < t < \infty, \\ u_x &= 0, & 0 < t < \infty, \\ u_x &= 0, & 0 < t < \infty, \\ u_x &= 0, & 0 < t < \infty, \\ u_x &= 0, & 0 < t < \infty, \\ u_x &= 0, & 0 < t < \infty, \\ u_x &= 0, & 0 < t < \infty, \\ u_x &= 0, & 0 < t < \infty, \\ u_x &= 0, & 0 < t < \infty, \\ u_x &= 0, & 0 < t < \infty, \\ u_x &= 0, & 0 < t < \infty, \\ u_x &= 0, & 0 < t < \infty, \\ u_x &= 0, & 0 < t < \infty, \\ u_x &= 0, & 0 < t < \infty, \\ u_x &= 0, & 0 < t < \infty, \\ u_x &= 0, & 0 < t < \infty, \\ u_x &= 0, & 0 < t < \infty, \\ u_x &= 0, & 0 < t < \infty, \\ u_x &= 0, & 0 < t < \infty, \\ u_x &= 0, & 0 < t < \infty, \\ u_x &= 0, & 0 < t < \infty, \\ u_x &= 0, & 0 < t < \infty, \\ u_x &= 0, & 0 < t < \infty, \\ u_x &= 0, & 0 < t < \infty, \\ u_x &= 0, & 0 < t < \infty, \\ u_x &= 0, & 0 < t < \infty, \\ u_x &= 0, & 0 < t < \infty, \\ u_x &= 0, & 0 < t < \infty, \\ u_x &= 0, & 0 < t < \infty, \\ u_x &= 0, & 0 < t < \infty, \\ u_x &= 0, & 0 < t < \infty, \\ u_x &= 0, & 0 < t < \infty, \\ u_x &= 0, & 0 < t < \infty, \\ u_x &= 0, & 0 < t < \infty, \\ u_x &= 0, & 0 < t < \infty, \\ u_x &= 0, & 0 < t < \infty, \\ u_x &= 0, & 0 < t < \infty, \\ u_x &= 0, & 0 < t < \infty, \\ u_x &= 0, & 0 < t \end{aligned}
\]

мають вигляд

\[
u(x, t) = = 0.24e^{-4t} \sin(2.02x) + 0.22e^{-24t} \sin(4.9x) - 0.03e^{-48.3t} \sin(7.98x) + \dots
\]

Фіг. На рисунку 7-5 показані графіки цього розв'язку для кількох точок
часу. Рекомендуємо читачеві відповісти на такі питання: Чи відповідає
тип рішення, яке ви отримуєте, з вашими інтуїтивними ідеями? Чи
виконуються граничні умови?

\section{ЗАУВАЖЕННЯ}\label{ux437ux430ux443ux432ux430ux436ux435ux43dux43dux44f-5}

Задача власних значень (7.5) є особливим випадком більш загальної
задачі, яку зазвичай називають проблемою Штурма-Ліувіля:

(7.10)

\[
\begin{aligned} & (\text{ОДУ}) \quad [p(x)y']' - q(x)y + \lambda r(x)y = 0, \quad 0 < x < 1, \\ & (\alpha_1 y(0) + \beta_1 y'(0) = 0, \\ & (\alpha_2 y(1) + \beta_1 y'(1) = 0. \end{aligned}
\]

Якщо розв'язати задачу з однорідними граничними умовами методом
розділення змінних, то функції \(X(x)\) завжди знаходяться як розв'язки
деякої задачі Штурма--Ліувіля. Очевидно, що задача власних значень (7.5)
є особливим випадком задачі (7.10).

Перелічимо основні властивості задачі (7.10), які були встановлені
Штурмом і Ліувілем за певних природних обмежень.

функції \(p(x)\), \(q(x)\) та \(r(x)\).

\begin{enumerate}
\def\labelenumi{\arabic{enumi}.}
\tightlist
\item
  Існує нескінченна послідовність власних значень, що задовольняють
  нерівності
\end{enumerate}

\[
\lambda_1 < \lambda_2 < \lambda_3 < \ldots < \lambda_n < \ldots \to \infty
\]

.

\begin{enumerate}
\def\labelenumi{\arabic{enumi}.}
\setcounter{enumi}{1}
\item
  Кожне власне значення \(\lambda_n\) відповідає одній (точній до
  константного множника) ненульовій власній функції \(y_n(x)\) .
\item
  Якщо \(y_n(x)\) і \(y_m(x)\) дві різні власні функції, що відповідають
  різним власним значенням \(\lambda_n\) і \(\lambda_m\) , то вони є
  ортогональними з вагою \(r(x)\) на сегменті {[}0, 1{]}, тобто вони
  задовольняють це відношення
\end{enumerate}

\[
\int_{0}^{1} r(x) y_{n}(x) y_{m}(x) dx = 0.
\]

Більше деталей про проблему Штурма-Ліувіля можна знайти в літературі,
наведеній наприкінці книги.

1) Див., наприклад, Тіхоног А. Н., Самарський А. А. Рівняння
математичної фізики. --- Москва: Наука, 1972. --- ред. Ед.

\section{ЗАВДАННЯ}\label{ux437ux430ux432ux434ux430ux43dux43dux44f-5}

\begin{enumerate}
\def\labelenumi{\arabic{enumi}.}
\tightlist
\item
  Розв'язати наступну задачу для рівняння теплопровідності методом
  розкладу
\end{enumerate}

(PPP)

\[
u_1 = u_{xx}, \quad 0 < x < 1, \quad 0 < t < \infty,
\]

(GI)

\[
\begin{cases} u(0, t) = 0, \\ u_x(1, t) = 0, \\ u(x, 0) = x, \quad 0 \le x \le 1. \end{cases}
\]

Чи відповідає рішення, яке ви знаходите, з інтонативними
представленнями? Що таке стаціонарне рішення?

\begin{enumerate}
\def\labelenumi{\arabic{enumi}.}
\setcounter{enumi}{1}
\tightlist
\item
  Знайдіть власні цінності та функції проблеми Ассолта-Ліувіля
\end{enumerate}

(ODY)

\[
X' + \lambda X = 0, \quad 0 < x < 1,
\]

(FY)

\[
\begin{cases} X(0) = 0, \\ X'(1) = 0. \end{cases}
\]

Задайте специфічні вирази для функцій \(p(x)\), \(q(x)\) та \(r(x)\) із
загальної задачі Штурма--Ліувіля для цього випадку.

\begin{enumerate}
\def\labelenumi{\arabic{enumi}.}
\setcounter{enumi}{2}
\tightlist
\item
  Розв'язати задачу теплопровідності для планки з термоізольованими
  межами
\end{enumerate}

(PPP)

\[
u_t = u_{xx}, \quad 0 < x < 1, \quad 0 < t < \infty,
\]

(GI)

\[
\begin{cases} u_x (0, t) = 0, \\ u_x (1, t) = 0, \end{cases} \quad 0 < t < \infty,
\]

(ГУ) \(u(x, 0) = x, \quad 0 \le x \le 1.\) Чи відповідає рішення, яке ви
отримуєте, з вашою інтерпретацією цієї проблеми? Що таке стаціонарне
рішення? Яке його фізичне значення?

\begin{enumerate}
\def\labelenumi{\arabic{enumi}.}
\setcounter{enumi}{3}
\tightlist
\item
  Знайдіть власні цінності та функції завдань
\end{enumerate}

(ОДА)

\[
X'' + \lambda X = 0, \quad 0 < x < 1,
\]

(GI)

\[
\begin{cases} X'(0) = 0, \\ X'(1) = 0. \end{cases}
\]

\#\# ПЕРЕТВОРЕННЯ СКЛАДНИХ РІВНЯНЬ У ПРОСТУ ФОРМУ

МЕТА ЛЕКЦІЇ Показати, як PPP для функції \(u(x, t)\) можна перетворити
на простіший PPP для нової функції \(w(x, t)\). Вибір правильної
трансформації зазвичай базується на ненавмисному намірі. У нашій лекції

\[
u_t = \alpha^2 u_{xx} - \beta u,
u_t = \alpha^2 u_{xx} - v u_x, \quad (v = \text{const})
\]

переходимо до найпростішого рівняння теплопровідності

\[
w_1 = \alpha^2 w_{xx}
\]

внаслідок відповідних перетворень

\[
u(x, t) = e^{-\beta t} w(x, t)
u(x, t) = e^{\epsilon t (1 - \delta t - 2) 2\alpha^2} w(x, t).
\]

Розв'язавши просте рівняння теплопровідності за функцією \(\omega(x,t)\)
, ми можемо знайти розв'язки початкових рівнянь, перетворивши функцію
\(\omega(x,t)\) відповідно до цих формул (звісно, ці перетворення слід
застосовувати і до початкових і граничних умов).

Після прочитання двох останніх лекцій читач може створити враження, що
метод розділення змінних застосовується \emph{лише} до рівняння

\[
u_1 = \alpha^2 u_{xx}
\]

Насправді рівняння теплопотоку є найпростішим, але не єдиним
параболічним рівнянням, яке можна розв'язати методом розділення змінних.
Як уже згадувалося, якщо рівняння є лінійним і однорідним, його можна
розв'язати, поділивши змінні. Наприклад, двовимірна задача
теплопровідності в колі описується рівнянням

\[
u_t = \alpha^2 \left[ u_{rr} + \frac{1}{r} u_r + \frac{1}{r^2} u_{\theta\theta} \right],
\]

І хоча це рівняння з \emph{змінними коефіцієнтами}, його можна звести до
трьох звичайних диференціальних рівнянь.

У цій лекції ми покажемо читачеві, що іноді, замість того, щоб
безпосередньо розв'язувати диференціальне рівняння в похідних, краще
перетворити його на простіше рівняння, яке можна розв'язати методом
розділення змінних або іншим методом. Наступний приклад ілюструє цей
підхід.

\section{Перетворення проблеми з перенесенням бічної поверхні у проблему
з термічно ізольованою бічною
поверхнею}\label{ux43fux435ux440ux435ux442ux432ux43eux440ux435ux43dux43dux44f-ux43fux440ux43eux431ux43bux435ux43cux438-ux437-ux43fux435ux440ux435ux43dux435ux441ux435ux43dux43dux44fux43c-ux431ux456ux447ux43dux43eux457-ux43fux43eux432ux435ux440ux445ux43dux456-ux443-ux43fux440ux43eux431ux43bux435ux43cux443-ux437-ux442ux435ux440ux43cux456ux447ux43dux43e-ux456ux437ux43eux43bux44cux43eux432ux430ux43dux43eux44e-ux431ux456ux447ux43dux43eux44e-ux43fux43eux432ux435ux440ux445ux43dux435ux44e}

Розглянемо наступне завдання:

(8.1)

\[
u_{t} = \alpha^{2}u_{xx} - \beta u, \quad 0 < x < 1, \quad 0 < t < \infty,
\begin{cases} u(0, t) = 0, \\ u(1, t) = 0, \\ u(x, 0) = \varphi(x), \quad 0 \le x \le 1, \end{cases}
\]

(8.1)

\[
(FY) \quad \begin{cases} u(0, t) = 0, \\ u(x, 0) = \varphi(x), \quad 0 \le x \le 1, \end{cases}
\]

де термін -\(\beta\) u описує потік тепла через бічну поверхню стрижня
(рис. 8.1).

\pandocbounded{\includegraphics[keepaspectratio]{images/_page_59_Figure_7.jpeg}}

\begin{verbatim}
Рис. 81. Задача теплопровідності, описана рівнянням u_t = $\alpha$^t u_{\lambda\lambda} - $\beta$ u : $\alpha$ — це втрата гепле через поверхню Боха; $\delta$ – дифузія вздовж стрижня.

Наша мета — ввести нову температурну w\left(x,\ t\right) замість u\left(x,\ t\right) , щоб рівняння для w\left(x,\ t\right) стало простішим за початкове рівняння
\end{verbatim}

\[
u_t = \alpha^2 u_{xx} - \beta u
\]

.

Це досить загальний метод розв'язання диференціальних рівнянь у
частинних похідних, але вибір відповідного перетворення базується на
інтуїтивних уявленнях про поведінку розв'язків початкового рівняння.
Наприклад, у задачі (8.1) температура \(u(x, t)\) у кожній точці \(x_0\)
змінюється внаслідок двох факторів:

\begin{verbatim}
- 1) дифузія тепла вздовж стрижня (через компонентний $\alpha$^2 u_{xx} ),
\end{verbatim}

\begin{itemize}
\item
  \begin{enumerate}
  \def\labelenumi{\arabic{enumi})}
  \setcounter{enumi}{1}
  \tightlist
  \item
    теплопередача через бічну поверхню стрижня (через член RI).
  \end{enumerate}

  Найважливішим елементом нашого міркування є наступне: якщо не було
  дифузії вздовж стрижня ( \(\alpha\) = 0 ),
\end{itemize}

тоді температура в кожній точці стрижня експоненційно знижується до нуля
згідно із законом

\[
u(x_0 t) = u(x_0, 0) e^{-\beta t}
\]

.

Виходячи з цього спостереження, запитаємо, чи можна розв'язок проблеми
(8.1) представити як утворення двох факторів

\[
(8.2) u(x, t) = e^{-\beta t}w(x, t)
\]

або:

\begin{verbatim}
\begin{pmatrix}
\end{verbatim}

Температура елемента з \textbackslash{}

Неізольовані боком

Поверхня

де \(w(x, t)\) описує розподіл температури, що виникає лише через
дифузію. Давайте подивимось, що станеться, якщо підставити цей вираз для
температури у задачу (8.1). Після виконання відповідних розрахунків
(читач може зробити їх сам), ми отримаємо

\begin{verbatim}
Нова задача майже точно така ж, як і оригінальна, але більше не має додатку 6u. Тому замість розв'язання задачі (8.1) ми розв'язуємо перетворену задачу (8.3), а потім множимо отриманий розв'язок w(x, t) на e^{-$\beta$ t} . У нашому випадку розв'язок задачі (8.3) вже було отримано раніше методом розділення змінних, тому ми можемо одразу записати
\end{verbatim}

(8.4)

\[
w(x, t) = \sum_{n=1}^{\infty} a_n e^{-(n\pi\alpha)^{\frac{n}{2}}t} \sin(n\pi x),
a_n = 2 \int_0^{\infty} \phi(\xi) \sin(n\pi \xi) d\xi,
\]

і, отже, розв'язок початкової задачі (8.1) має вигляд

\[
u(x, t) = e^{-\beta t}w(x, t).
\]

Нижче, у розділі «Нотатки», наведено ще один приклад проблеми,
розв'язаної цим методом.

\section{ПРИМІТКИ}\label{ux43fux440ux438ux43cux456ux442ux43aux438}

\begin{enumerate}
\def\labelenumi{\arabic{enumi}.}
\tightlist
\item
  Вирішити проблему
\end{enumerate}

\[
u_t = u_{xx} - u, \quad 0 < x < 1, \quad 0 < t < \infty,
\begin{cases} u(0, t) = 0, \\ u(1, t) = 0, \\ u(x, 0) = \sin(\pi x) + 0.5\sin(3\pi x) \end{cases}
\]

За методом, описаним вище, ми

\begin{itemize}
\tightlist
\item
  \begin{enumerate}
  \def\labelenumi{\alph{enumi})}
  \tightlist
  \item
    тимчасово відкинути конвективний член u у початковому рівнянні.
  \end{enumerate}
\item
  \begin{enumerate}
  \def\labelenumi{\alph{enumi})}
  \setcounter{enumi}{1}
  \tightlist
  \item
    записати розв'язок змішаної задачі без конвективного члена
  \end{enumerate}
\end{itemize}

\[
u(x, t) = e^{-\pi^2 t} \sin(\pi x) + 0.5e^{-(3\pi)^2 t} \sin(3\pi x),
\]

\begin{enumerate}
\def\labelenumi{\alph{enumi})}
\setcounter{enumi}{2}
\tightlist
\item
  помножимо це розв'язання на конвективний фактор
  \(e^{-\beta t}=e^{-t}\)отримаємо остаточне розв'язання початкової
  задачі
\end{enumerate}

\[
u(x, t) = e^{-t} \left[ e^{-\pi^2 t} \sin(\pi x) + 0.5e^{-(3\pi)^2 t} \sin(3\pi x) \right].
\]

\begin{enumerate}
\def\labelenumi{\arabic{enumi}.}
\setcounter{enumi}{1}
\tightlist
\item
  Рівняння конвективної дифузії
\end{enumerate}

\[
u_t = \alpha^2 u_{xx} - v u_x
\]

(коефіцієнт і є константою) також можна перетворити у вигляд

\[
w_t = \alpha^2 w_{xx}.
\]

Відповідне перетворення має вигляд

\[
u(x, t) = e^{v(x-vt/2)/2\alpha^2}w(x, t).
\]

Це перетворення підкреслює експоненціальний фактор, який повністю
зумовлений рухом середовища. Зверніть увагу, що цей множник описує
хвилю, що поширюється зліва направо зі швидкістю \upsilon/2 . Читач
матиме можливість скористатися цією трансформацією.

\section{ЗАВДАННЯ}\label{ux437ux430ux432ux434ux430ux43dux43dux44f-6}

\begin{enumerate}
\def\labelenumi{\arabic{enumi}.}
\item
  Розв'язати проблему конвективної дифузії

  \begin{array}{lll} (\mathrm{Y}\mathrm{H}\Pi) & u_t = u_{xx} - u_x, & 0 < x < 1, & 0 < t < \infty, \\ (\mathrm{\Gamma}\mathrm{Y}) & \left\{ \begin{array}{ll} u \ (0, \ t) = 0, & 0 < t < \infty, \\ u \ (1, \ t) = 0, & 0 < t < \infty, \\ u \ (x, \ 0) = e^{x/2}, & 0 \leqslant x \leqslant 1, \end{array}

   \Rightarrow. \\ \end{array}
\end{enumerate}

метод перетворення у простіше рівняння. Яке виглядає рішення? Можна дати
таку інтерпретацію цієї задачі: \(u(x, t)\) --- це концентрація певної
речовини у рухомому середовищі (середовище рухається зліва направо зі
швидкістю v = 1), початковий розподіл копконцентрації задається функцією
\(e^{x/2}\) , концентрація підтримується на межах на нулі (можна
припустити, що там є фільтри). Чи відповідає ваше рішення цій
інтерпретації?

\begin{enumerate}
\def\labelenumi{\arabic{enumi}.}
\setcounter{enumi}{1}
\item
  Решите залачу

  \begin{array}{llll} \text{(YUII)} & u_t = u_{xx} - u + x, & 0 < x < 1, & 0 < t < \infty, \\ \text{(FY)} & \left\{ \begin{array}{ll} u \ (0, \ t) = 0, \\ u \ (1, \ t) = 1, \\ u \ (x, \ 0) = 0, & 0 \le x \le 1, \end{array} \end{array}

  \Rightarrow. \textbackslash{} \textbackslash end\{array\}
\end{enumerate}

\begin{itemize}
\item
  а) преобразовав исоднородные граничные условия к однородным;
\item
  б) преобразовав исходное уравнение к повому уравнению, не содержащему
  члена -u;
\item
  в) решив получившуюся в результате эт іх преобразований задачу.
\item
  \begin{enumerate}
  \def\labelenumi{\arabic{enumi}.}
  \setcounter{enumi}{2}
  \tightlist
  \item
    Решите задачу
  \end{enumerate}

  \textbackslash begin\{array\}\{lll\} (\text{Y 4II}) \& u\_t =
  u\_\{\text{vx}\} - u, \& 0 \textless{} x \textless{} 1, \& 0
  \textless{} t \textless{} \infty, \textbackslash{} (\text{\Gamma Y})
  \&

  \begin{cases} u(0,\,t) = 0, \\ u(1,\,t) = 0, \\ u(x,\,0) = \sin{($\pi$ x)}, & 0 \leqslant x \leqslant 1, \end{cases}
\end{itemize}

применяя метод поділу змінних непосредственно к исходному уравнению и не
проводя никаких предварительных преобразований. Совпадает ли полученное
вами розв'язок с решением, полученным ранее після преобразования

\[
u(x, t) = e^{-t}w(x, t)
\]

?

\section{РОЗВ'ЯЗОК НЕОДНОРОДНЫХ УЧП МЕТОДОМ РАЗЛОЖЕНИЯ ПО СОБСТВЕННЫМ
ФУНКЦИЯМ.}\label{ux440ux43eux437ux432ux44fux437ux43eux43a-ux43dux435ux43eux434ux43dux43eux440ux43eux434ux43dux44bux445-ux443ux447ux43f-ux43cux435ux442ux43eux434ux43eux43c-ux440ux430ux437ux43bux43eux436ux435ux43dux438ux44f-ux43fux43e-ux441ux43eux431ux441ux442ux432ux435ux43dux43dux44bux43c-ux444ux443ux43dux43aux446ux438ux44fux43c.}

ЦЕЛЬ ЛЕКЦИИ: Научиться решать смешанные задачи вида

\begin{verbatim}
\begin{array}{lll} \text{(YUII)} & u_t = $\alpha$^2 u_{xx} + f(x,\ t), & 0 < x < 1, & 0 <t < \infty, \\ \text{(FY)} & \left\{ \begin{array}{ll} \alpha_1 u_x \left( 0,\ t \right) + \beta_1 u \left( 0,\ t \right) = 0, \\ \alpha_2 u_x \left( 1,\ t \right) + \beta_2 u \left( 1,\ t \right) = 0, \\ u \left( x,\ 0 \right) = q \left( x \right), & 0 \leqslant x \leqslant 1, \end{array}

\right. \\ \end{array}
\end{verbatim}

Розв'язок гетерогенного UBP можна знайти у вигляді серії

\[
u\left( x,\ t\right) =\sum_{n=1}^{\infty}T_{n}\left( t\right) X_{n}\left( x\right) ,
\]

де \(X_n\left(x\right)\) --- власні функції відповідної однорідної
задачі

\begin{verbatim}
\begin{array}{llll} (\text{YH}\Pi) & u_t = $\alpha$^2 u_{x_1}, & 0 < x < 1, & 0 < t < \infty, \\ (\text{\GammaV}) & \left\{ \begin{array}{ll} \alpha_1 u_x (0, \ t) + \beta_1 u \ (0, \ t) = 0, \\ \alpha_2 u_x \ (1, \ t) + \beta_2 u \ (1, \ t) = 0, \\ u \ (x, \ 0) = $\phi$ \ (x), & 0 \leqslant x \leqslant 1, \end{array} \end{array}

\right. \\ \end{array}
\end{verbatim}

А функції \(T_n\)(t) визначаються як розв'язки деяких звичайних
диференціальних рівнянь.

На лекції 6 ми ознайомилися з тим, як перетворювати неоднорідні граничні
умови на однорідні. На жаль, рівняння залишається гетерогенним, і ми
стоїмо перед завданням

\[
(9.1) \quad u_{t} = \alpha^{2} u_{xx} + f(x, t), \quad 0 < x < 1, \quad 0 < t < \infty,
(9.1) \quad (\Gamma Y) \quad
\begin{cases} \alpha_{1} u_{x}(0, t) + \beta_{1} u(0, t) = 0, \\ \alpha_{2} u_{x}(1, t) + \beta_{2} u(1, t) = 0, \\ u(x, 0) = \varphi(x), \quad 0 \le x \le 1. \end{cases}
\]

Мета нашої лекції --- опанувати метод розв'язання таких задач. Цей метод
аналогічний методу варіації довільних констант.

теорія звичайних диференціальних рівнянь називається методом розкладу за
власними функціями.

Ідея методу надзвичайно проста. Відомо, що розв'язок задачі (9.1) при
\(f(x, t)\) = 0 (так звана відповідна однорідна задача) записується як
ряд

\[
u\left(x,\ t\right)=\sum_{n=1}^{\infty}a_{n}e^{-\left(\lambda_{n}\alpha\right)^{2}}/X_{n}\left(x\right),
\]

де \(X_n\)(x) --- власні функції задачі Штурма-Ліувіля,

(9.2)

\[
X'' + \lambda^2 X = 0, \\ \alpha_1 X'(0) + \beta_1 X(0) = 0, \\ \alpha_2 X'(1) + \beta_2 X(1) = 0.
\]

Запитаймо себе, чи можна також шукати розв'язок гетерогенної проблеми
(9.1) у вигляді серії власних функцій, але більш загальної форми

\[
u(x, t) = \sum_{n=1}^{\infty} T_n(t) X_n(x).
\]

Існують певні фізичні підстави для такого припущення. Насправді терміни
\(f(x,t)\) описують джерело тепла, розташоване всередині стрижня. Це
означає, що залежність рішення від часу не буде виражена показниками

\[
e^{-(\lambda_n\alpha)^2 t}
\]

,

Це стосувалося однорідного завдання, але з деякими іншими функціями.

Щоб показати цей метод у дії і не ускладнювати подання деталями, ми
застосуємо його до розв'язання однієї дуже простої задачі.

\section{Построзв'язання методом демаркації за власними
функціями,}\label{ux43fux43eux441ux442ux440ux43eux437ux432ux44fux437ux430ux43dux43dux44f-ux43cux435ux442ux43eux434ux43eux43c-ux434ux435ux43cux430ux440ux43aux430ux446ux456ux457-ux437ux430-ux432ux43bux430ux441ux43dux438ux43cux438-ux444ux443ux43dux43aux446ux456ux44fux43cux438}

Розглянемо гетерогенну проблему

(9.3) (PY)

\[
\begin{aligned} u_t &= \alpha^2 u_{xx} + f(x, t), & 0 < x < 1, & 0 < t < \infty, \\ u(0, t) &= 0, & 0 < t < \infty, \\ u(1, t) &= 0, & 0 < t < \infty, \\ u(x, 0) &= \varphi(x), & 0 \leqslant x \leqslant 1. \end{aligned}
\]

Процес розв'язання задачі подається у вигляді послідовності наступних
кроків.

КРОК 1. Основна ідея методу полягає в розкладанні густини джерела
\(f(x, t)\) на ряд власними функціями

\[
f(x, t) = f_1(t) X_1(x) + f_2(t) X_2(x) + \dots + f_n(t) X_n(x) + \dots
\]

а також визначити реакції системи \(u_n\)\left(x,,t\right) на вплив
кожного компонента f\_\(n\left(t\right)\)X\_\(n\left(x\right)\) . Якщо
знайти всі відповіді, рішення проблеми виглядатиме так:

\[
u(x, t) = \sum_{n=1}^{\infty} u_n(x, t).
\]

Головна складність цього методу --- розклад густини джерела на
компоненти \(f_n\)(t),\(X_n\)(x) . Виявляється, що \(X_n\)(x) факторами
в цій задачі є власними векторами системи Штурма--Ліувіля, які виникають
при розв'язанні відповідної однорідної задачі методом розділення
змінних, а саме

\[
\begin{array}{c} u_t = \alpha^2 u_{xx} \quad \text{(отметим, что } f(x, \ t) = 0), \\ u(0, \ t) = 0, \\ u(1, \ t) = 0, \\ u(x, \ 0) = \varphi(x). \end{array}
\]

Завдання Штурма-Ліувіля в цьому випадку виглядає так:

\[
X'' + \lambda^2 X = 0,
\]

\(X(0)\) = 0, \(X(1)\) = 0, а його розв'язки є функціями

\[
X_n(x) = \sin(n\pi x), \quad n = 1, 2, 3 \dots
\]

Отже, розклад на густину джерела представлений як ряд

\[
(9.5) \ f(x,t) = f_1(t) \sin(\pi x) + f_2(t) \sin(2\pi x) + \dots + f_n(t) \sin(n\pi x) + \dots
\]

\mathcal{U} нарешті, щоб знайти функції \(f_n\)(t) , помножте обидві
сторони останнього відношення на \sin(m\pi x) і інтегруйте від 0 до 1 за
координатою x.

У результаті ми отримуємо

\[
\int_{0}^{1} f(x, t) \sin(m\pi x) dx = \sum_{n=1}^{\infty} f_{n}(t) \int_{0}^{1} \sin(m\pi x) \sin(n\pi x) dx =
= \frac{1}{2} f_{m}(t)
\]

або (після заміни t на p)

(9.6)

\[
f_{n}(t) = 2 \int_{0}^{1} f(x, t) \sin(n\pi x) dx.
\]

Останнє співвідношення встановлює співвідношення між коефіцієнтами
\(f_n(t)\) і часткою джерела тепла \(f(x, t)\).

КРОК 2. (Пошук \(u_n(x, t) = X_n(x) T_n(t)\) реакції на вхідний вплив
\(f_n(t) X_n(x)\) .)

Замінимо густину джерела на його розклад у рядок

\[
f(x, t) = \sum_{n=1}^{\infty} f_n(t) \sin(n\pi x)
\]

і намагатися знайти індивідуальні відповіді, тобто функції,
\(T_{\mathbf{n}}(t)\) у

\[
u(x, t) = \sum_{n=1}^{\infty} T_n(t) \sin(n\pi x).
\]

Оскільки розв'язок задачі має вигляд

\[
u(x, t) = \sum_{n=1}^{\infty} T_n(t) \sin(n\pi x),
\]

потім підставимо цей вираз у початкову задачу

\[
u_{t} = \alpha^{2} u_{xx} + \sum_{n=1}^{\infty} f_{n}(t) \sin(n\pi x),
u(0, t) = 0,
u(1, t) = 0,
u(x, 0) = \varphi(x),
\]

У результаті ми отримаємо

\[
\sum_{n=1}^{\infty} T'_n(t) \sin(n\pi x) = -\alpha^2 \sum_{n=1}^{\infty} (n\pi)^2 T_n(t) \sin(n\pi x) + + \sum_{n=1}^{\infty} f_n(t) \sin(\pi n\lambda),
\]

(9.7)

\[
\sum_{n=1}^{\infty} T_n(t) \sin 0 = 0
\]

(задовольняється ідентично), \sum\emph{\{n=1\}\^{}\{\infty\} \(T_n\)(t)
\sin (n\pi) = 0 (задовольняється ідентично), \sum}\{n=1\}\^{}\{\infty\}
\(T_n\)(0) \sin (n\pi x) = \phi(x) .

Переписування рівняння та початкової умови

(PPP)

\[
\sum_{n=1}^{\infty} [T'_n + (n\pi\alpha)^2 T_n - f_n(t)] \sin(n\pi x) = 0,
\]

(НУ)

\[
\sum_{n=1}^{\infty} T_n(0) \sin(n\pi x) = \varphi(x),
\]

можна побачити, що функції \(T_n\)(\ell) є розв'язками задачі Коші

(9.8)

\[
T_{n} + (n\pi\alpha)^{2} T_{n} = f_{n}(t),
T_{n}(0) = 2 \int_{0}^{1} \varphi(\xi) \sin(n\pi\xi) d\xi = a_{n}.
\]

Усі ці TAC легко розв'язати, а їхні рішення записуються у вигляді

\begin{enumerate}
\def\labelenumi{(\arabic{enumi})}
\setcounter{enumi}{98}
\tightlist
\item
\end{enumerate}

\[
T_n(t) = a_n e^{-(n\pi\alpha)^2 t} + \int_0^t e^{-(n\pi\alpha)^2 (t-\tau)} f_n(\tau) d\tau,
\]

З того місця, де легко отримати повне розв'язання проблеми (9.3):

(9.10)

\[
u(x, t) = \sum_{n=1}^{\infty} T_n(t) \sin(n\pi x) =
= \sum_{n=1}^{\infty} \left[ a_n e^{-(n\pi\alpha)^n t} \sin(n\pi x) \right] + \sum_{n=1}^{\infty} \left[ \sin(n\pi x) \int_0^t e^{-(n\pi\alpha)^n (t-\tau)} f_n(\tau) d\tau \right]
\]

Перехідна частина (обумовлена стаціонарною частиною (умовлена початковою
умовою) правою частиною \(f(x, t)\)).

Розв'язок (9.10) показує, що температурна реакція складається з двох
частин: перша частина зумовлена початковими умовами, а друга ---
густиною джерела \(f(x,t)\). Зверніть увагу, що термін «стаціонарна
частина» не є найкращим описом другої частини температурної реакції,
оскільки він не означає, що ця частина не залежить від часу
(стаціонарний стан може бути періодичним, якщо \(f(x, t)\) є періодичною
функцією часу t).

Проблема повністю вирішена. Перш ніж перейти до наступної теми, давайте
покажемо, як використовувати описаний вище метод для розв'язання

Конкретне завдання.

70

\section{Розв'язання конкретної задачі методом декомпозиції за власними
функціями}\label{ux440ux43eux437ux432ux44fux437ux430ux43dux43dux44f-ux43aux43eux43dux43aux440ux435ux442ux43dux43eux457-ux437ux430ux434ux430ux447ux456-ux43cux435ux442ux43eux434ux43eux43c-ux434ux435ux43aux43eux43cux43fux43eux437ux438ux446ux456ux457-ux437ux430-ux432ux43bux430ux441ux43dux438ux43cux438-ux444ux443ux43dux43aux446ux456ux44fux43cux438}

Розглянемо просту задачу

(9.11)

\[
\begin{aligned} u_t &= \alpha^2 u_{xx} + \sin{(3\pi x)}, & 0 < x < 1, \\ (\mathbf{FY}) & \begin{cases} u(0, t) &= 0, \\ u(1, t) &= 0, \end{cases} & 0 < t < \infty, \\ (\mathbf{HY}) & u(x, 0) &= \sin{(\pi x)}, & 0 \leq x \leq 1. \end{aligned}
\]

Наша мета --- обчислити коефіцієнти \(T_n\)(t) при розкладі

\[
u\left(x,\,t\right)=\sum_{n=1}^{\infty}\,T_{n}\left(t\right)\sin\left(n\pi x\right)
\]

(Власні функції цієї задачі \(X_n\)(x) збігатися з власними функціями
попереднього). Якщо підставити цю декомпозицію на початкову задачу,
отримаємо наступну задачу Коші для визначення функцій \(T_n\)(t) :

\[
T_n + (n\pi\alpha)^2 T_n = \int_n (t) = 2 \int_0^1 \sin(3\pi x) \sin(n\pi x) dx =
\begin{cases} 1 &n = 3, \\ 0 & n \neq 3, \end{cases}
T_n(0) = 2 \int_0^1 \sin(n\xi) \sin(n\pi\xi) d\xi =
\begin{cases} 1 &n = 1, \\ 0 & n \neq 1. \end{cases}
\]

Перепишем эти уравнения последовательно для n=1, 2\ldots{}

\[
\begin{array}{lll} (n=1) & T_1' + (\pi\alpha)^2 T_1 = 0 \\ T_1(0) = 1 & \Rightarrow T_1(t) = e^{-(\pi\alpha)^2 t}, \\ (n=2) & T_2' + (2\pi\alpha)^2 T_2 = 0 \T_2\  (0) = 0 & \Rightarrow T_2(t) = 0, \\ (n=3) & T_3' + (3\pi\alpha)^2 T_3 = 4 \\ T_3(0) = 0 & \Rightarrow T_3(t) = \frac{1}{(3\pi\alpha)^2} \left[1 - e^{-(3\ pi\alpha)^2 t}\right], \\ (n \geqslant 4) & T_n' + (n\pi\alpha)^2 T_n = 0 \\ T_n(0) = 0 & \Rightarrow T_n(0) = 0. \end{array}
\]

Отже, розв'язок задачі можна записати у вигляді

\[
u(x, t) = e^{-(\pi\alpha)^2 t} \sin(\pi x) + \frac{1}{(3\pi\alpha)^2} [1 - e^{-(3\pi\alpha)^2 t}] \sin(3\pi x).
\]

(9.12)

\section{ЗАУВАЖЕННЯ}\label{ux437ux430ux443ux432ux430ux436ux435ux43dux43dux44f-6}

\begin{enumerate}
\def\labelenumi{\arabic{enumi}.}
\item
  Метод розкладу власних функцій є одним із найпотужніших методів
  розв'язання неоднорідних диференціальних рівнянь у похідних похідних.
  Пізніше, коли ми ознайомимося з інтегральними перетвореннями,
  побачимо, що існують й інші методи розв'язання задач такого типу.
\item
  Власні функції \(X_n(x)\) , за допомогою якої виконується розклад,
  залежать від задачі, яку потрібно розв'язати, точніше від рівнянь і
  граничних умов. Читача просять розв'язати задачу 4 наприкінці лекції,
  щоб перевірити, чи дійсно він навчився знаходити власні функції
  \(X_n(x)\) .
\item
  Якщо читач пам'ятає основні факти з теорії звичайних диференціальних
  рівнянь, то він знає, що розв'язки, що відповідають неоднорідним
  членам вигляду
\end{enumerate}

\[
P_n(x) e^{\alpha x}
\begin{cases} \sin{(\beta x)}, \\ \cos{(\beta x)}, \end{cases}
\]

можно найти методом неопределенных коэффициентов. То же самое и здесь.
Задача (9.11) может быть решена таким же способом.

\section{ЗАДАЧИ}\label{ux437ux430ux434ux430ux447ux438}

\section{1. Решение
задачи}\label{ux440ux435ux448ux435ux43dux438ux435-ux437ux430ux434ux430ux447ux438}

\[
\begin{array}{ll} (\text{УЧП}) & u_t = u_{xx} + \sin{(3\pi x)}, & 0 < x < 1, \\ (\Gamma \text{У}) & \begin{cases} u \ (0, \ t) = 0, \\ u \ (1, \ t) = 0, \end{cases} & 0 < t < \infty \\ (\text{HV}) & u \ (x, \ 0) = \sin{(\pi x)}, \ 0 \leqslant x \leqslant 1 \end{array}
\]

дается формулой (9.12). Соответствует ли это решение вашим интуитивным
представлениям? Что представляет собой полученное решение?

\begin{enumerate}
\def\labelenumi{\arabic{enumi}.}
\setcounter{enumi}{1}
\tightlist
\item
  Решите задачу
\end{enumerate}

\[
\begin{array}{ll} (\text{YUII}) & u_t = u_{xx} + \sin{(\pi x)} + \sin{(2\pi x)}, & 0 < x < 1, & 0 <t < \infty, \\ (\text{\GammaY}) & \begin{cases} u\left(0,\ t\right) = 0, \\ u\left(1,\ t\right) = 0, \end{cases} & 0 < t < \infty, \\ (\text{HY}) & u\left(x,\ 0\right) = 0, & 0 \leqslant x \leqslant 1. \end{array}
\]

\section{3. Розв'яжіть
задачу}\label{ux440ux43eux437ux432ux44fux436ux456ux442ux44c-ux437ux430ux434ux430ux447ux443}

(PPP)

\[
u_t = u_{xx} + \sin(\pi x), \quad 0 < x < 1, \quad 0 < t < \infty,
\]

(GI)

\[
\begin{cases} u(0, t) = 0, \\ u(1, t) = 0, \\ u(x, 0) = 1, \quad 0 < x < 1, \end{cases}
\]

методом разложения по собственным функциям.

\begin{enumerate}
\def\labelenumi{\arabic{enumi}.}
\setcounter{enumi}{3}
\tightlist
\item
  Найдите решение задачи
\end{enumerate}

\[
\begin{array}{ll} (\text{YHII}) & u_t = u_{xx} + \sin{(\lambda_1 x)}, & 0 < x < 1, & 0 <t < \infty, \\ (\text{\GammaY}) & \begin{cases} u\left(0, \ t\right) = 0, \\ u_x\left(1, \ t\right) + u\left(1, \ t\right) = 0, \\ 0 < t < \infty, \end{cases} \\ (\text{HY}) & u\left(x, \ 0\right) = 0, & 0 \leqslant x \leqslant 1 \end{array}
\]

методом разложения по собственным функциям, если \lambda\_1 ---первый
корень уравнения \operatorname{tg} \lambda = -\lambda . Каковы
собственные функция этой задачи?

\section{5. Решите
задачу}\label{ux440ux435ux448ux438ux442ux435-ux437ux430ux434ux430ux447ux443}

\[
\begin{array}{ll} (\text{УЧП}) & u_t = u_{xx} & 0 < x < 1, \quad 0 < t < \infty, \\ (\Gamma \text{У}) & \begin{cases} u(0, t) = 0, \\ u(1, t) = \cos t, \end{cases} & 0 < t < \infty, \\ (\text{HY}) & u(x, 0), \quad 0 \leqslant x \leqslant 1, \end{array}
\]

Спочатку перетворивши граничні умови в нуль. Отримана задача
розв'язується методом розкладання за власними функціями.

\section{ІНТЕГРАЛЬНІ ПЕРЕТВОРЕННЯ (СИНУСОЇДАЛЬНІ ТА КОСИНУСОВІ
ПЕРЕТВОРЕННЯ)}\label{ux456ux43dux442ux435ux433ux440ux430ux43bux44cux43dux456-ux43fux435ux440ux435ux442ux432ux43eux440ux435ux43dux43dux44f-ux441ux438ux43dux443ux441ux43eux457ux434ux430ux43bux44cux43dux456-ux442ux430-ux43aux43eux441ux438ux43dux443ux441ux43eux432ux456-ux43fux435ux440ux435ux442ux432ux43eux440ux435ux43dux43dux44f}

МЕТА ЛЕКЦІЇ: Ввести концепцію інтегрального перетворення та показати, як
використання інтегральних перетворень дозволяє перейти від CPP з n
змінними до нового рівняння з n---1 змінними.

Лекція показує, що інтегральні перетворення можна інтерпретувати як
розклад вхідної дії на елементарні (частотний розклад), визначення
реакції системи на елементарні впливи та подальше підсумовування всіх
отриманих результатів.

У інтегральних перетвореннях операція диференціювання переходить у
операцію множення, і, відповідно, деякі часткові похідні переходять у
алгебраїчні вирази.

Для розв'язання задач дифузії на необмежених інтервалах вводять
синусоїдальні та косинусні перетворення. Розв'язок, отриманий у цьому
випадку, цікавий тим, що містить функцію, яка є додатковою до
ймовірнісного інтеграла.

Інтегральне перетворення зазвичай називається перетворенням, яке кожній
функції \(f(t)\) призначає нову функцію \(F(s)\) відповідно до формули

\[
F(s) = \int_{A}^{B} K(s, t) f(t) dt.
\]

Зверніть увагу, що початкова функція залежить від t, а отримана функція
--- від s. Функція \(K(s,t)\) називається ядром перетворення. Це
головний елемент, який відрізняє одну трансформацію від іншої. Зазвичай
ядро обирається так, щоб перетворення мало певні задані властивості.
Межі інтегрування також залежать від типу трансформації.

Голосною, що лежить в основі інтегральних перетворень, є здатність
позбутися часткових похідних однієї з незалежних змінних. Це означає, що
перетворене рівняння матиме на одну змінну менше, ніж початкове.
Наприклад, якщо застосувати інтегральне перетворення до рівняння

Щоб позбутися похідної часу, отримуємо звичайне диференціальне рівняння
з незалежною змінною X.

\[
u_{xx} + u_{cy} + u_{zz} = 0
\]

застосувати перетворення Фур'є до змінної x, тоді похідна \(u_{xx}\)
буде усунена, і в новому рівнянні залишаться лише дві незалежні змінні y
та z. Тоді отримуємо звичайне диференціальне рівняння з однією
незалежною змінною. Іншими словами, інтегральні перетворення спрощують
початкове завдання. Потім розв'язується перетворення задачі, і,
застосовуючи обернене перетворення до її розв'язку, отримується
розв'язок початкового перетворення. Схема описаного алгоритму показана
на рис. 10.1

\pandocbounded{\includegraphics[keepaspectratio]{images/_page_72_Figure_5.jpeg}}

РИС 10.1, Загальна схема застосування інтегрованих перетворень.

Фіг. На рисунку 10.1 видно, що кожне перетворення прямого інтеграла
пов'язане з оберненим перетворенням, яке відновлює початкову функцію з
трансформованої.

Прямі та обернені перетворення зазвичай називають парою перетворень.
Таблиця 1. Рисунок 10.1 показує кілька пар перетворень, які ми будемо
використовувати при розв'язуванні рівнянь із факторами.

\begin{verbatim}
Слід зазначити, що в цих перетвореннях використовуються два типи символів. Наприклад, у випадку перетворення Лапласа позначення \mathcal{L}[f] вказує, що функція f перетворена, тоді як позначення F(s) означає, що перетворення породжує функцію, яка залежить від s.
\end{verbatim}

У цій лекції ми не прагнемо вивчати всі пари перетворень. Ми
зосередимося лише на перетвореннях Фур'є синуса та косинуса (пари 1 і
2). Трохи пізніше ми вивчаємо іншу з цих пар. Ознайомляючись із
трансформаціями, ми дізнаємося про їхні переваги та недоліки,
дізнаємося, як

Таблиця 10.1. Деякі пари інтегральних перетворень

\begin{enumerate}
\def\labelenumi{\arabic{enumi}.}
\tightlist
\item
  Синус-перетворення Фур'є:
\end{enumerate}

\[
\mathcal{F}_s[f](\omega) = \frac{2}{\pi}\int_0^\infty f(t)\sin(\omega t)\,dt,
\mathcal{F}_s^{-1}[F](t) = \int_0^\infty F(\omega)\sin(\omega t)\,d\omega.
\]

\begin{enumerate}
\def\labelenumi{\arabic{enumi}.}
\setcounter{enumi}{1}
\tightlist
\item
  Косинус-перетворення Фур'є:
\end{enumerate}

\[
\mathcal{F}_c[f](\omega) = \frac{2}{\pi}\int_0^\infty f(t)\cos(\omega t)\,dt,
\mathcal{F}_c^{-1}[F](t) = \int_0^\infty F(\omega)\cos(\omega t)\,d\omega.
\]

\begin{enumerate}
\def\labelenumi{\arabic{enumi}.}
\setcounter{enumi}{2}
\tightlist
\item
  Перетворення Фур'є:
\end{enumerate}

\[
\mathcal{F}[f](\omega) = \frac{1}{\sqrt{2\pi}}\int_{-\infty}^{\infty} f(x)e^{-i\omega x}\,dx,
\mathcal{F}^{-1}[F](x) = \frac{1}{\sqrt{2\pi}}\int_{-\infty}^{\infty} F(\omega)e^{i\omega x}\,d\omega.
\]

\begin{enumerate}
\def\labelenumi{\arabic{enumi}.}
\setcounter{enumi}{3}
\tightlist
\item
  Скінченне синус-перетворення:
\end{enumerate}

\[
S_n = \frac{2}{L}\int_0^L f(x)\sin\!\left(\frac{n\pi x}{L}\right) dx,
f(x) = \sum_{n=1}^{\infty} S_n \sin\!\left(\frac{n\pi x}{L}\right).
\]

\begin{enumerate}
\def\labelenumi{\arabic{enumi}.}
\setcounter{enumi}{4}
\tightlist
\item
  Скінченне косинус-перетворення:
\end{enumerate}

\[
C_n = \frac{2}{L}\int_0^L f(x)\cos\!\left(\frac{n\pi x}{L}\right) dx,
f(x) = \frac{C_0}{2} + \sum_{n=1}^{\infty} C_n \cos\!\left(\frac{n\pi x}{L}\right).
\]

\begin{enumerate}
\def\labelenumi{\arabic{enumi}.}
\setcounter{enumi}{5}
\tightlist
\item
  Перетворення Лапласа:
\end{enumerate}

\[
\mathcal{L}[f](s) = \int_0^\infty e^{-st}f(t)\,dt,
\mathcal{L}^{-1}[F](t) = \frac{1}{2\pi i}\int_{\gamma-i\infty}^{\gamma+i\infty} e^{st}F(s)\,ds.
\]

\begin{enumerate}
\def\labelenumi{\arabic{enumi}.}
\setcounter{enumi}{6}
\tightlist
\item
  Перетворення Ханкеля:
\end{enumerate}

\[
H\{f\}(\xi) = F_n(\xi) = \int_0^\infty r J_n(\xi r) f(r)\,dr,
H^{-1}[F_n](r) = \int_0^\infty \xi J_n(\xi r) F_n(\xi)\,d\xi.
\]

Однак перед тим, як перейти до вивчення інтегральних перетворень, буде
корисно зрозуміти, що зазвичай називають спектром функції (або
спектральним розкладом функції).

\section{Спектральна
функція}\label{ux441ux43fux435ux43aux442ux440ux430ux43bux44cux43dux430-ux444ux443ux43dux43aux446ux456ux44f}

Інтегральні перетворення та спектр функції тісно пов'язані. Насправді
інтегральне перетворення можна інтерпретувати як розклад функції на
певний спектр компонент. Цей спектр залежить від обраного перетворення,
але кожна функція має певний спектр відносно цього перетворення Р и с

\pandocbounded{\includegraphics[keepaspectratio]{images/_page_74_Picture_7.jpeg}}

. 10 2 Розклад періодичної функції на ряд синусами та косинусами;
Послідовність квадратних імпульсів апроксимується синусами та косинусами

Наприклад, розглянемо розклад періодичної функції на ряд синусами та
косинусами (ряд Фур'є) 1:

\[
f(x) = \sum_{n=0}^{\infty} \left[ A_n \cos(nx) + B_n \sin(nx) \right]
\]

(див. рис. 10.2).

\begin{verbatim}
Тут коефіцієнти A_n і B_n позначають внесок \cos(nx) і \sin(nx) у функцію f(x), тоді як величина
\end{verbatim}

\[
V\overline{A_n^2+B_n^3}
\]

1) Фур'є-ряди детальніше розглядаються у лекції 11,

(називається спектром функції) показує частку амплітуди функції
\(f(x)\), що припадає на частоту n.

Наприклад, якщо функція \(f(x)\) --- це просто сума

синуси та косинуси

\[
f(x) = 1 + \sin x + \frac{1}{5}\sin(3x) + \cos x + \frac{1}{2}\cos(2x) + \frac{1}{4}\cos(4x),
\]

тоді його спектр (дискретний) виглядатиме так, як показано на рис. 10.3.

\begin{verbatim}
Зчитуючи значення значення \sqrt{A_n^2 + B_n^2} з графіка, можна визначити амплітудне значення компоненти функції f(x), яка має частоту n.
\end{verbatim}

\pandocbounded{\includegraphics[keepaspectratio]{images/_page_75_Figure_8.jpeg}}

Рис. 10 3. Дискретний спектр функції \(f(r)\).

Періодичні функції розкладаються на нескінченні ряди (вони мають
дискретний спектр), тоді як неперіодичні функції повинні мати
неперервний спектр. Якщо функція визначена лише на скінченному
інтервалі, її можна періодично розширювати на всю вісь і розкладати на
ряд Фур'є, який представлятиме нашу функцію в межах визначення області.

І хоча неперіодичну функцію не можна розкласти на нескінченний ряд
синусами та косинусами, можна спробувати записати неперервний аналог
ряду Фур'є:

\[
f(x) = \int_{-\infty}^{+\infty} [C(\omega)\cos(\omega x) + S(\omega)\sin(\omega x)] d\omega,
\]

де функції S(\omega) і C(\omega) слугують мірою спектральної густини
синуса і косинуса у функції \(f(x)\), а також значення

\[
V \overline{S^2(\omega) + C^2(\omega)}
\]

Він слугує мірою загального внеску частотного \omega у функцію \(f(x)\)
і називається спектром (неперервним) функції \(f(x)\). Виходячи з
наведеного вище інтуїтивного визначення спектра функції, давайте
проаналізуємо інтегральні перетворення «гвинт за гвинтом». Спершу
запишемо кілька властивостей цих перетворень, які ми будемо широко
використовувати.

\section{Синусоїдальні та косинусні перетворення
похідних}\label{ux441ux438ux43dux443ux441ux43eux457ux434ux430ux43bux44cux43dux456-ux442ux430-ux43aux43eux441ux438ux43dux443ux441ux43dux456-ux43fux435ux440ux435ux442ux432ux43eux440ux435ux43dux43dux44f-ux43fux43eux445ux456ux434ux43dux438ux445}

\begin{enumerate}
\def\labelenumi{\arabic{enumi}.}
\tightlist
\item
\end{enumerate}

\[
\mathcal{F}_s[f'] = -\omega \mathcal{F}_c[f]
\]

(доведено інтегруванням по частинах).

\begin{enumerate}
\def\labelenumi{\arabic{enumi}.}
\tightlist
\item
\end{enumerate}

\[
\mathcal{F}_s[f'] = -\omega \mathcal{F}_c[f]
\]

\begin{verbatim}
2. \mathcal{F}_s[f''] = \frac{2}{$\pi$} $\omega$ f(0) - $\omega$^2 \mathcal{F}_s[f] (10.1) 3.
\end{verbatim}

\[
\mathcal{F}_{\rho}[f'] = \frac{2}{\pi} f(0) + \omega \mathcal{F}_{s}[f]
\]

.

\begin{enumerate}
\def\labelenumi{\arabic{enumi}.}
\setcounter{enumi}{3}
\tightlist
\item
\end{enumerate}

\[
\mathcal{F}_{c}[f'') = -\frac{2}{\pi} f_{c}(0) - \omega^{2} \mathcal{F}_{c}[f].
\]

Деякі інші властивості прямих і обернених синусоїдальних і спеційних
перетворень наведені в таблиці наприкінці книги. Тепер покажемо, як
синусоїдальне перетворення можна використати для розв'язання важливої
змішаної задачі.

\section{Розв'язок задачі дифузії на напівпрямому синусоїдальному
метрі}\label{ux440ux43eux437ux432ux44fux437ux43eux43a-ux437ux430ux434ux430ux447ux456-ux434ux438ux444ux443ux437ux456ux457-ux43dux430-ux43dux430ux43fux456ux432ux43fux440ux44fux43cux43eux43cux443-ux441ux438ux43dux443ux441ux43eux457ux434ux430ux43bux44cux43dux43eux43cux443-ux43cux435ux442ux440ux456}

Задача, яку ми хочемо розв'язати, формулюється так:

(PPP)

\[
u_t = \alpha^2 u_{xx}
, 0 < x < \infty , 0 < t < \infty , (GU) u(0, t) = A , 0 < t < \infty , (NU) u(x, 0) = 0 , 0 \le x < \infty .
\]

\pandocbounded{\includegraphics[keepaspectratio]{images/_page_76_Figure_14.jpeg}}

\begin{verbatim}
Рис. 10.4. Задача дифузії в напівнескінченній області; u\left(0,\,t\right)= = A, u(x, 0) = 0.
\end{verbatim}

КРОК 1. Розберемо розв'язання проблеми на три досить прості кроки.
Діятимемо відповідно до наступного плану. На першому кроці виконуємо
перетворення синуса Фур'є на змінній x.

тоді отримуємо звичайне диференціальне рівняння для змінної t. Для
реалізації цього плану застосовуємо синусоїдальне перетворення до обох
частин рівняння з частковими похідними. У результаті ми отримуємо

\[
\mathcal{F}_s[u_t] = \alpha^2 \mathcal{F}_s[u_{xx}].
\]

Розглянемо кожну частину цього співвідношення.
\mathcal{F}\_s{[}\(u_t\){]} : часткову похідну \(u_t\) в цій задачі
можна назвати зовнішньою похідною, оскільки ми використовуємо
перетворення на x. У цьому випадку можна записати

\[
\mathcal{F}_{s}\left[u_{t}\right] = \frac{2}{\pi} \int_{0}^{\infty} u_{t}\left(x, t\right) \sin\left(\omega x\right) dx =
= \frac{\partial}{\partial t} \left[\frac{2}{\pi} \int_{0}^{\infty} u\left(x, t\right) \sin\left(\omega x\right) dx\right] =
= \frac{d}{dt} \mathcal{F}_{s}\left[u\right] =
= \frac{d}{dt} U\left(t\right).
\]

Тут ми взяли операцію диференціації поза знаком інтеграла, застосовуючи
відповідну теорему інтегрального числення. Зверніть увагу, що початкова
функція u залежить від змінних x і t, а також від її образу

\[
\mathcal{F}_{s}[u] = U(\omega, t)
\]

залежить від змінних \(\omega\) і t. Нова змінна \(\omega\) відіграє
роль параметра в новій задачі, і тому можна припустити, що образи
функцій залежать лише від однієї змінної t:

\[
\mathcal{F}_s[u] = U(t)
\]

. \({\mathscr F}_s[u_{xx}]\) : Для цього елемента дійсні наступні
рівняння:

\[
\begin{split} &\mathcal{F}_s\left[u_{xx}\right] = \frac{2}{\pi} \omega u \ (0, \t) - \omega^2 \mathcal{F}_s\left[u\right] = \\ &= \frac{2}{\pi} \omega u \ (0, \t) - \omega^2 U \ (t) = \\ &= \frac{2A\omega}{\pi} - \omega^2 U \ (t). \end{split}
\]

Зверніть увагу, що при виведенні відношень (10.1) ми розглядали функцію
f як функцію однієї змінної x. Тут випадок дещо інший, оскільки
\(u\left(x,\ t\right)\) залежить і від x, і від t. Однак ті ж формули
можна використовувати, якщо перетворення виконується однією змінною, а
решта вважаються константами.

У нашому випадку перетворення виконується змінною x, тобто змінну t
можна вважати сталою. Звернімо також увагу на той факт, що вже на цьому
етапі розв'язання задачі використовується гранична умова
\(u\left(0,\ t\right)=A\) . Підставляючи знайдені вирази у рівняння
\(u_{t}=\alpha^{2}u_{xx}\) , отримуємо звичайне диференціальне рівняння

\[
\frac{dU}{dt} = \alpha^2 \left[ -\omega^2 U(t) + \frac{2A\omega}{\pi} \right].
\]

Єдине, що існує --- це початкова умова функції \(U(t)\). Отримуємо його,
застосовуючи перетворення до початкової умови \(u(x, 0)\) = 0:

\[
\mathcal{F}_{s}[u(x, 0)] = U(0) = 0.
\]

Це завершує перший крок у побудові розв'язку методом інтегрального
перетворення: замість початкової змішаної задачі ми отримали задачу Коші
для звичайного диференціального рівняння:

(10.2)

\[
\begin{aligned} (O \Pi V) & \frac{dU}{di} + \omega^2 \alpha^2 U = \frac{2A\omega \alpha^2}{\pi}, \\ (H V) & U(0) = 0. \end{aligned}
\]

КРОК 2. Щоб розв'язати отриману задачу Коші, можна використати будь-який
елементарний метод розв'язання звичайних диференціальних рівнянь. У
результаті ми отримуємо

\[
U(t) = \frac{2A}{\pi\omega} (1 - e^{-\omega^2 \alpha^2 t}).
\]

Отже, перед нами маємо результат застосування синусоїдального
перетворення до бажаної функції \(u\left(x,\ t\right)\) . Останній крок
--- застосувати обернене перетворення до функції \(U\left(t\right)\) і
отримати потрібне рішення \(u\left(x,\ t\right)\) :

\[
u(x, t) = \mathcal{F}_{x}^{-1}[U].
\]

КРОК 3. Щоб отримати розв'язок, можна або безпосередньо обчислити
інтеграл, що входить до оберненого перетворення, або використати таблиці
перетворень. Використовуючи таблиці, ми знаходимо

\[
u(x, t) = A \operatorname{erfc}(x/2\alpha \sqrt{t}),
\]

де erfc(x) --- функція, доповнююча до ймовірнісного інтегралу, а
визначається відношенням

\[
\operatorname{erfc}(x) = \frac{2}{\sqrt{2\pi}} \int_{x}^{\infty} e^{-t^{2}} dt.
\]

Графік цієї функції показано на рис. 10.5.

Ймовірнісний інтеграл і додаткова функція до нього не виражаються через
елементарні функції. Вони добре популярні, і відповідні таблиці можна
знайти в довідниках з фізики та хімії.

\pandocbounded{\includegraphics[keepaspectratio]{images/_page_79_Figure_4.jpeg}}

Рис. 10.5. Графи функцій erf (x) та eric (x).

\section{Інтерпретація
рішення}\label{ux456ux43dux442ux435ux440ux43fux440ux435ux442ux430ux446ux456ux44f-ux440ux456ux448ux435ux43dux43dux44f}

Розв'язок

\[
u(x, t) = A \operatorname{erfc}[x/2\alpha V \overline{t}]
\]

дає багато їжі для роздумів. На рисунку 10.6 показані графіки
ймовірнісної комплементарної функції на різних значеннях часу t.

\pandocbounded{\includegraphics[keepaspectratio]{images/_page_79_Figure_10.jpeg}}

Рис. 10.6. Розподіл температури в напівнескінченному стрижні, на кінці
якого підтримується постійна температура.

Вони відрізняються один від одного лише характеристичною шкалою вздовж
осі x.

\section{ЗАВДАННЯ}\label{ux437ux430ux432ux434ux430ux43dux43dux44f-7}

\begin{itemize}
\tightlist
\item
  \begin{enumerate}
  \def\labelenumi{\arabic{enumi}.}
  \tightlist
  \item
    Довести ідентичності (10.1).
  \end{enumerate}
\item
  \begin{enumerate}
  \def\labelenumi{\arabic{enumi}.}
  \setcounter{enumi}{1}
  \tightlist
  \item
    Розв'яжіть задачу Коші (10.2).
  \end{enumerate}
\item
  \begin{enumerate}
  \def\labelenumi{\arabic{enumi}.}
  \setcounter{enumi}{2}
  \tightlist
  \item
    Використайте синусоїдальне або коспусне перетворення для розв'язання
  \end{enumerate}
\end{itemize}

\textbf{z}adachu

\[
\begin{array}{lll} (\text{УЧП}) & u_t = \alpha^2 u_{xx}, & 0 < x < \infty, \\ (\text{ГУ}) & u_x (0, \ t) = 0, & 0 < t < \infty, \\ (\text{НУ}) & u (x, \ 0) = H (1-x), & 0 \leqslant x < \infty, \end{array}
\]

де \(H(x)\) --- функція Хевісайду

\[
H(x) =
\begin{cases} 0, & x < 0, \\ 1, & 0 \le x. \end{cases}
\]

Входящая в начальное условие функция изображена на рисунке.

\pandocbounded{\includegraphics[keepaspectratio]{images/_page_80_Picture_7.jpeg}}

Как будет выглядеть график решения в различные моменты времени?

\begin{verbatim}
5. Найдите прообраз константы при синус-преобразовании, т. е. найдите f(x), если  F($\omega$) = c = \mathrm{const.}
\end{verbatim}

\section{РЯДЫ И ПРЕОБРАЗОВАНИЯ
ФУРЬЕ}\label{ux440ux44fux434ux44b-ux438-ux43fux440ux435ux43eux431ux440ux430ux437ux43eux432ux430ux43dux438ux44f-ux444ux443ux440ux44cux435}

ЦЕЛЬ ЛЕКЦИИ: Ввести понятие ряда Фурье и показать, каким образом
периодическую функцию \(f(x)\) можно представить в гиде ряда по синусам
и косинусам:

\[
f(x) = \frac{a_0}{2} + \sum_{n=1}^{\infty} [a_n \cos(n\pi x) + b_n \sin(n\pi x)].
\]

Если же функция \(f(x)\) непериодическая и определена на (-\infty,
+\infty) , то показать, как ряд Фурье переходит в \emph{преобразование
Фурье} и как функцию \(f(x)\) можно разложить в непрерывную совокупность
простейших функций. Эго разложение (интеграл Фурье) можно записать в
комплексной форме:

\[
f(x) = \frac{1}{\sqrt{2\pi}} \int_{-\infty}^{+\infty} \left[ \frac{1}{\sqrt{2\pi}} \int_{-\infty}^{+\infty} f(x) e^{-i\xi x} dx \right] e^{i\xi x} d\xi,
\]

откуда получаются прямое и обратное преобразования Фурье

\[
\mathcal{F}[f] = F(\xi) = \frac{1}{\sqrt{2\pi}} \int_{-\infty}^{+\infty} f(x) e^{-i\xi x} dx \qquad \text{(пряме перетворення Фур'є),}
\mathcal{F}^{-1}[F] = f(x) = \frac{1}{\sqrt{2\pi}} \int_{-x}^{+x} F(x) e^{i\xi x} d\xi \qquad \text{(обернене перетворення Фур'є)}.
\]

Важлива роль рядів Фур'є в теорії диференціальних рівнянь у похідних
пов'язана з тим, що періодичні функції, визначені в
\((-\infty, +\infty)\) або функції, визначені на скінченному інтерваті,
можна розкласти на ряд за допомогою сіпусів і косинусів. Це дозволяє
спростити початкове завдання. Наприклад, функція пилкоподібної форми

\[
f(x) = x, -L < x < L,
\]

\(f(x + 2L) = f(x)\) (условие периодичности),

изображенная на рис. 11.1, может быть представлена в виде ряда Фурье

(11.1)

\[
f(x) = \frac{a_0}{2} + \sum_{n=1}^{\infty} \left[ a_n \cos(n\pi x/L) + b_n \sin(n\pi x/L) \right],
\]

в котором коэффициенты Фурье \(a_n\) и \(b_n\) находятся по формулам
Эйлера

\[
a_n = \frac{1}{L} \int_{-L}^{L} f(x) \cos(n\pi x/L) dx = 0, \quad n = 0, 1, 2, \dots,
b_n = \frac{1}{L} \int_{-L}^{L} f(x) \sin(n\pi x/L) dx = -(-1)^n (2L/n\pi), \quad n = 1, 2, \dots.
\]

Эти ингегралы легко вычисляются. Чтобы вывести формулы Эйлера для
коэффициентов \(a_n\) и \(b_n\) , умножим обе части разложения (11.1) на
\sin{(n\pi x/L)} или \cos{(n\pi x/L)} и проинтегрируем получив- Р И С

\pandocbounded{\includegraphics[keepaspectratio]{images/_page_82_Figure_7.jpeg}}

. 11.1. Разложение пилообразной функции в ряд Фурье.

\begin{verbatim}
шееся соотношение от -L до L. Ортогональность функций  \{\sin(n$\pi$ x/L)\}  и  \{\cos(n$\pi$ x/L)\}  позволяет срязу получить формулы для коэффициентов  a_n  и  b_n  (см. задачу 6 в конце лекции). Для пилообразной функции разложение в ряд Фурье принимает вид
\end{verbatim}

(11.3)

\[
f(x) = \frac{2L}{\pi} \left[ \sin(\pi x/L) - \frac{1}{2} \sin(2\pi x/L) + \frac{1}{3} \sin(3\pi x/L) - \dots \right].
\]

Каждое слагаемое ряда принято называть гармоникой, частота каждой
следующей гармоники выше предыдущей. Частоты всех гармоник кратны
основной частоте, период которой совпадает

\begin{verbatim}
с периодом функции f(x). Один из недостатков рядов Фурье состоит в том, что в ряд Фурье можно разложить только периодическую функцию. Конечно, если функция определена на конечном интервале (например, f(x) = x при  0 \le x \le 1 ), ее можно периодически продолжить на всю числовую ось и затем использовать разложение (11.1).
\end{verbatim}

Тот факт, что сумма ряда Фурье определена и периодична вне интервала
{[}0, 1{]}, ничему не противоречит, поскольку мы интересуемся поведением
функции только внутри интервала. На самом деле для функции, определенной
внутри интервала, можно указать много различных разложений в ряд Фурье,
поскольку функцию можно продолжить за пределы интервала различными
способами. Сразу же отметим, что полученные ряды будут сходиться к
исходной функции с различными скоростями.

Читатель не должен думать, что всякую периодическую функцию можно
разложить в ряд Фурье. Пока мы знаем только, что если функцию \(f(x)\)
можно разложить в ряд Фурье (11.1), то коэффициенты \(a_n\) и \(b_n\)
находятся по формулам Эйлера (11.2). Более того, даже если функция
разложима в ряд Фурье, мы не знаем, можно ли ряд почленно
дифференцировать. И в самом деле, легко видеть, что производную функции
\(f(x)\) = x (пилообразной функции) нельзя найти почленным
дифференцированием ряда Фурье (11.3).

В литературе, перечисленной в конце книги, можно найти точные условия
разложимости функции \(f(x)\) в ряд Фурье, а также условия почленной
дифференцируемости получающегося ряда. Нам будет достаточно ознакомиться
с известной теоремой Дирихле.

\emph{Теорема Дирихле} (достаточные условия разложимости функции в ряд
Фурье):

Если \(f(x)\)---ограниченная периодическая функция, имеющая на каждом
периоде конечное число максимумов, минимумов и точек разрыва, то ряд
Фурье функции \(f(x)\) сходится к \(f(x)\) в каждой точке непрерывности
функции и к среднему значению пределов функции слева и справа в каждой
точке разрыва функции \(f(x)\).

\begin{verbatim}
Например, на рис. 11.1 ряд Фурье сходится к функции f(x) во всех точках, за исключением точк  x=\pm L, \pm 3L, \ldots  (точки разрыва). В каждой точке разрыва ряд сходится к нулю (среднее арифметическое значение функции в точках  \pm L  и  \pm L ).
\end{verbatim}

Теперь мы педголовлены для знакомства с преобразованием Фурье. Однако
целесообразно будет познакомиться сначала с понятием частотного спектра
периодической функции.

\section{Дискретный частотный спектр периодической
функции.}\label{ux434ux438ux441ux43aux440ux435ux442ux43dux44bux439-ux447ux430ux441ux442ux43eux442ux43dux44bux439-ux441ux43fux435ux43aux442ux440-ux43fux435ux440ux438ux43eux434ux438ux447ux435ux441ux43aux43eux439-ux444ux443ux43dux43aux446ux438ux438.}

Если \(f(x)\) --- периодическая функция, то разложение в ряд Фурье можно
интерпретировать как представление функции \(f(x)\) последовательностью
чисел \{\(c_n\)\} вида

\[
c_n = \sqrt{a_n^2 + b_n^2}, \quad n = 0, 1, 2, \dots,
\]

где коэффициенты \(c_n\) служат мерой вклада соответствующих частотных
компонент в функцию \(f(x)\). Например, пилообразная функция \(f(x)\)
разлагается в ряд Фурье

\[
f(x) = \frac{2L}{\pi} \left[ \sin(\pi x/L) - \frac{1}{2} \sin(2\pi x/L) + \frac{1}{3} \sin(3\pi x/L) - \dots \right]
\]

и, следовательно, частотный спектр \{\(c_n\)\} записывается в виде
\(c_n=2L/n\pi\)для n=1,~2,~\dots  и \(c_0\)=\textbar{}\(a_0\)\textbar=0
(см. рис. 11.2).

\pandocbounded{\includegraphics[keepaspectratio]{images/_page_84_Figure_8.jpeg}}

PNC. 11 2. Дискретный частотный спектр пилообразной функции a---
основная частота (период 2L); b---n-ая гармоника (период 2L/n).

Последовательность \{\(c_n\)\} напоминает спектральное разложение белого
света по цветам в спектроскове.

Введем теперь понятие преобразования Фурье.

\section{Преобразование
Фурье}\label{ux43fux440ux435ux43eux431ux440ux430ux437ux43eux432ux430ux43dux438ux435-ux444ux443ux440ux44cux435}

\begin{verbatim}
Как уже отмечалось рансе, непериодические функции, заданные на  (-\infty, +\infty) , не разлагаются в ряды Фурье Однако для некоторых таких функций можно построить аналог ряда Фурье. Не вдаваясь в детали доказательства, можно утверждать, что для непериодических функций, определенных на  (-\infty, \infty) , ряд Фурье
\end{verbatim}

\[
f(x_{l} = \frac{a_{0}}{2} + \sum_{n=1}^{\infty} \left[ a_{n} \cos(n\pi x/L) + b_{n} \sin(n\pi x/L) \right]
\]

переходит в интеграл Фирье (разложение по непрерывному спектру частот)

(11.4)

\[
f(x) = \int_{0}^{\infty} a(\xi) \cos(\xi x) d\xi + \int_{0}^{\infty} b(\xi) \sin(\xi x) d\xi,
\]

гле

(11.5)

\[
a(\xi) = \frac{1}{\pi} \int_{0}^{\infty} f(x) \cos(\xi x) dx
,  b(\xi) = \frac{1}{\pi} \int_{0}^{\infty} f(x) \sin(\xi x) dx .
\]

\begin{verbatim}
Из (11.4) видно, что интеграл Фурье осуществляет разложение функции f(x) в спектр по всем частотам  0 < \xi < \infty  (а не Р И С
\end{verbatim}

\pandocbounded{\includegraphics[keepaspectratio]{images/_page_85_Figure_7.jpeg}}

11.3. Частотные спектры некоторых функций.

только по частотам, кратным основной частоте, как это было у
периодических функций). По аналогии с рядами Фурье определим частотный
спектр по формуле

\[
C(\xi) = \sqrt{a^2(\xi) + b^2(\xi)}.
\]

Частотний спектр є мірою внеску частоти \xi у функ-

иню \(f(x)\).

\begin{verbatim}
На рис. 11.3 изображены некоторые функции f(x) и их спектры. Из рис. 11.3 видно, что функции, графики которых имеют угловые точки, обладают широким частотным спектром. Объясняется это тем, что для передачи излома графкка требуются высокочастотные спектральные компоненты. С другой стороны, если  f(x) = \sin \xi_0 x —простая периодическая функция, то ее частотный спектр обращается в нуль всюду, за исключением точки  \xi = \xi_0 .
\end{verbatim}

Теперь мы готовы к тому, чтобы познакомиться с экспоненциальной формой
преобразования Фурье (формулы 11.5 известны как синус- и
косинус-преобразования Фурье). Воспользуемся фор-

мулой Эйлера

\[
e^{i\theta} = \cos \theta + i \sin \theta
\]

и носле несложных преобразований приведем (11.4) к виду

(11.6)

\[
f(x) = \frac{1}{\sqrt{2\pi}} \int_{-\infty}^{+\infty} \left[ \frac{1}{\sqrt{2\pi}} \int_{-\infty}^{+\infty} f(x) e^{-i\xi x} dx \right] e^{i\xi x} d\xi,
\]

которое принято называть интегральным представлением Фурье. Отсюда
получаем пару интегральных преобразований

(11.7)

\[
\mathcal{F}[f] = F(\xi) = \frac{1}{\sqrt{2\pi}} \int_{-\infty}^{+\infty} f(x) e^{-i\xi x} dx
\]

(прямое преобразование Фурье),

\[
\mathcal{F}^{-1}[F] = f(x) = \frac{1}{\sqrt{2\pi}} \int_{-\infty}^{+\infty} F(\xi) e^{i\xi x} d\xi
\]

(обратное преобразование Фурье).

В следующей лекции мы остановимся на изучении свойств этих
преобразований.

\section{ЗАМЕЧАНИЯ}\label{ux437ux430ux43cux435ux447ux430ux43dux438ux44f}

\begin{verbatim}
1. Фурье-образ  F(\xi)  функции f(x) может оказаться комплекснозначной функцией. Например, для функции
\end{verbatim}

\[
f(x) =
\begin{cases} 0, & x \le 0, \\ e^{-x}, & x > 0, \end{cases}
\]

Фур'є-образ має вигляд

\[
F(\xi) = \frac{1}{\sqrt{2\pi}} \frac{1 - i\xi}{1 + \xi^2}.
\]

\begin{enumerate}
\def\labelenumi{\arabic{enumi}.}
\setcounter{enumi}{1}
\tightlist
\item
  Абсолютне значення зображення Фур'є \(|F(\xi)|\) називається частотним
  спектром функції \(f(x)\). Наприклад, частотний спектр функції
  \(f(x)\) з попереднього зауваження визначається формулою
\end{enumerate}

\textbar F(\xi)\textbar{} = \sqrt{\frac{1}{2$\pi$(1+\xi^2)}}

(зчитувач повинен мати змогу знайти модуль комплексного числа).

\begin{enumerate}
\def\labelenumi{\arabic{enumi}.}
\setcounter{enumi}{2}
\tightlist
\item
  Не всі функції піддаються перетворенню Фур'є (інтеграл (11.7) може не
  існувати). Наприклад, функції \(f(x)\) = c, \(\sin x\) , \(e^x\) , не
  мають \(x^2\) образу Фур'є. Лише ті функції, які мають тенденцію
  швидко нульовою на \(|x| \to \infty\) 1). У прикладних задачах ми
  застосовуємо перетворення Фур'є до температурних, хвильових та інших
  фізичних полів, які мають тенденцію до нуля на \(|x| \to \infty\) .
\end{enumerate}

\section{ЗАВДАННЯ}\label{ux437ux430ux432ux434ux430ux43dux43dux44f-8}

\begin{enumerate}
\def\labelenumi{\arabic{enumi}.}
\tightlist
\item
  Розташуйте функцію у рядку Фур'є над синусами
\end{enumerate}

\[
f(x) =
\begin{cases} -1, & -1 < x < 0, \\ 1, & 0 \le x < 1, \end{cases}
f(x+2) = f(x)
\]

(умова періодичності).

Побудуйте суму перших 2, 3 і 4 членів, щоб слідувати збіжності ряду до
\(f(x)\). Побудуйте графіки частотного спектра функції \(f(x)\).

\begin{itemize}
\tightlist
\item
  \begin{enumerate}
  \def\labelenumi{\arabic{enumi}.}
  \setcounter{enumi}{1}
  \tightlist
  \item
    Покажіть, що якщо розрізнити ряд Фур'є (11.3) функції пилки за
    членом, отримаємо послідовність, яка не збігається до похідної
    пилоподібної функції.
  \end{enumerate}
\item
  \begin{enumerate}
  \def\labelenumi{\arabic{enumi}.}
  \setcounter{enumi}{2}
  \tightlist
  \item
    Побудуйте частотний спектр наступних періодичних функцій:
  \end{enumerate}
\item
  \begin{enumerate}
  \def\labelenumi{\alph{enumi})}
  \tightlist
  \item
    \(f(x) = \sin x\) ,
  \end{enumerate}
\item
  \begin{enumerate}
  \def\labelenumi{\arabic{enumi})}
  \setcounter{enumi}{5}
  \tightlist
  \item
    \(f(x) = \sin x + \cos 2x,\) c)
    \(f(x) = \sin x + \cos x + 0.5 \sin 3x\) .
  \end{enumerate}
\end{itemize}

\begin{enumerate}
\def\labelenumi{\arabic{enumi}.}
\setcounter{enumi}{3}
\tightlist
\item
  Знайти \(F(\xi)\) зображення Фур'є та частотний спектр
  \(C(\xi) = |F(\xi)|\) функцію
\end{enumerate}

\[
f(x) =
\begin{cases} 1, & -1 < x < 1, \\ 0 & \text{в остальных точках.} \end{cases}
\]

Тут ми говоримо лише про перетворення Власіка-Фур'є. Ед.

\begin{verbatim}
5. Показати, що абсолютне значення функції F(\xi) = \frac{1}{1+i\xi} дорівнює
\end{verbatim}

\textbar F(\xi)\textbar{} = \sqrt{\frac{1}{1+\xi^2}}.

ПРИМІТКА. Спочатку помножте чисельник і знаменник на 1-i\xi , а потім
позбавтеся комплексної одиниці i у знаменнику. \^{}1 )

\begin{verbatim}
6. Перевірити ортогональність функцій \{\sin(n$\pi$ x/L)\} та \{\cos(n$\pi$ x/L)\} на сегменті [-L, L]:
\end{verbatim}

\[
\int_{-L}^{L} \sin(m\pi x/L) \sin(n\pi x/L) dx =
\begin{cases} 0, & m \neq n, \\ L, &m = n, \end{cases}
\int_{-L}^{L} \cos(m\pi x/L) \cos(n\pi x/L) dx =
\begin{cases} 0, & m \neq n, \\ L, &m = n, \end{cases}
\int_{-L}^{L} \sin(m\pi x/L) \cos(n\pi x/L) dx = 0 \text{ npn Bcex } m, n = 1, 2, \dots
\]

\section{Лекція 12}\label{ux43bux435ux43aux446ux456ux44f-12}

\section{ПЕРЕТВОРЕННЯ ФУР'Є ТА ЙОГО ЗАСТОСУВАННЯ ДО РОЗВ'ЯЗАННЯ
ДИФЕРЕНЦІАЛЬНИХ РІВНЯНЬ У ПОХІДНИХ
ПОХІДНИХ}\label{ux43fux435ux440ux435ux442ux432ux43eux440ux435ux43dux43dux44f-ux444ux443ux440ux454-ux442ux430-ux439ux43eux433ux43e-ux437ux430ux441ux442ux43eux441ux443ux432ux430ux43dux43dux44f-ux434ux43e-ux440ux43eux437ux432ux44fux437ux430ux43dux43dux44f-ux434ux438ux444ux435ux440ux435ux43dux446ux456ux430ux43bux44cux43dux438ux445-ux440ux456ux432ux43dux44fux43dux44c-ux443-ux43fux43eux445ux456ux434ux43dux438ux445-ux43fux43eux445ux456ux434ux43dux438ux445}

МЕТА ЛЕКЦІЇ: Проілюструвати важливі властивості перетворення Фур'є та
показати, як ці властивості можна використовувати для розв'язання PPP.
Зокрема, щоб показати, як під впливом перетворення Фур'є операція
диференціації замінюється множенням, а диференціальні рівняння
перетворюються на алгебранхічні рівняння. У цій лекції вводиться поняття
згортки.

\begin{verbatim}
<sup>2)</sup> Легше |F(\xi)| = \frac{|1|}{|1+i\xi|} = \frac{1}{\sqrt{1+\xi^2}} - \Pi_{PVM}, pcd. Якщо функція f(x) визначена в точці -\infty < x < \infty , то її чотиризображення визначається формулою
\end{verbatim}

(12.1)

\[
\mathcal{F}[f] = F(\xi) = \frac{1}{\sqrt{2\pi}} \int_{-\infty}^{+\infty} f(x) e^{-i\xi x} dx.
\]

Отже, якщо у нас є функція \(f(x)\), визначена на дійсній осі, ми
підставляємо її у формулу (12.1) і отримуємо нову функцію, F(\xi)
визначену для -\infty \textless{} \xi \textless{} \infty . Водночас не
слід забувати про конвенції існування раннього перетворення Фур'є (у
цьому випадку класичного).

\begin{verbatim}
Таблиця 1. 12.1 показує деякі функції та їхні фур'є-образи. Більш розгорнута таблиця функцій та їхніх фур'є-образів наведена в додатку. З наведеної вище таблиці видно, що образ Фур'є функції може виявитися комплекснозначною функцією аргументу \xi . У першому прикладі цієї таблиці функція F(\xi) містить як множник уявну одиницю i, і таким чином є комплекснозначною функцією екзистенційного аргументу ( \xi змінюється від -\infty до +\infty ). Іншими словами, аргумент \xi є реальним, а значення функції є складними.
\end{verbatim}

Перетворення Фур'є (як і більшість інших інтегральних перетворень)
виявилося дуже корисним, оскільки воно замінює операцію диференціації на
операцію множення і таким чином перетворює диференціальні рівняння на
алгебранхічні рівняння. Існує багато інших властивостей, які роблять
перетворення Фур'є зручним інструментом для розв'язання задач. Тут ми
перелічимо лише найважливіші з них.

Таблиця 12.1. Перетворення Фур'є деяких функцій

\begin{verbatim}
| Функція / (a) | \Gamma Перетворення Фур'є F ( \xi ) |
\end{verbatim}

\textbar----------------------------------------------------------------------------------\textbar-------------------------------------------------------------------------------------------\textbar{}

1.

\[
f(x) = \begin{cases} e^{-x}, & x \ge 0, \\ -e^{-x}, & x < 0. \end{cases}
\]

F(\xi) = -i \sqrt{\frac{2}{$\pi$}} \frac{\xi}{1 + \xi^2} (Xomnekches
funktsia). \textbar{}

2.

\[
f(x) = \begin{cases} 1, & -1 < x < 1, \\ 0, & x \le -1, x \ge 1 \end{cases}
\]

F(\xi) = \sqrt{\frac{\frac{2}{$\pi$}}{\frac{\sin \xi}{\xi}}}
(Beg.ectronnal (yul.ung). \textbar{}

\begin{verbatim}
| 3. f(x) = e^{-x^2} | F\left(\xi\right)=\frac{1}{\sqrt{2}}e^{-\sqrt{\xi}/2\right)^2} (основна синя функція). |
\end{verbatim}

\section{Властивості Властивості перетворення Фур'є Властивість 1 (Пара
перетворень
Фур'є)}\label{ux432ux43bux430ux441ux442ux438ux432ux43eux441ux442ux456-ux432ux43bux430ux441ux442ux438ux432ux43eux441ux442ux456-ux43fux435ux440ux435ux442ux432ux43eux440ux435ux43dux43dux44f-ux444ux443ux440ux454-ux432ux43bux430ux441ux442ux438ux432ux456ux441ux442ux44c-1-ux43fux430ux440ux430-ux43fux435ux440ux435ux442ux432ux43eux440ux435ux43dux44c-ux444ux443ux440ux454}

\begin{verbatim}
Після застосування перетворення Фур'є до функції f(x), визначеної для -\infty < x < \infty , отримуємо нову функцію F(\xi) за формулою
\end{verbatim}

\[
\mathscr{F}[f] = F(\xi) = \frac{1}{\sqrt{2\pi}} \int_{-\infty}^{+\infty} f(x) e^{-i\xi x} dx.
\]

Якщо застосувати до функції, F(\xi) визначеної для -\infty \textless{}
\xi \textless{} \infty , обернене перетворення Фур'є

\[
\mathscr{F}^{-1}[F] = f(x) = \frac{1}{\sqrt{2\pi}} \int_{-\infty}^{+\infty} F(\xi) e^{i\xi x} d\xi,
\]

\textbf{Тоді} отримуємо початкову функцію \(f(x)\).

\pandocbounded{\includegraphics[keepaspectratio]{images/_page_90_Figure_7.jpeg}}

Рис. 12.1. Графи функції та її фур'є-образу.

Наприклад, (див. рис. 12.1)

\[
e^{-|x|} \xrightarrow{\mathcal{F}} \sqrt{\frac{2}{\pi}} \xrightarrow{1 + \xi^2} \xrightarrow{\mathcal{F}^{-1}} e^{-|x|}.
\]

\section{Властивість 2 (Перетворення
лінійності)}\label{ux432ux43bux430ux441ux442ux438ux432ux456ux441ux442ux44c-2-ux43fux435ux440ux435ux442ux432ux43eux440ux435ux43dux43dux44f-ux43bux456ux43dux456ux439ux43dux43eux441ux442ux456}

Перетворення Фур'є --- це лінійне перетворення, тобто

\[
\mathscr{F}[af+bg]=a\mathscr{F}[f]+b\mathscr{F}[g].
\]

Це легко побачити. Читач може витратити кілька хвилин, щоб ознайомитися
з цією нерухомістю, яка використовується «на кожному кроці». Наприклад,
перетворення Фур'є виразу

\[
e^{-x} + 3e^{-x^2}
\]

має вигляд

\[
\mathscr{F}[e^{-c}] + 3\mathscr{F}[e^{-c^*}].
\]

\section{Властивість 3 (Часткова конвертація
деривативів)}\label{ux432ux43bux430ux441ux442ux438ux432ux456ux441ux442ux44c-3-ux447ux430ux441ux442ux43aux43eux432ux430-ux43aux43eux43dux432ux435ux440ux442ux430ux446ux456ux44f-ux434ux435ux440ux438ux432ux430ux442ux438ux432ux456ux432}

\begin{verbatim}
Якщо нас цікавить, як перетворення Фур'є діє на похідні, то перш за все слід звернути увагу на змінні, за якими диференційована функція. Наприклад, якщо перетворення Фур'є виконується на змінній \boldsymbol{x} (тобто змінна \boldsymbol{x} інтегрована у перетворення), і якщо часткова похідна функції \boldsymbol{u}(\boldsymbol{x},t) на тій самій змінній \boldsymbol{x} перетворена, то наступні формули перетворення похідних є дійсними
\end{verbatim}

\[
\begin{split} \mathcal{F}\left\{u_{x}\right\} &= \frac{1}{\sqrt{2\pi}} \int_{-\infty}^{+\infty} u_{x}\left(x,\,t\right) e^{-i\xi x} dx = i\xi \mathcal{F}\left[u\right], \\ \mathcal{F}\left\{u_{xx}\right\} &= \frac{1}{\sqrt{2\pi}} \int_{-\infty}^{+\infty} u_{xx}\left(x,\,t\right) \, \delta^{-i\xi x} dx = -\xi^{2} \mathcal{F}\left[u\right]. \end{split}
\]

С другой стороны, если подвергнуть преобразованию частную производную
\(u_t\)(x, t) (а переменная интегрирования в преобразовании, как и
прежде, будет x), получаем следующие формулы:

\[
\begin{aligned} \mathcal{F}\left[u_{t}\right] &= \frac{1}{\sqrt{2\pi}} \int_{-\infty}^{+\infty} u_{t}\left(x, \ t\right) e^{-i\frac{\pi}{2}x} dx = \frac{\partial}{\partial t} \mathcal{F}\left[u\right], \\ \mathcal{F}\left[u_{tt}\right] &= \frac{1}{\sqrt{2\pi}} \int_{-\infty}^{+\infty} u_{tt}\left(x, \ t\right) e^{-i\frac{\pi}{2}x} dx = \frac{\partial^{2}}{\partial t^{2}} \mathcal{F}\left[u\right]. \end{aligned}
\]

\section{Властивість 4 (Властивість
згортки)}\label{ux432ux43bux430ux441ux442ux438ux432ux456ux441ux442ux44c-4-ux432ux43bux430ux441ux442ux438ux432ux456ux441ux442ux44c-ux437ux433ux43eux440ux442ux43aux438}

Кожне інтегральне перетворення має так звану властивість згортки. Факт у
тому, що образ Фур'є добутку двох функцій загалом не відрізняється від
добутку образів Фур'є цих функцій, тобто

\[
\mathcal{F}[f(x)g(x)] \neq \mathcal{F}[f]\mathcal{F}[g].
\]

Однак у теорії перетворень існує операція, яка називається згорткою
функцій f і g (позначена як f*g), яка певною мірою виконує роль добутку,
а саме відношення

(12.2)

\[
\mathcal{F}[/*g] = \mathcal{F}[/]\mathcal{F}[g].
\]

Що це за загадкова операція під назвою згортка? Вона визначається
формулою

(12.3)

\[
(f*g)(x) = \frac{1}{\sqrt{2\pi}} \int_{-\infty}^{+\infty} f(x-\xi) g(\xi) d\xi.
\]

Неважко показати, що співвідношення (12,2) дійсно виконується. З
визначення видно, що згортка двох функцій \(f(x)\) і \(g(x)\) є новою
функцією (f*g)(x) аргументу x.

\section{Приклад обчислення згортки двох
функцій}\label{ux43fux440ux438ux43aux43bux430ux434-ux43eux431ux447ux438ux441ux43bux435ux43dux43dux44f-ux437ux433ux43eux440ux442ux43aux438-ux434ux432ux43eux445-ux444ux443ux43dux43aux446ux456ux439}

Нехай задано дві функції

\[
f(x) = x,
g(x) = e^{-x^2}.
\]

Тоді їхня згортка дорівнює

\[
(f \circ g)(x) = \frac{1}{\sqrt{2\pi}} \int_{-\infty}^{\infty} (x - \xi) e^{-\xi^2} d\xi = x/\sqrt{2}.
\]

Тут ми використали відому формулу

\[
\int_{-\infty}^{+\infty} e^{-\xi^2} d\xi = \sqrt{\pi}.
\]

Важливість концепції згортки (12.3) полягає в тому, що дуже часто,
розв'язуючи диференціальні рівняння, ми стикаємося з необхідністю
застосовувати обернене перетворення Фур'є до виразу, який є добутком
двох образів Фур'є, тобто потрібно знайти

(12.4)

\[
\mathcal{F}^{-1}\left\{\mathcal{F}\left[f\right]\mathcal{F}\left[g\right]\right\}.
\]

Застосовуючи обернене перетворення до обох частин (12.2), отримуємо

\[
f * g = \mathcal{J}^{-1} \{ \mathcal{F} [f] \mathcal{F} [g] \}.
\]

Отже, щоб визначити (12.4), потрібно знайти прототипи кожного з
факторів, тобто функцій f і g, а потім обчислити їхню згортку. Тепер ми
готові розв'язати одну важливу проблему: теорію диференціальних рівнянь
у похідних похідних.

\section{Розв'язання проблеми
Коші}\label{ux440ux43eux437ux432ux44fux437ux430ux43dux43dux44f-ux43fux440ux43eux431ux43bux435ux43cux438-ux43aux43eux448ux456}

Розглянемо проблему поширення тепла в \emph{нескінченному} стрижні, якщо
задано початкову температуру \(u(x, 0) = \varphi(x)\) .

Цю задачу зазвичай називають задачою Коші для рівняння теплопровідності
або задачою початкових умов

\[
\begin{array}{llll} (\text{УЧП}) & u_t = \alpha^2 u_{xx}, & -\infty < x < \infty, & 0 < t < \infty, \\ (\text{12.6}) & (\text{HV}) & u(x,0) = q(x), & -\infty < x < \infty. \end{array}
\]

Решение задачи можно разбить на три основных шага.

ШАГ 1 (Преобразование задачи).

Поскольку пространственная переменная x изменяется в пределах от
\(-\infty\) до \(+\infty\) , мы подвергнем уравнение и начальное условие
(12.6) преобразованию Фурье по переменной \(x^{(1)}\) (x---переменная
интегрирования в преобразовании). Сделав это, получим

\[
\mathcal{F}[u_t] = \alpha^2 \mathcal{F}[u_{xx}],
\mathcal{F}[u(x, 0)] = \mathcal{F}[\varphi(x)].
\]

Воспользуемся свойствами преобразования Фурье. Тогда

\[
\frac{\frac{dU(t)}{dt} = -\alpha^2 \xi^2 U(t), }{U(0) = \Phi(\xi)} = (\text{here } \Phi - \phi) \text{rje-image funcin } \phi),
\]

где \(U(t)\) = \mathcal{F}{[}\(u(x, t)\){]} . Читате ть должен обратить
внимание на то, что функция \(U(t)\) на салот доле зависит не только от
t, но и от \xi , но, поскольку в дифференциальном уравнении (12.7)
величина \xi  играет роль константы, мы не будем записывать ее в число
аргументов, т. е. будем считать, что U = \(U(t)\).

ШАГ 2 (Решение преобразованной задачи).

Поскольку новая переменная играет роль гонстанты в задаче (12.7), ее
решение легко находится:

(12.8)

\[
U(t) = \Phi(\xi) e^{-x^2 e^{tt}}.
\]

ШАГ 3 (Нахождение обратного преобразования). Искомое решение \(u(x, t)\)
находится по формуле

\[
u(x, t) = \mathcal{F}^{-1} [U(\xi, t)] =
= \mathcal{F}^{-1} [\Phi(\xi) e^{-\alpha \xi^2 t}].
\]

Предполагается, что преобразование Фурье этих функций существует. Прим.
ред.

Вот здесь нас выручит теорема о свертке (12.5). Используя эту теорему,
получаем

(12.9)

\[
u(x, t) = \mathcal{F}^{-1} \left[ \Phi(\xi) e^{-\alpha r \xi^{2} t} \right] =
= \mathcal{F}^{-1} \left[ \Phi(\xi) \right] * \mathcal{F}^{-1} \left[ e^{-\alpha^{2} \xi^{2} t} \right] =
= \Phi(x) * \left[ \frac{1}{2\alpha \sqrt{\pi t}} e^{-x^{2}/4\alpha^{2} t} \right] \text{ (знайдено в таблиці)} =
= \frac{1}{2\alpha \sqrt{\pi t}} \int_{-\infty}^{+\infty} \Phi(\xi) e^{-(x-\xi)^{2}/4\alpha^{2} t} d\xi.
\]

Дело сделано! Формула (12.9) дает решение нашей задачи. Прежде чем
поставить точку, давайте проанализируем полученный результат. Отметим,
что почынтегральная функция состои с

из двух сомножителей:

\begin{enumerate}
\def\labelenumi{\arabic{enumi}.}
\item
  Начальной температуры \(q(x)\),

  \begin{enumerate}
  \def\labelenumii{\arabic{enumii}.}
  \setcounter{enumii}{1}
  \tightlist
  \item
    Функции G(x, t) = \frac{1}{2$\alpha$ \sqrt{$\pi$ t}}
    e\textsuperscript{\{-(x-\xi)}2/4c\^{}2t\} , которая называется
    функцией Грина или функцией источника.
  \end{enumerate}
\end{enumerate}

Можно показать, что функция источника \(G(x, t)\) описывает отклик
системы на начальный температурный единичный импульс Р И С

\pandocbounded{\includegraphics[keepaspectratio]{images/_page_94_Figure_9.jpeg}}

\begin{verbatim}
. 12.2. Отклик G(x,t) на температурный импульс в точке  x=\xi , графчк G(x,t) похож на график нормального распределения, при этом время игриег роль среднеквадратического уклонения: a—функция G(x,t) при пебольших значениях t,  $\delta$ —функция G(x,t) сущестьенно уменьши...ась.

в ючке  x=\xi . Другими словами, функция  G(\kappa,t)  описывает распределение температуры в стержие в мочент времени t, если на точку  x=\xi  подействовал единичный тепловой импульс (см. рис. 12.2).

Теперь формуле (12.9) можно дать следующую интерпретацию: начальную температуру  u\left(x,\,0\right)=$\phi$\left(x\right)  можно представлять как конининуальное множество точечных импульсов каждой величины

q(\xi)  (в точке  x = \xi ). Каждый точечный импульс дает распределение температуры  $\phi$(\xi)  G(x, t). Результирующее распределение находится суммированием (интегрированием) температур точечных источников по формуле (12.9). Позже мы увидим, что это одно из проявлений общего принципа суперпозиции.
\end{verbatim}

\section{ЗАМЕЧАНИЕ}\label{ux437ux430ux43cux435ux447ux430ux43dux438ux435}

Главный недостаток преобразования Фурье в том, что его можно применять
не ко всем функциям. Например, функции вида

\[
f(x) = \text{const},
f(x) = e^x,
f(x) = \sin x
\]

не имеют фурье-образа, так как интеграт

\[
\mathcal{F}[f] = \frac{1}{\sqrt{2\pi}} \int_{-\pi}^{\pi} f(x) e^{-i\xi x} dx
\]

расходится. Фурье-образ имеют только те функции, которые достаточно
быстро сгремятся к нулю при \textbar x\textbar{} \to \infty . Кроме
того, нельзя применять преобразование Фурье по времени для решения задач
с начальными условиями, так как 0 \textless{} t \textless{}
\infty  (если не рассматривать продолжения функций на полуось в
соответствующем классе функций).

\section{ЗАДАЧИ}\label{ux437ux430ux434ux430ux447ux438-1}

\begin{enumerate}
\def\labelenumi{\arabic{enumi}.}
\tightlist
\item
  Найдите фурье-образ функции
\end{enumerate}

\[
f(x) =
\begin{cases} 0, & x \leq 0, \\ e^{-x}, & 0 < x, \end{cases}
\]

используя таблицы, приведенные в приложении.

\begin{itemize}
\tightlist
\item
  \begin{enumerate}
  \def\labelenumi{\arabic{enumi}.}
  \setcounter{enumi}{1}
  \tightlist
  \item
    Убедитесь в 10м, что прямое и обратное преобразования Фурье являются
    линейными преобразованиями.
  \end{enumerate}
\item
  \begin{enumerate}
  \def\labelenumi{\arabic{enumi}.}
  \setcounter{enumi}{2}
  \tightlist
  \item
    Используя преобразование Фурье, решите задачу Коши
  \end{enumerate}
\end{itemize}

\[
\begin{array}{lll} (\text{УЧП}) & u_t = \alpha^2 u_{xx}, & -\infty < x < \infty, \\ (\text{HV}) & u(x, 0) = e^{-x^2}, & -\infty < x < \infty. \end{array}
\]

\begin{enumerate}
\def\labelenumi{\arabic{enumi}.}
\setcounter{enumi}{3}
\tightlist
\item
  Проверьте справедливость соотношений
\end{enumerate}

\[
\begin{aligned} \mathcal{F}\left[u_{x}\right] &= i\xi \mathcal{F}\left[u\right], \\ \mathcal{F}\left[u_{xx}\right] &= -\xi^{2} \mathcal{F}\left[u\right]. \end{aligned}
\]

УКАЗАНИЕ, Воспользуйтесь интегрированием по частям.

\begin{enumerate}
\def\labelenumi{\arabic{enumi}.}
\setcounter{enumi}{4}
\tightlist
\item
  Перевірити, чи можна записати згортку двох функцій f і g у двох
  еквівалентних формах
\end{enumerate}

\[
(f * g)(x) = \frac{1}{\sqrt{2\pi}} \int_{-\infty}^{+\infty} f(\xi) g(x - \xi) d\xi
\]

или

\[
(f * g)(x) = \frac{1}{\sqrt{2\pi}} \int_{-\infty}^{+\infty} f(x-\xi) g(\xi) d\xi.
\]

Лекция 13

\section{ПРЕОБРАЗОВАНИЕ
ЛАПЛАСА}\label{ux43fux440ux435ux43eux431ux440ux430ux437ux43eux432ux430ux43dux438ux435-ux43bux430ux43fux43bux430ux441ux430}

ЦЕЛЬ ЛЕКЦИИ: Ввести понятие преобразования Лапласа

\[
\mathscr{L}[f] = F(s) = \int\limits_0^\infty f(t)\,e^{-st}dt
\]

(преобразование Лапласа),

познакомиться с формулой обращения преобразования Лапласа

\[
\mathscr{L}^{-1}[F] = f(t) = \frac{1}{2\pi t} \int_{c-t\infty}^{c+t\infty} F(s) e^{st} ds
\]

(обратное преобразование Лапласа)

и проиллюстрировать их свойства.

Преобразование Лапласа имеет важное преимущество перед преобразованием
Фурье, так как в подынтегральное выражение входит быстро затухающий
множитель \(e^{-st}\) , и поэтому преобразование Лапласа применимо к
более широкому классу функций по сравнению с преобразованием Фурье.
Кроме того, преобразование Лапласа имеет дело с функциями, определенными
на промежутке \([0, \infty)\) , а значит, особенно удобно проводить его
по переменной t.

После знакомства с основными свойствами преобразования Лапласа мы решим
несколько важных задач для уравнений с частными производными 113 всех
интегральных преобразований, рассматриваемых в этой книге,
преобразование Лапласа

(13 1)

\[
\mathscr{Z}[f] = \int_{0}^{\infty} f(t) e^{-st} dt,
\]

здається, це єдиний, з яким читач стикався раніше, адже це один із
потужних методів перетворення звичайних диференціальних рівнянь у
алгебранхічні рівняння. Крім того, перетворення Лапласа дозволяє
перетворити диференціальне рівняння в похідних у звичайне диференціальне
рівняння, і його можна виконати на будь-якій із тих змінних,
\(x, y, z, t \ldots\) які змінюються від 0 до \(\infty\) (найчастіше
перетворення Лапласа виконується на змінній t).

Якщо при застосуванні перетворення Лапласа до звичайної диференціальної
еквалізації отримують рівняння алгебрану, то у випадку часткових
диференціальних рівнянь перетворення Лапласа призводить або до нового
диференціального рівняння з меншою кількістю змінних, або до звичайного
диференціального рівняння. Перш ніж перейти до розв'язання дуже цікавих
задач, перелічимо деякі властивості перетворення Лапласа

\section{Свойства преобразования
Лапласа}\label{ux441ux432ux43eux439ux441ux442ux432ux430-ux43fux440ux435ux43eux431ux440ux430ux437ux43eux432ux430ux43dux438ux44f-ux43bux430ux43fux43bux430ux441ux430}

Свойство 1 (Прямое и обратное преобразования)

Преобразование Лапласа и формула его обращения имеют вид

\[
\mathscr{L}[f] = F(s) = \int_{0}^{\infty} f(t) \, e^{-st} \, dt \quad \text{(Перетворення Лапласа),}
\]

(13.2)

\[
\mathscr{L}^{-1}[F] = f(t) = \frac{1}{2\pi i} \int_{c-i\infty}^{c+i\infty} F(s) \, e^{st} \, ds \quad \text{(формула інверсії).}
\]

Преобразование Лапласа обладает важным преимуществом по сравнению с
преобразованием Фурье. Это преимущество обусловлено наличием затухающего
множителя \(e^{-st}\) в подынтегральном выражении, что обеспечивает
возможность применения преобразования Лапласа к более широкому классу
функций (множитель \(e^{t\xi t}\) в преобразовании Фурье не является
затухающим, поскольку модуль этого множителя равен единице). Более
точное представление об этом классе функций дает следующая теорема:

\section{Достаточные условия существования преобразования
Лапласа}\label{ux434ux43eux441ux442ux430ux442ux43eux447ux43dux44bux435-ux443ux441ux43bux43eux432ux438ux44f-ux441ux443ux449ux435ux441ux442ux432ux43eux432ux430ux43dux438ux44f-ux43fux440ux435ux43eux431ux440ux430ux437ux43eux432ux430ux43dux438ux44f-ux43bux430ux43fux43bux430ux441ux430}

Ести

\begin{enumerate}
\def\labelenumi{\arabic{enumi})}
\item
  для любого положительного числа A функция \(f(t)\) кусочнонепрерывна
  на интервале \(0 \le t \le A\) ;
\item
  существуют такие константы M, a и T, что \(|f(t)| \leq Me^{at}\) для
  всех t \textgreater{} T, то преобразование Лапласа
\end{enumerate}

\[
\mathscr{L}[f(\ell)] = F(s) = \int_{0}^{\infty} f(\ell) e^{-st} dt
\]

существует при всех s \textgreater{} a.

Приведем примеры некоторых функций и их преобразований Лапласа

У визначенні перетворення Лапласа мається на увазі, що s є дійсною
змінною, що належить інтервалу \((0, \infty)\) . Однак можливо (і часто
бажано) розширити це визначення на комплексні значення s

\[
\mathscr{L}^{-1}[F] = \frac{1}{2\pi i} \int_{\epsilon - i\infty}^{\epsilon + i\infty} F(s) e^{st} ds.
\]

Использование этой формулы подразумевает знакомство с техникой
интегрирования по контурам в комплексной плоскости и теорией вычетов Мы
не станем отвлекаться на изучение здесь этих вопросов, а при обращении
преобразования Лапласа будем пользоваться таблицей, приведенной в
приложении

\section{Свойство 2 (Преобразования частных
производных)}\label{ux441ux432ux43eux439ux441ux442ux432ux43e-2-ux43fux440ux435ux43eux431ux440ux430ux437ux43eux432ux430ux43dux438ux44f-ux447ux430ux441ux442ux43dux44bux445-ux43fux440ux43eux438ux437ux432ux43eux434ux43dux44bux445}

Пусть \(u(x, t)\) --- функция переменных, и мы хотим найти результаты
применения преобразования Лапласа к различным частным производным
\(u_t\) , \(u_{tt}\) , \(u_x\) , \(u_{\lambda x}\) . . . . Поскольку
преобразование Лапласа проводится по переменной t (t --- переменная
интегрирования), правила преобразования имеют вид

\[
\mathcal{L}\left[u_{t}\right] = \int_{0}^{\infty} u_{t}(x, t) e^{-st} dt = sU(x, s) - u(x, 0),
\mathcal{L}\left[u_{tt}\right] = \int_{0}^{\infty} u_{tt}(x, t) e^{-st} dt = s^{2}U(x, s) - su(x, 0) - u_{t}(x, 0),
\mathcal{L}\left[u_{x}\right] = \int_{0}^{\infty} u_{x}(x, t) e^{-st} dt = \frac{\partial U}{\partial x}(x, s),
\mathcal{L}\left[u_{xx}\right] = \int_{0}^{\infty} u_{xx}(x, s) e^{-st} dt = \frac{\partial^{2}U}{\partial x^{2}}(x, s),
\]

\begin{enumerate}
\def\labelenumi{\arabic{enumi}.}
\tightlist
\item
\end{enumerate}

\[
f(t) = 1
,  0 < t < \infty    (M = 1, a = 0)
\]

\pandocbounded{\includegraphics[keepaspectratio]{images/_page_99_Figure_3.jpeg}}

\[
F(s) = \int_{0}^{\infty} e^{-st} dt = \frac{1}{s}
\]

\pandocbounded{\includegraphics[keepaspectratio]{images/_page_99_Figure_5.jpeg}}

2

\[
f(t) = e^{2t}
,  0 < t < \infty
(M = 1,  n = 2 )
\]

\pandocbounded{\includegraphics[keepaspectratio]{images/_page_99_Figure_7.jpeg}}

\[
F(s) = \frac{1}{s-2}
,  s > 2
\]

a

\pandocbounded{\includegraphics[keepaspectratio]{images/_page_99_Figure_9.jpeg}}

\begin{enumerate}
\def\labelenumi{\arabic{enumi}.}
\setcounter{enumi}{2}
\tightlist
\item
\end{enumerate}

\[
f(t) = \sin(\omega t)
(M=1, \alpha=0)
\]

\pandocbounded{\includegraphics[keepaspectratio]{images/_page_99_Figure_11.jpeg}}

\pandocbounded{\includegraphics[keepaspectratio]{images/_page_99_Figure_12.jpeg}}

\pandocbounded{\includegraphics[keepaspectratio]{images/_page_99_Figure_13.jpeg}}

\begin{enumerate}
\def\labelenumi{\arabic{enumi}.}
\setcounter{enumi}{3}
\tightlist
\item
  \(f(t)\) = \(e^{t^2}\) (не имеет преобразования Лапласа)
\end{enumerate}

РИС. 13.1. Графики некоторых функций и их изображений по Лапласу.

\begin{verbatim}
где  U\left(x,\,s\right)=\mathcal{L}\left[u\left(x,\,t\right)\right] . Правила преобразования производных  u_{x}  и  u_{xy}  вытекают из теоремы о дифференцировании интеграла по параметру
\end{verbatim}

\[
\frac{d}{dx}\int_{a}^{b}f\left(x,\ y\right)dy=\int_{a}^{b}\frac{df}{dx}\left(x,\ y\right)dy,
\]

а формулы для преобразования производных \(u_t\) и \(u_{tt}\) можно
получить интегрированием по частям.

\section{Свойство 3 (Теорема о
свертке)}\label{ux441ux432ux43eux439ux441ux442ux432ux43e-3-ux442ux435ux43eux440ux435ux43cux430-ux43e-ux441ux432ux435ux440ux442ux43aux435}

Свертка играет здесь ту же роль, что и в теории преобразования Фурье, но
вводится несколько иначе.

\section{Определение конечной
свертки}\label{ux43eux43fux440ux435ux434ux435ux43bux435ux43dux438ux435-ux43aux43eux43dux435ux447ux43dux43eux439-ux441ux432ux435ux440ux442ux43aux438}

Конечная свертка двух функций f и g определяется формулами

\[
(f * g) (t) =
\begin{cases} \int_0^t f(\tau) g(t-\tau) d\tau, \\ \int_0^t f(t-\tau) g(\tau) d\tau \end{cases}
\]

(Можна довести, що ці інтеграли рівні). Тут інтегрування здійснюється
вздовж скінченного інтервалу від 0 до t, а не вздовж нескінченної (
\(\infty\) , \(+\infty\) ), як це було раніше з визначенням нескінченної
згортки.

Ось приклад скінченної згортки двох функцій. Нехай

\[
f(t) = t, \quad g(t) = t;
\]

тоді

\[
(f * g)(t) = \int_{0}^{t} \tau(t - \tau) d\tau = t^{3}/6.
\]

Як і у випадку нескінченної згортки, її найважливіша властивість
визначається формулою

(13.4)

\[
\mathscr{L}[f * g] = \mathscr{L}[f] \mathscr{L}[g]
\]

або еквівалентна формула

(13.5)

\[
\mathscr{L}^{-1}\left\{ \left[\mathscr{L}\left[f\right]\mathscr{L}\left[g\right]\right\} = f*g. \right.
\]

Ці формули дозволяють змінити перетворення Лапласа у випадках, коли
образ можна представити як добуток \(\mathscr{L}[f]\) і
\(\mathscr{L}[g]\) таких множників, для яких

Функції F і G легко знайти. Після пошуку функцій f і g можна знайти їхню
згортку. Наприклад,

\[
\mathcal{Z}^{-1}\left[\frac{1}{s} \cdot \frac{1}{s^2+1}\right] = \int_0^t \sin \tau \, d\tau = 1 - \cos t,
F(s) = \frac{1}{s} \xrightarrow{\mathcal{Z}^{-1}} f(t) = 1 \quad G(s) \xrightarrow{\mathcal{Z}^{-1}} g(t) = \sin t.
\]

Тепер ми готові до важливого змішаного виклику.

\section{Теплопровідність у напівнескінченному
середовищі}\label{ux442ux435ux43fux43bux43eux43fux440ux43eux432ux456ux434ux43dux456ux441ux442ux44c-ux443-ux43dux430ux43fux456ux432ux43dux435ux441ux43aux456ux43dux447ux435ux43dux43dux43eux43cux443-ux441ux435ux440ux435ux434ux43eux432ux438ux449ux456}

Розглянемо глибокий резервуар з рідиною, і нехай бічна поверхня
резервуара буде теплоізольованою. Нехай \(u_0\) буде початковою
температурою рідини, а температура повітря над рідиною ---

\pandocbounded{\includegraphics[keepaspectratio]{images/_page_101_Picture_7.jpeg}}

Рис. 13.2. Схематичне представлення проблеми теплопровідності: a ---
тепло тече всередину, якщо ви(0, t) \textless{} 0, і назовні, якщо ви(0,
t) \textgreater{} 0; b --- резервуар настільки глибокий, що гранична
умова внизу не впливає на розв'язок при значеннях x.

кістка дорівнює нулю (орієнтир температур) (див. рис. 13.2). Наше
завдання --- знайти температурне поле в рідині на різних глибинах і в
різний час. Іншими словами, нам потрібно розв'язати проблему

(13.6)

\[
\begin{array}{lll} (\text{УЧП}) & u_t = u_{xx}, & 0 < x < \infty, & 0 < t < \infty, \\ (\Gamma \text{У}) & u_x (0, t) - u (0, t) = 0, & 0 < t < \infty, \\ (\text{НУ}) & u (x, 0) = u_0, 0 < x < \infty. \end{array}
\]

Щоб розв'язати цю задачу, ми застосуємо перетворення Лапласа до змінної
t. Зверніть увагу, що можна також застосувати перетворення Лапласа до
змінної x, оскільки воно змінюється від 0 до \infty . Після застосування
перетворення

Ми отримуємо звичайне диференціальне рівняння за змінною X.

(13.7)

\[
(ODY) \quad sU(x) - u_0 = \frac{d^2U}{dx^2}, \quad 0 < x < \infty,
(\Gamma Y) \quad \frac{dU}{dx}(0) = U(0)
\]

(Ми перетворили лише рівняння та граничну умову, а не початкову умову.)
Маємо звичайне диференціальне рівняння другого порядку, але з однією
граничною умовою при x=0. Насправді існує ще одна гранична умова. Він
визначається фізичними міркуваннями і фіксується як \(U(x) \to 0\) ,
коли \(x \to +\infty\) . Зверніть увагу, що для спрощення позначення ми
опустили параметр s із аргументів функції \(U(x)\) всюди, оскільки
рівняння (13.7) є диференціальним рівнянням лише відносно x.

Щоб розв'язати задачу (13.7), запишіть загальне рішення (загальний
розв'язок однорідного + часткового розв'язку гетерогенного) TAC:

\[
U(x) = c_1 e^{V_{5x}} + c_2 e^{-V_{5x}} + u_0/s.
\]

Підставляючи цей вираз у граничну умову (13.7), можна визначити
константи \(c_1\) та \(c_2\) (одразу зрозуміло, що \(c_3 = 0\) ,

\pandocbounded{\includegraphics[keepaspectratio]{images/_page_102_Figure_7.jpeg}}

Рис. 13.3. Граф ймовірнісного інтеграла (eri) та його комплементарної
функції (eric).

Інакше температура буде зростати нескінченно зі збільшенням координати
X). Визначаючи сталу \(c_2\) з граничної умови при x=0, отримуємо
фінальний вираз для \(U\left(x\right)\) (13.8)

\[
U(x) = u_0 \left\{ \frac{e^{-V \widehat{s} x}}{s(V \widehat{s} + 1)} \right\} + \frac{1}{s}.
\]

Залишається зробити останній крок. Щоб визначити температурне поле,
потрібно обчислити \(u\left(x,\,t\right)\)

\[
u(x, t) = \mathcal{L}^{-1}[U(x, s)]
\]

(ми повернулися до написання всіх аргументів функції \(U(x, s)\)). Щоб
повернути перетворення Лапласа, ми використаємо таблиці, наведені в
додатку. У результаті ми отримуємо

(13.9)

\[
u(x, t) = u_0 - u_0 \left[ \text{eric} \left( \frac{x}{2} \sqrt{t} \right) + \text{eric} \left( \sqrt{t} + \frac{x}{2} \sqrt{t} \right) e^{x+t} \right],
\]

де

\[
\operatorname{eric}(x) = \frac{2}{\sqrt{\pi}} \int_{x}^{\infty} e^{-\xi^{2}} d\xi
\]

є функцією, доповнюючою до ймовірнісного інтеграла. Графіки цієї функції
показані на рис. 13.3. Аналіз цього рішення P H C

\pandocbounded{\includegraphics[keepaspectratio]{images/_page_103_Figure_7.jpeg}}

13.4. Температура всередині напівнескінченного середовища в різний час.

Потребує дуже малих зусиль. Фіг. (13.4) представлено результати
чисельного аналізу, отриманого на комп'ютері, оснащеному плоттером.

\section{ПРИМІТКИ}\label{ux43fux440ux438ux43cux456ux442ux43aux438-1}

\begin{itemize}
\tightlist
\item
  \begin{enumerate}
  \def\labelenumi{\arabic{enumi}.}
  \tightlist
  \item
    Перетворення Лапласа також можна застосувати до неоднорідного
    диференціального рівняння з частковими похідними (для застосування
    методу розділення змінних рівняння має бути однорідним), згідно з
    ним, загалом, його можна застосовувати до рівнянь із змінними
    коефіцієнтами (метод розділення змінних також може застосовуватися
    до змінних коефіцієнтів). Таблиця 1. 13.2 вказує сфери застосування
    кожного з двох методів.
  \end{enumerate}
\item
  \begin{enumerate}
  \def\labelenumi{\arabic{enumi}.}
  \setcounter{enumi}{1}
  \tightlist
  \item
    Для розв'язання змішаних і граничних задач також використовуються
    перетворення Ганкеля та Мелліна, які, однак, відрізняються від
    перетворення Лапласа в одному аспекті. Перетворення Лапласа замінює
    похідну операцією множення формули
  \end{enumerate}
\end{itemize}

\[
\mathscr{L}[y'] = s\mathscr{L}[y] - y(0),
\]

Таблиця 13.2 Порівняння методів розділення перетворення Лапласа та
змінних

{\def\LTcaptype{none} % do not increment counter
\begin{longtable}[]{@{}
  >{\raggedright\arraybackslash}p{(\linewidth - 4\tabcolsep) * \real{0.2907}}
  >{\raggedright\arraybackslash}p{(\linewidth - 4\tabcolsep) * \real{0.3140}}
  >{\raggedright\arraybackslash}p{(\linewidth - 4\tabcolsep) * \real{0.3953}}@{}}
\toprule\noalign{}
\begin{minipage}[b]{\linewidth}\raggedright
\end{minipage} & \begin{minipage}[b]{\linewidth}\raggedright
Метод
\end{minipage} & \begin{minipage}[b]{\linewidth}\raggedright
\end{minipage} \\
\midrule\noalign{}
\endhead
\bottomrule\noalign{}
\endlastfoot
& Перетворення Лапласа & Розділення \\
Гетерогенні PPP & Так & Так \\
Гетерогенні ГІ & Так & Так \\
Змінні ймовірності & Так & de' \\
Нелінійні рівняння & Так & Так \\
\end{longtable}
}

а перетворення Ганкеля і Мелліна замінюють диференціальний оператор на
множник. Наприклад, для перетворення Ганкеля, визначеного формулою

\[
H[y] = \int_{0}^{\infty} r J_{o}(\xi r) y(r) dr,
\]

справедливе співвідношення

\[
H\left[y''\left(r\right)+\frac{1}{r}\ y'\left(r\right)\right]=-\xi^{2}H\left[y\right].
\]

Таким чином можна розв'язати спеціальне диференціальне рівняння з
змінними коефіцієнтами (рівняння Бесселя). 3. Перетворення Лапласа на
змінній t можна інтерпретувати як проекцію площини xt на осі x, що
призводить до У ц ь о

\pandocbounded{\includegraphics[keepaspectratio]{images/_page_104_Figure_9.jpeg}}

му випадку початкові рівняння, граничні та початкові умови
перетворюються на нове диференціальне рівняння та нові граничні умови
(див. діаграму вище).

\section{ЗАВДАННЯ}\label{ux437ux430ux432ux434ux430ux43dux43dux44f-9}

\begin{enumerate}
\def\labelenumi{\arabic{enumi}.}
\tightlist
\item
  Переконайтеся, що наступна формула для перетворення часткової похідної
  \(u_t\) є істинною:
\end{enumerate}

\[
\mathscr{L}\left[u_{t}(x, t)\right] = sU(x, s) - u(x, 0).
\]

\begin{enumerate}
\def\labelenumi{\arabic{enumi}.}
\setcounter{enumi}{1}
\tightlist
\item
  Використати перетворення Лаціаса для розв'язання задачі Коші
\end{enumerate}

\[
\begin{array}{llll} (\text{УЧП}) & u_t = \alpha^2 u_{xx}, & -\infty < x < \infty, & 0 < t < \infty, \\ (\text{HV}) & u(x, 0) = \sin x, & -\infty < x < \infty. \end{array}
\]

\begin{enumerate}
\def\labelenumi{\arabic{enumi}.}
\setcounter{enumi}{2}
\tightlist
\item
  С помощью преобразования Лапласа по переменной t решите вадачу
\end{enumerate}

\[
\begin{array}{llll} (\text{УЧ\Pi}) & u_t = u_{x\,\text{v}}, & 0 < x < \infty, & 0 < t < \infty, \\ (\text{ГУ}) & u\left(0,\ t\right) = \sin t, & 0 < t < \infty, \\ (\text{HY}) & u\left(x,\ 0\right) = 0, & 0 \leqslant x < \infty. \end{array}
\]

Дайте фізичну інтерпретацію цієї зали.

\begin{enumerate}
\def\labelenumi{\arabic{enumi}.}
\setcounter{enumi}{3}
\tightlist
\item
  Розв'язати задачу з граничним значенням
\end{enumerate}

(ОДА)

\[
\frac{d^2U}{dx^2} - sU = A, \quad 0 < x < 1,
(\Gamma Y) \quad
\begin{cases} \frac{dU}{dx} (0) = 0 \\ U(1) = 0. \end{cases}
\]

Лекція 14

\section{ПРИНЦИП
ДЮАМЕЛЯ}\label{ux43fux440ux438ux43dux446ux438ux43f-ux434ux44eux430ux43cux435ux43bux44f}

МЕТА ЛЕКЦІЇ: Показати, як перетворення Лапласа можна використати для
виявлення цікавих особливостей розв'язання диференціальних рівнянь.
Зокрема, за допомогою алгебранічних перетворень образів Лапласа деякого
розв'язку рівняння часткових похідних можливо отримати відомий принцип
Дюамеля. Інтерпретація цього принципу в теорії звичайних диференціальних
рівнянь відома, але ми покажемо, як його використовувати для розв'язання
деяких змішаних задач для UPP.

технології.

Окрім того, що перетворення Лапласа є потужним методом розв'язання
часткових диференціальних рівнянь, воно дозволяє виявити фізичну
сутність отриманих розв'язків. За допомогою трансформації Лапласа ми
проілюструємо дуже важливий і цікавий принцип Дюамеля. Перш ніж перейти
до принципу Дюамеля, розглянемо проблему, яка дуже часто виникає у

\section{Теплопровідність у стрижні, кінці якого підтримуються при
заданій
температурі}\label{ux442ux435ux43fux43bux43eux43fux440ux43eux432ux456ux434ux43dux456ux441ux442ux44c-ux443-ux441ux442ux440ux438ux436ux43dux456-ux43aux456ux43dux446ux456-ux44fux43aux43eux433ux43e-ux43fux456ux434ux442ux440ux438ux43cux443ux44eux442ux44cux441ux44f-ux43fux440ux438-ux437ux430ux434ux430ux43dux456ux439-ux442ux435ux43cux43fux435ux440ux430ux442ux443ux440ux456}

Дуже часто необхідно знайти температуру всередині області, якщо відомо,
що на межі області вона змінюється з часом відповідно до певного закону
(граничні умови, залежні від часу).

\[
u(0,t)=0 \left(
\]

\begin{verbatim}
\begin{array}{c} \\ \\ \\ \\ \\ \\ \\ \\ \\ \\ \\ \\ \\ \\ \\ \\ \\ \\ \\
\end{verbatim}

Рис. 14.1. Часово-залежні граничні умови.

Наприклад, розглянемо термоізольований стержень, правий кінець якого
підтримується у \(f(t)\) (див. рисунок 14.1). Потрібно вирішити цю
проблему

(14.1)

\[
u_{t} = u_{xx}, \quad 0 < x < 1, \quad 0 < t < \infty,
(\Gamma Y) \quad \begin{cases} u(0, t) = 0, \\ u(1, t) = f(t), \end{cases} \quad 0 < t < \infty,
(HY) \quad u(x, 0) = 0, \quad 0 \le x \le 1.
\]

Сподіваємося, що розв'язання цієї задачі (14.1) буде легким, оскільки ми
знаємо розв'язок простішої задачі, де температура в кінці є сталою:

Насправді, якщо обидві задачі ((14.1) і (14.2)) розв'язуються одночасно
за допомогою перетворення Лапласа, то отримаємо визначний результат:
розв'язок задачі (14.1) виражається через розв'язання задачі (14.2).

Отже, розв'язуючи обидві задачі одночасно, ми отримуємо

\section{Легке завдання (14.2) (Постійні
солдати)}\label{ux43bux435ux433ux43aux435-ux437ux430ux432ux434ux430ux43dux43dux44f-14.2-ux43fux43eux441ux442ux456ux439ux43dux456-ux441ux43eux43bux434ux430ux442ux438}

Застосування перетворення -Лапласа до задачі (14.2)

\[
\frac{d^2W}{dx^2} - sW(x) = 0,
W(0) = 0,
W(1) = 1/s
\]

.

Пемеріль ODU

\[
W(x, s) = \frac{1}{s} \left[ \frac{\sinh(x \sqrt{s})}{\sinh(\sqrt{s})} \right].
\]

Пошук оберненого перетворення

\[
w(x, t) =
= x + \frac{2}{\pi} \sum_{n=1}^{\infty} \frac{(-1)^n}{n} e^{-(n\pi)^2 t} \times
\]

\(\times \sin(n\pi x)\) (розв'язання задачі з постійним ГІ).

\section{Важке завдання (14.1) (Солдати, що залежать від
часу)}\label{ux432ux430ux436ux43aux435-ux437ux430ux432ux434ux430ux43dux43dux44f-14.1-ux441ux43eux43bux434ux430ux442ux438-ux449ux43e-ux437ux430ux43bux435ux436ux430ux442ux44c-ux432ux456ux434-ux447ux430ux441ux443}

Застосування перетворення Лапласа до задачі (14.1)

\[
\frac{d^2U}{dx^2} - sU(x) = 0,
U\left( \mathbf{0}\right) =0,
U(1) = F(s)
\]

.

\[
U(x, s) = F(s) \left[ \frac{\sinh(x \sqrt{s})}{\sinh(\sqrt{s})} \right].
\]

Розділіть і помножте на z

\[
U(x, s) = F(s) \left\{ s \left[ \frac{\sinh(x \sqrt{s})}{s \sinh(s)} \right] \right\}.
\]

Використовуючи співвідношення \(\mathscr{L}[w_t] = sW - w(x, 0),\)
отримуємо

\[
U(x, s) = F(s) \mathcal{L}[w_t].
\]

Відповідно,

\[
u(x, t) = \mathcal{L}^{-1} \{ F(s) \mathcal{L}[w_t] \} =
= \mathcal{L}^{-1} [ F(s) ] * \mathcal{L}^{-1}[w_t] =
= f(t) * w_t(t) =
= \int_0^t w_t(x, t-\tau) f(\tau) d\tau =
\]

(після інтегрування частинами)

\[
= \int_{0}^{t} w(x, t-\tau) f'(\tau) d\tau + f(0) w(x, t).
\]

(Розв'язок задачі з часово-залежними ГІ виражається через розв'язок
задачі з постійними ГІ

розв'язано у задачах з константою

Отже, ми висловили розв'язок задач із граничними умовами, що змінюються
в часі, через розв'язання задачі з постійними граничними умовами.
Відповідні формули такі:

(14.3)

плутанина у наступних формулах. Записати диференціальне рівняння в
похідних у нових координатах (~, т), мы воспользуемся следующими
формулами перехода;

Наприклад, якщо dy/dx = B/2A = 3, то y - 3x = c, а отже
\(\xi(x, y) = y - 3x\) . Ця функція задовольняє умову
\([\xi_x/\xi_y] = -3\) і, відповідно, повертає коефіцієнт
\(\overline{A}\) до нуля.

Отже, половина роботи виконана: ми знайшли нову координатну
\(\xi = \xi(x, y)\) , яка забезпечила рівність \(\overline{A} = 0\) .
Залишається знайти \(\eta(x, y)\) таким чином, щоб перетворити
коефіцієнт \(\overline{B}\) на нуль.

Час зупинитися тут. Виявляється, що якщо вибрати \(B^2-4AC=0\) і \(\xi\)
так, щоб \(\overline{A}=0\) , то коефіцієнт \(\overline{B}\) автоматично
стане нулем 1). Тепер ми побачимо це.

Коефіцієнт \(\overline{B}\) задається виразом

\[
\overline{B}^2 - 4\overline{A}\overline{C} = (B^2 - 4AC) (\xi_x \eta_y - \xi_y \eta_x)^2 - \Pi \rho u u. \ \rho e \partial.
\overline{B} = 2 \left( V \overline{A} \xi_x + V \overline{C} \xi_y \right) \left( V \overline{A} \eta_x + V \overline{C} \eta_y \right).
\]

Функциональная диаграмма на рис. 15.5 делает эти формулы совсем
очевидными. Эта диаграмма особенно полезна при вычислении частных
производных функций и и и по переменным Р И С

\pandocbounded{\includegraphics[keepaspectratio]{images/_page_115_Picture_14.jpeg}}

15.5 Диаграмма, иллюстрирующая функциональную зависимость переменных

\begin{verbatim}
\xi   и   \tau  , поскольку   \xi   зависит и от x, и от t. Переменная   \tau   зависит только от t. Это все, что касается преобразований. Теперь подставим найденные выражения   u_t  ,   u_x   и   u_{vx}   в уравнение (15 2) и получим
\end{verbatim}

\[
\overline{B} = 2 \sqrt{\overline{A}} \left[ \sqrt{\overline{A}} \, \eta_x + \sqrt{\overline{C}} \, \eta_y \right].
\]

откуда

\[
u_{xy} + 2u_{xy} + u_{yy} = 0.
\]

.

\begin{verbatim}
Следовательно, новая задача с начальным условием в персменных   \xi   и   \tau   имеет вид
\end{verbatim}

(УЧП)

\[
\frac{dy}{dx} = -[\xi_x/\xi_y] = B/2A = 1.
\]

(НУ) u(\xi, 0) = 1 - H(\xi), \quad -\infty \textless{} \xi \textless{}
\infty.

\begin{verbatim}
Так как   \xi = x   при t = 0, то новые граничные условия совпадают со старыми. Эта задача была решена в лекции 12 с помощью интегрального преобразования Фурье, и ее решение записывается в виде
\end{verbatim}

\[
y = x + c
\]

где \(\phi\)(\(\beta\)) --- начальное условие. Следовательно, в нашем
случае

\[
\eta = y
\]

После подстановки

\begin{verbatim}
\begin{cases}
 \xi = y - x, \\
 \eta = y
 \end{cases}
\end{verbatim}

получаем окончательный результат

(15.3)

\[
\overline{C} = A \eta_x^2 + B \eta_x \eta_y + C \eta_y^2 = 1,
\overline{E} = A \eta_{xx} + B \eta_{xy} + C \eta_{yy} + D \eta_x + E \eta_y = 0,
\begin{cases} \frac{1}{2} \left[ 1 + \operatorname{erf} \left( \frac{-\xi}{2\sqrt{D\tau}} \right) \right], & \xi < 0, \\ \frac{1}{2} \operatorname{erfc} \left( \frac{\xi}{2\sqrt{D\tau}} \right), & 0 \leq \xi. \end{cases}
\overline{G} = G = 0.
\overline{A}u_{\xi\xi} + \overline{B}u_{\xi\eta} + \overline{C}u_{\eta\eta} + \overline{D}u_{\xi} + \overline{E}u_{\eta} + \overline{F}u = \overline{G}
\]

Графики этого решения для различных моментов времени изображены на рис.
15.6. Осталось сделать последний шаг --- записать решение нашей задачи в
координатах x и t. Подставляя

\[
u_{\eta} = f(\xi),
,   \tau = t
\]

в формулы (15.3), получаем

(15.4)

\[
u(\xi, \eta) = \eta f(\xi) + g(\xi),
u(x, y) = yf(y-x) + g(y-x)
\begin{cases} \frac{1}{2} \left[ 1 + \operatorname{erf}\left(\frac{Vt - x}{2\sqrt{Dt}}\right) \right], & Vt > x, \\ \frac{1}{2} \operatorname{erfc}\left(\frac{x - Vt}{2\sqrt{Dt}}\right), & Vt \leq x. \end{cases}
u_{xx} + 2u_{xy} + u_{yy} = 0.
u(x, y) = y \sin(y-x) + (y-x)^2
Au_{xx} + Bu_{xy} + Cu_{yy} + Du_x + Eu_y + Fu = G
\begin{array}{ll} u_t = u_{xx} - 2u_x, & -\infty < x < \infty, & 0 < t < \infty, \\ u(x, 0) = \sin x, & -\infty < x < \infty. \end{array}
u_{\xi\xi} + u_{\eta\eta} = \varphi(\xi, \eta, u, u_{\xi}, u_{\eta}).
\overline{A}u_{\xi\xi} + \overline{B}u_{\xi\eta} + \overline{C}u_{\eta\eta} + \overline{D}u_{\xi} + \overline{E}u_{\eta} + \overline{F}u = \overline{G}
u_{\xi\eta} = \psi(\xi, \eta, u, u_{\xi}, u_{\eta}).
\begin{array}{ll} u_t = u_{xx} - 2u_x, & -\infty < x < \infty, \ 0 < t < \infty, \\ u(x, \ 0) = e^x \sin x, & -\infty < x < \infty. \end{array}
\frac{dy}{dx} = \frac{B - \sqrt{B^2 - 4AC}}{2A}
\xi(x, y) = \text{const}, \quad \eta(x, y) = \text{const}.
\frac{dy}{dx} = -\sqrt{-4x^2} = -2ix,
\begin{array}{ll} (\text{YHI}) & u_t = -2u_x, \\ (\text{HY}) & u_t(x, 0) = e^{-x^2}. \end{array}
\]

також задаються комплексно-спряженим координатам

\# Конверсія 2

Наступне перетворення з \((\xi,\eta)\) на \((\alpha,\beta)\) виконується
за формулами

Проверьте ваш ответ.

\begin{enumerate}
\def\labelenumi{\arabic{enumi}.}
\setcounter{enumi}{3}
\tightlist
\item
  Решите задачу
\end{enumerate}

, \(\beta = \frac{\xi - \eta}{2i}\) .

Внаслідок застосування другого перетворення, рівняння

набуває остаточної форми

Проверьте ваше решение. Как выглядит решение в различные моменты
времени?

УКАЗАНИЕ. Переход к движущейся системе отсчета позволяет сразу отбросить
слагаемое \(Vu_x\) в уравнении конвективной диффузии. После решения
новой задачи

де \(\phi\) та \(\phi\) використовуються для позначення правих сторін
рівнянь у загальному вигляді. Замість того, щоб доводити, що ці дві
трансформації дійсно ведуть до бажаного результату, давайте розглянемо
простий приклад.

\section{Зведення рівняння y\^{}2u\_\{xx\} + x\^{}2u\_\{yy\} = 0 до його
канонічної
форми}\label{ux437ux432ux435ux434ux435ux43dux43dux44f-ux440ux456ux432ux43dux44fux43dux43dux44f-y2u_xx-x2u_yy-0-ux434ux43e-ux439ux43eux433ux43e-ux43aux430ux43dux43eux43dux456ux447ux43dux43eux457-ux444ux43eux440ux43cux438}

Розглянемо рівняння

у яких \(A=y^2\) , B=0, \(C=x^2\) , D=E=F=G=0. Дискримінантний
\(B^2-4AC\) дорівнює \(-4x^2y^2\) , і ми зведемо це рівняння до
канонічної форми в першому квадранті x\textgreater0, y\textgreater0.

КРОК 1 (Перше перетворення)

Запишемо характеристичні рівняння

\[
\frac{dy}{dx} = \frac{B + \sqrt{B^2 - 4AC}}{2A} = \frac{\sqrt{-4x^2y^2}}{2y^2} = i\frac{x}{y}
\]

сделайте замену переменных \(\xi = x - Vt\) , \(\tau = t\) . В нашем
случае можно в явном виде вычислить интеграл

\[
y^2 + ix^2 = \text{const},
\]

Такой интеграл уже встречался в лекции 12. Может быть, читателю будет
удобнее обратиться к таблицам преобразования Фурье.

\section{ОДНОМЕРНОЕ ВОЛНОВОЕ УРАВНЕНИЕ (ГИПЕРБОЛИЧЕСКИЕ
УРАВНЕНИЯ)}\label{ux43eux434ux43dux43eux43cux435ux440ux43dux43eux435-ux432ux43eux43bux43dux43eux432ux43eux435-ux443ux440ux430ux432ux43dux435ux43dux438ux435-ux433ux438ux43fux435ux440ux431ux43eux43bux438ux447ux435ux441ux43aux438ux435-ux443ux440ux430ux432ux43dux435ux43dux438ux44f}

ЦЕЛЬ ЛЕКЦИИ: Показать, что колебания сгруны описываются уравнением

Від

и что это уравнение --- следствие законов Ньютона. Рассматриваются и
некоторые другие виды волнового уравнения:

\eta(x, y) = y\^{}2 - ix\^{}2. (Неважливо, яка з двох функцій
називається \xi , а яка --- \eta ; ви можете поміняти їх місцями, якщо
хочете.) Після цього перетворення початкове рівняння набуває вигляду

Ми не звертатимемо уваги на це рівняння (яке є комплексним гіперболічним
рівнянням), а одразу перейдемо до другого перетворення.

КРОК 2. (Друга трансформація)

При здійсненні другої трансформації ми маємо

\[
\begin{split} u_{tt} &= \alpha^2 u_{xx} + F\left(x, \ t\right), \\ u_{tt} &= \alpha^2 u_{xx} - \beta u_t - \lambda u + F\left(x, \ t\right). \end{split}
\]

(реальна частина \(\xi\) і \(\eta\) ),
\(\beta = \frac{\xi - \eta}{2\iota} = x^2\) (уявна частина \(\xi\) і
\(\eta\) ).

Щоб спростити нотацію, найкраще перепозначити \((\alpha,\beta)\)
\((\xi,\eta)\) окремо і записати отримане \(\epsilon\) перетворення як

КРОК 3 (Отримання нового рівняння)

\begin{verbatim}
Нову канонічну форму можна отримати, обчисливши всі коефіцієнти $\overline{A}$, $\overline{B}$, $\overline{C}$, $\overline{D}$, $\overline{E}$, $\overline{F}$ і $\overline{G}$ у рівнянні
\end{verbatim}

До сих пор мы занкмались физическими процессами, которые описываются
одномерными параболическими уравнениями (диффузионные задачи). Теперь мы
приступим к изучению еще одного класса уравнений с частными
производными---гиперболических уравнений. Начнем с одномерного волнового
уравнения, которое описывает, в частности, поперечные колебания струны.

\section{Колебания
струны}\label{ux43aux43eux43bux435ux431ux430ux43dux438ux44f-ux441ux442ux440ux443ux43dux44b}

Рассмотрим малые колебания струны с закрепленными концами. Предположим,
что струна туго натянута, сделана из однородного материала и колеблется
в одной плоскости (рис. 16.1).

\pandocbounded{\includegraphics[keepaspectratio]{images/_page_119_Figure_10.jpeg}}

РИС. 16.1. Поперечные колебания струны: a --- смещение струны от
положения равновесия.

Действием сил тяжести на струну будем пренебрегать. Для построения
математической модели колебаний струны рассмотрим все силы, действующие
на небольшой участок струны (рис. 16.2). Оказывается, что волновое
уравнение---это не что иное, как закон движения Ньютона (изменение
импульса равно сумме действующих сил), примененный к струне. Глядя на
рис. 16.1, мы можем представить те силы, которые действуют на струну в
направлений, перпендикулярном оси х.

\pandocbounded{\includegraphics[keepaspectratio]{images/_page_120_Figure_2.jpeg}}

РИС. 16.2. Малый участок \((x, x + \Delta x)\) колеблющейся струны.

\begin{enumerate}
\def\labelenumi{\arabic{enumi}.}
\tightlist
\item
  Сила, обусловленная натяжением струны ( \(\alpha^2 u_{xx}\) ).
  Величина поперечной компоненты силы натяжения находится по формуле
\end{enumerate}

Сила =

за формулами (41.5) з \xi = y\^{}2 та \eta = x\^{}2 . Зробивши це,
отримуємо канонічну форму еліптичного рівняння

\(\simeq T [u_x(x + \Delta x, t) - u_x(x, t)].\) 2. Внешняя сила
\(F (x, t)\).

Внешняя сила \(F(x, t)\), приложенная к струне, может произвольным
образом зависеть от x и t. Некоторые примеры внешних сил приводятся
ниже:

а) гравитационная сила \(F(x, t)\) = -mg;

б) импульсы, распределенные вдоль струны и действующие на нее в
различные моменты времени:

в) силы, возникающие при действии звуковой волны на кожу барабана (позже
мы подробно изучим колебания мембраны, которые описываются волновым
уравнением).

\begin{enumerate}
\def\labelenumi{\arabic{enumi}.}
\setcounter{enumi}{2}
\item
  Силы трения, действующие на струну (--- \(\beta u_t\) ). Если струна
  колеблется в среде, то возникает сила сопротивления, которая
  пропорциональна скорости \(u_t\) .
\item
  Возвращающая сила (- үи).
\end{enumerate}

Эта сила направлена противоположно смещению струны. Если смещение u
положительно, то сила отрицательна.

Если применить теперь уравнения движения Ньютона к небольшому участку
струны, то получим

\section{ПРИМІТКИ}\label{ux43fux440ux438ux43cux456ux442ux43aux438-2}

\begin{enumerate}
\def\labelenumi{\arabic{enumi}.}
\tightlist
\item
  Аналогічно можна класифікувати лінійні рівняння другого порядку з
  \emph{трьох або більше змінних}, але для цього потрібен матричний
  аналіз. Наприклад, рівняння другого порядку з трьома змінними
\end{enumerate}

где \(\rho\) --- линейная плотность струны. Разделим обе части этого
сротношения на \(\Delta x\) и устремим \(\Delta x\) --- 0. В результате
получим хорошо известное телеграфное уравнение

(16.1)

згідно з цією класифікацією, буде параболічним, а рівняння

В этом уравнении следовало бы разделить \(\alpha\) , \(\beta\) ,
\(\gamma\) и \(F(x,t)\) на \(\rho\) , но мы ради простоты сохраним для
этих величин прежние обозначения. Искомое уравнение получено. Приведем
интерпретацию простейшего волнового уравнения

\begin{itemize}
\tightlist
\item
  Гіперболічно.
\end{itemize}

\begin{enumerate}
\def\labelenumi{\arabic{enumi}.}
\setcounter{enumi}{1}
\tightlist
\item
  Наш інтерес до класифікації диференціальних рівнянь у похідних
  пов'язаний із тим, що три основні типи рівнянь описують три різні типи
  фізичних явищ, і ми хотіли б отримати математичну класифікацію цих
  трьох класів фізичних явищ.
\end{enumerate}

\section{ЗАВДАННЯ}\label{ux437ux430ux432ux434ux430ux43dux43dux44f-10}

\begin{itemize}
\item
  \begin{enumerate}
  \def\labelenumi{\arabic{enumi}.}
  \tightlist
  \item
    Які з цих параболічних і еліптичних рівнянь записуються у канонічній
    формі:
  \end{enumerate}

  \begin{itemize}
  \tightlist
  \item
    \begin{enumerate}
    \def\labelenumi{\alph{enumi})}
    \tightlist
    \item
      u\_t = u\_\{xx\} hu ,

      \begin{itemize}
      \tightlist
      \item
        \begin{enumerate}
        \def\labelenumii{\arabic{enumii})}
        \setcounter{enumii}{5}
        \tightlist
        \item
          u\_\{xy\} + u\_\{xx\} + 3u = \sin x ,
        \end{enumerate}
      \item
        \mathbf{B}) ~\hat{u_{xx}} + 2\hat{u_{yy}} = 0, - \mathbf{r})
        ~u\_\{xx\} = \sin\^{}2 x ?
      \end{itemize}
    \end{enumerate}
  \end{itemize}
\item
  \begin{enumerate}
  \def\labelenumi{\arabic{enumi}.}
  \setcounter{enumi}{1}
  \tightlist
  \item
    Перетворимо параболічне рівняння \(u_{xx}\) + 2u\_\{xy\} +
    \(u_{yy}\) + u = 2 у канонічну форму.
  \end{enumerate}
\item
  \begin{enumerate}
  \def\labelenumi{\arabic{enumi}.}
  \setcounter{enumi}{2}
  \tightlist
  \item
    Перетворимо еліптичне рівняння
    \(u_{xx}\)+2u\_\{yy\}+x\textsuperscript{2u\_x=\$e}\{-x\^{}2/2\}\$ у
    форму каноха.
  \end{enumerate}
\end{itemize}

\section{МЕТОД МОНТЕ-КАРЛО
(ВСТУП)}\label{ux43cux435ux442ux43eux434-ux43cux43eux43dux442ux435-ux43aux430ux440ux43bux43e-ux432ux441ux442ux443ux43f}

МЕТА ЛЕКЦІЇ: Пояснити основні ідени, що лежать в основі методу
Монте-Карло, і показати, як цей метод можна використовувати для
розв'язання різних задач. Виявляється, що за допомогою азартних ігор
можна змоделювати (загалом, на комп'ютері) наближені розв'язки реальних
задач. Як простий приклад можна вказати обчислення інтеграла

основанную на интуитивных соображениях.

\section{Интуитивная интерпретация волнового
уравнения}\label{ux438ux43dux442ux443ux438ux442ux438ux432ux43dux430ux44f-ux438ux43dux442ux435ux440ux43fux440ux435ux442ux430ux446ux438ux44f-ux432ux43eux43bux43dux43eux432ux43eux433ux43e-ux443ux440ux430ux432ux43dux435ux43dux438ux44f}

У читателя может возникнуть вопрос: почему уравнение вида (16.2) должно
описывать нечто такое, что похоже на колебания скрипичной струны? Чтобы
разобраться в этом, необходимо понять,

\pandocbounded{\includegraphics[keepaspectratio]{images/_page_121_Picture_9.jpeg}}

РИС. 16 3. Интерпретация волнового уравнения
\(u_{tt} = \alpha^2 u_{xx}\)

что величина \(u_{tt}\) является вертикальным ускорением струны в точке
x. Следовательно, уравнение (16.2) можно интерпретировать следующим
образом: ускорение струны, обусловленное натяжением, в каждой точке тем
больше, чем больше вогнутость струны \(u_{xx}\) в данной точке
(коэффициентом пропорциональности служит константа \(\alpha^2 = T/\rho\)
) (см. рис. 133).

\section{ЗАМЕЧАНИЯ}\label{ux437ux430ux43cux435ux447ux430ux43dux438ux44f-1}

\begin{enumerate}
\def\labelenumi{\arabic{enumi}.}
\tightlist
\item
  Волновое уравнение \(u_{tt} = \alpha^2 u_{xx}\) описывает также
  продольные и крутильные колебания стержня. В продольных колебаниях
  смещения параллельны стержню, а через \(u\left(x,\,t\right)\)
  обозначается величина продольного смещения относительно равновесного
  положения. Такие колебания возникают, например, при ударе молотком по
  торцу стержня. Аналогично описываются крутиль-
\end{enumerate}

ные колебания

кидаючи дротик у одиничний квадрат \{(x, y): 0 \textless{} x \textless{}
1, 0 \textless{} y \textless{} 1\} Після приблизно 100 кидків відносна
частина дротиків, що опускається нижче кривої y = x\^{}2 , вважається
приблизним значенням інтеграла. Узагальнення такої гри (з випадковим
киданням дрогіка) завжди містить процедуру генерації послідовності
випадкових чисел.

У цій лекції ми коротко опишемо метод Монте-Карло, а в наступній лекції
покажемо, як його використовувати для розв'язання диференціальних
рівнянь у похідних

Методи Монте-Карло зазвичай називають групою методів розв'язання
детермінованих задач (тобто задач без випадковості), у яких суттєво
використовуються елементи випадковості. Загальна філософія етнічних
методів показана на рис. 42 1.

Ймовірнісна гра

\begin{verbatim}
Деторміноване завдання \rho — розв'язання задачі (наприклад, обчислення інтеграла, розв'язання UBP тощо).

Результат = \hat{p} Наближення p Рішення = p Рис. 42.1. Загальна теорія методу Монте-Карло.
\end{verbatim}

\section{Обчислення
інтеграла}\label{ux43eux431ux447ux438ux441ux43bux435ux43dux43dux44f-ux456ux43dux442ux435ux433ux440ux430ux43bux430}

Припустимо, ми хочемо обчислити інтеграл

Здесь \emph{k} --- модуль Юнга, характеризующий упругость материала.
Материалы, обладающие большим модулем Юнга, колеблются на более высоких
частотах. Звуковые волны представляют собой продольные волны.

\begin{enumerate}
\def\labelenumi{\arabic{enumi}.}
\setcounter{enumi}{1}
\tightlist
\item
  Если линейная плотность струны \(\rho(x)\) зависит от координаты, то
  волновое урагнение записывается в виде
\end{enumerate}

(детерміноване завдання). Щоб використати метод Монте-Карло, потрібно
винайти гру на удачу, яка дає приблизне значення цього інтеграла.

\begin{verbatim}
Звісно, можна придумати багато різних ігор. Остаточна версія залежить від точності розрахунків, простоти гри тощо. Очевидний варіант гри для обчислення інтегралу — кинути дротик у прямокутний R = \{(x,y): a \le x \le b, 0 \le y \le \max f(x)\} (див. рис. 42.2)
\end{verbatim}

\pandocbounded{\includegraphics[keepaspectratio]{images/_page_315_Figure_7.jpeg}}

Рис. 42 2. Обчислення інтегралу методом Монте-Карло.

Очевидно, якщо випадковим чином кинути близько 100 дротиків у
прямокутник R, що містить графік функції, можна отримати оцінку значення
інтеграла, помноживши площу прямокутника R на відносну кількість
дротиків під графіком функції. Отже, результат нашої гри

т. е. получается УЧП с переменными коэффициентами.

\begin{enumerate}
\def\labelenumi{\arabic{enumi}.}
\setcounter{enumi}{2}
\tightlist
\item
  Поскольку волновое уравнение \(u_{tt} = \alpha^2 u_{xx}\) содержит
  производную по времени второго порядка, для получения единственного
  решения при t\textgreater0 необходимо задать \(\partial sa\) начальных
  условия:
\end{enumerate}

--- це інтегральний бал 1.

Щоб виконати ці обчислення на комп'ютері, нам потрібно якимось чином
отримати послідовність випадкових точок (подивимось, як це буде зроблено
пізніше), тобто наказати комп'ютеру «кидати дротики». Блок-схема на рис.
42.3 показує, як комп'ютер вирішить цю проблему.

\pandocbounded{\includegraphics[keepaspectratio]{images/_page_316_Figure_2.jpeg}}

\begin{verbatim}
**РИС.** 42.3. Блок-схема розрахунку інтегрального \int_{a}^{b} f(x) dx методом Монте-Карло (100 кидків); M = \max_{a} f(x) .
\end{verbatim}

\section{Випадкові
числа}\label{ux432ux438ux43fux430ux434ux43aux43eux432ux456-ux447ux438ux441ux43bux430}

Перш ніж перейти до рівняння Монте-Карло з частковими похідними,
зупинимося на важливому питанні випадкових чисел. Щойно, обчислюючи
інтеграл, нам потрібно було побудувати послідовність випадкових точок
\(P_i = (x_i, y_i)\)всередині прямокутника R. Іншими словами,
x-координата точки P має бути випадковим числом з інтервалу {[}a, b{]},
а \(y_i\) від інтервалу {[}0, M{]}. Щоб знайти випадкові числа з цих
інтервалів, звертаємося до послідовності випадкових чисел, рівномірно
розподілених по сегменту {[}0, 1{]}. Тоді числа, рівномірно розподілені
на сегменті {[}a, b{]}, можна обчислити за формулою

(начальное смещение струны), \(u_t(x, 0) = g(x)\) (начальная скорость
струны).

Этим волновое уравнение отличается от уравнения теплопроводности, для
которого необходимо задать только одно начальное условие.

\begin{enumerate}
\def\labelenumi{\arabic{enumi}.}
\setcounter{enumi}{3}
\tightlist
\item
  С помощью волнового уравнения можно описать распределение
  электрического тока в проводе. Из законов Кирхгофа получается
  следующая система двух уравнений с частными производными первого
  порядка:
\end{enumerate}

(16.3)

Тепер, звичайно, виникає питання: як отримати послідовність випадкових
чисел, \{\(r_i\): i=1,2,\ldots\} рівномірно розподілених на {[}0,1{]}?
Найпоширенішим методом сьогодні є метод дедукцій (метод порівнянь).
Використовуючи цей метод, отримуємо послідовність випадкових цілих чисел
(наприклад, 2120, 1401, 177, 3013,\ldots), потім ліворуч від кожного
числа розміщується десяткова крапка, так що всі числа знаходяться між
нулем і одним (0,212; 0,1401; 0,0177; 0,3013; \ldots) .

Отже, щоб отримати послідовність випадкових цілих чисел (між 0 і P), ми
використаємо метод віднімання.

\section{Отримання випадкових чисел за допомогою методу
відрахування.}\label{ux43eux442ux440ux438ux43cux430ux43dux43dux44f-ux432ux438ux43fux430ux434ux43aux43eux432ux438ux445-ux447ux438ux441ux435ux43b-ux437ux430-ux434ux43eux43fux43eux43cux43eux433ux43eux44e-ux43cux435ux442ux43eux434ux443-ux432ux456ux434ux440ux430ux445ux443ux432ux430ux43dux43dux44f.}

\begin{enumerate}
\def\labelenumi{\arabic{enumi}.}
\item
  Як перше випадкове ціле число, обирайте будь-яке число, що знаходиться
  між 0 і P (число P обирається наперед).
\item
  Помножте це випадкове число на певне фіксоване число.
\end{enumerate}

Попередньо обраний множник M.

\begin{enumerate}
\def\labelenumi{\arabic{enumi}.}
\setcounter{enumi}{2}
\item
  Додайте до добутку деяке фіксоване ціле число \emph{K}, обране
  заздалегідь.
\item
  Поділити отриману суму на P і залишок
\end{enumerate}

Виберіть нове випадкове число.

\begin{verbatim}
Метод виведення можна записати у формі формули r_{i+1} \equiv (Mr_i + K) \bmod P , що означає, що для отримання наступного r_{i+1} потрібно взяти попередній r_i , помножити його на M, додати результат на K, поділити суму на P і взяти решту поділу. Наприклад, якщо P = 100, M = 37, K = 16, r_0 = 15 , то отримується наступна послідовність випадкових чисел:
\end{verbatim}

где

ж --- координата вдоль провода,

\emph{t} --- время,

\(i(x, t)\) --- распределение тока вдоль провода,

\(v\left(x,\ t\right)\) --- распределение потенциала вдоль провода,

С -- емкость провода на единицу длины,

G --- утечка на единицу длины,

Р --- сопротивление на единицу длины,

-- индуктивность провода на единицу длины.

Уравнения (16.3) принято называть системой телеграфных уравнений. В
такой форме они обычно не используются. Для преобразования этой системы
продифференцируем первое уравнение по x, второе продифференцируем по t,
умножим обе части на C и вычтем из первого. В результате получаем

\[
Mr_1 + K = 2643
\]

Воспользуемся теперь вторым уравнением из (16.3)

\[
r = 71
\]

окончательно получаем

(16.4)

Y

Это уравнение для тока, известное под названием \emph{телеграфного
уравнения}, является гиперболическим уравнением второго порядка (если
конечно ии C, ни L не равны нулю, в противном случае уравнение
становится параболическим).

Напряжение описывается точно таким же уравнением:

(16.5)

\[
I = \int_{0}^{1} \int_{0}^{1} \int_{0}^{1} \int_{0}^{1} e^{-(x^{2}+y^{2}+z^{2}+\omega^{2})} dx dy dz dw,
\]

Если G = R = 0, то уравнения (16.4)---(16.5) упрощаются:

(16.6)

\[
I = \int_{0}^{1} e^{\sin x} dx.
\]

\section{ЗАДАЧИ}\label{ux437ux430ux434ux430ux447ux438-2}

\begin{itemize}
\tightlist
\item
  \begin{enumerate}
  \def\labelenumi{\arabic{enumi}.}
  \tightlist
  \item
    Получите уравнение (16.5) для \emph{и} из системы двух уравнений
    первого порядка (16.3)
  \end{enumerate}
\item
  \begin{enumerate}
  \def\labelenumi{\arabic{enumi}.}
  \setcounter{enumi}{1}
  \tightlist
  \item
    Зная физический смысл каждого члена волнового уравнения, что вы
    можете сказать о поведении во времени решения следующей задачи:
  \end{enumerate}
\end{itemize}

(YYII)

\[
r_{i+1} \equiv (3r_i + 4) \mod 7, \quad r_0 = 0.
\]

, \(0 < x < 1\) , \(0 < t < \infty\) ,\\
(ГУ)

\[
I = \int_{0}^{1} \int_{0}^{1} \int_{0}^{1} e^{-(x^{2}+y^{2}+z^{2})} dx dy dz.
u(x,y) = g(x,y) =
\begin{cases} u(0, t) = 0, & 0 < t < \infty, \\ u(1, t) = 0, & 0 < t < \infty, \end{cases}
\]

После того как эта частная задача будет решена методом Монте-Карло, мы
обсудим, как решать этим методом более общие задачи.

Для иллюстрации метода Монте-Карло рассмотрим игру, которая называется
«Блуждающий пьяница». Чтобы играть в эту игру, необходим стол с
нанесенной на нем сеткой из линий, как показано на рис. 43.1.

Перейдем к описанию правил игры.

\section{Как играть в «Блуждающего
пьяницу»?}\label{ux43aux430ux43a-ux438ux433ux440ux430ux442ux44c-ux432-ux431ux43bux443ux436ux434ux430ux44eux449ux435ux433ux43e-ux43fux44cux44fux43dux438ux446ux443}

ШАГ 1. Блуждания пьяницы начинаются из произвольной точки сетки (в нашем
случае это точка A).

\pandocbounded{\includegraphics[keepaspectratio]{images/_page_321_Figure_4.jpeg}}

\begin{itemize}
\tightlist
\item
  А начальная точка - - внутренние точки сетки - □ -Конечная точка р;
\item
  \(g_i\) вознаграждение в m . \(p_I\)
\end{itemize}

РИС. 43.1. Игровое поле для игры «Блуждающий пьяница».

ШАГ 2. На каждом ходе игры пьяница случайным образом перебирается в одну
из четырех соседних точек сетки (в нашем случае соседними для A являются
точки B, C, D и E и вероятность попадания в каждую из этих точек равна
0,25).

Таблица 43.1

Вероятность окончания случайного блуждания в граничной точке \(p_i\) и
величина вознаграждения \(g_i\)

{\def\LTcaptype{none} % do not increment counter
\begin{longtable}[]{@{}
  >{\raggedright\arraybackslash}p{(\linewidth - 8\tabcolsep) * \real{0.2523}}
  >{\raggedright\arraybackslash}p{(\linewidth - 8\tabcolsep) * \real{0.4541}}
  >{\raggedright\arraybackslash}p{(\linewidth - 8\tabcolsep) * \real{0.2752}}
  >{\raggedright\arraybackslash}p{(\linewidth - 8\tabcolsep) * \real{0.0092}}
  >{\raggedright\arraybackslash}p{(\linewidth - 8\tabcolsep) * \real{0.0092}}@{}}
\toprule\noalign{}
\begin{minipage}[b]{\linewidth}\raggedright
Граничная точка р і
\end{minipage} & \begin{minipage}[b]{\linewidth}\raggedright
\(P_{\mathcal{A}}\left(p_{i}\right)\) --- относительная доля блужданий,
закончившихся в точке \(p_{i}\)
\end{minipage} & \begin{minipage}[b]{\linewidth}\raggedright
\(g_i\) --- величина вознаграждения при достижении точки \(p_i\)
\end{minipage} & \begin{minipage}[b]{\linewidth}\raggedright
\end{minipage} & \begin{minipage}[b]{\linewidth}\raggedright
\end{minipage} \\
\midrule\noalign{}
\endhead
\bottomrule\noalign{}
\endlastfoot
1234567891011 & 0,040,150,030,060,170,050,060,150,030,060,160,04 &
110000000 & & \\
\end{longtable}
}

ti

ШАГ 3. После перехода в соседнюю точку процесс блуждания возобновляется.
Так пьяница идет от точки к точке до тех пор, пока случайным образом не
окажется в граничной точке \(p_t\) . Здесь он останавливается, а мы
записываем номер точки \(p_t\) . На этом заканчивается одна случайная
прогулка.

ШАГ 4. Повторим щаги 1-3 достаточно много раз и найдем относительное
число прибытий пьяницы в каждую из граничных точек \(p_i\) . В табл.
43.1 приведены типичные результаты, полученные в результате 100
случайных прогулок.

ШАГ Б. Предположим, что пьяница получает вознаграждение
\(g_i = g\left(p_i\right)\) (величина граничного условия в точке \(p_i\)
, если он достиг точки \(p_i\) ) и предположим, что цель
игры---вычислить среднее вознаграждение \(R\left(A\right)\) для всех
случайных прогулок, начинающихся из точки A. Этот средний выигрыш
определяется по формуле

Игра завершается, когда величина \(R\left(A\right)\) найдена. В
соответствии с данными табл. 43.1 для \(R\left(A\right)\) получаем

Теперь объясним, зачем нам понадобилось играть в такую игру.

\section{Для чего играть в «Блуждающего
пьяницу»?}\label{ux434ux43bux44f-ux447ux435ux433ux43e-ux438ux433ux440ux430ux442ux44c-ux432-ux431ux43bux443ux436ux434ux430ux44eux449ux435ux433ux43e-ux43fux44cux44fux43dux438ux446ux443}

Оказывается, что среднее вознаграждение, о котором только что шла речь,
является приближенным решением задачи Дирихле в точке A. Это интересное
наблюдение основано на двух фактах.

\begin{enumerate}
\def\labelenumi{\arabic{enumi}.}
\item
  Предположим, что пъяница начал свою прогулку из точки, лежащей на
  еранице. Каждая такая прогулка немедленно заканчивается в той же
  точке, и пъяница получает вознаграждение \(g_i\) . Таким образом,
  среднее вознаграждение для каждой граничной точки равно \(g_i\) .
\item
  Теперь предположим, что прогулка нач нается из внутренней точки. Тогда
  ясно, что среднее вознаграждечие для точки \(R(A)\) будет средним
  арифметическим от средних вознаграждений для четырех соседних точек
\end{enumerate}

\[
\begin{cases} u(x, 0) = \sin(\pi x), & 0 \le x \le 1, \\ u_t(x, 0) = 0. & 0 \le x \le 1, \end{cases}
\]

Зададимся еще раз вопросом: почему же среднее значение вознаграждения
\(R(A)\) аппроксимирует решение задачи Дирихле в точке A? Мы видели, что
величина \(R(A)\) удовлетворяет двум

уравнениям

(во внутренних точках), \(R(A)\) = \(g_i\) (в граничных точках).

Если \(g_i\) ---это значения функции \(g(x, y)\) из граничного условия в
граничных точках \(p_i\) , то два наших уравнения точно совпадают с
двумя уравнениями, которые были получены при решении задачи Дирихле
методом конечных разностей. То есть величина \(R(A)\) соответствует
величине \(u_{i,f}\) в разностных уравнениях

(i, j)---внутренняя точка, \(u_{ij} = g_{ij},\) \(g_{ij}\) ---значение
решения в граничной точке (i, j).

Следовательно, величина \(R(A)\) действительно аппроксимирует решение
уравнения с частными производными в точке A.

Теперь можно считать, что игра в «Блуждающего пьяницу» состоит из трех
шагов.

\section{Решение уравнения Лапласа методом
Монте-Карло}\label{ux440ux435ux448ux435ux43dux438ux435-ux443ux440ux430ux432ux43dux435ux43dux438ux44f-ux43bux430ux43fux43bux430ux441ux430-ux43cux435ux442ux43eux434ux43eux43c-ux43cux43eux43dux442ux435-ux43aux430ux440ux43bux43e}

\begin{itemize}
\tightlist
\item
  Приведенные ниже правила позволяют получить решение в одной внутренней
  точке квадрата.
\item
  ШАГ 1. Совершим некоторое число случайных прогулок, начинающихся в
  точке А и заканчивающихся в одной из граничных точек. Проследим за
  тем, сколько раз прогулки заканчивались в каждой граничной точке.
\item
  ШАГ 2. После того как все прогулки завершены, для каждой граничной
  точки вычисляем относительное число прогулок, завершившихся в этой
  точке. Обозначим эти величины \(P_A(p_i)\) .

  \begin{itemize}
  \tightlist
  \item
    ШАГ 3. Вычисляем приближенное решение \(u(A)\) по формуле
    \(u(A) = g_1 P_A(p_1) + g_2 P_A(p_2) + \dots + g_N P_A(p_N)\) ,
  \end{itemize}
\end{itemize}

где \(g_i\) --- значение функции g в точке \(p_i\) , а N --- число
граничных точек.

Игру «Блуждающий пьяница» можно модернизировать, чтобы решить более
сложную задачу. Ниже приводится пример такой модернизации.

Решение задачи Дирихле с переменными коэффициентами.

Рассмотрим эллиптическую краевую вадачу в квадрате

Дайте физическую интерпретацию этой задачи.

\begin{enumerate}
\def\labelenumi{\arabic{enumi}.}
\setcounter{enumi}{3}
\tightlist
\item
  Для уравнения \(u_{tt} = u_{xx}\)найдите все решения вида
\end{enumerate}

\(u_{xx} + (\sin x) u_{yy} = 0\) , \(0 < x < \pi\) , \(0 < y < \pi\) ,
\((\Gamma Y)\) \(u(x, y) = g(x, y)\) на границе квадрата.

Чтобы решить эту задачу, заменим \(u_{xx}\) , \(u_{yy}\) и \(\sin x\)
следующим образом:

Будет ли решением сумма двух решений такого вида?

\section{ФОРМУЛА
ДАЛАМБЕРА}\label{ux444ux43eux440ux43cux443ux43bux430-ux434ux430ux43bux430ux43cux431ux435ux440ux430}

ЦЕЛЬ ЛЕКЦИИ: Построить решение задачи Коши для волнового уравнения

и \(_{yy}=u_{l+1,\ j}-2u_{l,\ j}+u_{l-1,\ j},\) (центральные разностные
производные),

и подставим их в уравнение с частными производными. Соответствующая
сетка показана на рис. 43.2. Разрешив получившееся уравнение
относительно \(u_{i,l}\) , получаем

Посмотрим внимательно на это уравнение. Коэффициенты при
\(u_{i+1,\ j}\),~\(u_i\),~j+i,~\(u_l\),~j-1 н \(u_{i-1}\),~j
положительны, и их сумма равна единице. Другими словами, решение
\(u_i\),~j является взвешенным средним решений в четырех соседних
точках. Следовательно, можно Р И С

\pandocbounded{\includegraphics[keepaspectratio]{images/_page_324_Picture_9.jpeg}}

. 43.2. Узлы сетки для случайного маршруга.

модифицировать «Блуждающего пьяницу» так, чтобы вероятности перехода в
соседние точки были не 0,25, а равнялись бы коэффициенту в
соответствующем члене. Другими словами, если пьяница находится в точке
(i, j), то он переходит в точку

Эта задача описывает движение бесконечной струны с заданными начальными
условиями. Она была решена в 1750 году французским математиком
Даламбером. Формула Даламбера

з імовірністю \(\frac{1}{2(1+\sin x_j)}\) у точку \((i, j-1)\), з
імовірністю \(\frac{1}{2(1+\sin x_j)}\) у точку \((i+1, j)\), з
імовірністю \(\frac{\sin x_j}{2(1+\sin x_j)}\) у точку \((i-1, j)\) і з
імовірністю \(\frac{\sin x_j}{2(1+\sin x_j)}\).

В остальном игра не меняется. Модификация для других задач может
оказаться более хитроумной, но основная идея остается той же. Читатель
может заняться разработкой своих игр для решения других задач.
Параболический случай рассмотрен в разделе «Задачи».

\section{ЗАМЕЧАНИЯ}\label{ux437ux430ux43cux435ux447ux430ux43dux438ux44f-2}

\begin{verbatim}
1. Отметим, что если относительные частоты попадания  P_A(p_l)  известны, то можно легко найти величину  u\left(A\right)  при других граничных условиях  g_l . Для этого достаточно подставить  g_l  в формулу
\end{verbatim}

позволяет легко найти решение, если известны начальные условия. Кроме
того, она позволяет дать интересную физическую интерпретацию решения на
языке бегущих волн.

Как, наверное, помнит читатель, в параболническом случае мы сначала
решали диффузионные задачи на конечном отрезке (методом разделения
переменных), а уже затем перешли к неограниченным промежуткам (---
\(\infty < x < \infty\) ), где применяли преобразование Фурье. В
гиперболическом случае мы поступим наоборот. Начнем с решения
одномерного волнового уравнения на всей действительной оси, т. е. с
задачи Коши

(17.1)

а не разыгрывать вновь случайные прогулки.

\begin{enumerate}
\def\labelenumi{\arabic{enumi}.}
\setcounter{enumi}{1}
\tightlist
\item
  Во многих случаях исследователя интересует решение уравнения с
  частными производными в одной точке. Если граница области сложна и
  уравнение зависит от трех или четырех независимых переменных, метод
  Монте-Карло может оказаться удобным. Методы Монте-Карло возникли при
  решении очень трудных задач диффузии нейтронов. Эти задачи невозможно
  решать аналитически.
\end{enumerate}

\section{ЗАДАЧИ}\label{ux437ux430ux434ux430ux447ux438-3}

\begin{enumerate}
\def\labelenumi{\arabic{enumi}.}
\tightlist
\item
  Построить блок-схему игры «Блуждающий пьяница» для решения задачи
\end{enumerate}

(УЧП)

, \(0 < x < 1\) , \(0 < y < 1\) , (ГУ) \(u(x, y) = g(x, y)\) на границе,

\begin{itemize}
\tightlist
\item
  во внутренних узлах сетки, если число узлов сетки произвольно. 2.
  Напишите программу для решения задачи по блок-схеме из задачи 1.
\item
  \begin{enumerate}
  \def\labelenumi{\arabic{enumi}.}
  \setcounter{enumi}{2}
  \tightlist
  \item
    Как модифицировать игру, чтобы решить задачу
  \end{enumerate}
\end{itemize}

(УЧП)

содержащей только начальные условия. Эту задачу можно было бы решить
либо с помощью преобразования Фурье (по переменной t), либо с помощью
преобразования Лапласа (по переменной t). Однако мы воспользуемся другим
методом (методом канонических координат), который, как мы надеемся,
заинтересует читателя. Этот метод базируется на тех же идеях, что и
метод перехода к движущейся системе отсчета, с которым мы познакомились
в лекции 15. Итак, приступим к решению задачи (17.1).

\section{Решение одномерного волнового уравнения. Формула
Даламбера}\label{ux440ux435ux448ux435ux43dux438ux435-ux43eux434ux43dux43eux43cux435ux440ux43dux43eux433ux43e-ux432ux43eux43bux43dux43eux432ux43eux433ux43e-ux443ux440ux430ux432ux43dux435ux43dux438ux44f.-ux444ux43eux440ux43cux443ux43bux430-ux434ux430ux43bux430ux43cux431ux435ux440ux430}

Решение задачи (17.1) разобьем на несколько шагов:

\begin{verbatim}
ШАГ 1 (Замена координат (x, t) новыми каноническими координатами   (\xi, \eta)  ).
\end{verbatim}

Для решения задачи (17.1) воспользуемся тем, что если заменить две
независимые переменные x и t новыми пространственновременными
координатами

, \(0 < x < 1\) , \(0 < y < 1\) , (ГУ) \(u(x, y) = g(x, y)\) на границе?

Как выглядит в этой игре случайная прогулка?

\begin{enumerate}
\def\labelenumi{\arabic{enumi}.}
\setcounter{enumi}{3}
\tightlist
\item
  Можете ли вы так модифицировать игру, чтобы решить задачу
\end{enumerate}

\[
\begin{array}{l} u_{xx} + u_{yy} = 0, \quad 0 < x < 1, \quad 0 < y < 1, \\ u(x, 1) = 0, \quad 0 < x < 1, \quad 0 < y < 1, \\ u(x, 0) = 0, \quad 0 < x < 1, \quad 0 < y < 1, \\ u(0, y) = 1, \quad \frac{\partial u}{\partial x}(1, y) = 0 \end{array}
\]

то уравнение

\begin{verbatim}
\begin{array}{lll} (\text{УЧП}) & u_t = $\alpha$^2 u_{xx} - $\beta$ u_x - $\gamma$ u, \quad 0 < x < 1, \quad 0 < t < \infty \\ (\text{ГУ}) & \left\{ \begin{array}{ll} u\left(0,\ t\right) = f\left(t\right), & 0 < t < \infty, \\ u\left(1,\ t\right) = g\left(t\right), & 0 < t < \infty, \\ u\left(x,\ 0\right) = $\varphi$\left(x\right), & 0 \leqslant x \leqslant 1. \end{array} \end{array}
\end{verbatim}

преобразуется в уравнение

ПРИМІТКА. Замінимо це рівняння на скінченну різницю Кренка--Ніколсона,
виразимо \(u_{i+1,j}\) через п'ять значень \(u_{i+1,j-1}\) ,
\(u_{i+1,j+1}\) , \(u_{i,j+1}\) , \(u_{i,j+1}\) , \(u_{i,j+1}\) у
суміжних точках.

\section{КАЛЬКУЛЮС ВАРІАЦІЙ (РІВНЯННЯ
ЕЙЛЕРА--ЛАГРАНЖА)}\label{ux43aux430ux43bux44cux43aux443ux43bux44eux441-ux432ux430ux440ux456ux430ux446ux456ux439-ux440ux456ux432ux43dux44fux43dux43dux44f-ux435ux439ux43bux435ux440ux430ux43bux430ux433ux440ux430ux43dux436ux430}

МЕТА ЛЕКЦІЇ: Ввести поняття функціоналу (функцію з функції) і пояснити,
як функціонали природно виникають у фізиці. Типовий функціонал є
інтегралом

.

В этом легко убедиться, если воспользоваться формулами преобразования
координат:

(17.2)

що є функцією функції y (припускається, що інтегративна функція
F\left(x,,y,,y'\right) задана). Приклад функціональності:

Подставляя эти выражения для производных в волновое уравнение, получим

Ми покажемо, як знайти функцію \$\(\overline{y}\)(x)\$ , яка мінімізує
функціональність J{[}y{]}. Виявляється, що мінімізуюча функція має
\(\overline{y}\) задовольняти так зване рівняння Ейлера-Лагранжа. Це
рівняння виконує ту ж роль, що й необхідна умова мінімуму

.. Первый шаг сделан.

ШАГ 2 (Решение преобразованного уравнения).

Преобразованное уравнение легко решить двумя последовательными
интегрированиями (сначала по переменной ξ, а затем по η). После
интегрирования по ξ получаем

функції \(f(x)\) у точці x у диференціальному численні.

Варіаційне числення тісно пов'язане з диференціальними рівняннями, але,
на жаль, більшість студентів його не вивчає. Ця лекція (і наступна) є
вступом до варіаційного числення та показує, як можна розв'язувати
диференціальні рівняння в похідних на основі варіаційних принципів.

Калькулюс варіацій виник одночасно з математичним аналізом у зв'язку з
розв'язанням задач максимізації та мінімізації функцій функцій (які
називаються функціоналами). Першою задачею варіаційного числення була
задача брахістохрона, сформульована Йоганном Бернуллі у 1696 році.

Бернуллі показав, що час сходження записується у формі

--- произвольная функция от \(\eta\) .

Интегрирование последнего соотношения по \textbf{η} приводит к общему
решению

і, отже, є функцією T\{y\} з функції. Оскільки багато функціоналів
впорядковані таким чином, ми зосередимося на вивченні загальної
функціональності

(44.1)

где \(\phi(\eta)\) --- первообразная функция \(\phi_1(\eta)\) , а
\(\psi(\xi)\) --- произвольная функция. Итак, общее решение уравнения

Тепер сформульуємо головну мету наших лекцій: знайти функцію, яка дає
мінімум (або максимум) функціоналу (44.1). Стратегія пошуку буде такою
ж, як при пошуку мінімуму функції в диференціальному численні. Там ми
знайшли критичні точки функції за умовою f'(x) = 0. У варіаційному
численні все буде трохи складніше, оскільки аргумент більше не є
числовою змінною, а функцією. Однак загальний підхід залишається
незмінним: ми обчислимо так звану функціональну похідну функції \(y(x)\)
і прирівняємо її до нуля. Це нове рівняння буде аналогічним

записывается в виде

(17.3)

з диференціального числення, але тепер це буде звичайне диференціальне
рівняння, відоме як рівняння Ейлера-Лагранжа. Нам просто потрібно знайти
це рівняння і розв'язати з ним деякі проблеми.

Мінімізація функціональності

где \(\phi(\eta)\) и \(\psi(\xi)\) --- произвольные функции своих
аргументов. Папример, читатель может проверить, что функции

Розглянемо задачу пошуку функції \(y(x)\), яка мінімізує
функціональність

\(u(\xi, \eta) = 1/\eta + tg \xi,\) \(u(\xi, \eta) = \eta^2 + e\)

являются решениями уравнения \(u_{\mathrm{ }nh}=0\) . Второй шаг сделан.

\(\text{IIIA}\Gamma\) 3 (Возвращение к старым координатам x и t).

Для нахождения общего решения, т. е. всех решений уравнения
\(u_{tt} = c^2 u_{xx}\) , подставим

у класі гладких функцій, які задовольняють граничні умови

, \(y(b) = B\) (див. рис. 44.2).

Нехай існує \(y = \overline{y}(x)\) потрібна функція, і розглянемо малу
варіацію функції \(\overline{y}\) , тобто функцію
\(\overline{y} + \varepsilon \eta(x)\) , де \(\varepsilon\) --- мале
число, а \(\eta(x)\) --- гладка функція, яка задовольняє

граничні умови \(\eta(a) = \eta(b) = 0\) (див. рис. 44.2). Очевидно, що
якщо обчислити інтеграл J з близької функції \(y + \varepsilon \eta\) ,
то функціонал буде збільшуватися, тобто

в решение

за всі \(\epsilon\) . Іншими словами, графік
\(\phi\left(\epsilon\right)=J\left[\overline{y}+\epsilon\eta\right]\) як
функція, \(\epsilon\) буде виглядати як рис. 44.3.

\pandocbounded{\includegraphics[keepaspectratio]{images/_page_329_Figure_5.jpeg}}

Рис. 44.1. Задача брахістохрана (варіаційне числення починалося з цієї
задачі).

\pandocbounded{\includegraphics[keepaspectratio]{images/_page_329_Figure_7.jpeg}}

Рис. 44.2. Варіація функції: a --- варіація функції
\(\overline{y} = \overline{y}(x) + \epsilon \eta(x)\) ; \(\delta\) є
мінімізуючою функцією \(\overline{y}(x)\) .

З Фіг. 44.3 Очевидно, що в цьому випадку необхідно обчислити похідну від

В результате получаем

(17.4)

Po E, поставити \(\varepsilon = 0\) і прирівняти отриманий вираз до
нуля, тобто

Это общее решение волнового уравнения. С физической точки эрения оно
интересно тем, что представляет сумму двух бегущих волн Р И С

\pandocbounded{\includegraphics[keepaspectratio]{images/_page_126_Figure_14.jpeg}}

. 17.1. Волна \(u(x, t)\) = \(e^{-(x-ct)^2}\) движется слева направо

произвольной формы, движущихся в прогивоположных направлениях со
скоростью c.~Например, функции

\[
= \int_{0}^{d} \left[\frac{\partial F}{\partial \overline{y}} \eta\left(x\right) + \frac{\partial F}{\partial \overline{y'}} \eta'\left(x\right)\right] dx = 0.
\]

(волна, Лвижущаяся слева направо),

\[
\frac{d\varphi(e)}{de} \equiv \int_{a}^{b} \left\{ \frac{\partial F}{\partial \overline{y}} - \frac{\partial}{\partial x} \left[ \frac{\partial F}{\partial \overline{y'}} \right] \right\} \, \eta(x) \, dx = 0.
\]

(волна, движущаяся справа налево),

\[
\frac{\partial F}{\partial \overline{y}} - \frac{d}{dx} \left[ \frac{\partial F}{\partial \overline{y'}} \right] = 0
\]

(две волны, движущиеся в противоположных направлениях)

являются типичными решениями волиового уравнения. На рисунке 17.1
изображена простейшая бегущая волна.

ШАГ 4 (Учет начальных условий).

Напомним читателю общий метод решения задачи Коши в теории обыкновенных
дифференциальных уравнений. Сначала находят общее решение уравнения, а
затем оно подставляется в начальные условия для того, чтобы найти
конкретные значения произвольных постоянных. У нас ситуация аналогична.
Для решения задачи Коши для волнового уравнения мы подставим общее
решение

\[
\frac{\partial F}{\partial y} - \frac{d}{dx} \left[ \frac{\partial F}{\partial y'} \right] = 0
\]

(содержащее две произвольные функции) в начальные условия

\[
J[y] = \int_{0}^{1} [y^{2} + y'^{2}] dx.
\]

\(u_1\)(x, 0) = \(g(x)\), чтобы найти конкретные выражения для
произвольных функций ф и ф. После подстановки получаем

(17.5)

\[
F_y - \frac{d}{dx} F_{y'} = 0,
\]

Если теперь проинтегрировать второе уравнение (17.5), то получится
алгебранческая система двух уравнений относительно неизвестных функций
\varphi(x) и \psi(x) . В самом деле, после интегрирования в пределах от
\(x_0\) до x получаем

(17.6)

\[
\begin{aligned} F_y &= 2y, \ F_{y'} &= 2y'. \end{aligned}
\]

Решаем теперь уравнение (17.6) совместно с первым уравнением (17.5) и
получаем следующие выражения для функций \(\varphi\)(x) и \psi(x) :

(17.7)

\[
2y - \frac{d}{dx}(2y') = 0,
y'' - y = 0.
\]

Следовательно, общее решение задачи (17.1) дается формулой

(17.8)

\[
\overline{y}(x) = 0.42e^x - 0.42e^{-x}
\]

Цель достигнута. Мы получили общее решение вадачи (17.1). Решение вида
(17.8) принято называть формулой Даламбера. Читателю предлагается
самостоятельно проследить за расстановкой пределов интегрирования, чтобы
убедиться в том, что они меняются от x-ct до x+ct. Задача полностью
решена.

Прежде чем завершить лекцию, приведем несколько примеров, показывающих,
как применять формулу Даламбера в конкретных валачах.

\section{Примеры применения формулы
Даламбера}\label{ux43fux440ux438ux43cux435ux440ux44b-ux43fux440ux438ux43cux435ux43dux435ux43dux438ux44f-ux444ux43eux440ux43cux443ux43bux44b-ux434ux430ux43bux430ux43cux431ux435ux440ux430}

\section{1. Движение синусондальной
волны}\label{ux434ux432ux438ux436ux435ux43dux438ux435-ux441ux438ux43dux443ux441ux43eux43dux434ux430ux43bux44cux43dux43eux439-ux432ux43eux43bux43dux44b}

Рассмотрим начальные условия вида

\[
J[y] = \int_a^b y \sqrt{1 + y'^2} \, dx,
\]

\(u_1(x, 0) = 0.\) По формуле Даламбера получаем решение

\[
\overline{y}(x) = \alpha \operatorname{ch}[(x-\beta)/\alpha],
\]

Его можно интерпретировать следующим образом. Начальное смещение струны
\(u(x, 0) = \sin x\) делится на две одинаковые части

\[
\overline{y}(x) = 0.42e^x - 0.42e^{-x}
\]

H \(\frac{\sin x}{2}\) .

Каждая из частей распространяется со скоростью с в виде бегущих волн.
Одна из волн перемещается слева направо, а вторая---в противоположном
направлении. Читателю предлагается подумать над тем, как выглядит в этом
случае результирующая волна.

\section{2. Распространение простейшего прямоугольного
импульса}\label{ux440ux430ux441ux43fux440ux43eux441ux442ux440ux430ux43dux435ux43dux438ux435-ux43fux440ux43eux441ux442ux435ux439ux448ux435ux433ux43e-ux43fux440ux44fux43cux43eux443ux433ux43eux43bux44cux43dux43eux433ux43e-ux438ux43cux43fux443ux43bux44cux441ux430}

В этом случае начальные условия задаются следующим образом:

\[
J[y] = \int_{0}^{1} [y^{2} + y'^{2}] dx.
\]

\(u_t(x, 0) = 0\) (в остальных точках).

Поскольку начальное смещение расшепляется на две полуволны, движущиеся в
противоположных направлениях, результирующее движение будет таким, как
показано на рис. 17.2.

\pandocbounded{\includegraphics[keepaspectratio]{images/_page_129_Figure_2.jpeg}}

РИС, 17.2. Начальное возмущение распадается на две бегущие волны.

\section{3. Задана начальная
скорость}\label{ux437ux430ux434ux430ux43dux430-ux43dux430ux447ux430ux43bux44cux43dux430ux44f-ux441ux43aux43eux440ux43eux441ux442ux44c}

Предположим теперь, что в начальный момент струна находится в положении
равновесия. Придадим струне начальную скорость (как в фортепиано) вида
\(\sin x\) , т. е.

\[
\int\limits_{t_1}^{t_2}
J\left[u\right] = \iint\limits_{D} F\left(x,\ y,\ u,\ u_{x},\ u_{y}\right) \, dx \, dy.
\]

Решение находим по формуле Даламбера

\[
F_{u} - \frac{\partial}{\partial x} F_{ux} - \frac{\partial}{\partial y} F_{uy} = 0.
\]

Оно представляет сумму двух бегущих косинусоидальных воли. Читатель
должен спросить себя, кажется ли ему это решение разумным.

На этом лекция 17 заканчивается. В следующей лекции мы покажем, как
формула Даламбера помогает дать интерпретацию в плоскости переменных x и
t.

\section{ЗАМЕЧАНИЯ}\label{ux437ux430ux43cux435ux447ux430ux43dux438ux44f-3}

\begin{itemize}
\item
  \begin{enumerate}
  \def\labelenumi{\arabic{enumi}.}
  \tightlist
  \item
    Общее решение уравнения с частными производными второго порядка
    содержит две произвольные функции, а общее решение обыкновенного
    дифференциального уравнения содержит две произвольные постоянные.
    Значит, у уравнения с частными производными больше решений, чем у
    обыкновенного дифференциального уравнения.
  \end{enumerate}
\item
  2 Замена переменных для приведения уравнения с частными производными к
  более простому виду является одним из общих методов теории. Новые
  координаты (ξ, η) принято называть каноническими координатами. Мы
  будем обращаться к ним и в дальнейшем, особенно при изучении
  гиперболических уравнений.
\item
  3 Метод нахождения общего решения уравнения с частными производными и
  подстановки этого решения в начальные условия не является общим
  методом решения уравнений с частными производными. Решение, которое
  рассматривается в данной лекции,--- единственный пример использования
  этого метода. Обычно невозможно найти общее решение уравнения с
  частными производными, но даже если удается это сделать, оказывается
  очень сложным подставить это решение в краевые условия.
\end{itemize}

\section{ЗАДАЧН}\label{ux437ux430ux434ux430ux447ux43d}

\begin{enumerate}
\def\labelenumi{\arabic{enumi}.}
\item
  Проверьте, что решение Даламбера (17-8) удовлетворяет начальным
  условиям задачи (17.1).

  \begin{enumerate}
  \def\labelenumii{\arabic{enumii}.}
  \setcounter{enumii}{1}
  \tightlist
  \item
    Подставьте (17.7) в общее решение u(x, t) = \psi(x - ct) + \psi(x +
    \epsilon t) и получите формулу Даламбера.
  \end{enumerate}
\item
  Найдите решение задачи Коши:
\end{enumerate}

\[
J[y] = \int_{0}^{1} \sqrt{1 + y'^{2}} dx.
\]

Как выглядит это решение в различные моменты времени? 4. Найдите решение
задачи Коши:

(УЧП)

\[
\int_{t_{1}}^{t_{2}} [KE - PE] dt = \frac{1}{2} \int_{t_{1}}^{t_{2}} [m\dot{y}^{2} - ky^{2}] dt
\]

(НУ)

\[
J[y] = \int_{0}^{\pi/2} [y'^{2} - y^{2}] dx
\]

Постройте график решения в различные моменты времени. 5. Решая совместно
уравнения (17.6) и первое уравнение из (17.5),

получите формулы (17.7).

\section{ФОРМУЛА ДАЛАМБЕРА
(ПРОДОЛЖЕНИЕ)}\label{ux444ux43eux440ux43cux443ux43bux430-ux434ux430ux43bux430ux43cux431ux435ux440ux430-ux43fux440ux43eux434ux43eux43bux436ux435ux43dux438ux435}

ЦЕЛЬ ЛЕКЦИИ! Показать, как с помощью формулы Даламбера можно найти
решение волнового уравнения на полуограниченной прямой

(УЧП)

\[
J[u] = \iint_D F[x, y, u, u_x, u_y) \, dx \, dy,
\]

(ГУ) \(u(0, t)\) = 0, \quad 0 \textless{} t \textless{} \infty,

(НУ)

\[
F_{u} - \frac{\partial}{\partial x} F_{uv} - \frac{\partial}{\partial y} F_{uv} = 0.
\]

Дополнительно дается интерпретация формулы Даламбера в координатной
плоскости x, t.

На прошлой лекции было показано, что если заданы начальное смещение
струны \(u(x, 0)\) = \(f(x)\) и начальная скорость \(u_t\)(x, 0) =
\(g(x)\) , то выражение

(18.i)

\[
u_{xx} + u_{yy} = f
\]

описывает смещение струны в произвольный момент времени. В этой лекции
мы познакомим читателя с одной интересной интерпретацией этой формулы в
плоскости переменных x и t (пространственновременная или фазовая
плоскость) и покажем, как изменить формулу Даламбера, чтобы получить
решение задачи о колебаниях полубесконечной струны.

Начнем с интерпретации формулы (18.1) в плоскости переменных x и t.

\section{Пространственно-временная интерпретация формулы
Даламбера}\label{ux43fux440ux43eux441ux442ux440ux430ux43dux441ux442ux432ux435ux43dux43dux43e-ux432ux440ux435ux43cux435ux43dux43dux430ux44f-ux438ux43dux442ux435ux440ux43fux440ux435ux442ux430ux446ux438ux44f-ux444ux43eux440ux43cux443ux43bux44b-ux434ux430ux43bux430ux43cux431ux435ux440ux430}

На предыдущей лекции мы показали, что решение задачи Коши

(18.2)

\[
J[u] = \iint_{D} [u_{x}^{2} + u_{y}^{2} + 2uf] dx dy.
J[u] = \int_{0}^{1} \int_{0}^{1} \left[ u_{1}^{2} + u_{y}^{2} \right] dx dy
u(x, t) = \frac{1}{2} [f(x-ct) + f(x+ct)] + \frac{1}{2c} \int_{x-ct}^{x+ct} g(\xi) d\xi.
u_{xx} + u_{yy} = 0
u(x, 0) = f(x),
u_{xx} + u_{yy} = 0
u(x, t) = \frac{1}{2} [f(x-ct) + f(x+ct)].
J[u] = \int_{0}^{1} \int_{0}^{1} \left[ u_{x}^{2} + u_{y}^{2} \right] dx dy
x-ct=x_0-ct_0
\]

(PPP)

\[
(Y \text{ ЧП}) \qquad u_{tt} = c^2 u_{xx}, \quad -\infty < x < \infty, \quad 0 < t < \infty,
J[u] = \iint_{D} [u] + u_{\eta} + 2uf ] dx dy.
(HY) \qquad
u_{xx} + u_{yy} = f
J[u] = \int_{0}^{1} \int_{0}^{1} [u' + u' + 2uf] dx dy,
J[u] = \int_{0}^{1} \int_{0}^{1} \left[ u_{x}^{2} + u_{y}^{2} + 2uf \right] dx dy,
\]

На рис. 18.2 решение изображено в плоскости переменных x, t, а не u, x,
как в предыдущей лекции.

\pandocbounded{\includegraphics[keepaspectratio]{images/_page_133_Figure_5.jpeg}}

РИС. 18.2. Решение задачи Коши (18.2) в xt-плоскости a --- передний
фронт волны, движущейся налево; b --- задний фронт волны; b --- переднии
фронт волны, движущейся направо

Теперь займемся интерпретацией формулы Даламбера, когда начальное
смещение равно нулю, а начальная скорость произвольна.

СЛУЧАЙ 2 (Начальное смещение равно нулю, начальная скорость
произвольна).

Рассмотрим начальные условия вида

\[
u_n(x, y) = a_1 \varphi_1(x, y) + a_2 \varphi_2(x, y) + \ldots + a_n \varphi_n(x, y),
\]

\(u_t(x, 0) = g(x).\) Решение в этом случае имеет вид

\[
\varphi_1(x, y) = xy(1-x)(1-y)
\]

Значение величины u в точке \((x_{\mathrm{ }e},\ t_{\mathrm{ }e})\)
можно интерпретировать как интеграл от начальной скорости в пределах от
\(x_{\mathrm{ }e}-ct_{\mathrm{ }e}\) до Р И С

\pandocbounded{\includegraphics[keepaspectratio]{images/_page_134_Figure_2.jpeg}}

. 18.3. Интерпретация задачи с начальной скоростью в x1-плоскости.

\(x_{\bullet}+ct_{\bullet}\) (рис. 18.3). В качестве примера на рис.
18.4 приводится графическая интерпретация решения задачи

(13.4)

\[
u_{1}(x, y) = a_{1}xy(1-x)(1-y),
u_{3}(x, y) = xy(1-x)(1-y)[a_{1} + a_{2}x + a_{3}y],
\vdots
u(x, t) = \frac{1}{2c} \int_{x-ct}^{x+ct} g(\xi) d\xi =
J[u_n] = \int_0^1 \int_0^1 \left\{ \left[ \sum_{j=1}^n a_j \frac{\partial \varphi_j}{\partial x} \right]^2 + \left[ \sum_{j=1}^n a_j \frac{\partial \varphi_j}{\partial y} \right]^2 + 2f \sum_{j=1}^n a_j \varphi_j \right\} dx dy
= \frac{1}{2c} \int_{x-ct}^{x+ct} 0 d\xi, \quad (x, t) \in \text{ области 1,}
\frac{\partial J\left[u_{n}\right]}{\partial a_{1}} = 2 \int_{0}^{1} \int_{0}^{1} \left\{ \sum_{j=1}^{n} \left[ \frac{\partial \varphi_{j}}{\partial x} \frac{\partial \varphi_{j}}{\partial x} + \frac{\partial \varphi_{j}}{\partial y} \frac{\partial \varphi_{1}}{\partial y} \frac{\partial \}{\partial y} \right] a_{j} + f \varphi_{1} \right\} dx dy = 0,
= \frac{1}{2c} \int_{x-ct}^{x+ct} d\xi, \quad (x, t) \in \text{ области 2,}
\]

Все это представляется изрядно сложным, но если переписать все эти
уравнения в матричной форме, то получится система линейных уравнений

\[
= \frac{1}{2c} \int_{x-ct}^{-1} d\xi, \quad (x, t) \in \text{ области 3,}
\]

где A = (\(A_{IJ}\)) матрица размером n \times n , элементы которой
вычисляются по формулам

(45.3)

\[
= \frac{1}{2c} \int_{x-ct}^{1} d\xi, \quad (x, t) \in \text{ области 4,}
\]

b = (\(b_t\)) --- вектор с компонентами

\[
= \frac{1}{2c} \int_{x-ct}^{x+ct} 0 d\xi, \quad (x, t) \in \text{ области 5,}
\]

а a=(\(a_i\)) --- неизвестный вектор, компоненты которого являются
коэффициентами в приближенном решении
\(a_n\)(x,y)=\(a_1\)\varphi\_1(x,y)+\(a_2\)\varphi\_2(x,y)+\ldots+\(a_n\)\varphi\_n(x,y)
.

\begin{verbatim}
ШАГ 3. Решаем линейную систему Aa = b относительно коэффициентов  a_1, a_2, \ldots, a_n  и получаем приближенную минимизирующую функцию
\end{verbatim}

\[
= \frac{1}{2c} \int_{x-ct}^{x+ct} d\xi, \quad (x, t) \in \text{ области 6.}
\]

а эначит, и приближенное решение задачи Дирихле.

\section{ЗАМЕЧАНИЯ}\label{ux437ux430ux43cux435ux447ux430ux43dux438ux44f-4}

\begin{enumerate}
\def\labelenumi{\arabic{enumi}.}
\tightlist
\item
  В этой лекции методом Ритца минимизировался функционал, вависящий от
  функции двух переменных. Точно так же можно минимизировать функционалы
  вида
\end{enumerate}

\[
y(0) = 0, \quad y(1) = 1.
\begin{array}{llll} (\text{Y}4\Pi) & u_{tt} = c^2 u_{xx}, & 0 < x < \infty, & 0 < t < \infty, \\ (\Gamma \text{Y}) & u(0, t) = 0, & 0 < t < \infty, \\ (\text{H}\text{Y}) & \begin{cases} u(x, 0) = f(x), \\ u_t(x, 0) = g(x). \end{cases} & 0 < x < \infty, \\ \end{array}
\varphi_i(0) = 0,
J[u] = \int_{0}^{1} \int_{0}^{1} [u^{2} + u^{2}_{y}] dx dy,
\]

\section{Решение задачи для полубесконечной струны с помощью формулы
Даламбера}\label{ux440ux435ux448ux435ux43dux438ux435-ux437ux430ux434ux430ux447ux438-ux434ux43bux44f-ux43fux43eux43bux443ux431ux435ux441ux43aux43eux43dux435ux447ux43dux43eux439-ux441ux442ux440ux443ux43dux44b-ux441-ux43fux43eux43cux43eux449ux44cux44e-ux444ux43eux440ux43cux443ux43bux44b-ux434ux430ux43bux430ux43cux431ux435ux440ux430}

Наша задача---найти волновые движения полубесконечной струны, левый
конец которой закреплен, при заданных начальных условиях. Для решения
задачи (18.6) мы воспользуемся тем же приемом, что и при отыскании
общего решения волнового уравнения на всей прямой. Подставляя общее
решение

\[
u_{xx} + u_{yy} = 1
\]

в начальные условия, находим (как в предыд, щей лекции)

(18.7)

\[
J[y] = \int_{0}^{1} [y^{2} + {y'}^{2}] dx, \ y(0) = 0, \ y(1) = 1,
z(x) = (1 - x) y(x),
\]

``Теперь перед нами проблема, которой не было в задаче о бесконечной
струне. Решение u\left(x,\t\right) должно быть определено всюду только
внутри первого квадранта (x\textgreater0,,t\textgreater0) плоскости
переменных x и t. Значит, мы должны уметь вычислить функции

\[
u_{xx} + u_{yy} = 0
\]

при всех -\infty \textless{} x-ct \textless{} \infty  , \psi(x+ct) при
всех 0 \textless{} x+ct \textless{} \infty  .

\begin{verbatim}
К сожалению, первая из формул (18.7) позволяет вычислять   $\varphi$(x-ct)   только при   x-ct\geqslant 0  , поскольку в начальных данных функции f(x) и g(x) определены только для положительных значений аргумента.

Если   x-ct \geqslant 0  , то после подстановки (18.7) в общее решение
\end{verbatim}

\[
J[u] = \int_{0}^{1} \int_{0}^{1} [u_{x}^{2} + u_{y}^{2}] dx dx.
\]

получаем

\[
u_{xx} + u_{yy} = \sin(\pi x)
u(x, y) = \left[Ae^{\pi y} + Be^{-\pi y} - \frac{1}{n^2}\right] \sin(\pi x),
\begin{split} u\left(x,\,t\right) &= \phi\left(x-ct\right) + \psi\left(x+ct\right) = \\ &= \frac{1}{2} \left[ f\left(x-ct\right) + f\left(x+ct\right) \right] + \frac{1}{2c} \int\limits_{x-ct}^{x+ct} g\left(\xi\right) \, d\xi. \end{split}
\]

\Дельта u + u\^{}2 = 0

\[
\Delta u = 0
\]

Что же делать, если x \textless{} ct? Пришла пора воспользоваться
граничным условием \(u(0,t)\)=0. С помощью этого граничного условия мы
доопределяем функцию \varphi(x-ct) при x \textless{} ct. Для этого
подставим общее решение \(u(x,t)\)=\varphi(x-ct)+\psi(x+ct) в граничное
условие \(u(0,t)\)=0. В результате получим \varphi(-ct)=-\psi(ct) ,

откуда

\[
\Delta u = 0
\]

Подстановка этого выражения для ф в общее решение дает

\[
\Delta u + \varepsilon u^2 = 0
\]

Комбинируя решения для x \textless{} ct и x \textgreater{} ct,
окончательно получаем (см. рис. 18.6)

\Дельта u + u\^{}2 = 0

\[
\Delta u = 0 \rightarrow \Delta u + \varepsilon u^2 = 0 \rightarrow \Delta u + u^2 = 0
\begin{cases} \frac{1}{2} \left[ f(x - ct) + f(x + ct) \right] + \frac{1}{2c} \int_{x - ct}^{x + ct} g(\xi) d\xi, & x \ge ct, \\ \frac{1}{2} \left[ f(x + ct) - f(ct - x) \right] + \frac{1}{2c} \int_{ct - x}^{x + ct} g(\xi) d\xi, & x < ct. \end{cases}
u_0 \rightarrow u = \sum_{k=0}^{\infty} \varepsilon^k u_k \rightarrow u = \sum_{k=0}^{\infty} u_k
\Delta u + cu^2 = 0,
\]

Этим завершается наша лекция. Интерпретация решения (18.8) будет
приведена в замечаниях.

\pandocbounded{\includegraphics[keepaspectratio]{images/_page_137_Figure_9.jpeg}}

\begin{verbatim}
РИС. 18.6. Интерпретация задачи Коши для полуограниченной струны в xt-плоскости a—возмущение отражается от границы;   $\delta$  —возмущение распространяется вдоль характеристики из точки (  x_0+ct_0  , 0).
\end{verbatim}

\section{ЗАМЕЧАНИЯ}\label{ux437ux430ux43cux435ux447ux430ux43dux438ux44f-5}

\begin{verbatim}
1. Решение (18.8) полностью соогветствует нашим представлениям о полубесконечной струне с граничным условием u(0, t) = 0. При   x \geqslant ct   оно совпадает с формулой Даламбера для бесконечной струны. Если x < ct, то формула модифицируется таким образом, что в нее включается волна, отраженная от границы (при отражении знак волны меняется на противоположный).
\end{verbatim}

\begin{enumerate}
\def\labelenumi{\arabic{enumi}.}
\setcounter{enumi}{1}
\tightlist
\item
  Если граничное условие \(u(0, t)\) = 0 изменить на другое, то и
  решение (18.8), естественно, изменится. Можно получить в явном виде
  решения для граничных условий
\end{enumerate}

\[
\Delta u + \varepsilon u^2 = 0
\]

или

\[
(46.1) u = u_0 + \varepsilon u_1 + \varepsilon^2 u_2 + \dots,
\]

Дополнительную информацию по этим вопросам читатель может найти в
литературе, указанной в конце книги.

\begin{enumerate}
\def\labelenumi{\arabic{enumi}.}
\setcounter{enumi}{2}
\tightlist
\item
  Прямые линии
\end{enumerate}

\Дельта u + u\^{}2 = 0

x - ct = \text{const} называются характеристиками. Возмущение
распространяется вдоль характеристик. Понятие характеристики тесно
связано с уравнениями гиперболического типа.

\section{ЗАДАЧИ}\label{ux437ux430ux434ux430ux447ux438-4}

\begin{enumerate}
\def\labelenumi{\arabic{enumi}.}
\tightlist
\item
  Решите смещанную задачу для полубесконечной струны
\end{enumerate}

\[
\Delta u + u^2 = 0,
\Delta u = 0, \quad 0 < r < 1, \quad 0 \le \theta < 2\pi,
\]

Постройте графики решения для разл чных моментов времени.

\begin{enumerate}
\def\labelenumi{\arabic{enumi}.}
\setcounter{enumi}{1}
\tightlist
\item
  Для решения предыдущей задачи можно воспользоваться следующим
  алгоритмом:
\end{enumerate}

а) продолжить начальные условия на всю действительную ось по формулам

имеет только вулевое решение. - Прим. ред.

\begin{enumerate}
\def\labelenumi{\arabic{enumi})}
\setcounter{enumi}{1}
\item
  Здесь и ниже речь идет об отыскании одного из решений (второе таким
  методом не находится).--- Прим ред.
\item
  Эта «ракета», во-первых, не при всех граничных условиях долетает до
  є=1 (где при опредсленных вещественных граничных условиях открывается
  мир комплексных решений) и, во-вторых, ни при каких в не долетает до
  второго вещественного решения, когда оно существует.--- Прим. ред.
\end{enumerate}

ции \(u_0\) , \(u_1\) , \(u_2\) , \ldots, входящие в степенной ряд
(46.1). В оставшейся части лекции мы покажем, как этим методом можно
решить различные краевые задачи.

\begin{verbatim}
# Решение нелинейного уравнения $\Delta$ u + u^2 = 0 методом возмущений
\end{verbatim}

Предположим, что нам нужно решить следующую нелинейную задачу Дирихле:

(46.2)

\(u_t(x, 0) = 0, -\infty < x < 0.\) б) усреднить две волны, движущиеся
налево и направо (как это было в предыдущей лекции);

\textbf{в}) рассмотреть полученное решение при \(x \ge 0\) .

Используя этот алгоритм, постройте графики решения для различных
моментов времени t, если начальное условие задано графически.

\pandocbounded{\includegraphics[keepaspectratio]{images/_page_139_Figure_3.jpeg}}

\begin{enumerate}
\def\labelenumi{\arabic{enumi}.}
\setcounter{enumi}{2}
\tightlist
\item
  Решите смешанную задачу для полубесконечной струны
\end{enumerate}

Розглянемо цю нелінійну задачу як збурену лінійну задачу:

(46.3)

Розв'язок якої \(u_0(r, \theta) = r \cos \theta\) відомий. Як змінити
\(u_0(r, \theta)\) , щоб отримати рішення проблеми (46.2)? Як уже
згадувалося, розглянемо клас задач \(\Delta u + \epsilon u^2 = 0\) з
граничними умовами і шукаємо розв'язок кожної задачі з цього класу у
вигляді

(46.4)

Воспользуйтесь методом, аналогичным рассмотренному в лекции. Дайте
физическую интерпретацию этой задачи.

\begin{enumerate}
\def\labelenumi{\arabic{enumi}.}
\setcounter{enumi}{3}
\tightlist
\item
  Предположим, что колебания струны описываются уравнением
  \(u_{tt} = u_{xx}\),начальное смещение изображено на рисунке, а
  начальная скорость равна нулю. Изобразите решение этой задачи в
  xt-плоскости. Обратите внимание на то, что начальные условия в этой
  задаче разрывны.
\end{enumerate}

\pandocbounded{\includegraphics[keepaspectratio]{images/_page_139_Figure_8.jpeg}}

\section{ВОЛНОВОЕ УРАВНЕНИЕ И ГРАНИЧНЫЕ
УСЛОВИЯ}\label{ux432ux43eux43bux43dux43eux432ux43eux435-ux443ux440ux430ux432ux43dux435ux43dux438ux435-ux438-ux433ux440ux430ux43dux438ux447ux43dux44bux435-ux443ux441ux43bux43eux432ux438ux44f}

ЦЕЛЬ ЛЕКЦИИ: Познакомыть читателя с основными типами граничных условий
для волнового уравнения

Потім знаходимо розв'язок задачі (46.2), поклавши
\(\varepsilon = 1^{\circ}\) ).

Вираз (46.4), який претендує на розв'язок, буде замінений на обурену
проблему

(PPP)

на конечном отрезке.

В дальнейшем будут рассматриваться три основных типа граничных условий.

\begin{enumerate}
\def\labelenumi{\arabic{enumi}.}
\tightlist
\item
  Заданный режим (граничные условия 1-го рода)
\end{enumerate}

, (GI) \(u(1, \theta) = \cos \theta\) .

У результаті ми отримуємо

\(u(L, t)\) = \(g_2\)(t). 2. Заданные силы (граничные условия 2-го рода)

\[
u_0 (1, \theta) + \varepsilon u_1 (1, \theta) + \varepsilon^2 u_2 (1, \theta) + \ldots = \cos \theta.
\]

\(u_x\)(L, t) = \(g_2\)(t). 3. Упругое закрепление (граничные условия
3-го рода)

\[
P_{0}
\]

\(u_x\)(L, t) - \gamma\_2 \(u(L, t)\) = \(g_2\)(t). Могут встречаться и
произвольные комбинации этих трех типоз граничных условий. В лекции
будет показан физический смысл этих граничных условий.

Итак, пока что мы рассмотрели единственный вид волнового
движения---поперечные колебания струны. Читатель должен понимать, что
это только один из видов волновых движений.

Ниже приведен перечень основных типов волн:

\begin{itemize}
\tightlist
\item
  \begin{enumerate}
  \def\labelenumi{\arabic{enumi}.}
  \tightlist
  \item
    Звуковые волны (продольные волны).
  \end{enumerate}
\item
  \begin{enumerate}
  \def\labelenumi{\arabic{enumi}.}
  \setcounter{enumi}{1}
  \tightlist
  \item
    Электромагнитные волны, в том числе световые.
  \end{enumerate}
\item
  \begin{enumerate}
  \def\labelenumi{\arabic{enumi}.}
  \setcounter{enumi}{2}
  \tightlist
  \item
    Колебания твердых тел (продольные, поперечные и крутильные).
  \end{enumerate}
\item
  \begin{enumerate}
  \def\labelenumi{\arabic{enumi}.}
  \setcounter{enumi}{3}
  \tightlist
  \item
    Волны вероятности в квантовой механике.
  \end{enumerate}
\item
  \begin{enumerate}
  \def\labelenumi{\arabic{enumi}.}
  \setcounter{enumi}{4}
  \tightlist
  \item
    Волны на воде (поперечные волны).
  \end{enumerate}
\item
  \begin{enumerate}
  \def\labelenumi{\arabic{enumi}.}
  \setcounter{enumi}{5}
  \tightlist
  \item
    Колебания струны (поперечные волны).
  \end{enumerate}
\end{itemize}

Цель нашей лекции --- обсудить различные типы граничных условий, которые
возникают при решении физических задач волнового характера. Мы
остановимся только на одномерных задачах и линейных граничных условиях.
Обычно различают граничные условия трех типов

\begin{enumerate}
\def\labelenumi{\arabic{enumi}.}
\tightlist
\item
  Заданные режимы в граничных точках (граничные условия 1-го рода)
\end{enumerate}

\[
\qquad u_{\theta}(r, \theta) = r \cos \theta,
\]

\(u(L, t)\) = \(g_2\)(t). 2. Заданные силы в граничных точках (граничные
условия 2-го рода)

\[
\begin{cases} \Delta u_{1} = -u_{0}^{2}, \\ u_{1}(1, \theta) = 0, \\ u_{2}(1, \theta). \end{cases}
\]

\(u_x\)(L, t) = \(g_2\)(t). 3. Упругое закрепление в граничных точках
(граничные условия 3-го рода)

\[
\begin{cases} \Delta u_{2} = -2u_{0}u_{1}, \\ u_{2}(1, \theta). \end{cases}
\]

\(u_x\)(L, t) - \gamma\_2 \(u(L, t)\) = \(g_2\)(t). Мы начнем с
граничных условий 1-го рода.

\begin{enumerate}
\def\labelenumi{\arabic{enumi}.}
\tightlist
\item
  Заданные режимы в граничных точках.
\end{enumerate}

Нам предстоит решить задачу вида

(19.1)

\[
P_{1}
\]

Щоб розв'язати \emph{неоднорідне рівняння}, ми також звернемося до
методу, який використовувався в теорії звичайних диференціальних рівнянь
і складається з наступного:

\begin{itemize}
\item
  \begin{enumerate}
  \def\labelenumi{\arabic{enumi})}
  \tightlist
  \item
    знайти загальний розв'язок \(y_n\) однорідного рівняння;
  \end{enumerate}
\item
  \begin{enumerate}
  \def\labelenumi{\arabic{enumi})}
  \setcounter{enumi}{1}
  \tightlist
  \item
    знайти часткове розв'язок \(y_p\) гетерогенного рівняння;
  \end{enumerate}
\item
  \begin{enumerate}
  \def\labelenumi{\arabic{enumi})}
  \setcounter{enumi}{2}
  \tightlist
  \item
    Підставити \(y_h\) + \(y_p\) у початкову умову та визначити
    константи.
  \end{enumerate}

  Ми використаємо цей метод для розв'язання нашої проблеми. У нашій
  задачі P\_1 загальний розв'язок однорідного рівняння \(\Delta\) u=0
  має вигляд із розділеними змінними
\end{itemize}

\[
\begin{aligned} u_{tt} &= c^2 u_{xx}, & 0 < x < 1, & 0 < t < \infty, \\ u(0, t) &= g_1(t), & 0 < t < \infty, \\ u(1, t) &= g_2(t), & 0 < t < \infty, \\ (HY) &\begin{cases} u(x, 0) &= f(x), \\ u_t(x, 0) &= g(x), \end{cases} & 0 \leqslant x \leqslant 1, \end{aligned}
\]

Часткове (лише одне) розв'язок неоднорідного рівняння

Ми розглядаємо у формі

\[
u_x(0, t) = 0, \quad u_x(L, t) = 0.
\]

Правильні частини форми \(r^n\) , \(r^n \cos{(n\theta)}\) ,
\(r^n \sin{(n\theta)}\) ведуть до рішень \(Ar^{n+2}\) ,
\(Br^{n+2} \cos{(n\theta)}\) \(Cr^{n+2} \sin{(n\theta)}\) відповідно.

Підставляючи функції \(u_p(r, \theta)\) у гетерогенне рівняння,
отримуємо A = -1/32, B = -1/24.

Отже,

\[
u_{x}(1, t) = \frac{1}{k} v(t) - (k - \text{модуль Юнга}).
\]

Щоб розв'язати задачу, \(P_1\) має зробити останній крок --- підставити
загальне розв'язання \(u(r, \theta) = u_h(r, \theta) + u_p(r, \theta)\)
на граничні умови \(u(1, \theta) = 0\) і знайти коефіцієнти \(a_n\) і
\(b_n\) .

Після цього ми отримуємо

\[
u_{x}(0, t) = \frac{h}{T} u(0, t),
\]

Отже, \(a_0=1/32\) , \(a_2=1/24\) , і всі інші \(a_n\) і \(b_n\)
дорівнюють нулю. Отже, рішення проблеми \(P_1\) знайдено:

\[
u_{x}(L, t) = -\frac{h}{T} u(L, t).
\]

Функція \(u_1(r, \theta)\) називається пертурбацією першого порядку до
\(u_4(r, \theta)\) . Додавши \(u_0\) і \(u_1\) , отримуємо новий підхід
до розв'язання задачі (46.2):

(46.5)

\[
u_{x}(0, t) - \frac{h}{T}u(0, t) = 0,
\]

Щоб знайти збурення наступного порядку \(u_{z}(r, \theta)\) , потрібно
розв'язати задачу

\[
u_{x}(L, t) + \frac{h}{T}u(L, t) = 0.
\begin{cases} \Delta u_{1} = -2u_{0}u_{1}, \\ u_{2}(1, \theta) = 0, \end{cases}
\]

де \(u_{\phi}(r,\theta)\) і \(u_{1}(r,0)\) вже відомі, раніше знайдені
функції. Звісно, виконувати всі ці алгебраїчні перетворення без допомоги
комп'ютера надзвичайно складно. На щастя, вираз (46.5) є майже точним
рішенням нашої проблеми. Насправді, якщо підставити цей вираз зліва від
\(\Delta u + u^{2}\) нашого нелінійного рівняння, отримаємо вираз, який
майже всюди відрізняється від нуля в межах кола 0 \textless{} r
\textless{} 1.

Окрім розв'язання нелінійних задач, теорія збурень може
використовуватися для розв'язання задач у областях з неправильними
межами (звісно, якщо нерегулярність не дуже велика). Давайте розглянемо
простий приклад.

\section{Приклад задачі з порушеною
межею}\label{ux43fux440ux438ux43aux43bux430ux434-ux437ux430ux434ux430ux447ux456-ux437-ux43fux43eux440ux443ux448ux435ux43dux43eux44e-ux43cux435ux436ux435ux44e}

Порушення можна порушити не лише рівняння з похідними похідними. Можливо
знайти розв'язок рівняння Лапласа всередині деформованого кола, збуривши
розв'язок рівняння Лапласа всередині кругової області. Наприклад,
припустимо, ми хочемо знайти потенціал у площі, обмеженій кривою
\(r=1+\frac{1}{4}\sin\theta\) (у полярних координатах). Нехай потенціал
u на межі заданий. Тобто, нам потрібно вирішити проблему

(46.6)

\[
u \left( 1 + \frac{1}{4} \sin \theta, \ \theta \right) = \cos \theta, \quad 0 \le \theta \le 2\pi.
\begin{aligned} u_x(0, T) &= \frac{h}{T} [u(0, t) - \theta_1(t)], \\ u_x(L, T) &= -\frac{h}{T} [u(L, t) - \theta_2(t)]. \end{aligned}
u\left(1+\frac{1}{4}\sin\theta,\ \theta\right)=\cos\theta
u(1+\epsilon\sin\theta, \theta) = \cos\theta, \quad 0 \le \epsilon \le 1/4.
\]

Этим мы завершаем обсуждение наиболее важных граничных условий,
связанных с гиперболическими задачами. Следующие несколько лекций мы
посвятим решению задач, содержащих полобные граничные условия.

\section{ЗАМЕЧАНИЯ}\label{ux437ux430ux43cux435ux447ux430ux43dux438ux44f-6}

\begin{enumerate}
\def\labelenumi{\arabic{enumi}.}
\tightlist
\item
  В лекции не рассматривалось граничное условие еще одного типа, когда
  струна подвергается действию силы, пропорциональной скорости и
  направленной в противоположном направ-
\end{enumerate}

лении. Граничное условие такого типа (на левом конце) записывается в
виде

\[
f(x+h) = f(x) + f'(x)h + f''(x)\frac{h^2}{2!} + \dots,
\]

\begin{enumerate}
\def\labelenumi{\arabic{enumi}.}
\setcounter{enumi}{1}
\tightlist
\item
  Нелинейное упругое закрепление левого конца приводит к граничному
  условию вида
\end{enumerate}

\[
u(1+\varepsilon\sin\theta,\ \theta)=u(1,\ \theta)+u_r(1,\ \theta)(\varepsilon\sin\theta)+u_{rr}(1,\ \theta)\frac{(\varepsilon\sin\theta)^2}{21}+\ldots
\]

где \(\varphi\)(u) --- произвольная функция величины u. Например,
граничное условие

\[
\Delta u = 0, \quad 0 < r < 1 + \varepsilon \sin \theta,
\]

говорит о том, что возвращающая сила на левом конце пропорциональна кубу
смещения (а не первой степени и, как это будет согласно закону Гука в
линейном случае).

\begin{enumerate}
\def\labelenumi{\arabic{enumi}.}
\setcounter{enumi}{2}
\tightlist
\item
  Если к нижнему концу пружины, совершающей продольные колебания,
  прикрепить массу \emph{m}, то граничное условие следует задавать в
  виде
\end{enumerate}

Коли всі розрахунки будуть виконані, ми внесемо \(\varepsilon=1/4\) у
кінцеві результати. Звісно, здається, що цю задачу не можна легко
розв'язати, але її можна звести до послідовності задач, у кожній з яких
визначена одна з функцій \(u_{\mathrm{ }e}\) , \(u_{\mathrm{ }e}\) ,
\ldots, що включені у розклад

\section{ЗАДАЧИ}\label{ux437ux430ux434ux430ux447ux438-5}

Опираясь на интуитивные представления о граничных условиях различных
типов, изобразить в общих чертах решение задачи

Після підстави ряду в (467) отримуємо таку послідовність задач для
визначення функцій \(u_a\) , \(u_1\) , \ldots;

\[
\begin{cases} \Delta u_{0} = 0, & 0 < r < 1 \text{ (внутри круга),} \\ u_{0}(1, \theta) = \cos \theta, & u_{0}(r, \theta) = r \cos \theta, \end{cases}
\]

\begin{verbatim}
\begin{array}{ll} (\mathrm{Y} \, \Pi) & u_{tt} = c^2 u_{xx}, \quad 0 < x < 1, \quad 0 < t < \infty, \\ (\mathrm{\Gamma} \, \mathrm{Y}) & \left\{ \begin{array}{ll} u \, (0, \, t) = 0, \\ u \, (1, \, t) = \sin t, \end{array}
\end{verbatim}

\[
\begin{cases} \Delta u_{1} = 0, & 0 < r < 1 \text{ (внутри круга),} \\ u_{1}(1, \theta) = -\sin \theta \frac{\partial u_{0}(1, \theta)}{\partial r} = -\sin \theta \cos \theta \end{cases}
u = u_0 + \frac{1}{4}u_1 + \left(\frac{1}{4}\right)^2 u_1 + \cdots
u_u + \frac{1}{4} u_t
\begin{array}{ll} u \, (x, \, 0) = 0, \\ u_{t} \, (x, \, 0) = 0, \end{array}
u_t = (1+x) u_{xx}, \\ u(x, 0) = q(x), \quad -\infty < x < \infty,
\]

\begin{verbatim}
\right. & 0 \leqslant x \leqslant 1, \end{array}
\end{verbatim}

\[
u_1 = (1 + \varepsilon_1) u_{xx}
u = u_1 + \varepsilon u_1 + \varepsilon^2 u_2 + \cdots
P_{0}
\]

в различные моменты времени. Могли бы вы угадать решение этой задачи?

\begin{enumerate}
\def\labelenumi{\arabic{enumi}.}
\setcounter{enumi}{2}
\tightlist
\item
  Во что перейдет граничное условие
\end{enumerate}

Отметим, что в первых двух задачах коэффициенты постоянны. Читатель
должен понимать, что параметр в должен быть достаточно мал, ниаче
бескопечный ряд может оказаться расходящимся

\section{ЗАДАЧИ}\label{ux437ux430ux434ux430ux447ux438-6}

\begin{verbatim}
- 1. Подставьте разложение (46.8) в задачу (46.7) и получите последовательность задач  P_0,\ P_1,\ P_2,\ \dots . 2. Покажите, что нелинейную задачу
\end{verbatim}

npit

\begin{itemize}
\tightlist
\item
  \begin{enumerate}
  \def\labelenumi{\alph{enumi})}
  \tightlist
  \item
    \(h \rightarrow \infty\) ,
  \end{enumerate}
\item
  \begin{enumerate}
  \def\labelenumi{\arabic{enumi})}
  \setcounter{enumi}{5}
  \tightlist
  \item
    \(h \rightarrow 0\) .
  \end{enumerate}
\end{itemize}

Согласуется ли результат с вашей интунцией?

\begin{enumerate}
\def\labelenumi{\arabic{enumi}.}
\setcounter{enumi}{3}
\tightlist
\item
  Изобразить в общих чертах решение задачи
\end{enumerate}

(УЧП)

можна звести до послідовності лінійних задач \(P_{o}\) , \(P_{i}\) ,
\(P_{z}\) .

\begin{itemize}
\tightlist
\item
  \begin{enumerate}
  \def\labelenumi{\arabic{enumi}.}
  \setcounter{enumi}{2}
  \tightlist
  \item
    Замініть нелінійну задачу 2 на приблизне розв'язання (46,5) і
    оцініть її точність.
  \end{enumerate}
\item
  \begin{enumerate}
  \def\labelenumi{\arabic{enumi}.}
  \setcounter{enumi}{3}
  \tightlist
  \item
    Розв'яжіть задачу \(P_1\) з прикладу з збуреною межою і перевірте,
    наскільки добре є наближене розв'язання
  \end{enumerate}
\end{itemize}

(ГУ)

задовольняє відношення

\(u\left(1 + \frac{1}{4}\sin\theta, \theta\right) = \cos\theta.\) \#
Лекція 47

\section{РОЗВ'ЯЗАННЯ ДИФЕРЕНЦІАЛЬНИХ РІВНЯНЬ У ПОХІДНИХ ПОХІДНИХ ЗА
ДОПОМОГОЮ МЕТОДУ КОНФОРМНИХ
ВІДОБРАЖЕНЬ}\label{ux440ux43eux437ux432ux44fux437ux430ux43dux43dux44f-ux434ux438ux444ux435ux440ux435ux43dux446ux456ux430ux43bux44cux43dux438ux445-ux440ux456ux432ux43dux44fux43dux44c-ux443-ux43fux43eux445ux456ux434ux43dux438ux445-ux43fux43eux445ux456ux434ux43dux438ux445-ux437ux430-ux434ux43eux43fux43eux43cux43eux433ux43eux44e-ux43cux435ux442ux43eux434ux443-ux43aux43eux43dux444ux43eux440ux43cux43dux438ux445-ux432ux456ux434ux43eux431ux440ux430ux436ux435ux43dux44c}

МЕТА ЛЕКЦІЇ: Показати, як деякі двовимірні граничні задачі можна
перетворити на інші, простіші задачі за допомогою конформного
відображення. Припустимо, наприклад, що ми хочемо розв'язати рівняння
Лапласа

\[
\begin{cases} u(0, t) = 0, \\ u_x(1, t) = -u(1, t), \\ u(x, 0) = x, \\ u_t(x, 0) = 0. \end{cases}
\]

з певною граничною умовою в області комплексної форми (площина змінних x
і y). Цю задачу можна перетворити на нову, в якій потрібно знайти
розв'язок рівняння Лапласа

у простішій області (площина змінних і і і).

\pandocbounded{\includegraphics[keepaspectratio]{images/_page_353_Figure_9.jpeg}}

\begin{verbatim}
У конформному відображенні рівняння Лапласа \varphi_{xx}+\varphi_{yy}=0 у координатній площині (x,y) повертається у рівняння Лапласа \varphi_{nn}\dashv -\varphi_{vo}=0 у координатній площині (u,v). Інші рівняння, такі як \varphi_{xx}+2\varphi_{xy}+\varphi_{yy}=0 , не мають такої властивості. Після розв'язання рівняння Лапласа $\varphi$(u,v) у простій області (

наприклад, у колі, напівплощині, квадраті) знайдено, достатньо підставити вирази U=U(X, Y), V=V(X, Y) у цей розв'язок, і ми отримаємо розв'язок нашої задачі $\phi$[u(x, y), v(x, y)] , вираженого через початкові змінні.
\end{verbatim}

У цій лекції ми розглянемо:

\begin{enumerate}
\def\labelenumi{\arabic{enumi})}
\item
  про властивості конформних відображень,
\item
  як їх будувати (для деяких районів),
\item
  як розв'язувати задачі методом конформних відображень.
\end{enumerate}

Однією з головних труднощів, що виникають при розв'язанні задач меж, є
складна форма кордонів. Навіть межі відносно простої форми часто роблять
задачу дуже складною для розв'язання. Звісно, можна спробувати
розв'язати задачу за допомогою методів теорії збурень, але теорія
збурень працює добре лише у випадках, коли межа близька до простої.

Однак існує спосіб розв'язати двовимірне рівняння Лапласа в областях з
комплексною межею. Цей метод базується на конформних відображеннях.

Перш ніж почати розв'язувати задачі цим методом, давайте трохи
ознайомимося з поняттям конформного відображення та комплексних змінних
загалом.

\section{Комплексні функції та конформні
відображення}\label{ux43aux43eux43cux43fux43bux435ux43aux441ux43dux456-ux444ux443ux43dux43aux446ux456ux457-ux442ux430-ux43aux43eux43dux444ux43eux440ux43cux43dux456-ux432ux456ux434ux43eux431ux440ux430ux436ux435ux43dux43dux44f}

У цій лекції ми ознайомимося лише з тими поняттями з теорії функцій
комплексної змінної, які нам знадобляться пізніше. Представимо
комплексне число z = x + iy як точку на площині комплексного xy (рис.
47.1).

Щоб ввести поняття конформного відображення, спершу ознайомимося з
поняттям функції комплексної змінної

.

Тут z --- це комплексна змінна (з деякої області площини z), а w ---
нова комплексна змінна, яка отримає значення відповідно до формули w =
\(f(z)\) у новій комплексній площині w.

Наприклад, комплексна функція, визначена в \emph{першому квадранті}
комплексної z-площини, \(w=z^2\) відображає цю область у всю півплощину
v\textgreater0 у комплексній площині w.

Таблиця 47 2

{\def\LTcaptype{none} % do not increment counter
\begin{longtable}[]{@{}
  >{\raggedright\arraybackslash}p{(\linewidth - 4\tabcolsep) * \real{0.2342}}
  >{\raggedright\arraybackslash}p{(\linewidth - 4\tabcolsep) * \real{0.2252}}
  >{\raggedright\arraybackslash}p{(\linewidth - 4\tabcolsep) * \real{0.5405}}@{}}
\toprule\noalign{}
\begin{minipage}[b]{\linewidth}\raggedright
Площина d
\end{minipage} & \begin{minipage}[b]{\linewidth}\raggedright
\(z \longrightarrow z^2\)
\end{minipage} & \begin{minipage}[b]{\linewidth}\raggedright
Площина S
\end{minipage} \\
\midrule\noalign{}
\endhead
\bottomrule\noalign{}
\endlastfoot
0 & & 0 \\
t1 → t & & -1(2) \\
Вісь x Перший квадрант & & Позитивна дійсна піввісь верхньої
півплощини \\
\end{longtable}
}

Таблиця 47.2 показує відповідність, яка встановлюється шляхом
відображення \(w=z^2\) між деякими точками на площинах z і w.

\pandocbounded{\includegraphics[keepaspectratio]{images/_page_355_Figure_5.jpeg}}

Рис. 47.1. Площина комплексної змінної \(\boldsymbol{z}\) та
некольорових формул бігунів Р И С У

\pandocbounded{\includegraphics[keepaspectratio]{images/_page_355_Figure_7.jpeg}}

НОК 47.2. Відображення першого квадранта площини \({\bf z}\) у верхню
півплощину \({\bf w}\) .

Інакше кажучи, якщо ми говоримо про комплексісні функції змінної w =
\(f(z)\), то можна розглянути питання, як криві площини z перетворюються
на криві в площині

Це питання не ставиться в теорії функцій дійсної змінної, з якою читач
уже добре знайомий.

Перш ніж рухатися далі, давайте розглянемо, як \(w=z^2\) складне
відображення типів у еквівалентній дійсній формі. Дійсна форма показує,
як початкові координати (x,y) у площині z перетворюються на нові
координати (u,v) у площині w \(w=z^2\) .

\[
\begin{cases} u(x, 0) = x, \\ u_t(x, 0) = 0. \end{cases}
\]

І давайте окремо прирівняємо реальні та уявні частини в цьому
співвідношенні. У результаті ми отримуємо

\[
\begin{array}{ll} u=x^2-y^2,\\ v=2xy, \end{array}
\]

(действительная форма отображения \(w=z^2\) ).

Эта форма понадобится нам позже. В этом примере она интересна тем, что
показывает, как гиперболы в плоскости z переходят в координатные линии u
= const и v = const в плоскости w.

Теперь мы перейдем к изучению одного частного вида функций комплексного
переменного, которые осуществляют конформные отображения.

\section{Определение конформного
отображения}\label{ux43eux43fux440ux435ux434ux435ux43bux435ux43dux438ux435-ux43aux43eux43dux444ux43eux440ux43cux43dux43eux433ux43e-ux43eux442ux43eux431ux440ux430ux436ux435ux43dux438ux44f}

Отображение w=\(f(z)\) комплексиой плоскости z на комплексную плоскость
w называется конформным в точке \(z_{\bullet}\) плоскости z, если
производная \(f'(z_{\bullet}) \neq 0\) . Отображение \(f(z)\) называется
конформным в области D, если \(f'(z) \neq 0\) в каждой точке области D.

Например, отображение \(f(z) = z^2\) конформно всюду, за исключением
точки z = 0, поскольку \(z'(z) = 2z \neq 0\) для всех \(z \neq 0\) . C
другой стороны, отображение \(e^z\) конформно во всей плоскости z, так
как всюду \(f'(z) = e^z \neq 0\) . Какая же польза от конформных
отображений? Ответ состоит в том, что, решая уравнение Лапласа
\(\phi_{xx} + \phi_{yy} = 0\) внутри какой-то области плоскости
переменных x и y, мы можем рассматривать эту плоскость как плоскость
комплексного переменного z. Рассмотрим комплексное отображение w =
\(f(z)\) плоскости z на плоскость w, в которой введены координаты (u,
v). При отображении w = \(f(z)\) уравнение Лапласа
\(\varphi_{xx} + \varphi_{yy} = 0\) преобразуется в некоторое новое
уравнение с частными производными, зависящее от новых координат и и v.
Замечательным является следующее: если отображение w = \(f(z)\) является
конформным в области, где задано уравнение \(\phi_{xy} + \phi_{uy} = 0\)
, то новое уравнение также будет уравнением Лапласа в координатах и и v.
To есть уравнение \(\phi_{xx} + \phi_{uu} = 0\) перейдет снова в
уравнение \(\varphi_{nn} + \varphi_{nn} = 0\) . Так мы приходим к идее
найти такое конформное отображение, которое переводит область со сложной
границей в область с простой границей (напомним, что отображение
\(w=z^*\) переводит границу первого квадранта плоскости z в
действительную ось плоскости w).

Приведем несколько примеров, в которых будет показано, как можно решать
уравнения с частными производными с помощью конформных отображений.

Уравнение Лапласа в верхней полуплоскости.

Предположим, что нам необходимо решить следующую задачу Дирихле в
верхней полувлоскости (рис. 47.3):

(47.1)

можно найти стандартным методом разделения переменных. При этом
оказывается, что решение u\left(x,,t\right) , представленное в виде ряда

Один із способів розв'язання цієї проблеми --- застосувати перетворення
Фур'є до змінної x, але набагато краще конформно відобразити верхню
половину площини у площину (u, v). Читачеві не буде складно побачити, що
конформне відображення

является суперпозицией простейших колебаний. Пространственная часть
\(X_n\)(x) каждого элементарного колебания является решением
(собственной функцией) некоторой краевой задачи Штурма---Лиувилля.

\begin{verbatim}
Итак, мы научились решать волиовое уравнение   u_{II}=c^2u_{xx}   на всей прямой   -\infty < x < \infty  . Формула Даламбера представляет решение в виде суммы двух бегущих волн, распространяющихся в противоположные стороны. Если же рассматривать это же уравнение на ограниченном промежутке 0 < x < L, то бегущих волн уже не станет, так как они будут взаимодействовать с границами. Вместо них возникают другие волны, которые называют стоячими волнами. Посмотрим, например, что случится, если гитарную струну, закрепленную в точках x=0, x=L, привести в движение. Для ответа на этот вопрос нам необходимо решить следующую смешанную задачу гиперболического типа:
\end{verbatim}

Вона має такі властивості (див. RNS. 47.4).

\begin{enumerate}
\def\labelenumi{\arabic{enumi}.}
\tightlist
\item
  Вона відображає верхню півплощину площини z у смуги
  \(-\infty < u < \infty\) , \(0 < v < \pi\) , у площині w.
\end{enumerate}

\pandocbounded{\includegraphics[keepaspectratio]{images/_page_357_Figure_11.jpeg}}

РИС 47.3 Задача граничного значення у верхній половині площини.

\begin{itemize}
\tightlist
\item
  \begin{enumerate}
  \def\labelenumi{\arabic{enumi}.}
  \setcounter{enumi}{1}
  \tightlist
  \item
    Ребро прямої y=0, -1 \textless{} x \textless{} 1, у площині z (де
    потенціал дорівнює 1), переходить у пряму \(v=\pi,\) -
    \(\infty < u < \infty\) , у площині w. - 3. Два промені
    x\textgreater1, y=0, і x\textless-1, y=0, у площині z проходять
    відповідно у позитивний і негативний
  \end{enumerate}
\end{itemize}

Читач повинен бути впевнений, що може підтвердити ці твердження.

\pandocbounded{\includegraphics[keepaspectratio]{images/_page_358_Figure_2.jpeg}}

Рис. 47.4. Перетворення складної задачі на легку за допомогою
конформного відображення.

Важливість цих властивостей полягає в тому, що вони перетворюють
початкову вадаху (47.1) на дуже просту задачу

(47.2)

Розв'язок очевидний: \varphi(u, v)\(= \frac{1}{\pi}v\).

\begin{verbatim}
Тепер, щоб знайти розв'язок початкової задачі, потрібно пояснити v через x і y і підставити цей вираз на q(u,v)=v'$\pi$ . Після цього ми отримуємо
\end{verbatim}

Оказывается, в этом случае бегущие волны отражаются от границ таким
образом, что результирующие колебания становятся не бегущими, а
сохраняющими форму в одном положении, т. е.

\pandocbounded{\includegraphics[keepaspectratio]{images/_page_147_Picture_5.jpeg}}

РИС. 20.1. Три типичных стоячих волны \(X(x)\) \(T(t)\).

превращаются в стоячие волны Несколько примеров стоячих волн показано на
рис. 20.1. Если известны профили этих стоячих волн \(X_n(x)\) и характер
колебаний каждой волны \(T_n(t)\) , то решение задачи о колебаниях
гитарной струны представляет собой суперпозицию простейших колебаний
\(X_n(x)\,T_n(t)\) и запи-

сывается в виде

\[
= \ln\left[\frac{z-1}{z+1}\right] + i \arg\left[\frac{z-1}{z+1}\right]
\]

Единственное, что остается сделать нам,--- это подобрать коэффициенты
\(c_n\) так, чтобы при t=0 удовлетворялись начальные условия
\(u\left(x,\ 0\right)=f\left(x\right)\) и
\(u_t\left(x,\ 0\right)=g\left(x\right)\) .

Задачу о колебаниях гитарной струны мы решим методом разделения
переменных.

\section{Решение задачи о колебаниях ограниченной струны методом
разделения
переменных}\label{ux440ux435ux448ux435ux43dux438ux435-ux437ux430ux434ux430ux447ux438-ux43e-ux43aux43eux43bux435ux431ux430ux43dux438ux44fux445-ux43eux433ux440ux430ux43dux438ux447ux435ux43dux43dux43eux439-ux441ux442ux440ux443ux43dux44b-ux43cux435ux442ux43eux434ux43eux43c-ux440ux430ux437ux434ux435ux43bux435ux43dux438ux44f-ux43fux435ux440ux435ux43cux435ux43dux43dux44bux445}

Чтобы решить смешанную задачу

(УЧП)

\[
v = \arg\left[\frac{z-1}{z+1}\right] = \arg\left[\frac{x+iy-1}{x+1+iy}\right] =
= \arg\left[\frac{2y}{x^2+y^2-1}\right]
\varphi(x, y) = \frac{1}{\pi} \arctan\left[\frac{2y}{x^2 + y^2 - 1}\right].
\begin{cases} u(0, t) = 0, \\ u(L, t) = 0, \end{cases}
x^{2} + y^{2} = 1,
\]

якщо потенціал на внутрішньому колі дорівнює одинці, а на зовнішньому
колі дорівнює двом (рис. 47.5). Іншими словами, нам потрібно розв'язати
проблему

(47.3)

Проблема полягає в тому, щоб знайти конформне відображення, яке
перетворює цю область на щось простіше, де нашу проблему легко
розв'язати. У цьому випадку очевидно, що ця проста ділянка має бути
кільцем. Відповідна експозиція Р И С У

\pandocbounded{\includegraphics[keepaspectratio]{images/_page_359_Figure_10.jpeg}}

НОК 47.5. Конформне відображення на кільце.

Це вже не так очевидно, але існують довідники, де можна знайти дисплей
із вказаними властивостями. Одним із найзагальних є «Довідник конформних
відображень»

G. Kober 1), яка містить сотні конформних відображень між різними
регіонами. Якщо читач витратить трохи часу і краще ознайомиться з
конформними таблицями відображення, він знайде відповідне конформне
відображення для задачі (47.3). У нашому конкретному випадку це буде
образ виду

(47.4)

де s=-0,146, а t=-6,85. У своїй еквівалентній дійсній формі це буде

(47.5)

\[
\begin{cases} u(x, 0) = f(x), \\ u_t(x, 0) = g(x), \end{cases}
v = \gamma [y(x-t)],
\gamma = 2t/[(x-t)^2 + y^2].
\]

найдем сначала стоячне волны, т. е. решения вида

\[
\begin{aligned} \varphi_{uu} + \varphi_{vv} &= 0, & \text{в кольче,} \\ (\Gamma Y) & & & & & & & \\ (\Gamma Y) & & & & & & \\ (\Gamma Y) & & & & & \\ & & & & & \\ (\Gamma Y) & & & & & \\ & & & & \\ & & & & \\ & & & & \\ & & & & \\ & & & \\ & & & \\ & & & \\ & & & \\ & & & \\ & & & \\ & & \\ & & \\ & & \\ & & \\ & & \\ & & \\ & & \\ & & \\ & & \\ & & \\ & & \\ & & \\ & & \\ & & \\ & & \\ & & \\ & & \\ & & \\ & & \\ & & \\ & & \\ & & \\ & & \\ & & \\ & & \\ & & \\ & & \\ & & \\ & & \\ & & \\ & & \\ & & \\ & & \\ & & \\ & & \\ & & \\ & & \\ & & \\ & & \\ & & \\ & & \\ & & \\ & & \\ & & \\ & & \\ & & \\ & & \\ & & \\ & & \\ & & \\ & & \\ & & \\ & & \\ & & \\ & & \\ & & \\ & & \\ & & \\ & & \\ & & \\ & & \\ & & \\ & & \\ & & \\ & & \\ & & \\ & & \\ & & \\ & & \\ & & \\ & & \\ & & \\ & & \\ & & \\ & & \\ & & \\ & & \\ & & \\ & & \\ & & \\ & & \\ & & \\ & & \\ & & \\ & & \\ & & \\ & & \\ & & \\ & & \\ & & \\ & & \\ & & \\ & & \\ & & \\ & & \\ & & \\ & & \\ & & \\ & & \\ & & \\ & & \\ & & \\ & & \\ & & \\ & & \\ & & \\ & & \\ & & \\ & & \\ & & \\ & & \\ & & \\ & & \\ & & \\ & & \\ & & \\ & & \\ & & \\ & & \\ & & \\ & & \\ & & \\ & & \\ & & \\ & & \\ & & \\ & & \\ & & \\ & & \\ & & \\ & & \\ & & \\ & & \\ & & \\ & & \\ & & \\ & & \\ & & \\ & & \\ & & \\ & & \\ & & \\ & & \\ & & \\ & & \\ & & \\ & & \\ & & \\ & & \\ & & \\ & & \\ & & \\ & & \\ & & \\ & & \\ & & \\ & & \\ & & \\ & & \\ & & \\ & & \\ & & \\ & & \\ & & \\ & & \\ & & \\ & & \\ & & \\ & & \\ & & \\ & & \\ & & \\ & & \\ & & \\ & & \\ & & \\ & & \\ & & \\ & & \\ & & \\ & & \\ & & \\ & & \\ & & \\ & & \\ & & \\ & & \\ & & \\ & & \\ & & \\ & & \\ & & \\ & & \\ & & \\ & & \\ & & \\ & & \\ & & \\ & & \\ & & \\ & & \\ & & \\ & & \\ & & \\ & & \\ & & \\ & & \\ & & \\ & & \\ & & \\ & & \\ & & \\ & & \\ & & \\ & & \\ & & \\ & & \\ & & \\ & & \\ & & \\ & & \\ & & \\ & & \\ & & \\ & & \\ & & \\ & & \\ & & \\ & & \\ & & \\ & & \\ & & \\ & & \\ & & \\ & & \\ & & \\ & & \\ & & \\ & & \\ & & \\ & & \\ & & \\ & & \\ & & \\ & & \\ & & \\ & & \\ & & \\ & & \\ & & \\ & & \\ & & \\ & & \\ & & \\ & & \\ & & \\ & & \\ & & \\ & & \\ & & \\ & & \\ & & \\ & & \\ & & \\ & & \\ & & \\ & & \\ & & \\ & & \\ & & \\ & & \\ & & \\ & & \\ & & \\ & & \\ & & \\ & & \\ & & \\ & & \\ & & \\ & & \\ & & \\ & & \\ & & \\ & & \\ & & \\ & & \\ & & \\ & & \\ & & \\ & & \\ & & \\ & & \\ & & \\ & & \\ & & \\ & & \\ & & \\ & & \\ & & \\ & & \\ & & \\ & & \end{aligned}
\]

Подставив это выражение в уравнение (20.1) и разделив переменные,
получим два обыкновенных дифференциальных уравнения

\[
\varphi(r) = a \ln r + b, \quad r = u^2 + v^2.
\varphi(u, v) = 0.57 \ln(u^2 + v^2) + 1
\]

где константа \lambda   --- пока что любое действительное число ---
\infty \textless{} \lambda \textless{} \infty   .

\begin{verbatim}
Результаты исследования решений этих двух уравнений при различных значениях   \lambda   приведены на рис 20.2.

Теперь мы должны *опбросить* те решения, которые либо неограничены при   t \to \infty  , либо становятся тождественно равными нулю, если потребовать выполнения граничных условий u(0, t) = u(L, t) = 0.

В качестве упражнения читателю предлагается проверить, что нетривиальные и ограниченные решения можно получить только при отрицательных значениях параметра   \lambda  . Таким образом, наша задача— найти такие значения констант A, B, C, D и отрица-

тельной константы   \lambda = -$\beta$^2  , чтобы выражение
\end{verbatim}

(20.2)

\[
\varphi(x, y) = 0.57 \ln(u^2 + v^2) + 1,
\]

удовлетворяло граничным условиям. Если это удастся, то мы будем
располагать множеством простейших колебаний струны, каждое из которых
удовлетворяет уравнению с частными произ- Р И С

\pandocbounded{\includegraphics[keepaspectratio]{images/_page_149_Figure_5.jpeg}}

20.2. Построение стоячих воли при раздичных значениях параметра λ.

водными и граничным условиям. Затем останется только найти такую
линейную комбинацию этих колебаний, чтобы сумма удовлетворяла начальным
условиям при t=0.

\begin{verbatim}
Подстановка (20.2) в граничные условия u(0, t) = u(L, t) = 0 дает   u(0, t) = X(0) T(t) = D[A \sin($\alpha$$\beta$ t) + B \cos($\alpha$$\beta$ t)] = 0 \Rightarrow D = 0  , (20.3)   u(L, t) = X(L) T(t) = C \sin($\beta$ L) [A \sin($\alpha$$\beta$ t) + B \cos($\alpha$$\beta$ t)] = 0 \Rightarrow \sin($\beta$ L) = 0  .

Другими словами, константа разделения   $\beta$   (мы можем вместо   \lambda   искать константу   $\beta$  ) должна удовлетворять уравнению   \sin{($\beta$ L)} = 0  , откуда
\end{verbatim}

\[
w = \ln \left[ \frac{z-1}{z+1} \right].
\]

Заметим, что если во втором уравнении (20.3) положить C=0, то получится
тривиальное решение X\left(x\right)\$\$T\left(t\right)\equiv0 . Таким
образом, мы нашли последовательность элементарных колебаний струны
(будем помечать их индексом n)

\[
z = re^{i\theta}
\]

(20.4)

или

\begin{verbatim}
\begin{aligned} (\text{Y} \forall \Pi) & \qquad \phi_{\text{vx}} + \phi_{yy} = 0, \quad 0 < x < \infty, \quad 0 < y < \infty, \\ \left\{ \begin{array}{l} $\phi$(x, \ 0) = \begin{cases} 1, & 0 < x < 1, \\ 0, & 1 \leqslant x < \infty, \end{cases} \\ $\phi$(0, \ y) = \begin{cases} 1, & 0 < y < 1, \\ 0, & 1 \leqslant y < \infty. \end{cases} \end{aligned}
\end{verbatim}

где \(a_n\) , \(b_n\) , \(R_n\) и \delta\_n ---произвольные постоянные.
Заметим, что элементарные колебания удовлетворяют граничным условиям.
Читателю должно быть ясно, что каждое элементарное колебание
представляет собой стоячую волну. В стоячей волне все точки колеблются с
одной и той же частотой, а форма волны подобна волнам, изображенным на
рис. 20.3.

\pandocbounded{\includegraphics[keepaspectratio]{images/_page_150_Picture_5.jpeg}}

РИС. 20.3. Стоячие волны \(u_n\)(x, t) = \(X_n\)(x) \(T_n\)(t) .

В нашей задаче уравнение и граничные условия линейны и однородны,
поэтому произвольная сумма элементарных стоячих волн также является
решением волнового уравнения и удовлетворяет граничным условиям. Мы
должны теперь найти такую суперпозицию стоячих волн, чтобы она
удовлетворяла также и начальным условиям, тогда получим решение нашей
задачи. Подстановка суммы

Зверніть увагу на зображення.

\pandocbounded{\includegraphics[keepaspectratio]{images/_page_362_Figure_1.jpeg}}

\pandocbounded{\includegraphics[keepaspectratio]{images/_page_362_Figure_2.jpeg}}

\begin{enumerate}
\def\labelenumi{\arabic{enumi}.}
\setcounter{enumi}{3}
\tightlist
\item
  Розв'язати змішану задачу Діріхле-Неймана з кутовим порошком за
  допомогою розчину під кутом 45°
\end{enumerate}

\pandocbounded{\includegraphics[keepaspectratio]{images/_page_362_Figure_5.jpeg}}

ПРИМІТКА. Комплексна функція \(w = \ln z = \ln |z| + i \arg(z)\)
відображає промінь \(\theta = c_1\) у площині z у пряму \(v = c_1\) , а
коло \(r = c_2\) у пряму \(u = \ln c_2\) у площині w (див. рисунки).

\pandocbounded{\includegraphics[keepaspectratio]{images/_page_362_Figure_7.jpeg}}

\section{ТАБЛИЦІ ІНТЕГРАЛЬНИХ
ПЕРЕТВОРЕНЬ}\label{ux442ux430ux431ux43bux438ux446ux456-ux456ux43dux442ux435ux433ux440ux430ux43bux44cux43dux438ux445-ux43fux435ux440ux435ux442ux432ux43eux440ux435ux43dux44c}

Таблиця A: Експоненціальне перетворення Фур'є,

Таблиця B: Перетворення синуса Фур'є.

Таблиця St Cosine Фур'є-перетворення.

Таблиця Dr скінченне перетворення Фур'є, таблиця та скінченне косинусне
перетворення Фур'є.

Таблиця D: Трансформація Лапласа.

\section{\texorpdfstring{Прийняті позначення \(\delta\left(x\right)\)
---
дельта-функція,}{Прийняті позначення \textbackslash delta\textbackslash left(x\textbackslash right) --- дельта-функція,}}\label{ux43fux440ux438ux439ux43dux44fux442ux456-ux43fux43eux437ux43dux430ux447ux435ux43dux43dux44f-deltaleftxright-ux434ux435ux43bux44cux442ux430-ux444ux443ux43dux43aux446ux456ux44f}

в начальные условия

\[
\begin{cases} 0, x < a, \\ 1, x \geqslant a, \end{cases}
\]

\(u_t\)(x, \theta) = \(g(x)\) приводит к двум уравнениям

\[
H(x-a) =
\]

Функція Хевісайду,

\[
H(a-x) = \begin{cases} 1, & x \leqslant a, \\ 0, & x > a, \end{cases}
\]

дзеркальною функцією Хевісайду.

ТАБЛИЦЯ A. Експоненціальне перетворення Фур'є

Воспользовавшись соотношениями ортогональности

1. f'(x) \textbar{} iωF (ω) \textbar{}

\textbar------------------------------\textbar-------------------------------------------------\textbar{}

\begin{verbatim}
| 2. f''(x) | -$\omega$^2 F($\omega$) |

| 3. fn (x), (n-та похідна) | (i$\omega$)^n F($\omega$) |

| 4. f(ax), a > 0 | \frac{1}{a} F\left(\frac{$\omega$}{a}\right) |

| 5. f(x-a) | e^{-ia$\omega$}F($\omega$) |

| 6. e^{-a^2x^2} | \frac{1}{a\sqrt{2}}e^{-$\omega$^2/4a^2} |

| 7. e^{-a_1 x} | \sqrt{\frac{2}{$\pi$}} \frac{a}{a^2 + $\omega$^2} |
\end{verbatim}

\begin{enumerate}
\def\labelenumi{\arabic{enumi}.}
\setcounter{enumi}{19}
\item
  гріх (ax)
\item
\item
\end{enumerate}

\[
\begin{cases} 0, & m \neq n, \\ L/2, & m = n, \end{cases}
\]

\begin{enumerate}
\def\labelenumi{\arabic{enumi}.}
\setcounter{enumi}{9}
\tightlist
\item
\item
\end{enumerate}

находим коэффициенты \(a_n\) и \(b_n\)

(20.5)

\begin{enumerate}
\def\labelenumi{\arabic{enumi}.}
\setcounter{enumi}{11}
\tightlist
\item
\end{enumerate}

Итак, задача решена. Ее решение записывается в виде

(20.6)

\begin{enumerate}
\def\labelenumi{\arabic{enumi}.}
\setcounter{enumi}{12}
\tightlist
\item
\end{enumerate}

а коэффициенты \(a_n\) и \(b_n\) находятся по формулам (20.5). Прежде
чем закончить лекцию, сделаем несколько замечаний.

\section{ЗАМЕЧАНИЯ}\label{ux437ux430ux43cux435ux447ux430ux43dux438ux44f-7}

\begin{enumerate}
\def\labelenumi{\arabic{enumi}.}
\item
  Если начальная \emph{скорость} струны равна нулю, то решение (20.6)
  принимает вид
\item
\end{enumerate}

и его можно интерпретировать следующим образом. Предположим, что мы
разложили начальное смещение струны

\begin{enumerate}
\def\labelenumi{\arabic{enumi}.}
\setcounter{enumi}{14}
\tightlist
\item
\end{enumerate}

в ряд по синусам

\begin{enumerate}
\def\labelenumi{\arabic{enumi}.}
\setcounter{enumi}{15}
\tightlist
\item
\end{enumerate}

Тогда каждый член такого разложения должен совершать колебания

\begin{enumerate}
\def\labelenumi{\arabic{enumi}.}
\setcounter{enumi}{16}
\tightlist
\item
\end{enumerate}

Если теперь просуммировать все эти колебания, то получится решение нашей
задачи. Например, предположим, что

начальное смещение струны задается функцией
\(f(x) = \sin(\pi x/L) + 0.5 \sin(3\pi x/ L) + 0,25 \sin(5\pi x/L)\) .

Полный отклик системы на начальное условие такого вида равен сумме
откликов на каждое слагаемое, т. е.

\begin{enumerate}
\def\labelenumi{\arabic{enumi}.}
\setcounter{enumi}{17}
\item
\item
  Слагаемое с номером n в разложении (20.6) называется \textbf{п-й модой
  колебаний} или \textbf{п-й гармоникой}. Используя тригонометрические
  формулы, можно записать эти гармоники в виде
\item
\end{enumerate}

где \(R_n\) и \(\delta_n\) --- новые произвольные постоянные (амплитуда
и фазовый угол). Эта новая форма представления n-й моды оказывается
более полезной при анализе колебаний. Частота колебаний од (рад/с) и-й
моды находится по формуле

\begin{enumerate}
\def\labelenumi{\arabic{enumi}.}
\setcounter{enumi}{9}
\tightlist
\item
\end{enumerate}

где Т --- натяжение струны, р --- линейная плотность струны. Отметим,
также, что частота n-й гармоники в n раз больше основной частоты (n=1).
Это свойство частоты характерно не для всех типов колебаний. Приятный
звук скривичной или гитарной струны обусловлен кратностью частогы
обертонов основной частоте. В барабане частоты гармоник высокого порядка
не кратны основной частоте колебаний.

\section{ЗАДАЧИ}\label{ux437ux430ux434ux430ux447ux438-7}

\begin{enumerate}
\def\labelenumi{\arabic{enumi}.}
\item
  Найдите решение задачи (20.1), если начальные условия задаются
  формулами
\item
\end{enumerate}

\(u_t(x, 0) = 0.\) Постройте графики решения для различных моментов
времени. Является ли это решение периодическим по времени? Каков его
период?

\begin{enumerate}
\def\labelenumi{\arabic{enumi}.}
\setcounter{enumi}{1}
\item
  Найдите решение задачи о колебаниях струны (20.1) при начальных
  условиях вида
\item
\end{enumerate}

\(u_t(x, 0) = \sin(3\pi x/L).\) Как выглядит график решения в различные
моменты времени?

\begin{itemize}
\tightlist
\item
  \begin{enumerate}
  \def\labelenumi{\arabic{enumi}.}
  \setcounter{enumi}{2}
  \tightlist
  \item
    Покажите, что если на рис. 20.2 величина \(\lambda \geqslant 0\) ,
    то решения \(X(x)\) \(T(t)\) будут либо не ограничены, либо
    тождественно равны нулю.
  \end{enumerate}
\item
  Найдите решение задачи о колебаниях струны при начальных условиях вида
\end{itemize}

\begin{enumerate}
\def\labelenumi{\arabic{enumi}.}
\setcounter{enumi}{11}
\tightlist
\item
\end{enumerate}

\(u_1(x, 0) = (3\pi \alpha/L)\sin(3\pi x/L).\) 5. Гитарная струна длины
L=1 оттягивается за среднюю точку на величину h (см. рис.) Начальное
положение струны можно задать в виде

\begin{enumerate}
\def\labelenumi{\arabic{enumi}.}
\setcounter{enumi}{12}
\tightlist
\item
\item
\item
\item
\end{enumerate}

\[
\begin{array}{llll} (\mathrm{Y} \Pi \Pi) & u_{tt} = \alpha^2 u_{xx} - \beta u_t, & 0 < x < 1, & 0 < t < \infty, \\ (\Gamma \mathrm{Y}) & \begin{cases} u \ (0, \ t) = 0, \\ u \ (1, \ t) = 0, \end{cases} & 0 < t < \infty, \\ (\mathrm{H} \mathrm{Y}) & \begin{cases} u \ (x, \ 0) = f \ (x), \\ u_1 \ (x, \ 0) = 0. \end{cases} & 0 \leqslant x \leqslant 1, \end{array}
\]

\begin{enumerate}
\def\labelenumi{\arabic{enumi}.}
\setcounter{enumi}{16}
\tightlist
\item
\item
\item
\end{enumerate}

\[
\begin{array}{ll} (\mathrm{Y} \mathrm{H} \Pi) & u_{tt} = \alpha^2 u_{xx} + Kx, \quad 0 < x < 1, \ 0 < t < \infty, \\ (\mathrm{\Gamma} \mathrm{Y}) & \begin{cases} u \ (0, \ t) = 0, \\ u \ (1, \ t) = 0, \end{cases} & 0 < t < \infty, \\ (\mathrm{H} \mathrm{Y}) & \begin{cases} u \ (x, \ 0) = f \ (x), \\ u_t \ (x, \ 0) = 0. \end{cases} & 0 \leqslant x \leqslant 1, \end{array}
\]

\begin{enumerate}
\def\labelenumi{\arabic{enumi}.}
\setcounter{enumi}{18}
\tightlist
\item
\item
\end{enumerate}

\section{КОЛЕБАНИЯ БАЛКИ (УРАВНЕНИЕ С ЧАСТНЫМИ ПРОИЗВОДНЫМИ ЧЕТВЕРТОГО
ПОРЯДКА)}\label{ux43aux43eux43bux435ux431ux430ux43dux438ux44f-ux431ux430ux43bux43aux438-ux443ux440ux430ux432ux43dux435ux43dux438ux435-ux441-ux447ux430ux441ux442ux43dux44bux43cux438-ux43fux440ux43eux438ux437ux432ux43eux434ux43dux44bux43cux438-ux447ux435ux442ux432ux435ux440ux442ux43eux433ux43e-ux43fux43eux440ux44fux434ux43aux430}

ЦЕЛЬ ЛЕКЦИИ: Показать, каким образом задача о колебании балки приводит к
уравнению четвертого порядка и как эта задача решается в случае свободно
опирающихся концов. Проводится сравнение колебаний балки с колебаниями
скрипичной струны.

\pandocbounded{\includegraphics[keepaspectratio]{images/_page_154_Figure_5.jpeg}}

Главное отличие поперечных колебаний тонкой балки от поперечных
колебаний струны в том, что балка оказывает сопротивление изгибу. Не
вдаваясь в механику тонких балок, мы примем, что учет сопротивления
изгиба приводит (вместо волнового уравнения) к уравнению четвертого
порядка

(21.1)

\section{\texorpdfstring{ТАБЛИЦЯ B. Перетворення Спнуса-Фур'є i
\sqrt{`\frac{\pi}{2}}`{[}\(\delta\)(\(\omega\)+a) -
\(\delta\)(\(\omega\)-a){]}}{ТАБЛИЦЯ B. Перетворення Спнуса-Фур'є i `{[}\textbackslash delta(\textbackslash omega+a) - \textbackslash delta(\textbackslash omega-a){]}}}\label{ux442ux430ux431ux43bux438ux446ux44f-b.-ux43fux435ux440ux435ux442ux432ux43eux440ux435ux43dux43dux44f-ux441ux43fux43dux443ux441ux430-ux444ux443ux440ux454-i-deltaomegaa---deltaomega-a}

rge \(\alpha^2 = K/\rho\) ,

K---модуль сдвига (чем больше K, тем жестче балка, тем выше частота
колебаний),

р --- линейная плотность балки (масса/ед. длины).

Поскольку читатель впервые встречается в нашей книге с уравнением выше
чем второго порядка, будет полезно решить типовую задачу о колебаниях
балки. Позже мы познакомимся и с другими задачами теории балок.

\section{Свободно опирающаяся
балка}\label{ux441ux432ux43eux431ux43eux434ux43dux43e-ux43eux43fux438ux440ux430ux44eux449ux430ux44fux441ux44f-ux431ux430ux43bux43aux430}

Рассмотрим малые колебания тонкой балки, концы которой свободно
опираются на две опоры. Говоря «свободно опираются», мы подразумеваем,
что концы балки не перемещаются, но на-

клоны балки в концевых точках могут изменяться (концы балки вакреплены с
помощью штифтовых устройств, рис. 21.1).

Совершенно ясно, что на концах балки должны выполняться граничные
условия

\[
F(\omega) = \frac{?}{\pi} \int_{0}^{\infty} f(x) \sin(\omega x) dx
\]

\(u(1, t) = 0,\) но совсем не очевидно, что на концах балки должны
удовлетворяться еще два граничных условия:

\[
0 < \omega < x
\]

\(u_{xx}(1, t) = 0.\) Используя теорию тонких балок, можно показать, что
изгибающий момент в балке пропорционален величине \(u_{xx}\) , а
изгибающий момент в свободно опирающемся конце должен быть ра- Р И С

\pandocbounded{\includegraphics[keepaspectratio]{images/_page_155_Figure_8.jpeg}}

.21.1. Свободно оппрающаяся балка.

вен нулю. Следовательно, изображенная на рис. 21.1 колеблющаяся балка
описывается следующей смешанной задачей (для простоты величина о принята
равной единице):

(21.2)

f'\,'(0)

\begin{enumerate}
\def\labelenumi{\arabic{enumi}.}
\setcounter{enumi}{1}
\tightlist
\item
  \(f(\alpha x)\)
\end{enumerate}

\[
\begin{aligned} & (\text{Y} \cdot \Pi) \quad u_{tt} = -u_{xxxx}, \quad 0 < x < 1, \quad 0 < t < \infty, \\ & \begin{cases} u \cdot (0, \ t) = 0, \\ u_{xx} \cdot (0, \ t) = 0, \\ u \cdot (1, \ t) = 0, \\ u_{xx} \cdot (1, \ t) = 0, \end{cases} \quad 0 < t < \infty, \\ & \begin{cases} u \cdot (x, \ 0) = f(x), \\ u_{1}(x, \ 0) = g(x). \end{cases} \quad 0 \leqslant x \leqslant 1, \end{aligned}
\]

\begin{enumerate}
\def\labelenumi{\arabic{enumi}.}
\setcounter{enumi}{2}
\item
  \(e^{-\alpha x}\)
\item
\end{enumerate}

Для решения этой задачи воспользуемся методом разделения переменных.
Будем искать только периодические решения, т. е. колебания вида

(21.3)

\[
[2/\pi\omega]^{1/2}
\]

Отметим, что выбор решения вида (21.3) на самом деле означает, что
константу разделения в методе разделения переменных мы взяли
отрицательной.

Подставим теперь (21.3) в уравнение колебаний балки и получим уравнение
для \(X\left(x\right)\)

\[
\frac{2}{\pi\omega} [1-\cos(\omega a)]
\]

общее решение которого записывается в виде

1

Для определения коэффициентов C, D, E и F подставим общее решение в
граничные условия

\[
e^{-1}a + \omega
\frac{1}{2}e^{-\alpha}\sin\omega
\begin{split} u\left(0,\ t\right) &= 0 \Rightarrow X\left(0\right)T\left(t\right) = 0 \Rightarrow X\left(0\right) = 0 \Rightarrow C + E = 0 \\ u_{xx}\left(0,\ t\right) &= 0 \Rightarrow X''\left(0\right)T\left(t\right) = 0 \Rightarrow X''\left(0\right) = 0 \Rightarrow -C + E = 0 \\ u\left(1,\ t\right) &= 0 \Rightarrow D\sin\sqrt{\omega} + F\sin\sqrt{\omega} = 0, \\ u_{xx}\left(1,\ t\right) &= 0 \Rightarrow -D\sin\sqrt{\omega} + F\sin\sqrt{\omega} = 0. \end{split}
\frac{1-e^{-a\omega}}{\omega}
\frac{2}{\pi}F''(\omega)
F \sinh \sqrt{\omega} = 0,
f(x) = \int_{0}^{\infty} F(\omega) \cos(\omega x) d\omega \qquad \qquad F(x) = \frac{2}{\pi} \int_{0}^{\infty} f(x) \cos(\omega x) dx
D \sin \sqrt{\omega} = 0.
\]

\begin{enumerate}
\def\labelenumi{\arabic{enumi}.}
\tightlist
\item
\end{enumerate}

\[
F = 0
\]

\section{ТАБЛИЦЯ D. Скінченне перетворення синуса
Фур'є}\label{ux442ux430ux431ux43bux438ux446ux44f-d.-ux441ux43aux456ux43dux447ux435ux43dux43dux435-ux43fux435ux440ux435ux442ux432ux43eux440ux435ux43dux43dux44f-ux441ux438ux43dux443ux441ux430-ux444ux443ux440ux454}

\begin{verbatim}
У скінченному перетворенні синуса Фур'є функція f(x), визначена в інтервалі 0 \le x \le $\pi$ , перетворюється у послідовність S_n , n = 1, 2 ... За формулою
\end{verbatim}

\[
\omega_n = (n\pi)^2,
\]

Коли йдеться про таблиці, зручно розглядати змінну x як належну до
інтервалу {[}0, l{]}. Будь-який інший чисельний інтервал {[}a, b{]}
можна перетворити на {[}0, \(\pi\){]} за допомогою формули

\[
u_n(x, t) = X_n(x) T_n(t) = [a_n \sin(n\pi)^2 t + b_n \cos(n\pi)^2 t] \sin(n\pi x).
\]

У скінченному синусоїдальному перетворенні немає нічого загадкового:
елементи послідовності \mathcal{S}\_n збігаються з коефіцієнтами в
\sin(nx) у розкладі функції \(f(x)\) у ряд Фур'є, тобто

\[
u(x, t) = \sum_{n=1}^{\infty} \left[ a_n \sin(n\pi)^2 t + b_n \cos(n\pi)^2 t \right] \sin(n\pi x)
f(x) = \sum_{n=1}^{\infty} S_n \sin(nx)
u(x, 0) = f(x) = \sum_{n=1}^{\infty} b_n \sin(n\pi x),
0 \le x \le \pi
u_t(x, 0) = g(x) = \sum_{n=1}^{\infty} (n\pi)^2 a_n \sin(n\pi x).
\]

\begin{enumerate}
\def\labelenumi{\arabic{enumi}.}
\tightlist
\item
\end{enumerate}

\[
a_n = \frac{2}{(n\pi)^2} \int_0^1 g(x) \sin(n\pi x) dx,
\]

\begin{enumerate}
\def\labelenumi{\arabic{enumi}.}
\setcounter{enumi}{8}
\tightlist
\item
\end{enumerate}

\[
b_n = 2 \int_0^1 f(x) \sin(n\pi x) dx.
\]

\(\frac{n}{n^2 + a^2}\){[}1 - (-1)\^{}n \(e^{ax}\){]}
10.\(\frac{\sin a (\pi - x)}{\sin a\pi} \frac{2n}{\pi (n^2 + a^2)}\)ТАБЛИЦЯ
E. Скінченне перетворення Косніуса Фур'є

\begin{verbatim}
У скінченному косинусному перетворенні Фур'є функція f(x), визначена на інтервалі 0 \le x \le $\pi$ , перетворюється у послідовність C_n , n=0,\ 1,\ 2,\ \ldots , за формулою
\end{verbatim}

\[
u(x, 0) = \sin(\pi x) + 0.5 \sin(3\pi x),
\]

Коли йдеться про таблиці, зручно розглядати змінну x як таку, що
належить до інтервалу {[}0,\(\pi\){]} . Будь-який інший числовий
інтервал {[}a,b{]} можна перетворити на {[}0,\(\pi\){]} за допомогою
формули

\[
a_n=0
\]

У скінченному косинусному перетворенні немає нічого загадкового:
елементи послідовності \(C_n\) збігаються з коефіцієнтами, коли
\cos{(nx)} у косинусному розкладі функції \(f(x)\) у ряд Фур'є, тобто

\[
u(x, t) = \cos(\pi^2 t) \sin(\pi x) + 0.5 \cos(9\pi^2 t) \sin(3\pi x)
f(x) = \frac{C_0}{2} + \sum_{n=1}^{\infty} C_n \cos(nx) \qquad C_n = \frac{2}{\pi} \int_{0}^{\pi} f(x) \cos nx \, dx
u(x, t) = \cos(\pi t) \sin(\pi x) + 0.5 \cos(3\pi t) \sin(3\pi x)
\]

\begin{enumerate}
\def\labelenumi{\arabic{enumi}.}
\tightlist
\item
\end{enumerate}

\[
u(0, t) = 0,
-n^{2}C_{n} - \frac{2}{\pi} \{f'(0) - (-1)^{n} f'(n)\}
u(x, t) = \sum_{n=1}^{\infty} X_n(x) \left[ a_n \sin(\omega_n t) + b_n \cos(\omega_n t) \right],
\begin{cases} 2, &n=0 \\ 0, &n=1, 2, \dots \end{cases}
(\text{УЧП}) \quad u_{tt} + u_{xxxx} = 0, \quad 0 < x < 1, \quad 0 < t < \infty,
\begin{cases} 1, & n=m \\ 0, & n\neq m \end{cases}
(\text{ГУ}) \quad
\]

x

\[
\begin{cases} u(0, t) &= 0, \\ u_x(0, t) &= 0, \\ u_{xx}(1, t) &= 0, \\ u_{xxx}(1, t) &= 0, \end{cases}
(\text{HV}) \quad
\begin{cases} u(x, 0) &= f(x), \\ u_t(x, 0) &= g(x). \end{cases}
\quad 0 \leqslant x \leqslant 1,
\]

6 M 601

УКАЗАНИЕ. Хотя собственные функции \(X_n(x)\) в этой задаче уже не
обычные синусы, согласно общей теории Штурма --- Ли вилля они будут
ортогональны на отрезке {[}0, 1{]}.

\begin{enumerate}
\def\labelenumi{\arabic{enumi}.}
\setcounter{enumi}{1}
\tightlist
\item
  Найдите решение для свободно опирающейся на оба конца балки при
  начальных условиях
\end{enumerate}

\[
u(x, 0) = \sin(\pi x),
\]

\(u_t(x, 0) = \sin(\pi x).\) \(0 \le x \le 1,\)

\begin{enumerate}
\def\labelenumi{\arabic{enumi}.}
\setcounter{enumi}{2}
\tightlist
\item
  Решите задачу 2 с начальными условиями
\end{enumerate}

\[
\begin{array}{l} u(x, 0) = 1 - x^2, \\ u_t(x, 0) = 0. \end{array}
\quad 0 \le x \le 1,
\]

\begin{enumerate}
\def\labelenumi{\arabic{enumi}.}
\setcounter{enumi}{3}
\tightlist
\item
  Пусть левый конец (x=0) балки жестко закреплен в стену, а правый конец
  (x=1) свободно опирается, как показано на рис. 21.3. Решите задачу для
  балки с такими граничными условиями и объясните, как найти собственные
  частоты колебаний в этом случае. Знать собственные частоты балок очень
  важно, чтобы избежать резонанса с внешними воздействиями.
\end{enumerate}

Лекция 22

\section{ПЕРЕХОД К БЕЗРАЗМЕРНЫМ
ПЕРЕМЕННЫМ}\label{ux43fux435ux440ux435ux445ux43eux434-ux43a-ux431ux435ux437ux440ux430ux437ux43cux435ux440ux43dux44bux43c-ux43fux435ux440ux435ux43cux435ux43dux43dux44bux43c}

ЦЕЛЬ ЛЕКЦИИ: Показать, каким образом краевые, начальные и другие виды
физических задач можно записать в безразмерном виде. В конкретных
областях знания (физике, химии, биологии или экономике) уравнения
записываются по-разному.

Переход к безразмерным переменным позволяет привести уравнения к одному
и тому же виду. Именно по этой причине при изучении уравнений с частными
производными в математике стараются отвлекаться от физического смысла
входящих в них параметров. Вот почему химик, физик или биолог должен
преобразовать свое уравнение к той форме, которая принята в нашей книге.

Основная идея введения новых (безразмерных) переменных состоит в том,
что после перехода к безразмерным переменным вадача становится чисто
математической и уже не содержит

карактерных физических констант. Именно таким способом многие различные
уравнения физики, биологии и химии, содержащие всякие нюансы, связанные
с физическими параметрами, приводятся к одной и той же простой форме
(см. рис. 22.1).

Познакомимся с этим преобразованием на простом примере.

\pandocbounded{\includegraphics[keepaspectratio]{images/_page_161_Figure_4.jpeg}}

РИС. 22.1. Некоторые задачи, приводящиеся к одной и той же безразмерной
форме.

\section{Приведение диффузионной задачи к безразмерному
виду}\label{ux43fux440ux438ux432ux435ux434ux435ux43dux438ux435-ux434ux438ux444ux444ux443ux437ux438ux43eux43dux43dux43eux439-ux437ux430ux434ux430ux447ux438-ux43a-ux431ux435ux437ux440ux430ux437ux43cux435ux440ux43dux43eux43cux443-ux432ux438ux434ux443}

\begin{verbatim}
Начнем со смешанной задачи для отрезка с начальным условием на температуру  u\left(x,\,0\right)=\sin\left($\pi$ x/L\right)  и граничными температурами на концах  T_1  и  T_2 . Другими словами, мы решаем задачу
\end{verbatim}

(22.1)

\begin{verbatim}
\begin{array}{ccccccccccccccccccccccccccccccccccc
\end{verbatim}

(см. рис. 22.2).

\pandocbounded{\includegraphics[keepaspectratio]{images/_page_161_Figure_10.jpeg}}

РИС. 22.2. Задача теплопроводности в плоскости переменных (х, т).

Наша цель --- дать новую эквивалентную (22.1) постановку задачи таким
образом, чтобы

в новой задаче не было никаких физических параметров (вроде α);

\begin{enumerate}
\def\labelenumi{\arabic{enumi})}
\setcounter{enumi}{1}
\item
  начальные и граничные условия стали проще.

  Для осуществления этой цели, мы введем три новых безразмерных
  переменных U, \xi   и \tau  , чтобы заменить прежние u, x и t по схеме
\end{enumerate}

\[
u \to U
\]

(безразмерная температура), x \to \xi   (безразмерная длина), t
\to \tau   (безразмерное время).

Мы проведем все три преобразования одновременно.

\begin{verbatim}
# Преобразование зависимой переменной  u \rightarrow U
\end{verbatim}

Определим функцию \(U(x, t)\) по формуле

\[
U(x, t) = \frac{u(x, t) - T_1}{T_2 - T_1}
\]

.

\begin{verbatim}
Ясно, что новая температура   U\left(x,\,t\right)   безразмерна, поскольку мы делим градусы Цельсия на градусы Цельсия. Новые граничные условия для   U\left(x,\,t\right)   при x=0 и x=L примут вид   U\left(0,\,t\right)=0   и   U\left(L,\,t\right)=1  . Давайте поставим задачу для новой функции   U\left(x,\,t\right)  . После несложных преобразований задача (22.1) переходит в новую задачу
\end{verbatim}

задачу

\begin{verbatim}
\begin{array}{ll} \text{(УЧП)} & U_t = $\alpha$^2 U_{xx}, \quad 0 < x < L, \quad 0 < t < \infty, \\ U_t = 0, \quad 0 < t < \infty, \\ U_t = 0, \quad 0 < t < \infty, \\ U_t = 0, \quad 0 < t < \infty, \\ U_t = 0, \quad 0 < t < \infty, \\ U_t = 0, \quad 0 < t < \infty, \\ U_t = 0, \quad 0 < t < \infty, \\ U_t = 0, \quad 0 < t < \infty, \\ U_t = 0, \quad 0 < t < \infty, \\ U_t = 0, \quad 0 < t < \infty, \\ U_t = 0, \quad 0 < t < \infty, \\ U_t = 0, \quad 0 < t < \infty, \\ U_t = 0, \quad 0 < t < \infty, \\ U_t = 0, \quad 0 < t < \infty, \\ U_t = 0, \quad 0 < t < \infty, \\ U_t = 0, \quad 0 < t < \infty, \\ U_t = 0, \quad 0 < t < \infty, \\ U_t = 0, \quad 0 < t < \infty, \\ U_t = 0, \quad 0 < t < \infty, \\ U_t = 0, \quad 0 < t < \infty, \\ U_t = 0, \quad 0 < t < \infty, \\ U_t = 0, \quad 0 < t < \infty, \\ U_t = 0, \quad 0 < t < \infty, \\ U_t = 0, \quad 0 < t < \infty, \\ U_t = 0, \quad 0 < t < \infty, \\ U_t = 0, \quad 0 < t < \infty, \\ U_t = 0, \quad 0 < t < \infty, \\ U_t = 0, \quad 0 < t < \infty, \\ U_t = 0, \quad 0 < t < \infty, \\ U_t = 0, \quad 0 < t < \infty, \\ U_t = 0, \quad 0 < t < \infty, \\ U_t = 0, \quad 0 < t < \infty, \\ U_t = 0, \quad 0 < t < \infty, \\ U_t = 0, \quad 0 < t < \infty, \\ U_t = 0, \quad 0 < t < \infty, \\ U_t = 0, \quad 0 < t < \infty, \\ U_t = 0, \quad 0 < t < \infty, \\ U_t = 0, \quad 0 < t < \infty, \\ U_t = 0, \quad 0 < t < \infty, \\ U_t = 0, \quad 0 < t < \infty, \\ U_t = 0, \quad 0 < t < \infty, \\ U_t = 0, \quad 0 < t < \infty, \\ U_t = 0, \quad 0 < t < \infty, \\ U_t = 0, \quad 0 < t < \infty, \\ U_t = 0, \quad 0 < t < \infty, \\ U_t = 0, \quad 0 < t < \infty, \\ U_t = 0, \quad 0 < t < \infty, \\ U_t = 0, \quad 0 < t < \infty, \\ U_t = 0, \quad 0 < t < \infty, \\ U_t = 0, \quad 0 < t < \infty, \\ U_t = 0, \quad 0 < t < \infty, \\ U_t = 0, \quad 0 < t < \infty, \\ U_t = 0, \quad 0 < t < \infty, \\ U_t = 0, \quad 0 < t < \infty, \\ U_t = 0, \quad 0 < t < \infty, \\ U_t = 0, \quad 0 < t < \infty, \\ U_t = 0, \quad 0 < t < \infty, \\ U_t = 0, \quad 0 < t < \infty, \\ U_t = 0, \quad 0 < t < \infty, \\ U_t = 0, \quad 0 < t < \infty, \\ U_t = 0, \quad 0 < t < \infty, \\ U_t = 0, \quad 0 < t < \infty, \\ U_t = 0, \quad 0 < t < \infty, \\ U_t = 0, \quad 0 < t < \infty, \\ U_t = 0, \quad 0 < t < \infty, \\ U_t = 0, \quad 0 < t < \infty, \\ U_t = 0, \quad 0 < t < \infty, \\ U_t = 0, \quad 0 < t < \infty, \\ U_t = 0, \quad 0 < t < \infty, \\ U_t = 0, \quad 0 < t < \infty, \\ U_t = 0, \quad 0 < t < \infty, \\ U_t = 0, \quad 0 < t < \infty, \\ U_t = 0, \quad 0 < t < \infty, \\ U_t = 0, \quad 0 < t < \infty, \\ U_t = 0, \quad 0 < t < \infty, \\ U_t = 0, \quad 0 < t < \infty, \\ U_t = 0, \quad 0 < t < \infty, \\ U_t = 0, \quad 0 < t < \infty, \\ U_t = 0, \quad 0 < t < \infty, \\ U_t = 0, \quad 0 < t < \infty, \\ U_t = 0, \quad 0 < t < \infty, \\ U_t = 0, \quad 0 < t < \infty, \\ U_t = 0, \quad 0 < t < \infty, \\ U_t
\end{verbatim}

Если угодно, то здесь можно остановиться, решить задачу для U\left(x,,
t\right) и затем найти u\left(x,,t\right) по формуле

\[
u(x, t) = T_1 + (T_2 - T_1) U(x, t).
\]

Однако мы пойдем дальше и займемся преобразованием независимых
переменных x и t.

\begin{verbatim}
# Преобразование пространственной координаты  x \rightarrow \xi

Представляется совершенно очевидным, как следует выбрать безразмерную переменную   \xi  . Поскольку   0 \leqslant x \leqslant L  , полагаем
\end{verbatim}

\[
\xi = x/L
\]

.

После вычисления производных

\[
U_x = U_{\xi} \xi_x = \frac{1}{L} U_{\xi},
U_{xx} = \frac{1}{L^2} U_{\xi\xi}
\]

становится очевидной постановка следующей задачи (в переменных U,
\(\xi\) , t):

\[
\begin{array}{ll} (\text{YM}\Pi) & U_t = (\alpha/L)^2 \, U_{\xi\xi}, \quad 0 < \xi < 1, \quad 0 < t < \infty, \\ (22.3) & (\Gamma \text{Y}) & \begin{cases} U(0,\ t) = 0, \\ U(1,\ t) = 1, \\ U(\xi,\ 0) = \frac{\sin\left(\pi\xi\right) - T_1}{T_2 - T_1}, \quad 0 \leqslant \xi \leqslant 1. \end{cases} \end{array}
\]

Мы проделали уже две трети пути. Последний шаг связан с введением
безразмерного времени \(\tau\) так, чтобы исключить коэффициент
\([\alpha/L]^2\) из дифференциального уравнения.

\section{\texorpdfstring{Преобразование времени
\(t \rightarrow \tau\)}{Преобразование времени t \textbackslash rightarrow \textbackslash tau}}\label{ux43fux440ux435ux43eux431ux440ux430ux437ux43eux432ux430ux43dux438ux435-ux432ux440ux435ux43cux435ux43dux438-t-rightarrow-tau}

Как выбрать безразмерное время? Это не так очевидно, как в случае первых
двух переменных. Однако, поскольку наша цель --- исключить константу
\([\alpha/L]^2\) из уравнения (22.3), мы поступим следующим образом.

\begin{itemize}
\item
  \begin{enumerate}
  \def\labelenumi{\arabic{enumi}.}
  \tightlist
  \item
    Попытаемся провести замену вида \(\tau = ct\) , где c неизвестная
    пока константа.
  \end{enumerate}

  \begin{itemize}
  \tightlist
  \item
    \begin{enumerate}
    \def\labelenumi{\arabic{enumi}.}
    \setcounter{enumi}{1}
    \tightlist
    \item
      Вычислим \(u_t = u_{\tau} \tau_t = c u_{\tau}\) .
    \end{enumerate}
  \item
    \begin{enumerate}
    \def\labelenumi{\arabic{enumi}.}
    \setcounter{enumi}{2}
    \tightlist
    \item
      Подставим эту производную в уравнение (22.3) и получим
    \end{enumerate}
  \end{itemize}
\end{itemize}

\[
cu_{\mathbf{t}} = [\alpha/L]^2 u_{\xi\xi}.
\]

Следовательно, нужно выбрать \(c = [\alpha/L]^2\) , а значит, новое
время должно вводиться по формуле

\[
\tau = [\alpha/L]^2 t.
\]

Применяя это преобразование к предыдущей задаче (22.3), можно полностью
перейти к безразмерным переменным U, \(\xi\) и \(\tau_i\)

\begin{verbatim}
\begin{array}{cccccccccccccccccccccccccccccccccccc
\end{verbatim}

Новая безразмерная система обладает следующими свойствами:

\begin{enumerate}
\def\labelenumi{\arabic{enumi}.}
\item
  Уравнение не содержит параметров.
\item
  Граничные условия имеют простой вид.
\item
  Начальное условие изменилось несущественно (т. е. оно содержит только
  известные функции).
\item
  Задача стала проще и компактнее по сравнению с исходной. Решение этой
  задачи можно построить раз и навсегда, поэтому если читатель
  преобразует свою исходную задачу (22.1) к безразмерному виду (22.4) и
  найдет ее решение \(U(\xi, \tau)\) в учебнике или научном журнале, то
  он может найти решение исходной задачи (22.1) простым пересчетом
\end{enumerate}

\[
u(x, t) = T_1 + (T_2 - T_1) U(x/L, \alpha^2 t/L^2).
\]

Этим завершается наше обсуждение проблемы преобразования задач к
безразмерному виду. Не существует универсальных правил, по которым
следует вводить новые переменные, здесь следует опираться на физическую
интуицию и пробовать различные возможности.

Мы закончим эту лекцию простым примером преобразования задачи к
безразмерному виду, решения новой задачи и затем возвращения к исходным
переменным.

\section{Пример преобразования гиперболической задачи к безразмерному
виду}\label{ux43fux440ux438ux43cux435ux440-ux43fux440ux435ux43eux431ux440ux430ux437ux43eux432ux430ux43dux438ux44f-ux433ux438ux43fux435ux440ux431ux43eux43bux438ux447ux435ux441ux43aux43eux439-ux437ux430ux434ux430ux447ux438-ux43a-ux431ux435ux437ux440ux430ux437ux43cux435ux440ux43dux43eux43cux443-ux432ux438ux434ux443}

Рассмотрим колеблющуюся струну

(22.5)

\[
\begin{array}{llll} (\text{УЧП}) & u_{tt} = \alpha^2 u_{xx}, & 0 < x < L, & 0 < t < \infty, \\ u(0, t) = 0, & 0 < t < \infty, \\ u(L, t) = 0, & 0 < t < \infty, \\ u(X, 0) = \sin(\pi x/L) + (0.5)\sin(3\pi x/L), \\ u_t(x, 0) = 0. & \end{array}
\]

С помощью преобразования независимых переменных к безразмерному виду
(преобразовывать u нет необходимости)

\[
\xi = x/L
\]

, \(\tau = [\alpha/L] t\)

получаем новую задачу

\[
(\text{YYII}) \qquad u_{\tau\tau} = u_{\xi\xi}, \quad 0 < \xi < 1, \quad 0 < \tau < \infty, 1
(\text{\GammaY}) \quad
\begin{cases} u(0, \tau) = 0, \\ u(1, \tau) = 0, \\ u(\xi, 0) = \sin(\pi\xi) + 0.5\sin(3\pi\xi), \\ u_{\xi}(\xi, 0) = 0, \end{cases}
\quad 0 < \xi \leqslant 1,
(\text{HY}) \quad
\begin{cases} u(\xi, 0) = \sin(\pi\xi) + 0.5\sin(3\pi\xi), \\ u_{\xi}(\xi, 0) = 0, \end{cases}
\quad 0 \leqslant \xi \leqslant 1,
\]

решением которой является функция

\[
u(\xi, \tau) = \cos(\pi \tau) \sin(\pi \xi) + 0.5 \cos(3\pi \tau) \sin(3\pi \xi).
\]

Если вернуться обратно к координатам x и t, то решение задачи (22.5)
принимает вид

\[
u(x, t) = \cos(\pi \alpha t/L)\sin(\pi x/L) + 0.5\cos(3\pi \alpha t/L)\sin(3\pi x/L).
\]

\section{ЗАМЕЧАНИЯ}\label{ux437ux430ux43cux435ux447ux430ux43dux438ux44f-8}

\begin{itemize}
\tightlist
\item
  \begin{enumerate}
  \def\labelenumi{\arabic{enumi}.}
  \tightlist
  \item
    Анализ размерности особенно важен в численных методах, поскольку
    большинство программ написаны для решения общих математических задач
    и не подразумевают использование большого числа физических
    параметров. Желающие воспользоваться этими программами должны
    преобразовать свою задачу к соответствующему виду, решить задачу на
    ЭВМ, а затем привести численные результаты к нужной размерности.
  \end{enumerate}
\item
  \begin{enumerate}
  \def\labelenumi{\arabic{enumi}.}
  \setcounter{enumi}{1}
  \tightlist
  \item
    Анализ размерностей позволяет математикам работать с уравнением в
    частных производных, исключив из него множество параметров и
    констант, которые не затрагивают математического существа задачи.
  \end{enumerate}
\item
  \begin{enumerate}
  \def\labelenumi{\arabic{enumi}.}
  \setcounter{enumi}{2}
  \tightlist
  \item
    Не всегда обязательно преобразовывать все переменные к безразмерному
    виду, иногда достаточно преобразовать одну или две переменных.
  \end{enumerate}
\end{itemize}

\section{ЗАДАЧИ}\label{ux437ux430ux434ux430ux447ux438-8}

\begin{enumerate}
\def\labelenumi{\arabic{enumi}.}
\tightlist
\item
  С помощью замены переменных
\end{enumerate}

\[
\xi = x/L
\]

, \(\tau = [\alpha/L] t\)

преобразуйте задачу (22.5) к безразмерному виду (22.6).

\begin{enumerate}
\def\labelenumi{\arabic{enumi}.}
\setcounter{enumi}{1}
\item
  Найдите безразмерную форму для задачи

  \begin{array}{lll} (\text{Y} \, \Pi) & u_t = $\alpha$^2 u_{xx}, & 0 < x < L, \\ (\Gamma \, \text{Y}) & \left\{ \begin{array}{ll} u \, (0, \, t) = T_1, \\ u \, (L, \, t) = 0, \\ u \, (x, \, 0) = T_2, & 0 \leqslant x \leqslant L. \end{array} \end{array}

  \right. \textbackslash{} \textbackslash end\{array\}
\item
  С помощью замены
\end{enumerate}

\[
U(x, t) = \frac{u(x, t) - T_1}{T_2 - T_1}
\]

преобразуйте задачу (22.1) в задачу (22.2).

\begin{verbatim}
4. Почему замена переменной   \tau = $\alpha$ t   приводит к исключению параметра   $\alpha$   из волнового уравнения   u_{tt} = $\alpha$^2 u_{xx}  ? В чем физический смысл этой замены? Напоминаем вам, что   $\alpha$  —скорость
\end{verbatim}

распространения волн. При поиске новых координат в большинстве случаев
главную роль играет интуиция.

Как выбрать новую пространственную координату ξ, чтобы исключить и из
уравнения

\[
u_t + vu_x = 0?
\]

Лекция 23

\section{КЛАССИФИКАЦИЯ УРАВНЕНИЙ С ЧАСТНЫМИ ПРОИЗВОДНЫМИ (КАНОНИЧЕСКАЯ
ФОРМА ГИПЕРБОЛИЧЕСКОГО
УРАВНЕНИЯ)}\label{ux43aux43bux430ux441ux441ux438ux444ux438ux43aux430ux446ux438ux44f-ux443ux440ux430ux432ux43dux435ux43dux438ux439-ux441-ux447ux430ux441ux442ux43dux44bux43cux438-ux43fux440ux43eux438ux437ux432ux43eux434ux43dux44bux43cux438-ux43aux430ux43dux43eux43dux438ux447ux435ux441ux43aux430ux44f-ux444ux43eux440ux43cux430-ux433ux438ux43fux435ux440ux431ux43eux43bux438ux447ux435ux441ux43aux43eux433ux43e-ux443ux440ux430ux432ux43dux435ux43dux438ux44f}

ЦЕЛЬ ЛЕКЦИИ: Показать, что линейное уравнение второго порядка с двумя
независимыми переменными

\[
Au_{xx} + Bu_{xy} + Cu_{yy} + Du_x + Eu_y + Fu = G
\]

(A, B, C, D, E, F и G являются функциями x и y или могут быть
константами) относится к одному из следующих типов:

\begin{itemize}
\item
  \begin{enumerate}
  \def\labelenumi{\arabic{enumi})}
  \tightlist
  \item
    гиперболическому (если B\^{}2 4AC \textgreater{} 0 ),
  \end{enumerate}
\item
  \begin{enumerate}
  \def\labelenumi{\arabic{enumi})}
  \setcounter{enumi}{1}
  \tightlist
  \item
    параболическому (если B\^{}2 4AC = ),
  \end{enumerate}
\item
  \begin{enumerate}
  \def\labelenumi{\arabic{enumi})}
  \setcounter{enumi}{2}
  \tightlist
  \item
    эллиптическому (если 0 B\^{}2 4AC \textless{} 0 ).
  \end{enumerate}

  Показать, какие новые переменные \xi = \xi(x, y) и \eta = \eta(x, y)
  можно ввести вместо x и y, чтобы упростить исходное уравнение. В новых
  переменных \xi   и \eta   уравнение с частными производными приводится
  к одному из следующих видов (в вависимости от того, будет ли величина
  B\^{}2 - 4AC положительна, равна нулю или отрицательна):
\item
  \begin{enumerate}
  \def\labelenumi{\arabic{enumi}.}
  \tightlist
  \item
  \end{enumerate}
\end{itemize}

\[
\begin{cases} u_{\xi\xi} u_{\eta\eta} = \Psi(\xi,  \eta, u, u_{\xi}, u_{\eta}) і \text{(дві канонічні форми гіперболічного рівняння),} \end{cases}
\]

\begin{itemize}
\item
  \begin{enumerate}
  \def\labelenumi{\arabic{enumi}.}
  \setcounter{enumi}{1}
  \tightlist
  \item
    \(u_{\eta\eta} = \Phi(\xi, \eta, u, u_{\xi}, u_{\eta}) (\)каноническая
    форма параболического уравнения),
  \end{enumerate}

  \begin{itemize}
  \tightlist
  \item
    \begin{enumerate}
    \def\labelenumi{\arabic{enumi}.}
    \setcounter{enumi}{2}
    \tightlist
    \item
      u\_\{\xi\xi\} + u\_\{\eta\eta\} = \Phi(\xi, \eta, u, u\_\{\xi\},
      u\_\{\eta\}) (каноническая форма эллиптического уравнения),
    \end{enumerate}
  \end{itemize}

  где \Phi   и \Psi  ---функции от частных производных первого порядка,
  зависимой переменной u и новых независимых переменных \xi   и \eta  .
  Конкретный вид функций \Phi   и \Psi   зависит от исходного уравнения.
\end{itemize}

У читателя может возникнуть мысль, что глава, посвященная классификации
уравнений, должна размещаться в самом начале книги. Вероятно, это
правильно, и многие книги действительно начинаются с этого материала.
Однако правда и то, что многие студенты не могут с энтузиазмом изучать
нечто такое, о чем у них нет никакого представления. Именно по этой
причине мы отложили изучение темы, посвященной классификации уравнений с
частными производными, до настоящего момента.

Наша задача - классифицировать все уравнения вида

(23 1)

\[
Au_{xx} + Bu_{xy} + Cu_{yy} + Du_x + Eu_y + Fu = G,
\]

где A, B, C, D, E, F и G---вообще говоря, функции {[}x и y. Согласно
предлагаемой схеме классификации уравнение называется

\begin{enumerate}
\def\labelenumi{\arabic{enumi})}
\item
  гиперболическим в точке (\(x_0\), \(y_0\)) , если B\^{}2(\(x_0\),
  \(y_0\)) - 4A(\(x_0\), \(y_0\)) \помножити на \помножити на C(\(x_0\),
  \(y_0\)) \textgreater{} 0 2) параболнческим в точке (\(x_0\), \(y_0\))
  , если B\^{}2(\(x_0\), \(y_0\)) - 4A(\(x_0\), \(y_0\)) \помножити на
  C(\(x_0\), \(y_0\)) = 0 ,
\item
  эллиптическим в точке (\(x_0\), \(y_0\)) , если B\^{}2(\(x_0\),
  \(y_0\)) - 4A(\(x_0\), \(y_0\)) \помножити на \помножити на C(\(x_0\),
  \(y_0\)) \textless{} 0 и в зависимости от типа его можно привести к
  одной из канонических форм. Для того чтобы было легче разобраться в
  классификационной схеме, мы приведем четыре примера гиперболических,
  параболических и эллиптических уравнений.
\end{enumerate}

\section{Примеры гиперболических, параболических и эллиптических
уравнений
10}\label{ux43fux440ux438ux43cux435ux440ux44b-ux433ux438ux43fux435ux440ux431ux43eux43bux438ux447ux435ux441ux43aux438ux445-ux43fux430ux440ux430ux431ux43eux43bux438ux447ux435ux441ux43aux438ux445-ux438-ux44dux43bux43bux438ux43fux442ux438ux447ux435ux441ux43aux438ux445-ux443ux440ux430ux432ux43dux435ux43dux438ux439-10}

\begin{enumerate}
\def\labelenumi{\arabic{enumi}.}
\tightlist
\item
  У равнение теплопроводности \(u_t = u_{xx}\)--- линейное уравнение
  второго порядка вида (23.1) с коэффициентами
\end{enumerate}

\[
A=1
,   B=0  ,   C=0  ,   D=0  ,   E=-1  ,   F=0  ,   G=0  ,
\]

так что B\^{}2-4AC=0 при всех x и t. Следовательно, это уравнение будет
параболическим при всех x и t. В общем уравнении (23 1) переменная y
играет роль времени t. Результат не изменится, если в общем уравнении
поменять местами x и y.

\begin{enumerate}
\def\labelenumi{\arabic{enumi}.}
\setcounter{enumi}{1}
\tightlist
\item
  Волновое уравнение \(u_{1t} = u_{xx}\)также принадлежит классу
  уравнений вида (23.1) с коэффициентами
\end{enumerate}

\[
A=1
,   B=0  ,   C=-1  ,   D=E=F=G=0  .
\]

Так как B\^{}2 - 4AC = 4 для всех x и t, то это уравнение является
еиперболическим при всех x и t.

1) Следует указать, что приведенные в примерах уравнения являются
частными случаями более общих параболических, гиперболических в
эллиптических уравнений, а основные их свойства характерны и для общего
случая.

\begin{enumerate}
\def\labelenumi{\arabic{enumi}.}
\setcounter{enumi}{2}
\item
  У равнение Лапласа \(u_{xx}\) + \(u_{yy} = 0\)является эллиптическим
  при всех x и y, так как B\^{}2 - 4AC = - 4 \textless{} 0 .

  \begin{enumerate}
  \def\labelenumii{\arabic{enumii}.}
  \setcounter{enumii}{3}
  \tightlist
  \item
    Линейное уравнение xu\_\{xx\} + u\_\{yy\} = \sin x с переменными
    коэффициентами также принадлежит классу уравнений вида (23.1), но
    для него B\^{}2-4AC=-4x и, следовательно, уравнение будет
    эллиптическим при x \textgreater{} 0,
  \end{enumerate}
\end{enumerate}

лараболическим при x=0, гиперболическим при x \textless{} 0.

Этот пример иллюстрирует тот факт, что тип уравнения о переменными
коэффициентами может изменяться от точки к точке.

Читатель должен обратить внимание на следующее: тип уравнения (23.1)
определяется только коэффициентами при вторых првизводных и никак не
зависит ни от коэффициентов при первых производных и самой функции, ни
от свободного члена.

\begin{verbatim}
Обратимся теперь к главной теме нашей лекции — приведению гиперболических уравнений к каноническому виду. Известно, что если в данной области пространства уравнение является гиперболическим, то вместо x и y можно ввести координаты   \xi   и   \eta  (характеристические переменные) таким образом, что уравнение примет простейшую форму
\end{verbatim}

(23.2)

\[
u_{\xi\eta} = \Phi(\xi, \eta, u, u_{\xi}, u_{\eta}).
\]

Это уравнение содержит только одну производную второго порядка
\(u_{\xi\eta}\) , а функция \Phi(\xi,) \eta,u,\(u_{\xi}\),\(u_{\eta}\))
зависит только от новых независимых переменных § и η, зависимой
переменной и и первых производных из и ип. Конкретный вид функции Ф
зависит, конечно, от исходного уравнения и формул перехода к новым
координатам § и п. Научимся отыскивать эти функции.

\section{Каноническая форма гиперболического
уравнения}\label{ux43aux430ux43dux43eux43dux438ux447ux435ux441ux43aux430ux44f-ux444ux43eux440ux43cux430-ux433ux438ux43fux435ux440ux431ux43eux43bux438ux447ux435ux441ux43aux43eux433ux43e-ux443ux440ux430ux432ux43dux435ux43dux438ux44f}

Рассмотрим уравнение общего вида

(23.3)

\[
Au_{xx} + Bu_{xy} + Cu_{yy} + Du_x + Eu_y + Fu = G
\]

и пусть в интересующей нас области B\^{}2 - 4AC \textgreater{} 0 . Нам
нужно ввести новые координаты 1)

\[
\xi = \xi(x, y),
\eta = \eta(x, y)
\]

так, чтобы в уравнении (23.3) осталась только одна частная производная
второго порядка \(u_{\xi\eta}\) (не следует пытаться искать такую замену
переменных, с помощью которой гиперболическое уравне-

1) Необходимо, чтобы это преобразование координат было локально обратимо
и дважды дифференцируемо. --- Прим. ред.

ние можно было бы привести и параболической или заливитической
канонической форме, это невозможно сделать).

Вычислим сначала частные производные

\[
\begin{array}{c} u_x = u_\xi \xi_x + u_\eta \eta_x, \\ u_y = u_\xi \xi_y + u_\eta \eta_y, \\ (23.4) \ u_{xx} = u_{\xi\xi} \xi_x^2 + 2u_{\xi\eta} \xi_x \eta_x + u_{\eta\eta} \eta_x^2 + u_{\xi} \xi_{xx} + u_{\eta} \eta_{xx}, \\ u_{xy} = u_{\xi\xi} \xi_x^2 + u_{\xi\eta} (\xi_x \eta_y + \xi_y \eta_x) + u_{\eta\eta} \eta_x \eta_y + u_{\xi} \xi_{xy} + u_{\eta} \eta_{xy}, \\ u_{yy} = u_{\xi\xi} \xi_y^2 + 2u_{\xi\eta} \xi_y \eta_y + u_{\eta\eta} \eta_y^2 + u_{\xi} \xi_{yy} + u_{\eta} \eta_{yy}. \end{array}
\]

Подставляя эти соотношения в исходное уравнение (23.3), после простых,
но громоздких вычислений получаем

(23.5)

\[
\overline{A}u_{\xi\xi} + \overline{B}u_{\xi\eta} + \overline{C}u_{\eta\eta} + \overline{D}u_{\xi} + \overline{E}u_{\eta} + \overline{F}u = \overline{G},
\overline{A} = A\xi_{x}^{2} + B\xi_{x}\xi_{y} + C\xi_{y}^{2},
\overline{B} = 2A\xi_{x}\eta_{y} + B(\xi_{x}\eta_{y} + \xi_{y}\eta_{x}) + 2C\xi_{y}\eta_{y},
\overline{C} = A\eta_{x}^{2} + B\eta_{x}\eta_{y} + C\eta_{y}^{2},
\overline{D} = A\xi_{xx} + B\xi_{xy} + C\xi_{yy} + D\xi_{x} + E\xi_{y},
\overline{E} = A\eta_{xx} + B\eta_{xy} + C\eta_{yy} + D\eta_{x} + E\eta_{y},
\overline{F} = F,
\overline{G} = G.
\]

Читателю предоставится возможность проделать эти выкладки при решении
задач.

\begin{verbatim}
Наш следующий шаг—выбрать функции   \xi = \xi(x, y)   и   \eta = \eta(x, y)   таким образом, чтобы коэффициенты   $\overline{A}$   и   $\overline{C}$   обратились в нуль. Это позволит нам привести исходное уравиение к канонической форме. Итак, необходимо потребовать, чтобы
\end{verbatim}

\[
\bar{A} = A\xi_x^2 + B\xi_x\xi_y + C\xi_y^2 = 0,
\]

\(\bar{C}\) = A\eta\_x\^{}2 + B\eta\_x\eta\_y + C\eta\_y\^{}2 = 0. Эти
уравнения можно представить в виде

\[
\begin{split} &A\left[\xi_{x}/\xi_{y}\right]^{2}+B\left[\xi_{x}/\xi_{y}\right]+C=0,\\ &A\left[\eta_{x}/\eta_{y}\right]^{2}+B\left[\eta_{x}/\eta_{y}\right]^{2}+C=0. \end{split}
\]

Разрешив их относительно {[}\xi\_x/\xi\_y{]} и {[}\eta\_x/\eta\_y{]} ,
получаем

(23.7)

\[
\begin{bmatrix} \xi_x/\xi_y \end{bmatrix}
= \frac{-B + \sqrt{B^2 - 4AC}}{2A}
\]

(характеристические

\[
[\eta_x/\eta_y] = \frac{-B - \sqrt{B^2 - 4AC}}{2A}
\]

уравнения).

Kaждое из квадратных уравнений для \([\xi_x/\xi_y]\) и
\([\eta_x/\eta_y]\) имеет два кория, но мы оставляем только по одному,
но так чтобы они были различны.

\pandocbounded{\includegraphics[keepaspectratio]{images/_page_170_Figure_3.jpeg}}

РИС. 23.1. Характеристики \(\xi(x, y) = c\) и \(\eta(x, y) = c\) .

Задача свелась к нахождению двух функций \(\xi(x, y)\) и \(\eta(x, y)\)
, таких, чтобы отношения \([\xi_x/\xi_y]\) и \([\eta_x/\eta_y]\)
удовлетворяли уравнению (23.7). Найти такие функции очень просто 1),
если обратить внимание на рис. 23.1. Чтобы понять, каким образом можно
это сделать, рассмотрим простое уравнение

\[
u_{xx} - 4u_{yy} + u_x = 0.
\]

Его характеристики определяются из уравнений

\[
\begin{vmatrix} \frac{dy}{dx} = -[\xi_x/\xi_y] = \frac{B - \sqrt{B^2 - 4AC}}{2A} = -2, \\ \frac{dy}{dx} = -[\eta_x/\eta_y] = \frac{B^{\frac{1}{2}} + \sqrt{B^2 - 4AC}}{2A} = 2. \end{vmatrix}
\]

Решая эти уравнения относительно у, получаем

\[
y = -2x + c_1,
y = 2x + c_2.
\]

Чтобы найти \xi   и \eta  , разрешим полученные уравнения относительно
констант \(c_1\) и \(c_2\) , оставляя их в правой части и перенося
остальные слагаемые в левую.

Выражения, стоящие в левых частях, примем за искомые функции § и пр

\[
\xi = y + 2x = c_1,
\eta = y - 2x = c_2.
\]

1) См., например, Тихонов А. Н., Самарский А. А. Уравнения
математической физики. --- М.: Наука, 1972 и др. издания. --- Прим. ред.

Очевидно, введенные таким образом функции \(\xi\) и \(\eta\)
удовлетворяют уравнениям характеристик. Эти новые координаты изображены
на рис. 23.2. На этом мы заканчиваем обсуждение вопроса о том, как найти
новые координаты. Осталось последнее --- определить новый вид исходного
уравнения.

\pandocbounded{\includegraphics[keepaspectratio]{images/_page_171_Figure_3.jpeg}}

РИС. 23.2. Характеристические координаты для уравнения
\(u_{xx} - 4u_{uu} + u_x = 0\) .

Это совсем просто: чтобы найти каноническую форму уравнения с частными
производными, подставим новые координаты \(\xi(x, y)\) и \(\eta(x,) y)\)
в уравнение

\[
\overline{A}u_{\xi\xi} + \overline{B}u_{\xi\eta} + \overline{C}u_{\eta\eta} + \overline{D}u_{\xi} + \overline{E}u_{\eta} + \overline{F}u = \overline{G}
\]

,

где \(\overline{A}\) , \(\overline{B}\) , \(\overline{C}\) ,
\(\overline{D}\) , \(\overline{E}\) , \(\overline{F}\) и
\(\overline{G}\) определены формулами (23.6).

Прежде чем закончить лекцию, давайте на конкретном примере посмотрим,
как «работает» общий метод приведения к каноническому виду.

Приведение гиперболического уравнения \(y^2u_{xx}-x^2u_{yy}=0\) к
канонической форме

Рассмотрим уравнение

\[
y^2u_{xx}-x^2u_{yy}=0
\]

, \(x>\) , \(0 y>0\) ,

которое является гиперболическим в первом квадранте. Найдем новые
координаты таким образом, чтобы привести исходное уравнение к
канонической форме.

ШАГ 1. Решаем уравнения характеристик

\[
\frac{dy}{dx} = \frac{B + \sqrt{B^2 - 4AC}}{2A} = \frac{x}{y}
\]

(напомним, что этот шаг

\[
\frac{dy}{dx} = \frac{B - \sqrt{B^2 - 4AC}}{2A} = -\frac{x}{y}.
\]

обеспечивает \(\overline{A} = \overline{C} = 0\) ),

\pandocbounded{\includegraphics[keepaspectratio]{images/_page_172_Figure_2.jpeg}}

РИС, 23.3. Новые карактеристические координаты.

Проинтегрировав эти обыкновенные дифференциальные уравнения методом
разделения переменных, получаем два неявных соотнощения (при желании
можно было бы явно выразить у через х)

\[
y^2 - x^2 = \text{const},
y^2 + x^2 = \text{const}.
\]

Следовательно, новые координаты следует вводить по формулам

\[
\xi = y^2 - x^2
\]

,

\[
\eta = y^2 + x^2  .
\]

Эти новые координаты изображены на рис. 23.3. Теперь следует получить
новое уравнение. Вычислим все коэффициенты

\begin{verbatim}
\widetilde{A} = 0   (так и должно быть; мы специально выбирали   \xi \xi   и   \eta   так, чтобы этот коэффициент обращался в нуль).
\end{verbatim}

\[
\vec{B} = 2A\xi_x\eta_x + B(\xi_x\eta_y + \xi_y\eta_x) + 2C\xi_y\eta_y = -16x^2y^2
\]

,

\begin{verbatim}
$\overline{C}$ = 0   (по той же причине, что и   $\overline{A}$ = 0  ),
\end{verbatim}

\[
\bar{D} = A\xi_{xx} + B\xi_{xy} + C\xi_{yy} + D\xi_x + E\xi_y = -2(x^2 + y^2)
\widetilde{E} = A\eta_{xx} + B\eta_{xy} + C\eta_{yy} + D\eta_x + E\eta_y = 2(y^2 - x^2),
\tilde{F} = F = 0
\]

,

\[
\overline{G} = G = 0
\]

,

подставим их в уравнение

\[
\overline{A}u_{\xi\xi} + \overline{B}u_{\xi\eta} + \overline{C}u_{\eta\eta} + \overline{D}u + \overline{E}u_{\eta} + \overline{F}u = \overline{G}
\]

и получим

\[
u_{\xi\eta} = \frac{-(x^2 + y^2) u_{\xi} + (y^2 - x^2) u_{\eta}}{8x^2y^3}.
\]

ШАГ 2. Выражая, наконец, x и y через \xi   и \eta  , окончательно
получаем

(23.8)

\[
u_{\xi\eta} = \frac{\eta u_{\xi} - \xi u_{\eta}}{2(\xi^2 - \eta^2)}.
\]

Задача решена. Теперь мы знаем, что нужно сделать, чтобы найти новые
координаты и привести уравнение к канонической форме.

\section{ЗАМЕЧАНИЯ}\label{ux437ux430ux43cux435ux447ux430ux43dux438ux44f-9}

\begin{enumerate}
\def\labelenumi{\arabic{enumi}.}
\tightlist
\item
  На самом деле для гиперболических уравнений существует две
  канонические формы. Вторую можно получить из первов заменой переменных
  вида
\end{enumerate}

\[
\alpha = \alpha (\xi, \eta) = \xi + \eta,
\]

\(\beta\) = \(\beta\) (\xi, \eta) = \xi - \eta. Воспользуемся этой
заменой для преобразования уравнения (23.8). Так как

\[
\begin{array}{l} u_{\xi} = u_{\alpha}\alpha_{\xi} + u_{\beta}\beta_{\xi} = u_{\alpha} + u_{\beta}. \\ u_{\eta} = u_{\alpha}\alpha_{\eta} + u_{\beta}\beta_{\eta} = u_{\alpha} - u_{\beta}, \\ u_{\xi\eta} = u_{\alpha\alpha}\alpha_{\eta} + u_{\alpha\beta}\beta_{\eta} + u_{\beta\alpha}\alpha_{\eta} + u_{\beta\beta}\beta_{\eta} = u_{\alpha\alpha} - u_{\beta\beta}, \end{array}
\]

TO

(23 9)

\[
u_{\alpha\alpha} - u_{\beta\beta} = \frac{-\beta u_{\alpha} - \alpha u_{\beta}}{2\alpha\beta}.
\]

При желании можно выразить \(\alpha\) и \(\beta\) через исходные
координаты x и y. В результате получаем

\[
\alpha = \xi + \eta = (y^2 - x^2) + (y^2 + x^2) = 2y^3,
\beta = \xi - \eta = (y^2 - x^2) - (y^2 + x^2) = -2x^2.
\]

\begin{itemize}
\tightlist
\item
  Читатель может спросить, зачем нужно классифицировать и приводить к
  канонической форме уравнения с частными производными.

  \begin{itemize}
  \tightlist
  \item
    а) Разбиение уравнений на классы гиперболических, параболических и
    эллиптических соответствует разбиению физических процессов на три
    основных класса: волновых, диффузионных и стационарных.
    Математические особенности решений уравнений этих трех типов
    совершенно различны.
  \item
    б) В большинстве работ, посвященных решению гиперболических задач,
    предполагается, что исходное уравнение записано в канонической
    форме, т. е. в виде
  \end{itemize}
\end{itemize}

\[
u_{\xi\xi} - u_{\eta\eta} = \Psi(\xi, \, \eta, \, u_{\xi}, \, u, u_{\eta}).
\]

Если у нас есть уравнение и мы хотим получить его решение, мы должны
преобразовать его к канонической форме и воспользоваться известными
результатами.

в) Для численного решения канонического гиперболического уравнения
создано много различных программ для ЭВМ.

Функция \(\Phi\) ( \(\xi\) , \(\xi \eta\) , u, \(u_{\xi}\) ,
\(u_{\eta}\) ) вводится в \(\Theta\) ВМ в виде подпрограммы, так что
сначала необходимо привести уравнение к канонической форме. После того
как решение в новых координатах найдено, его всегда можно преобразовать
обратно к исходным координатам.

\section{ЗАДАЧИ}\label{ux437ux430ux434ux430ux447ux438-9}

\begin{itemize}
\tightlist
\item
  \begin{enumerate}
  \def\labelenumi{\arabic{enumi}.}
  \tightlist
  \item
    Выясните, являются ли приведенные ниже уравнения гыперболическими,
    параболическими или эллиптическими:
  \end{enumerate}
\item
  : \textbf{a)} \(u_{xx} u_{xy} = 0\) ,

  \begin{itemize}
  \tightlist
  \item
    \begin{enumerate}
    \def\labelenumi{\arabic{enumi})}
    \setcounter{enumi}{5}
    \tightlist
    \item
      \(u_{tt} = u_{xx} + u_x + hu\) ,
    \end{enumerate}
  \item
    \textbf{B)} \(u_{xx} + 3u_{yy} = \sin x\) ,
  \item
    \begin{enumerate}
    \def\labelenumi{\alph{enumi})}
    \setcounter{enumi}{17}
    \tightlist
    \item
      \(u_{xx} + u_{yy} = f(x, y)\) ,
    \end{enumerate}
  \item
    \(\mathbf{D} ) \ u_{rr} + \frac{1}{r} u_r + \frac{1}{r^2} u_{\theta \theta} = f(r, \ \theta).\)
  \end{itemize}
\item
  \begin{enumerate}
  \def\labelenumi{\arabic{enumi}.}
  \setcounter{enumi}{1}
  \tightlist
  \item
    Получите соотношения (23.4), (23.5) и (23.6).
  \end{enumerate}
\item
  \begin{enumerate}
  \def\labelenumi{\arabic{enumi}.}
  \setcounter{enumi}{2}
  \tightlist
  \item
    Убедитесь в том, что уравнение
  \end{enumerate}
\end{itemize}

\[
3u_{xx} + 7u_{xy} + 2u_{yy} = 0
\]

гиперболического типа при всех x и y и найдите
\(x a p a \kappa m e p u c\) - m u + e c \kappa u + e c \kappa u + e c
\kappa u + e c \kappa u + e c \kappa u + e c \kappa u + e c \kappa u + e
c \kappa u + e c \kappa u + e c \kappa u + e c \kappa u + e c \kappa u +
e c \kappa u + e c \kappa u + e c \kappa u + e c \kappa u + e c \kappa u
+ e c \kappa u + e c \kappa u + e e c \kappa u + e c \kappa u + e c
\kappa u + e c \kappa u + e c \kappa u + e c \kappa u + e c \kappa u + e
c \kappa u + e c \kappa u + e c \kappa u + e c \kappa u + e c \kappa u +
e c \kappa u + e c \kappa u + e c \kappa u + e c \kappa u + e c \kappa u
+ e c \kappa u + e c \kappa u + e c \kap u + e c \kap

pa u + e c \kappa u + e c \kappa u + e c \kappa u + e c \kappa u + e c
\kappa u + e c \kappa u + e c \kappa u + e c \kappa u + e c \kappa u + e
c \kappa u + e c \kappa u + e c \kappa u + e c \kappa u + e c \kappa u +
e c \kappa u + e c \kappa u + e c \kappa u + e c \kappa u + e c \kappa u
+ e c \kappa u + e c \kappa u + e c \kappa u + e e c \kappa u + e c
\kappa u + e c \kappa u + e c \kappa u + e c \kappa u + e c \kappa u + e
c \kappa u + e c \kappa u + e c \kappa u + e c \kappa u + e c \kappa u +
e c \kappa u + e c \kappa u + e c \kappa u + e c \kappa u + e c \kappa u
+ e c \kappa u + e c \kappa u + e c \kappa u + e c \kappa u + e c
\kappa u + e c \kappa u + e c \kappa u + e c \kappa u + e c \kappa u + e
c \kappa u + e c \kappa u + e c \kappa u + e c \kappa u + e c \kappa u +
e c \kappa u + e c \kappa u + e c \kappa u + e c \kappa u + e c \kappa u
+ e c \kappa u + e c \kappa u + e c \kappa u + e c \kappa u + e c
\kappa u + e u + e c \kappa u + e c \kappa u + e c \kappa u + e c
\kappa u + e c \kappa u + e c \kappa u + e c \kappa u + e c \kappa u + e
c \kappa u + e c \kappa u + e c \kappa u + e c \kappa u + e c \kappa u +
e c \kappa u + e c \kappa u + e c \kappa u + e c \kappa u + e c \kappa u
+ e c \kappa u + e c ~kappa u + e c \kappa u + e c \kappa u + e c
\kappa u + e c \kappa u + e c \kappa u + e c \kappa u + e c \kappa u + e
c \kappa u + e c \kappa u + e c \kappa u + e c \kappa u + e c \kappa u +
e c \kappa u + e c \kappa u + e c \kappa u + e c \kappa u + e c \kappa u
+ e c \kappa u + e e c \kappa u + e c \kappa u + e c \kappa u + e c
\kappa u + e c \kappa u + e c \kappa u + e c \kappa u + e c \kappa u + e
c \kappa u + e c \kappa u + e c \kappa u + e c \kappa u + e c \kappa u +
e c \kappa u + e c \kappa u + e c \kappa u + e c \kappa u + e c \kappa u
+ e c \kappa u + e c \kappa u + e c \kappa u + e c \kappa u + e c
\kappa u + e c \kappa u + e c \kappa u + e c \kappa u + e c \kappa u + e
c \kappa u + e c \kappa u + e c \kappa u + e c \kappa u + e c \kappa u +
e c \kappa u + e c \kappa u + e c \kappa u + e c \kappa u + e c \kappa u
+ e c \kappa u + e c \kappa u + e c \kappa u + e u + e c \kappa u + e c
\kappa u + e c \kappa u + e c \kappa u + e c \kappa u + e c \kappa u + e
c \kappa u + e c \kappa u + e c \kappa u + e c \kappa u + e c \kappa u +
e c \kappa u + e c \kappa u + e c \kappa u + e c \kappa u + e c \kappa u
+ e c \kappa u + e c \kappa u + e c \kappa u + e c ~kappa u + e c
\kappa u + e c \kappa u + e c \kappa u + e c \kappa u + e c \kappa u + e
c \kappa u + e c \kappa u + e c \kappa u + e c \kappa u + e c \kappa u +
e c \kappa u + e c \kappa u + e c \kappa u + e c

\begin{enumerate}
\def\labelenumi{\arabic{enumi}.}
\setcounter{enumi}{3}
\tightlist
\item
  Приведите уравнение из предыдущей задачи к каноническому виду
\end{enumerate}

\[
u_{\xi\eta} = \Phi(\xi, \eta, u, u_1, u_2).
\]

\begin{enumerate}
\def\labelenumi{\arabic{enumi}.}
\setcounter{enumi}{4}
\tightlist
\item
  Продолжите вадачу 4 и приведите уравнение ко второй канонической форме
\end{enumerate}

\[
u_{\alpha\alpha} - u_{\beta\beta} = \Psi(\alpha, \beta, u, u_{\alpha}, u_{\beta}).
\]

\begin{enumerate}
\def\labelenumi{\arabic{enumi}.}
\setcounter{enumi}{5}
\tightlist
\item
  Найдите характеристики уравнения
\end{enumerate}

\[
u_{xx} + 4u_{xy} = 0.
\]

Преобразуйте уравнение к новым координатам, решите и, вернувшись к
координатам x и y, получите решение исходной вадачи.

\section{ВОЛНОВОЕ УРАВНЕНИЕ В СВОБОДНОМ ПРОСТРАНСТВЕ (ДВУМЕРНЫЕ И
ТРЕХМЕРНЫЕ
ЗАДАЧИ)}\label{ux432ux43eux43bux43dux43eux432ux43eux435-ux443ux440ux430ux432ux43dux435ux43dux438ux435-ux432-ux441ux432ux43eux431ux43eux434ux43dux43eux43c-ux43fux440ux43eux441ux442ux440ux430ux43dux441ux442ux432ux435-ux434ux432ux443ux43cux435ux440ux43dux44bux435-ux438-ux442ux440ux435ux445ux43cux435ux440ux43dux44bux435-ux437ux430ux434ux430ux447ux438}

ЦЕЛЬ ЛЕКЦИИ: Решить задачу с начальными условиями

(УЧП)

\[
u_{tt} = e^{2} (u_{xx} + u_{yy} + u_{zz}),
\begin{cases} -\infty < x < \infty, \\ -\infty < y < \infty, \\ -\infty < z < \infty, \end{cases}
\]

(НУ)

\[
\begin{cases} u(x, y, z, 0) = \varphi(x, y, z), \\ u_{t}(x, y, z, 0) = \psi(x, y, z) \end{cases}
\]

в трехмерном пространстве и показать, что полученное решение
удовлетворяет \emph{принципу Гюйгенса}. Для решения соответствующей
двумерной задачи

(УЧП)

\[
u_{tt} = c^2 (u_{xx} + u_{yy}), \quad
\begin{cases} -\infty < x < \infty, \\ -\infty < y < \infty, \end{cases}
\]

(НУ)

\[
\begin{cases} u(x, y, 0) = \varphi(x, y), \\ u_t(x, y, 0) = \psi(x, y) \end{cases}
\]

используется метод спуска и демонстрируется, что двумерные решения не
удовлетворяют принципу Гюйгенса. И наконец, с помощью метода спуска
показано, что решение одномерного уравнения выражается формулой
Даламбера, с которой мы познакомились ранее.

В предыдущих лекциях мы рассматривали задачу с начальными условиями для
бесконечной колеблющейся струны и показали, что ее решение дается
формулой Даламбера. Читателю должно быть ясно, что в трехмерном
пространстве одномерное волновое уравнение описывает плоские волны.
Например, звуковые или электромагнитные волны на достаточно больших
расстояниях от источников можно считать плоскими волнами и,
следовательно, описывать одномерным уравнением. Существует следующая
терминология:

\begin{itemize}
\tightlist
\item
  \begin{enumerate}
  \def\labelenumi{\arabic{enumi})}
  \tightlist
  \item
    одномерные волны называются плоскими волнами,
  \end{enumerate}
\item
  \begin{enumerate}
  \def\labelenumi{\arabic{enumi})}
  \setcounter{enumi}{1}
  \tightlist
  \item
    двумерные волны называются цилиндрическими волнами.
  \end{enumerate}
\end{itemize}

\begin{enumerate}
\def\labelenumi{\arabic{enumi})}
\setcounter{enumi}{2}
\tightlist
\item
  трехмерные волны называются сферическими волнами.
\end{enumerate}

Другими словами, одномерное волновое уравнение может описывать либо
плоские волны в пространствах большего числа измерений, либо колебания
струны. Задача лекции --- обобщить формулу Даламбера на случай двух и
трех измерений.

\section{Волны в трехмерном
пространстве}\label{ux432ux43eux43bux43dux44b-ux432-ux442ux440ux435ux445ux43cux435ux440ux43dux43eux43c-ux43fux440ux43eux441ux442ux440ux430ux43dux441ux442ux432ux435}

Рассмотрим сферические волны в трехмерном пространстве при ваданных
начальных условиях, т. е. решим задачу с начальными условиями

(24.1)

\[
\begin{cases} -\infty < x < \infty, \\ -\infty < y < \infty, \\ -\infty < z < \infty, \end{cases}
\]

(24.1)

\[
\begin{cases} u(x, y, z, 0) = \varphi(x, y, z), \\ u_{t}(x, y, z, 0) = \psi(x, y, z). \end{cases}
P_{2}
\begin{aligned} u_{tt} &= c^2 \Delta u, \\ (HY) & \begin{cases} u(x, y, z, 0) = 0, \\ u_t(x, y, z, 0) = \psi(x, y, z), \end{cases} \end{aligned}
\]

где экаком \(\Delta\) обозначен дифференциальный оператор

Цю задачу можна розв'язати методом \Phi  Юр'єра \^{}\{1\} ), і
виявляється, що розв'язок можна представити у вигляді

\[
P_{1}
\]

где \(\overline{\psi}\) --- \emph{среднее} значение начального
распределения \psi  *по сфере* радиуса ct с центром в точке (x, y, z),
т. е.

\(+ ct \cos \varphi\) ) \((ct)^2 \sin \varphi d\varphi d\theta\) .

Аргументы функции \(\psi\) пробегают по поверхности сферы, если величины
\(\theta\) и \(\phi\) изменяются в пределах \([0, 2\pi]\) и \([0, \pi]\)
соответственно (рис. 24.1). Эту формулу можно интерпретировать сле-

См Владимиров В. С. Уравнения математической физики. --- М.: Наука, 1971
и др. издания. --- Прим ред.

дующим образом: каждая точка пространства излучает расходящуюся (со
скоростью c) сферическую волну, через t секунд точка с координатами (x,
y, z) окажется под воздействием начального Р И С

\pandocbounded{\includegraphics[keepaspectratio]{images/_page_177_Picture_3.jpeg}}

. 24 1. Решение является средним значением начатьного распределения на
сфере.

возмущения, сосредоточенного на сфере раднуса ct с центром в заданной
точке (см. рис. 24.2).

Формула (24.3) позволяет рассчитать на ЭВМ решение задачи для
большинства начальных условий. Может быть, читателю будет интересно
найти решение для некоторых простых функций ф.

\pandocbounded{\includegraphics[keepaspectratio]{images/_page_177_Picture_7.jpeg}}

РИС. 24.2. Начальное возмущение \(\psi\) распространяется из каждой
точки во все стороны.

Чтобы получить полное решение, рассмотрим вторую половину задачи

(24.4)

\[
\begin{aligned} (\text{YAII}) & u_{tt} = c^2 \Delta u, \\ (\text{HY}) & \begin{cases} u = 0, \\ u_t = \varphi, \end{cases} \end{aligned}
u_h(r, \theta) = \sum_{n=0}^{\infty} r^n [a_n \cos(n\theta) + b_n \sin(n\theta)].
\Delta u = -\frac{r^2}{2} - \frac{r^3}{2} \cos{(2\theta)}
u = \frac{\partial}{\partial t} [t\overline{\varphi}].
u_p(r, \theta) = Ar^4 + Br^4 \cos(2\theta).
u(x, t) = \frac{1}{2c} \int_{x-ct}^{x+ct} \varphi(s) ds.
u_p(r, \theta) = -\frac{r^4}{32} - \frac{r^4}{24} \cos(2\theta).
u_t(x, t) = \frac{1}{2} [\varphi(x + ct) + \varphi(x - ct)].
\sum_{n=0}^{\infty} [a_n \cos(n\theta) + b_n \sin(n\theta)] - \frac{1}{32} - \frac{1}{24} \cos(2\theta) = 0.
(\text{УЧП}) \quad u_{tt} = c^2 \Delta u, \quad (x, y, z) \in \mathbb{R}^3,
u_1(r, \theta) = \frac{1}{32} + \frac{1}{24}r^2\cos(2\theta) - \frac{r^4}{24}\cos(2\theta) - \frac{r^4}{32} = \frac{(r^4 - 1)}{32} - \frac{(r^4 - r^2)}{24}\cos(2\theta).
(\text{HV}) \quad
u = u_0 + u_1 = r \cos \theta - \frac{(r^4 - 1)}{32} - \frac{(r^4 - r^2)}{24} \cos (2\theta).
\begin{cases} u = \varphi, \\ u_t = \psi \end{cases}
P_{2}
\Delta u = 0, \quad 0 < r < 1 + \frac{1}{4} \sin \theta,
\]

записывается в виде

\[
u \left( 1 + \frac{1}{4} \sin \theta, \ \theta \right) = \cos \theta, \quad 0 \le \theta \le 2\pi.
\]

где \(\overline{\)\phi\(}\) и \(\overline{\psi}\) --- средние значения
функций \(\phi\) и \psi  на сфере раднусом ct с центром в точке (x, y,
z).

\begin{verbatim}
Цей розв'язок зазвичай називають формулою Пуассона для хвильового рівняння у тривимірному просторі. Це природне узагальнення формули д'Аламбера на тривимірний випадок. Найважливішим у формулі Пуассона є те, що обидва інтеграли в  $\bar{$\phi$}$  і  $\bar{\psi}$  взяті з поверхні сфери. Цей факт дозволяє нам надати таку важливу інтерпретацію рішення.

У момент  t=t_1  розв'язок u у (x,y,z) залежить лише від початкових розподілів сфери з радіусом  ct_1  центрованим на (x,y,z) (рисунок 24.3). Тепер припустимо, що початкові розподіли  $\phi$  і  \psi  дорівнюють нулю всюди, крім малої сфери (див. рисунок 24.3). З часом радіус P И С
\end{verbatim}

\pandocbounded{\includegraphics[keepaspectratio]{images/_page_179_Picture_3.jpeg}}

. 24.3. Схема поясияет, как начальное возмущение влияет на точку (к, ч,
г).

\begin{verbatim}
Сфера, що оточує точку (x, y, z), зростає зі швидкістю c, так що через  t_{\hat{z}}  секунди ця сфера починає перетинати область початкового збурення, і тому розв'язок u(x, y, z, t) стає ненульовим. У  t_{\hat{z}} < t < t_{\hat{a}}  розв'язок U у (x, y, z) залишається ненульовим, оскільки сфера продовжує перетинати область початкового збурення, але на  t=t_{\hat{a}}  розв'язок знову стає нульовим. Іншими словами, розв'язок, що поширюється з початкової області збурення, має чітко визначену задню кромку. Цей загальний принцип відомий як принцип Гюйгенса для тривимірного простору. Виявляється, що передній край хвилі завжди чітко окреслений, але задній край чітко окреслений лише в просторах розмірності 3, 5, 7.... Ми вже знаємо з формули д'Аламбера, що початкові умови
\end{verbatim}

\pandocbounded{\includegraphics[keepaspectratio]{images/_page_349_Figure_5.jpeg}}

\[
u\left(1+\frac{1}{4}\sin\theta,\ \theta\right)=\cos\theta
\]

приводят, вообще говоря, к нерезкому заднему фронту (поскольку формула
Даламбера подразумевает интегрирование функции \psi(x) в пределах от
x-ct до x+ct).

Покажем теперь, что принцип Гюйгенса неприменим к цилиндрическим волнам.
Такого рода решения описывают волны на воде от точечного источника,
когда задний фронт волны не резкий, а постепенно затухает до нуля.

\section{Двумерное волновое
уравнение}\label{ux434ux432ux443ux43cux435ux440ux43dux43eux435-ux432ux43eux43bux43dux43eux432ux43eux435-ux443ux440ux430ux432ux43dux435ux43dux438ux435}

Чтобы решить двумерную волновую задачу

(УЧП)

\[
u(1+\epsilon\sin\theta, \theta) = \cos\theta, \quad 0 \le \epsilon \le 1/4.
f(x+h) = f(x) + f'(x)h + f''(x)\frac{h^2}{2!} + \dots,
\begin{cases} -\infty < x < \infty, \\ -\infty < y < \infty, \end{cases}
\]

\pandocbounded{\includegraphics[keepaspectratio]{images/_page_350_Figure_2.jpeg}}

\[
u(1+\varepsilon\sin\theta,\ \theta)=u(1,\ \theta)+u_r(1,\ \theta)(\varepsilon\sin\theta)+u_{rr}(1,\ \theta)\frac{(\varepsilon\sin\theta)^2}{21}+\ldots
\]

(НУ)

\[
\begin{cases} u(x, y, 0) = \varphi(x, y), \\ u_t(x, y, 0) = \psi(x, y), \end{cases}
\]

мы рассмотрим решение трехмерной задачи при начальных условиях,
зависящих только от двух переменных x и y.

Сделав это, получим, что трехмерная формула

\[
\Delta u = 0, \quad 0 < r < 1 + \varepsilon \sin \theta,
\]

для решения и будет описывать цилиндрические волны и, следовательно,
даст решение двумерной задачи. Этот метод называется методом спуска.
После всех вычислений 1) (отнюдь не тривиальных) получаем

\[
u(1, \theta) + u_r(1, \theta) (\varepsilon \sin \theta) + u_{rr}(1, \theta) \frac{(\varepsilon \sin \theta)^2}{2!} + \dots = \cos \theta.
\]

Перед нами решение двумерного волнового уравнения, и хотя, вероятнее
всего, его практическое применение подразумевает использование ЭВМ,
интересно проанализировать его с точки зрения принципа Гюйгенса. В этом
решении оба интеграла от начальных возмущений берутся по внутренности
(слово «внутренность» играет здесь решающую роль) круга радиусом сt е
центром в точке (x, y). Другими словами, если проанализировать это
решение так же, как мы анализировали решение в трехмерном случае, мы
увидим, что начальное возмущение распространяется с резко очерченным
передним фронтом и размытым задвим фронтом. Таким образом, принцип
Гюйгенса не выполняется в пространстве двух измерений.

И наконец, предположим, что начальные возмущения \(\phi\) и \(\phi\)
зависят только от одной переменной. В этом случае возникают плоские
волны, и после применения метода спуска получаем

\&lt;sup\textgreater2) См., например, Тихонов А. Н., Самарский А. А.
Уравнения математической физики. --- М.: Наука, 1972 и др. издания. ---
Прим. ред.

хорошо известную формулу Даламбера

\[
(46.8) u = u_0 + \varepsilon u_1 + \varepsilon^2 u_2 + \dots.
\]

Отметим, что и в этом случае применение метода спуска требует проведения
нетривиальных вычислений. В формуле Даламбера начальное смещение ф
приводит к образованию резкого заднего фронта, а начальная скорость ф
его размывает. Другими словами, одномерный случай необычен тем, что
принцип Гюйгенса выполняется для начальных смещений и не выполняется для
начальных скоростей. Можно утверждать, что в целом принцип Гюйгенса в
одномерном случае не выполняется.

\section{ЗАМЕЧАНИЯ}\label{ux437ux430ux43cux435ux447ux430ux43dux438ux44f-10}

Мы не стали подробно излагать метод спуска, поскольку не котели
останавливаться на вопросе о том, как можно понизить размерность
входящих в формулы интегралов. Все эти вычисления можно найти в учебной
литературе, список которой приведен в конце книги. Общая идея
заключается в том, что при решении задач в пространстве невысокой
размерности можно воспользоваться известным решением аналогичной задачи
в пространстве более высокой размерности, а затем упростить его, считая,
что начальные и граничные условия не зависят от части переменных.
Читателю должно быть ясно, что метод спуска можно применять не только к
рассмотренной задаче.

\section{ЗАДАЧИ}\label{ux437ux430ux434ux430ux447ux438-10}

\begin{enumerate}
\def\labelenumi{\arabic{enumi}.}
\tightlist
\item
  Покажите, что в одномерном случае решение задачи
\end{enumerate}

\[
P_{0}
\begin{cases} \Delta u_{0} = 0, & 0 < r < 1 \text{ (внутри круга),} \\ u_{0}(1, \theta) = \cos \theta, & u_{0}(r, \theta) = r \cos \theta, \end{cases}
\begin{cases} u(x, \ 0) = \varphi(x), \\ u_t(x, \ 0) = 0 \end{cases}
P_{1}
\begin{cases} \Delta u_{1} = 0, & 0 < r < 1 \text{ (внутри круга),} \\ u_{1}(1, \theta) = -\sin \theta \frac{\partial u_{0}(1, \theta)}{\partial r} = -\sin \theta \cos \theta \end{cases}
\begin{aligned} & (\forall \mathsf{Y}\Pi) \quad u_{tt} = c^2 u_{xx}, \\ & (\mathsf{H} \mathsf{Y}) \quad \begin{cases} u\left(x, \ 0\right) = 0, \\ u_t\left(x, \ 0\right) = \phi\left(x\right). \end{cases} \end{aligned}
u = u_0 + \frac{1}{4}u_1 + \left(\frac{1}{4}\right)^2 u_1 + \cdots
u_u + \frac{1}{4} u_t
\]

\begin{enumerate}
\def\labelenumi{\arabic{enumi}.}
\setcounter{enumi}{1}
\tightlist
\item
  Результаты задачи 1 используйте для решения следующей задачи:
\end{enumerate}

(УЧП)

\[
u_t = (1+x) u_{xx}, \\ u(x, 0) = q(x), \quad -\infty < x < \infty,
\]

(НУ)

\[
u_1 = (1 + \varepsilon_1) u_{xx}
u = u_1 + \varepsilon u_1 + \varepsilon^2 u_2 + \cdots
\begin{cases} u(x, 0) = x, \\ u_t(x, 0) = 0. \end{cases}
P_{0}
\]

\begin{enumerate}
\def\labelenumi{\arabic{enumi}.}
\setcounter{enumi}{2}
\tightlist
\item
  Проиллюстрируйте и прокомментируйте решение трехмерной задачи
\end{enumerate}

(УЧП)

\[
,  (x,  y, z) \in R^3  ,  u = 0   u_t =
g(x) = e^{-x^2}.
\]

\begin{enumerate}
\def\labelenumi{\arabic{enumi}.}
\setcounter{enumi}{3}
\tightlist
\item
  Решите аналогичную (см. задачу 3) задачу для двумерного уравнения
\end{enumerate}

(УЧП)

(НУ)

Зверніть увагу, що в перших двох задачах коефіцієнти є сталими. Читач
повинен розуміти, що параметр b має бути достатньо малим, а необережний
ряд може відрізнятися

\section{ЗАДАЧИ}\label{ux437ux430ux434ux430ux447ux438-11}

\begin{verbatim}
- 1. Подставьте разложение (46.8) в задачу (46.7) и получите последовательность задач P_0,\ P_1,\ P_2,\ \dots . 2. Покажите, что нелинейную задачу
\end{verbatim}

\begin{enumerate}
\def\labelenumi{\arabic{enumi}.}
\setcounter{enumi}{4}
\tightlist
\item
  Решите аналогичную (см. задачу 3) задачу для одномерных плоских волн
\end{enumerate}

(УЧП)

(НУ)

\begin{itemize}
\tightlist
\item
  \begin{enumerate}
  \def\labelenumi{\arabic{enumi}.}
  \setcounter{enumi}{5}
  \tightlist
  \item
    Какова физическая природа того факта, что принцип Гюйгенса не
    справедлив в двумерном случае?
  \end{enumerate}
\item
  \begin{enumerate}
  \def\labelenumi{\arabic{enumi}.}
  \setcounter{enumi}{6}
  \tightlist
  \item
    Воспользуйтесь формулой Лейбница
  \end{enumerate}
\end{itemize}

\[
u(r, \theta) = u_0(r, \theta) + \frac{1}{4}u_1(r, \theta)
\]

и продифференцируйте интеграл

\[
\Delta u = 0,
\]

по времени t.

\section{КОНЕЧНЫЕ ПРЕОБРАЗОВАНИЯ ФУРЬЕ (СИНУС- И
КОСИНУС-ПРЕОБРАЗОВАНИЯ)}\label{ux43aux43eux43dux435ux447ux43dux44bux435-ux43fux440ux435ux43eux431ux440ux430ux437ux43eux432ux430ux43dux438ux44f-ux444ux443ux440ux44cux435-ux441ux438ux43dux443ux441--ux438-ux43aux43eux441ux438ux43dux443ux441-ux43fux440ux435ux43eux431ux440ux430ux437ux43eux432ux430ux43dux438ux44f}

ЦЕЛЬ ЛЕКЦИИ: Ввести два новых интегральных преобразования (конечные
спиус- и косинус-преобразования):

\[
q_{xx} + \varphi_{uu} = 0
\]

(конечное синус-преобразование),

\[
q_{uu} + q_{vv} = 0
\]

(конечное косинус-преобразование),

\pandocbounded{\includegraphics[keepaspectratio]{images/_page_353_Figure_9.jpeg}}

(обратное синус-преобразование),

\[
w = f(z)
\]

(обратное косинус-преобразование),

и показать, как с помощью этих преобразований решать краевые задачи (в
частности, неоднородные). Ранее мы познакомились с регулярными
преобразованиями Фурье и Лапласа и увидели, как применение этих
преобразований к уравнению с частными производными сводит задачу к
решению обыкновениого дифференциального уравнения. Обычное
преобразование Фурье подразумевает, что преобразуемая переменная
изменяется от \(-\infty\) до \(+\infty\) , и, следовательно, оно
применяется для решения задач в свободном пространстве (без границ). В
этой лекции мы покажем, как решать краевые задачи (с границами),
преобразуя ограниченные переменные (с таким преобразованием мы
встречаемся впервые).

Давайте на время забудем о наших задачах, а займемся самими
преобразованнями. Дадим определение преобразований, приведем формулы
обращения и основные свойства. Для чего нужны преобразования и как ими
пользоваться, мы узнаем несколько позже. Впрочем, если говорить кратко,
то суть интегральных преобразований в разложении входящих в задачу
функций по

частотам, решении семейства задач для каждой частоты и после-

дующем суммировании всех полученных результатов.

Начнем с функции \(f(x)\), определенной на интервале {[}0, L{]}.
Конечные синус- и косинус-преобразования этой функции определяются
формулами

Конечное синус-преобразование
\(S[f] = S_n = \frac{2}{L} \int_{\sigma}^{L} f(x) \sin(n\pi x/L) dx\)
\(n = 1, 2, \dots\)
\(C[f] = C_n = \frac{2}{L} \int_{\sigma}^{L} f(x) \cos(n\pi x/L) dx\)
\(n = 0, 1, \dots\)

Конечное косинус-преобразование

Читатель, наверное, заметил, что эти преобразования не отличаются от
формул для определения коэффициентов Фурье при разложении их в ряд по
синусам или косинусам. Формулы обращения для этих преобразований---это
обычные ряды Фурье по синусам и косинусам.

Обратное синус-преобразование

\pandocbounded{\includegraphics[keepaspectratio]{images/_page_355_Figure_5.jpeg}}

\pandocbounded{\includegraphics[keepaspectratio]{images/_page_355_Figure_7.jpeg}}

Обратное косинус-преобразование

Отметим, что в обратном косинус-преобразовании суммирование начинается с
n=0, а в синус-преобразовании --- с n=1.

\section{Примеры
синус-преобразования}\label{ux43fux440ux438ux43cux435ux440ux44b-ux441ux438ux43dux443ux441-ux43fux440ux435ux43eux431ux440ux430ux437ux43eux432ux430ux43dux438ux44f}

Пусть

\[
u + iv = (x + iy)^2 = x^2 - y^2 + 2ixy
\]

тогда

\begin{verbatim}
\left\{
\end{verbatim}

График функции \(f(x)\) и ее преобразования изображены на рис. 25.1.
Обратное преобразование записывается в виде

Чи знаєте ви, як виглядає графік цієї функції з сегментом {[}0, 1{]}?
Подумай про це. Зверніть увагу, що синусоїдальне перетворення функції
\(f(x)\) визначається лише для додатних цілих чисел n.~Іншими словами,
скінченний синус і косинус перетворюють функції у чисельні
послідовності.

\pandocbounded{\includegraphics[keepaspectratio]{images/_page_185_Figure_3.jpeg}}

РИС. 25.1. Графики функции \(f(x)\) = 1 и ее преобразования.

\section{Свойства
преобразований}\label{ux441ux432ux43eux439ux441ux442ux432ux430-ux43fux440ux435ux43eux431ux440ux430ux437ux43eux432ux430ux43dux438ux439}

Прежде чем приступить к решению задач, мы должны получить некоторые
свойства этих преобразований.

\begin{verbatim}
Якщо you(x, t) — функція  \partial $\theta$ yx  змінних, і ми виконуємо перетворення змінної x, то
\end{verbatim}

(действительная форма отображения \(w=z^2\) ).

Эта форма понадобится нам позже. В этом примере она интересна тем, что
показывает, как гиперболы в плоскости z переходят в координатные линии u
= const и v = const в плоскости w.

Теперь мы перейдем к изучению одного частного вида функций комплексного
переменного, которые осуществляют конформные отображения.

\section{Определение конформного
отображения}\label{ux43eux43fux440ux435ux434ux435ux43bux435ux43dux438ux435-ux43aux43eux43dux444ux43eux440ux43cux43dux43eux433ux43e-ux43eux442ux43eux431ux440ux430ux436ux435ux43dux438ux44f-1}

Отображение w=\(f(z)\) комплексиой плоскости z на комплексную плоскость
w называется конформным в точке \(z_{\bullet}\) плоскости z, если
производная \(f'(z_{\bullet}) \neq 0\) . Отображение \(f(z)\) называется
конформным в области D, если \(f'(z) \neq 0\) в каждой точке области D.

Например, отображение \(f(z) = z^2\) конформно всюду, за исключением
точки z = 0, поскольку \(z'(z) = 2z \neq 0\) для всех \(z \neq 0\) . C
другой стороны, отображение \(e^z\) конформно во всей плоскости z, так
как всюду \(f'(z) = e^z \neq 0\) . Какая же польза от конформных
отображений? Ответ состоит в том, что, решая уравнение Лапласа
\(\phi_{xx} + \phi_{yy} = 0\) внутри какой-то области плоскости
переменных x и y, мы можем рассматривать эту плоскость как плоскость
комплексного переменного z. Рассмотрим комплексное отображение w =
\(f(z)\) плоскости z на плоскость w, в которой введены координаты (u,
v). При отображении w = \(f(z)\) уравнение Лапласа
\(\varphi_{xx} + \varphi_{yy} = 0\) преобразуется в некоторое новое
уравнение с частными производными, зависящее от новых координат и и v.
Замечательным является следующее: если отображение w = \(f(z)\) является
конформным в области, где задано уравнение \(\phi_{xy} + \phi_{uy} = 0\)
, то новое уравнение также будет уравнением Лапласа в координатах и и v.
To есть уравнение \(\phi_{xx} + \phi_{uu} = 0\) перейдет снова в
уравнение \(\varphi_{nn} + \varphi_{nn} = 0\) . Так мы приходим к идее
найти такое конформное отображение, которое переводит область со сложной
границей в область с простой границей (напомним, что отображение
\(w=z^*\) переводит границу первого квадранта плоскости z в
действительную ось плоскости w).

Приведем несколько примеров, в которых будет показано, как можно решать
уравнения с частными производными с помощью конформных отображений.

Уравнение Лапласа в верхней полуплоскости.

Предположим, что нам необходимо решить следующую задачу Дирихле в
верхней полувлоскости (рис. 47.3):

(47.1)

(Зверніть увагу, що перетворення було виконано на змінній x, і отримана
послідовність залежить лише від часу t.)

А как быть с производными? Ниже приводятся несколько полезных формул

\begin{verbatim}
\begin{aligned} (Y & \Pi) & \varphi_{xx} + \varphi_{yy} = 0, & -\infty < x < \infty, & 0 < y < \infty, \\ (47.1) & $\varphi$(x, 0) = \begin{cases} 0, & |x| > 1, \\ 1, & |x| \le 1. \end{cases}
\end{verbatim}

\[
w = \ln\left\{\frac{z-1}{z+1}\right\}
\begin{split} S\left[u_{t}\right] &= \frac{dS\left[u\right]}{dt}, \quad S\left[u_{tt}\right] = \frac{d^{2}S\left[u\right]}{dt^{2}}, \\ S\left[u_{xx}\right] &= -\left[n\pi/L\right]^{2}S\left[u\right] + \frac{2n\pi}{L^{2}}\left[u\left(0,\ t\right) + (-1)^{n+2}u\left(L,\ t\right)\right], \\ C\left[u_{xx}\right] &= -\left[n\pi/L\right]^{2}C\left[u\right] - \frac{2}{L}\left[u_{x}\left(0,\ t\right) + (-1)^{n+1}u_{x}\left(L,\ t\right)\right]. \end{split}
\]

\pandocbounded{\includegraphics[keepaspectratio]{images/_page_357_Figure_11.jpeg}}

\pandocbounded{\includegraphics[keepaspectratio]{images/_page_358_Figure_2.jpeg}}

\[
(\text{УЧП}) \quad u_{tt} = u_{xx} + \sin(\pi x), \quad 0 < x < 1, \quad 0 < t < \infty,
\begin{cases} (\text{УЧП}) & \varphi_{uu} + \varphi_{vv} = 0, \quad -\infty < u < \infty, \quad 0 < v < \pi, \\ (\text{47.2}) & \begin{cases} \varphi(u, 0) = 0, \\ \varphi(u, \pi) = 1. \end{cases} \end{cases}
(\text{ГУ}) \quad
w = u + iv = \ln\left[\frac{z-1}{z+1}\right] =
\begin{cases} u(0, t) = 0, \\ u(1, t) = 0, \end{cases}
= \ln\left[\frac{z-1}{z+1}\right] + i \arg\left[\frac{z-1}{z+1}\right]
v = \arg\left[\frac{z-1}{z+1}\right] = \arg\left[\frac{x+iy-1}{x+1+iy}\right] =
= \arg\left[\frac{x^2+y^2-1+i2y}{(x+1)^2+y^2}\right] =
= \arg\left[\frac{2y}{x^2+y^2-1}\right]
\begin{cases} u(x, 0) = 1, \\ u_t(x, 0) = 0, \end{cases}
\varphi(x, y) = \frac{1}{\pi} \arctan\left[\frac{2y}{x^2 + y^2 - 1}\right].
x^{2} + y^{2} = 1,
\]

Решение задачи выполним за три следующих шага:

ШАГ І (Определение подходящего преобразования.)

Поскольку переменная x меняется от 0 до 1, мы воспользуемся конечным
преобразованием, а именно синус-преобразованием; позже станет ясно
почему. Мы могли бы решить эту задачу с помощью преобразования Лапласа
по переменной t (по трудоемкости этот метод не отличается от конечного
синус-преобразования).

ШАГ 2 (Выполнение преобразования).

Для удобства будем использовать обозначение \(S_n(t) = S[u]\) . Применим
синус-преобразование к исходному уравнению

\[
(x-1)^{2} + y^{2} = 9,
\]

и воспользуемся тождеством для синус-преобразования, получим

\[
\begin{aligned} & (\text{УЧП}) & \phi_{xx} + \phi_{yy} = 0, & \text{внутри } D, \\ & (\text{ГУ}) & \begin{cases} \varphi(x, y) = 1 & \text{на } x^2 + y^2 = 1, \\ \varphi(x, y) = 2 & \text{на } (x - 1)^2 + y^2 = 9. \end{cases} \end{aligned}
\begin{split} \frac{d^2S_n\left(t\right)}{dt^2} &= -\left(\pi n\right)^2S_n\left(t\right) + 2n\pi\left[u\left(0,\ t\right) + \left(-1\right)^{n+1}u\left(1,\ t\right)\right] + D_n\left(t\right) = \\ &= -\left(n\pi\right)^2S_n\left(t\right) + D_n\left(t\right), \end{split}
\]

\pandocbounded{\includegraphics[keepaspectratio]{images/_page_359_Figure_10.jpeg}}

\[
w = 2t \left[ \frac{z-s}{z-t} \right],
\]

где

\[
u = \gamma [(x-s)(x-t)-\gamma^2],
v = \gamma [y(x-t)],
\gamma = 2t/[(x-t)^2 + y^2].
\]

Если теперь преобразовать начальные условия краевой задачи, то получим
начальные условия для обыкновенного дифференциального уравнения

\[
\begin{aligned} \varphi_{uu} + \varphi_{vv} &= 0, & \text{в кольче,} \\ (\Gamma Y) & & & & & & & \\ (\Gamma Y) & & & & & & \\ (\Gamma Y) & & & & & \\ & & & & & \\ (\Gamma Y) & & & & & \\ & & & & \\ & & & & \\ & & & & \\ & & & & \\ & & & \\ & & & \\ & & & \\ & & & \\ & & & \\ & & & \\ & & \\ & & \\ & & \\ & & \\ & & \\ & & \\ & & \\ & & \\ & & \\ & & \\ & & \\ & & \\ & & \\ & & \\ & & \\ & & \\ & & \\ & & \\ & & \\ & & \\ & & \\ & & \\ & & \\ & & \\ & & \\ & & \\ & & \\ & & \\ & & \\ & & \\ & & \\ & & \\ & & \\ & & \\ & & \\ & & \\ & & \\ & & \\ & & \\ & & \\ & & \\ & & \\ & & \\ & & \\ & & \\ & & \\ & & \\ & & \\ & & \\ & & \\ & & \\ & & \\ & & \\ & & \\ & & \\ & & \\ & & \\ & & \\ & & \\ & & \\ & & \\ & & \\ & & \\ & & \\ & & \\ & & \\ & & \\ & & \\ & & \\ & & \\ & & \\ & & \\ & & \\ & & \\ & & \\ & & \\ & & \\ & & \\ & & \\ & & \\ & & \\ & & \\ & & \\ & & \\ & & \\ & & \\ & & \\ & & \\ & & \\ & & \\ & & \\ & & \\ & & \\ & & \\ & & \\ & & \\ & & \\ & & \\ & & \\ & & \\ & & \\ & & \\ & & \\ & & \\ & & \\ & & \\ & & \\ & & \\ & & \\ & & \\ & & \\ & & \\ & & \\ & & \\ & & \\ & & \\ & & \\ & & \\ & & \\ & & \\ & & \\ & & \\ & & \\ & & \\ & & \\ & & \\ & & \\ & & \\ & & \\ & & \\ & & \\ & & \\ & & \\ & & \\ & & \\ & & \\ & & \\ & & \\ & & \\ & & \\ & & \\ & & \\ & & \\ & & \\ & & \\ & & \\ & & \\ & & \\ & & \\ & & \\ & & \\ & & \\ & & \\ & & \\ & & \\ & & \\ & & \\ & & \\ & & \\ & & \\ & & \\ & & \\ & & \\ & & \\ & & \\ & & \\ & & \\ & & \\ & & \\ & & \\ & & \\ & & \\ & & \\ & & \\ & & \\ & & \\ & & \\ & & \\ & & \\ & & \\ & & \\ & & \\ & & \\ & & \\ & & \\ & & \\ & & \\ & & \\ & & \\ & & \\ & & \\ & & \\ & & \\ & & \\ & & \\ & & \\ & & \\ & & \\ & & \\ & & \\ & & \\ & & \\ & & \\ & & \\ & & \\ & & \\ & & \\ & & \\ & & \\ & & \\ & & \\ & & \\ & & \\ & & \\ & & \\ & & \\ & & \\ & & \\ & & \\ & & \\ & & \\ & & \\ & & \\ & & \\ & & \\ & & \\ & & \\ & & \\ & & \\ & & \\ & & \\ & & \\ & & \\ & & \\ & & \\ & & \\ & & \\ & & \\ & & \\ & & \\ & & \\ & & \\ & & \\ & & \\ & & \\ & & \\ & & \\ & & \\ & & \\ & & \\ & & \\ & & \\ & & \\ & & \\ & & \\ & & \\ & & \\ & & \\ & & \\ & & \\ & & \\ & & \\ & & \\ & & \\ & & \\ & & \\ & & \\ & & \\ & & \\ & & \\ & & \\ & & \\ & & \\ & & \\ & & \\ & & \\ & & \\ & & \\ & & \\ & & \\ & & \\ & & \\ & & \\ & & \\ & & \\ & & \\ & & \\ & & \\ & & \\ & & \\ & & \\ & & \\ & & \\ & & \\ & & \\ & & \\ & & \\ & & \\ & & \\ & & \\ & & \\ & & \\ & & \end{aligned}
\begin{cases} 4/n\pi, & n = 1, 3, \dots, \\ 0, & n = 2, 4, \dots, \end{cases}
\varphi(r) = a \ln r + b, \quad r = u^2 + v^2.
\varphi(u, v) = 0.57 \ln(u^2 + v^2) + 1
\varphi(x, y) = 0.57 \ln(u^2 + v^2) + 1,
\]

Итак, решим теперь новое семейство задач Кощи:

(ОДУ)

\[
w = \ln \left[ \frac{z-1}{z+1} \right].
z = re^{i\theta}
\begin{cases} 1, & n = 1, \\ 0, & n = 2, 3, \dots, \end{cases}
\]

\begin{verbatim}
\begin{aligned} (\text{Y} \forall \Pi) & \qquad \phi_{\text{vx}} + \phi_{yy} = 0, \quad 0 < x < \infty, \quad 0 < y < \infty, \\ \left\{ \begin{array}{l} $\phi$(x, \ 0) = \begin{cases} 1, & 0 < x < 1, \\ 0, & 1 \leqslant x < \infty, \end{cases} \\ $\phi$(0, \ y) = \begin{cases} 1, & 0 < y < 1, \\ 0, & 1 \leqslant y < \infty. \end{cases} \end{aligned}
\end{verbatim}

\[
\right.
\]

\pandocbounded{\includegraphics[keepaspectratio]{images/_page_362_Figure_1.jpeg}}

\begin{verbatim}
\begin{cases} S_n(0) = \begin{cases} 4/n$\pi$, & n = 1, 3, \dots, \\ 0, & n = 2, 4, \dots, \end{cases}
\end{verbatim}

\pandocbounded{\includegraphics[keepaspectratio]{images/_page_362_Figure_2.jpeg}}

\pandocbounded{\includegraphics[keepaspectratio]{images/_page_362_Figure_5.jpeg}}

\pandocbounded{\includegraphics[keepaspectratio]{images/_page_362_Figure_7.jpeg}}

\[
\begin{cases} \frac{dS_n(0)}{dt} = 0. \end{cases}
H(x-a) =
\begin{cases} 0, x < a, \\ 1, x \geqslant a, \end{cases}
\]

В результате получаем

,

где

\[
H(x-a) =
\begin{cases} 0, & x < a, \\ 1, & x \geqslant a, \end{cases}
\begin{cases} 0, & n = 2, 4, \dots, \\ (4/n\pi)\cos(n\pi t), & n = 3, 5, 7, \dots \end{cases}
f(x) = \mathcal{F}^{-1}[F] = \frac{1}{\sqrt{2\pi}} \int_{-\infty}^{+\infty} F(\omega) e^{-i\omega x} d\omega \qquad F(\omega) = \mathcal{F}[f] = \frac{1}{\sqrt{2\pi}} \int_{-\infty}^{+\infty} f(x) e^{-i\omega x} dx
\]

Следовательно, решение нашей задачи можно представить в виде

\[
\begin{cases} 1, |x| < a \\ 0, |x| > a \end{cases}
\]

\section{ЗАМЕЧ
АНИЯ}\label{ux437ux430ux43cux435ux447-ux430ux43dux438ux44f}

\begin{enumerate}
\def\labelenumi{\arabic{enumi}.}
\item
  Для решения задачи методом конечного синус- или косинуспреобразования
  краевые условия при x=0 и x=L должны иметь вид
\item
\end{enumerate}

\[
\begin{cases} 1, |x| < 1 \\ 0, |x| > 1 \end{cases}
\begin{cases} u(0, t) = f(t) \\ u(L, t) = g(t) \end{cases}
\delta(x-a)
\]

\(f(x)\) * \(g(x)\)

\begin{verbatim}
\begin{array}{ll} (\mathrm{Y} \mathrm{U} \Pi) & u_t = u_{xx}, \quad 0 < x < 1, \quad 0 < t < \infty, \\ (\mathrm{\Gamma} \mathrm{Y} & \left\{ \begin{array}{ll} u_x \left( 0, \ t \right) = 0 \\ u_x \left( 1, \ t \right) = 0, \end{array} \end{array}
\end{verbatim}

\[
(1+x^2)^{-1}
xe^{-a+x^4}, a > 0
\]

\begin{enumerate}
\def\labelenumi{\arabic{enumi}.}
\setcounter{enumi}{1}
\tightlist
\item
  Решите в общем виде задачу
\end{enumerate}

\(H(x+a)\) - \(H(x-a)\)

\[
\frac{a}{x^2+a^2}
\]

\begin{verbatim}
\right. & 0 < t < \infty, \\ (\text{НУ}) & u(x, 0) = 0, \quad 0 \leqslant x \leqslant 1. \end{array}
\end{verbatim}

\[
\frac{2ax}{(x^2+a^2)^2}
\begin{cases} \cos(ax), |x| > \pi/2a \\ 0, |x| > \pi/2a \end{cases}
\]

\section{ПРИНЦИП СУПЕРПОЗИЦИИ--- ОСНОВА ТЕОРИИ ЛИНЕЙНЫХ
СИСТЕМ}\label{ux43fux440ux438ux43dux446ux438ux43f-ux441ux443ux43fux435ux440ux43fux43eux437ux438ux446ux438ux438-ux43eux441ux43dux43eux432ux430-ux442ux435ux43eux440ux438ux438-ux43bux438ux43dux435ux439ux43dux44bux445-ux441ux438ux441ux442ux435ux43c}

ЦЕЛЬ ЛЕКЦИИ: Познакомить читателя с принципом суперпозиции и показать,
как с его помощью исходную задачу можно разложить на подзадачи, затем
решить сразу все подзадачи и, сложив полученные результаты, найти
решение исходной задачи. Мы хотим также показать, что два основных
метода решения линейных уравнений---метод разделения переменных и метод
интегральных преобразований---используют принцип суперпозиции.

Предположим, что инженеру нужно найти отклик линейной системы и на
входное воздействие f.~Общий подход к решению задачи таков:

\begin{itemize}
\tightlist
\item
  \begin{enumerate}
  \def\labelenumi{\arabic{enumi})}
  \tightlist
  \item
    разложить f на элементарные воздействия \(f = \sum f_k\) ;
  \end{enumerate}
\item
  \begin{enumerate}
  \def\labelenumi{\arabic{enumi})}
  \setcounter{enumi}{1}
  \tightlist
  \item
    найти отклик системы \(u_k\) на воздействие \(f_k\) ;
  \end{enumerate}
\end{itemize}

\begin{enumerate}
\def\labelenumi{\arabic{enumi})}
\setcounter{enumi}{2}
\tightlist
\item
  сложить отклики \(u = \sum_{i=1}^{n} u_{i}\) .
\end{enumerate}

Если система линейная, то полученная сумма и является откликом системы
на входное воздействие f.~В этом состоит принцип суперпоэиции (см. рнс.
26.1).

\pandocbounded{\includegraphics[keepaspectratio]{images/_page_189_Picture_10.jpeg}}

РИС, 26.1. Основная идея принципа суперпоэнции; L---дифференциальный
оператор; f---вход; u---выход.

Эту основную идею можно использовать при решении сложных смешанных
задач, разбивая их на простые подзадачи, ре-

шая каждую из них и затем суммируя полученные результаты (конечно,
дифференциальное уравнение и граничные условия должны быть линейны).

\section{Разложение смешанной задачи на две более
простые}\label{ux440ux430ux437ux43bux43eux436ux435ux43dux438ux435-ux441ux43cux435ux448ux430ux43dux43dux43eux439-ux437ux430ux434ux430ux447ux438-ux43dux430-ux434ux432ux435-ux431ux43eux43bux435ux435-ux43fux440ux43eux441ux442ux44bux435}

Предположим, нам необходимо решить линейную задачу (обозначим ее P)

(УЧП)

\begin{enumerate}
\def\labelenumi{\arabic{enumi}.}
\setcounter{enumi}{17}
\tightlist
\item
\end{enumerate}

(Р) (ГУ)

\[
\begin{cases} 1-|x|, |x| < 1 \\ 0, |x| > 1 \end{cases}
\begin{cases} u(0, t) = 0, \\ u(1, t) = 0, \\ u(x, 0) = \sin(2\pi x), \quad 0 \le x \le 1. \end{cases}
\cos(ax)
\frac{2}{\pi} \frac{\sin a\omega}{\omega}
\]

Перед нами неоднородное уравнение теплопроводности, так что метод
разделения переменных неприменим. Можпо, конечно, воспользоваться
конечным синус-преобразованием по переменной x или преобразованием
Лапласа по переменной t, но вместо этого мы рассмотрим две подзадачи

\[
\frac{2}{\pi} \frac{\sin a\omega}{(\omega^2+a^2)^2}
\frac{\pi}{2} e^{-a+\omega}
\begin{cases} u(0, t) = 0, \\ u(1, t) = 0, \\ u(1, t) = 0, \end{cases}
\frac{\pi}{2} e^{-a+\omega}
\frac{\pi}{2} e^{-a+\omega}
\frac{\pi}{2} e^{-a+\omega}
\frac{\pi}{2} e^{-a+\omega}
\]

Каждая из этих задач легко решается и в то же время ясно, что
\mathit{suma} решений задач \(P_1\) и \(P_2\) является решением исходной
вадачи, т. е.

\[
\frac{2ax}{(x^2+a^2)^2}
\]

Решение задачи \(P_{i}\) Розв'язання задачі \(P_i\) Розв'язання задачі
\(P_i\) Можна показати, що в загальному випадку розв'язок задачі

\[
\frac{\cos(ax), |x| > \pi/2a}{0, |x| > \pi/2a}
\frac{1-|x|, |x| < 1}{0, |x| > 1}
\begin{cases} u_t = u_{xx} + f(x, t), \\ u(0, t) = 0, \\ u(1, t) = 0, \\ u(x, 0) = \varphi(x) \end{cases}
\cos(ax)
\cos(ax)
\frac{\pi}{2} [\delta(\omega+a) + \delta(\omega-a)]
\begin{cases} u_t = u_{xx} + f(x, t), \\ u(0, t) = 0, \\ u(1, t) = 0, \\ u(x, 0) = 0 \end{cases}
f(x) = \int_{0}^{\infty} F(\omega) \sin(\omega x) d\omega
F(\omega) = \frac{?}{\pi} \int_{0}^{\infty} f(x) \sin(\omega x) dx
0 < x < \infty
\begin{cases} u_t = u_{xx}, \\ u(0, t) = 0, \\ u(1, t) = 0, \\ u(x, 0) = \varphi(x). \end{cases}
0 < \omega < x
\]

f'\,'(0)

\[
-\omega^2 F(\omega) + \frac{2}{\pi} \omega f(0)
\frac{1}{\alpha} F\left(\frac{\omega}{\alpha}\right)
\]

проследим, как оно решается методом конечного синус-преобразования. Вот
что мы должны сделать: разложить еходное воздействие \(f(x, t)\) на
компоненты, найти отклик \(U_n\) на каждую компоненту, а затем сложить
эти отклики. Математически это может быть и не совсем очевидно, но
давайте внимательно проследим за ходом решения. Разложим входящие в
уравнение с частными производными члены в ряды Фурье по синусам:

\[
\frac{2\omega}{\pi (\alpha^2 + \omega^2)}
x^{-1/2}
[2/\pi\omega]^{1/2}
\]

where

\[
\frac{2}{\pi\omega} [1-\cos(\omega a)]
\]

1

\[
e^{-1}a + \omega
\]

Зверніть увагу, по-перше, що коефіцієнти \(A_n\) , \(B_n\) \(F_n\)
дійсно є функціями часу t, оскільки ми застосували синусоїдальне
перетворення до функції двох змінних x, t, u, а по-друге, ми розклали
вхідний ефект \(f(x,t)\) на прості компоненти \(F_n\)(t) . Тепер
знайдіть відповіді \(U_n\)(t) для кожного компонента \(F_n\)(t) , потім
додайте всі відповіді \(U_n\)(t) і отримайте розв'язок \(u(x,t)\).

Щоб знайти відповіді \(U_n\)(t) , необхідно робити невеликі розрахунки з
коефіцієнтами \(A_n\)(t) і \(B_n\)(t) , щоб у вирази інтегрування була
включена лише функція u, а не її похідні \(u_i\) і \(u_{xx}\) .
Інтегруючи на частини, ми отримуємо

\[
\frac{1}{2}e^{-\alpha}\sin\omega
\frac{1-e^{-a\omega}}{\omega}
\begin{split} A_n(t) &= 2\int_0^1 u_t \sin(n\pi x) \, dx = \frac{d}{dt} \bigg[ 2\int_0^1 u(x, t) \sin(n\pi x) \, dx \bigg] = \frac{dU_n(t)}{dt} \,, \\ B_n(t) &= 2\int_0^1 u_{xx} \sin(n\pi x) \, dx = -(n\pi)^2 \, U_n(t) \,+ \\ &\quad + 2n\pi \, \big[ u(0, t) + (-1)^{n+1} \, u(1, t) \big], \end{split}
\frac{2}{\pi}F''(\omega)
f(x) = \int_{0}^{\infty} F(\omega) \cos(\omega x) d\omega \qquad \qquad F(x) = \frac{2}{\pi} \int_{0}^{\infty} f(x) \cos(\omega x) dx
0 < x < \infty \qquad \qquad 0 < \omega < \infty
\begin{cases} u(0, t) = 0, \\ u(1, t) = 0 \end{cases}
\]

j'\,'(x)

\[
S_n = \frac{2}{\pi} \int_0^{\pi} f(x) \sin(nx) dx.
\]

в последнее соотношение, получаем

\[
y = \frac{\pi}{b - a} (x - a).
\]

и, следовательно, (26.1) принимает вид

\[
f(x) = \sum_{n=1}^{\infty} S_n \sin(nx).
\]

Поскольку это разложение является тождеством по x, его коэффициенты
должны быть равны нулю, \tau  . е.

\[
f(x) = \sum_{n=1}^{\infty} S_n \sin(nx)
\]

.

Таким образом, мы получили связь между входом и выходом системы, между
коэффициентами \(F_n\) и \(U_n\) . Прежде чем решать уравнения для
\(U_n\)(t) , обратимся к начальному условию

\[
S_n = \frac{2}{\pi} \int_{0}^{\pi} f(x) \sin(nx) dx
\]

Если разложить \(u(x, 0)\) в ряд Фурье по синусам и приравнять его нулю:

\[
0 \le x \le \pi
\]

то определяются начальные условия для функции U\_\(n\left(t\right)\) ,
т. е.

\[
n = 1, 2, ...
\]

, n = 1, 2, \dots  Мы получили разложение нашей смешанной задачи на ряд
простых задач об отклике системы на элементарные входные воздействия:

(OДУ)

f'\,'(x)

(HУ) \(U_n\)(0) = 0, n = 1, 2, \ldots{} Каждое из этих уравнений легко
решается методом интегрирующего множителя (или с помощью преобразования
Лапласа). В любом случае получаем

\[
\frac{\pi}{2} e^{ax}
\]

Таким образом, мы получили отклики \(U_n\)(t) на простейшие входные
воздействия \(F_n\)(t) . Для получения решения исходной задачи осталось
совершить последний шаг---сложить все отклики

\[
C_n = \frac{2}{\pi} \int_0^{\pi} f(x) \cos(nx) dx.
\]

Читатель, должно быть, заметил, что каждый отклик \(U_n\)(t) умножается
на весовой множитель \sin(n\pi x) . Конечно, в разложении функции
\(f(x, t)\) на компоненты \(F_n\)(t) присутствуют те же весовые
множители.

\section{ЗАМЕЧАНИЕ}\label{ux437ux430ux43cux435ux447ux430ux43dux438ux435-1}

Використання скінченного синусоїдального перетворення дозволяє записати
розв'язок у вигляді ряду, але для більшості інших інтегральних
перетворень розв'язки представлені як інтеграли (неперервний розклад).
Читача можна попросити дати інтерпретацію проблеми

\[
y = \frac{\pi}{b-a}(x-a).
f(x) = \frac{C_0}{2} + \sum_{n=1}^{\infty} C_n \cos(nx).
f(x) = \frac{C_0}{2} + \sum_{n=1}^{\infty} C_n \cos(nx) \qquad C_n = \frac{2}{\pi} \int_{0}^{\pi} f(x) \cos nx \, dx
0 \le x \le \pi \qquad n = 0, 1, 2, \dots
\]

\begin{verbatim}
\begin{array}{ll} (\mathrm{Y} + \Pi) & u_{tt} = u_{xx} + \sin{(3$\pi$ x)}, & 0 < x < 1, & 0 < t < \infty, \\ (\mathrm{\Gamma} \mathrm{Y}) & \left\{ \begin{array}{ll} u\left(0, \ t\right) = 0, \\ u\left(1, \ t\right) = 0, \\ \end{array}
\end{verbatim}

f'\,'(x)

\[
-n^{2}C_{n} - \frac{2}{\pi} \{f'(0) - (-1)^{n} f'(n)\}
\begin{array}{ll} u\left(x, \ 0\right) = \sin{(\pi x)}, \\ u_{t}\left(x, \ 0\right) = 0, \end{array}
\]

Для каждой из подзадач выберите метод решения по своему BKYCY,

\begin{enumerate}
\def\labelenumi{\arabic{enumi}.}
\setcounter{enumi}{2}
\tightlist
\item
  Предположим, что \(u_1\) и \(u_2\) --- решения следующих уравнений:
\end{enumerate}

\begin{enumerate}
\def\labelenumi{\alph{enumi})}
\tightlist
\item
\end{enumerate}

;\\
b) \(u_t = u_{xx} + e^t\) ;

\begin{enumerate}
\def\labelenumi{\Alph{enumi})}
\setcounter{enumi}{1}
\tightlist
\item
  \begin{enumerate}
  \def\labelenumii{\arabic{enumii}.}
  \setcounter{enumii}{4}
  \tightlist
  \item
    \cos(mx), m=1, 2, \dots
  \end{enumerate}
\end{enumerate}

;\\
F) \(u_t = u_{xx} + u^2\) .

\mathbf{\Pi} для яких рівнянь функція \(u_1\) + \(u_2\) також буде
розв'язком? Які висновки ви можете зробити зі своїх відповідей?

\begin{enumerate}
\def\labelenumi{\arabic{enumi}.}
\setcounter{enumi}{3}
\tightlist
\item
  Найдите четыре смещанные задачи, сумма решений которых дает решение
  следующей задачи:
\end{enumerate}

\[
\begin{aligned} u_t &= u_{xx} + f(x, t), & 0 < x < 1, & 0 < t < \infty, \\ & \begin{cases} u(0, t) = g_1(t), \\ u(1, t) = g_2(t), \end{cases} & 0 < t < \infty, \\ & u(x, 0) = \varphi(x), & 0 \le x \le 1. \end{aligned}
\]

\begin{enumerate}
\def\labelenumi{\arabic{enumi}.}
\setcounter{enumi}{4}
\tightlist
\item
  Решите задачу Коши
\end{enumerate}

(OAY)

x

(HY) \(U_n\)(0) = 0. Можете ли вы проверить полученное решение?

\begin{enumerate}
\def\labelenumi{\arabic{enumi}.}
\setcounter{enumi}{5}
\item
  Предположим, что функции \(u_1\) и \(u_2\) удовлетворяют линейным
  однородным граничным условиям

  \textbackslash begin\{cases\} u\_x(1, t) + h\_2 u(1, t) = 0. Будет ли
  функция u\_i + u\_i удовлетворять этим условиям? , .
\end{enumerate}

Лекция 27

\section{УРАВНЕНИЯ ПЕРВОГО ПОРЯДКА (МЕТОД
ХАРАКТЕРИСТИК)}\label{ux443ux440ux430ux432ux43dux435ux43dux438ux44f-ux43fux435ux440ux432ux43eux433ux43e-ux43fux43eux440ux44fux434ux43aux430-ux43cux435ux442ux43eux434-ux445ux430ux440ux430ux43aux442ux435ux440ux438ux441ux442ux438ux43a}

ЦЕЛЬ ЛЕКЦИИ: Ввести понятие уравнения с частными производными первого
порядка (до сих пор мы рассматривали уравнения второго порядка) и
познакомиться с важным методом решения задач с начальными
условиями---методом характеристик. Задача, которой мы теперь займемся,
имеет вид

(УЧП)

\_ \textbar, -\infty \textless{} x \textless{} \infty  , 0 \textless{} t
\textless{} \infty  ,

\[
(HV)   u(x, 0) = \varphi(x), \quad -\infty < x < \infty.
\]

Сразу же отметим, что впервые мы будем решать задачу с переменными
коэффициентами. Оказывается, если от координат (x, t) перейти к новым,
характеристическим координатам (s, \tau) , то наше УЧП превратится в
обыкновенное дифференциальное уравнение. Затем можно решить обыкновенное
дифференциальное уравнение, т. е. найти u(s, \tau) , а на последнем шаге
выразить s и t через x и t и получить \(u(x, t)\).

Напомним читателю, что, когда мы решали уравнение дифф

\[
\Delta u = u_{xx} + u_{yy},   -\infty < x < \infty  ,   0 < t < \infty
\]

константа \alpha\^{}2 играла роль коэффициента диффузии, а v---была
скоростью среды. Очевидно, что если \alpha=0 (т. е. диффузии нет), то
решение будет смещаться как целое вдоль оси x со скоростью v (у нас
осталась только конвекция). Другими словами, если начальное условие
\(u(x,0)\)=\varphi(x) , то соответствующее решение уравнения

\[
\Delta u = u_{rr} + \frac{1}{r}u_r + \frac{1}{r^2}u_{\theta\theta}
\]

будет иметь вид \(u(x, t)\) = \varphi(x-vt) .

Это дает нам основания думать, что решение уравнения первого порядка

\[
\Delta u = u_{xx} + u_{yy} + u_{zz}
\]

переносится вдоль потока со скоростью

\[
\Delta u = u_{rr} + \frac{1}{r}u_r + \frac{1}{r^2}u_{\theta\theta} + u_{zz}.
\]

Конечно, если a и b---константы, то решение представляет собой бегущую с
постоянной скоростью волну. Если же \(a(x, t)\) и \(b(x, t)\)

\pandocbounded{\includegraphics[keepaspectratio]{images/_page_196_Figure_10.jpeg}}

РИС. 27.1. Начальное значение в точке x оказывает влияние на решение
только вдоль характеристики в плоскости переменных (x,t). Параметр s
изменяется вдоль характеристик от нуля до бесконечности, начиная с линии
начального условия; \tau   постоянна вдоль каждой характеристики;
\tau   изменяется вдоль линии начального условия.

зависят от x и t, то скорость потока изменяется как вдоль потока, так и
во времени (читатель сможет увидеть, что начальная кривая сильно
искажается). Здесь так много аналогий с конвекцией!

Вернемся к нашей основной задаче:

(УЧП)

\[
\Delta u = u_{rr} + \frac{2}{r}u_r + \frac{1}{r^2}u_{\theta\theta} + \frac{\cot\theta}{r^2}
\]

(НУ) \(u(x, 0)\) = \varphi(x), -\infty \textless{} x \textless{} \infty.

Решение этого линейного уравнения первого порядка основывается на
следующем физическом факте: начальное условие в некоторой точке x
переносится в tx-плоскости вдоль линии, которая

называется характеристи-кой (см. рис. 27.1)

Этим наше уразнение отличается от других уравнений (таких, как уравнение
теплопроводности \(u_t = u_{xx} )\),для которых начальное значение в
точке x оказывает влияние на решение во всех точках пространства и во
все моменты времени. Если читатель помнит, то начальное сме- Р И С

\pandocbounded{\includegraphics[keepaspectratio]{images/_page_197_Figure_5.jpeg}}

. 27.2. Начальное возмущение \(u\left(x,0\right)\) в точке порождает две
волны. Начальное возмущение струны согласно волновому уравнению
распространяется вдоль двух характеристик.

щение скрипичной струны в точке x оказывает влияние на решение вдоль
\partial syx линий в tx-плоскости (соответствующих двум бегущим волнам)
(см. рис. 27.2).

Естественно, напрашивается идея ввести две новые координаты s и
\tau   (вместо x и t) так, чтобы

\begin{itemize}
\tightlist
\item
  s изменялась вдоль характеристических кривых,
\item
  т изменялась вдоль начальной кривой (желательно вдоль линии t=0).
\end{itemize}

Рассмотрим сначала новую координату s. Если s выбрать так, что она
удовлетворяет приведенному выше условию, то уравнение

\[
\Delta u = u_{rr} + \frac{2}{r}u_r + \frac{1}{r^2}u_{\theta\theta} + \frac{\cot\theta}{r^2}u_0 + \frac{1}{r^2\sin\theta}u_{\phi\phi}
\]

превратится в обыкновенное дифференциальное уравнение К о н е

\pandocbounded{\includegraphics[keepaspectratio]{images/_page_370_Figure_1.jpeg}}

чно, остается открытым вопрос, как найти эти характеристики. Ответ
прост: мы выберем характеристики \{{[}\(x(s)\), \(t(s)\){]}: 0
\textless{} s \textless{} \infty\} так, чтобы они удовлетворяли системе

\pandocbounded{\includegraphics[keepaspectratio]{images/_page_370_Picture_2.jpeg}}

\pandocbounded{\includegraphics[keepaspectratio]{images/_page_370_Figure_3.jpeg}}

\[
\begin{cases} \frac{dx}{ds} = a(x, t), \\ \frac{dt}{ds} = b(x, t). \end{cases}
-\infty < x < +\infty
-\infty < y < +\infty
\]

Ясно, что в этом случае

\[
-\infty < z < +\infty
\]

Другими словами, на характеристиках \{{[}\(x(s)\), \(t(s)\){]}: 0
\textless{} s \textless{} \infty\} уравнение с частными производными
становится обыкновенным дифференциальным уравнением. Поясним эти идеи на
конкретном примере.

\section{ПРИМЕР}\label{ux43fux440ux438ux43cux435ux440}

Припустимо, нам потрібно розв'язати наступну задачу Коші з сталими
коефіцієнтами:

\begin{verbatim}
\begin{array}{lll}
\end{verbatim}

\[
\begin{array}{llll} (\forall \Pi \Pi) & u_x + u_t + 2u = 0, & -\infty < x < \infty, & 0 < t < \infty \\ (\forall \Psi) & u(x, 0) = \sin x, & -\infty < x < \infty. \end{array}
-\infty & < x < +\infty & r \geqslant 0, \ 0 \leqslant \theta \leqslant 2\pi, \\
\frac{dx}{ds} = 1, \quad \frac{dt}{ds} = 1, \quad 0 < s < \infty.
-\infty & < y < +\infty & -\alpha \leqslant z \leqslant \alpha, \\
x(s) = s + c_1, \quad t(s) = s + c_2.
x(0) = \tau,
-\infty & < z < +\infty & x = r \cos \theta, \\
x = s + \tau
x-t=\tau
y = r \sin \theta, \\
\frac{du}{ds} + 2u = 0, \quad 0 < s < \infty,
u(0) = \sin \tau.
z = z.
u\left(s,\,\tau\right)=\sin\tau e^{-2s}
\]

\begin{verbatim}
\end{array}
\end{verbatim}

\[
x = s + \tau
s=t
r \geqslant 0
u(x, t) = \sin(x-t)e^{-xt}
\Delta u = 0
(\forall \Pi) \quad a(x, t) u_x + b(x, t) u_t + c(x, t) u = 0, \quad -\infty < x < \infty, \\ 0 < t < \infty,
\Delta u + \lambda^2 u = 0
u(x, 0) = f(x), \quad -\infty < x < \infty.
\Delta u = k
\frac{dx}{ds} = a(x, t), \qquad \frac{dt}{ds} = b(x, t).
\Delta u + k(E - V)u = 0
x = x (s, \tau),
t = t (s, \tau).
\frac{du}{ds} + c[x(s, \tau), t(s, \tau)]u = 0, 0 < s < \infty,
(\text{УЧП})
\frac{dx}{ds} = x
\frac{dt}{ds} = 1
x = \tau e^s,
u_{tt} = c^2 \Delta u
t = s.
u_{tt} = c^2 \Delta u - hu_t
\frac{du}{ds} + su = 0, \quad 0 < s < \infty,
u_t = \alpha^2 u_{xx}
u(0) = F(\tau).
u_t = \alpha^2 u_{xx} - hu_x
u(s, \tau) = F(\tau) e^{-s} \ell^2
u_t = \alpha^2 u_{xx} - ku
s=t
u_t = \alpha^2 u_{xx} + f(x, t)
u(x, t) = F(xe^{-t}) e^{-t^2/2}
\]

Уравнение одномерной диффузии Уравнение конвективной диффузии

Уравнение теплопроводности с поглощением Уравнение теплопроводности с
источником

\section{НЕЛИНЕЙНЫЕ УРАВНЕНИЯ ПЕРВОГО ПОРЯДКА (ЗАКОНЫ
СОХРАНЕНИЯ)}\label{ux43dux435ux43bux438ux43dux435ux439ux43dux44bux435-ux443ux440ux430ux432ux43dux435ux43dux438ux44f-ux43fux435ux440ux432ux43eux433ux43e-ux43fux43eux440ux44fux434ux43aux430-ux437ux430ux43aux43eux43dux44b-ux441ux43eux445ux440ux430ux43dux435ux43dux438ux44f}

ЦЕЛЬ ЛЕКЦИИ: Ввести понятие нелинейных уравнений первого порядка и
показать, как некоторые из них (так называемые законы сохранения) можно
использовать для описания физических явлений. Например, закон сохранения

\[
u_i + g(u)u_x = 0
\]

с начальным условием

\[
u(x, 0) = \varphi(x)
\]

используется для описания потока автомобилей на автостраде. Этот пример
показывает, что уравнения с частными производными применимы не только в
физике, биологии или технике. Более того, это нелинейное уравнение имеет
разрывные решения, которые описывают ударные волны, распространяющиеся в
потоке автомобилей вдоль дороги.

Одно из наиболее важных уравнений в частных производных --- это закон
сохранения

\[
u_t + f_x = 0.
\]

Этот закон утверждает, что увеличение некоторой физической величины
\(u_t\) равно изменению потока--- \(f_x\) этой величины через поперечное
сечение (поток, направленный слева направо, считается положительным). В
динамике жидкости величина \(u(x,t)\) может обозначать плотность
жидкости в точке x, а величина \(f(x,t)\)---поток жидкости (т. е.
количество жидкости, протекающей через поперечное сечение в точке x в
момент времени t). Мы не станем в этой лекции рассматривать динамику
невязкой несжимаемой жидкости. Вместо этого мы покажем, как закон
сохранения можно использовать для прогноза динамики дорожного движения
(нас будет интересовать поток автомобилей, а не поток молекул воды).

Сначала займемся выводом закона сохранения.

\section{Вывод закона
сохранения}\label{ux432ux44bux432ux43eux434-ux437ux430ux43aux43eux43dux430-ux441ux43eux445ux440ux430ux43dux435ux43dux438ux44f}

Предположим, что автомобили движутся слева направо по скоростной
автостраде, у которой нет боковых съездов и въездов. Обозначим

\(u(x, t)\) --- плотность автомобилей в точке x (число автомобилей на
единицу длины),

\(f(x, t)\) --- ноток автомобилей в точке x (число автомобилей,
проходящих через точку x в одну минуту).

Совершенно очевидно, что на участке дороги {[}a, b{]} изменение числа
автомобилей (за единицу времени) дается следующими двумя выражениями:

Изменение числа автомобилей на
\([a, b] = \frac{d}{dt} \int_a^b u(x, t) dx\)

И

Изменение числа автомобилей на
\([a,b] = f(a,t) - f(b,t) = -\int_{0}^{b} \frac{\partial f}{\partial x}(x,t)dx\)
.

Последнее соотношение --- следствие основной теоремы интегрального
исчисления. Приравнивая эти два интеграла, получаем

\[
\int_{a}^{b} \frac{\partial u}{\partial t}(x, t) dx = -\int_{a}^{b} \frac{\partial f}{\partial x}(x, t) dx.
\]

Поскольку промежуток {[}a,b{]} произволен, делаем вывод, что
подынтегральные функции равны, т е. выполняется закон сохранения

\[
u_t + f_x = 0.
\]

\section{Применение закона сохранения к задаче о дорожном
движении}\label{ux43fux440ux438ux43cux435ux43dux435ux43dux438ux435-ux437ux430ux43aux43eux43dux430-ux441ux43eux445ux440ux430ux43dux435ux43dux438ux44f-ux43a-ux437ux430ux434ux430ux447ux435-ux43e-ux434ux43eux440ux43eux436ux43dux43eux43c-ux434ux432ux438ux436ux435ux43dux438ux438}

Уравнение \(u_t + f_x = 0\) по форме очень простое. Правда, в нем две
неизвестные функции! Вопрос в том, как воспользоваться этим уравнением и
что это нам даст5

В задачах дорожного движения пользуются экспериментально найденной
зависимостью потока автомобилей \(f(u)\) от плотности автомобилей в
данной точке u. Представляется очевидным, что если плотность u растет,
то растет и поток f (по крайней мере в точке). Типичная модель дорожного
движения задается формулой (см. рис. 28.1)

\[
f(u) = Au (1-u).
\]

Используются и другие модели:

\begin{itemize}
\tightlist
\item
  \begin{enumerate}
  \def\labelenumi{\arabic{enumi})}
  \tightlist
  \item
    \(f(u)\) = ku (линейный рост потока),
  \end{enumerate}
\item
  \begin{enumerate}
  \def\labelenumi{\arabic{enumi})}
  \setcounter{enumi}{1}
  \tightlist
  \item
    \(f(u) = u^2\) (квадратичный рост потока).
  \end{enumerate}
\end{itemize}

\pandocbounded{\includegraphics[keepaspectratio]{images/_page_205_Figure_5.jpeg}}

РИС. 28 1. Типичная зависимость потока от плотности.

Теперь мы подготовлены к тому, чтобы применить уравнение \(u_i+f_x=0\) к
задаче дорожного движения. По правилу дифференцирования сложной функции
\(f_x=\frac{df}{du}\,u_x\) ; значит, закон сохранения можно переписать в
виде

\[
u_t + \frac{df}{du} u_x = u_t + g(u) u_x = 0.
\]

Пусть, например, зависимость потока от плотности квадратичная, т. е.
\(f(u) = u^{2}\) , тогда закон сохранения принимает следующую форму:

\[
u_t + 2uu_x = 0.
\]

Если начальная плотность распределения автомобилей была
\(u(x, 0) = \varphi(x)\) , то, чтобы найти плотность распределения в
произвольный момент времени, необходимо решить задачу с начальным
условием (задачу Коши):

\[
\begin{array}{llll} (\text{УЧП}) & u_t + 2uu_x = 0, & -\infty < x < \infty, & 0 < t < \infty, \\ (\text{HV}) & u\left(x, \ 0\right) = \varphi\left(x\right), & -\infty < x < \infty. \end{array}
\]

Займемся теперь решением нелинейной задачи Коши.

\section{Нелинейная задача с начальным условием (задача
Коши)}\label{ux43dux435ux43bux438ux43dux435ux439ux43dux430ux44f-ux437ux430ux434ux430ux447ux430-ux441-ux43dux430ux447ux430ux43bux44cux43dux44bux43c-ux443ux441ux43bux43eux432ux438ux435ux43c-ux437ux430ux434ux430ux447ux430-ux43aux43eux448ux438}

Рассмотрим задачу

\[
\begin{array}{llll} (\mathsf{Y}\mathsf{P}\Pi) & u_t + g\left(u\right)u_x = 0, & -\infty < x < \infty, & 0 < t < \infty, \\ (\mathsf{H}\mathsf{Y}) & u\left(x, \ 0\right) = \varphi\left(x\right), & -\infty < x < \infty. \end{array}
\]

Вспомним предыдущую лекцию, где мы изучали уравнение конвективного
переноса:

\[
u_t + vu_x = 0
\]

.

У цьому рівнянні \(u(x, t)\) --- це концентрація матерії в потоку, що
рухається зі швидкістю v

\[
u_1 + g(u)u_x = 0
\]

и уравнением переноса. Будем считать, что частина воды в гочке \(x_0\)
движется со скоростью \(g(u)\) (вдоль по течению или против него). Тогда
через t секунд координата частины будет

\[
x = x_0 + g(u) t
\]

(уравнение характеристик).

Пам'ятайте, що концентрація \(u (x, t)\) не змінюється вздовж
характеристики. Якщо ми знаємо початкову концентрацію \(u(x_a, 0)\) ,

\pandocbounded{\includegraphics[keepaspectratio]{images/_page_206_Figure_8.jpeg}}

Рис. 28.2. Характеристики уравнения \(u_t+g\) (u) \(u_x=0\) ; a---общий
вид характеристики, выходящей из т. ( \(x_0\) , 0), вдоль характеристики
решение остается постоянным и равным величине \(u (\) x\_0 \(, 0)\);
b---пример характеристики, начинающейся из т. (2, 0), уравнения
\(u_t+3uu_x=0\) с \(u (2, 0)\)=1.

то уравнение характеристик принимает вид \(x = x_0 + g[u(x_0, 0)] t\)
(уравнение характеристики, начинающейся в точке \((x_0, 0)\) ).

Пусть, например, мы решаем задачу

(УЧП)

\[
u_t + 3uu_x = 0
\]

,\\
(НУ) u(x, 0) =

\[
u_t = Du_{xx} - Vu_x.
\]

Чтобы найти характеристику, выходящую из точки, например (2, 0), мы
должны написать

\[
x = 2 + g[u(2, 0)]t =
\]

= \(2 + g[1]t =\) = \(2 + 3t\) .

Можна стверджувати, що розв'язок нашої задачі \(u(x, t)\) дорівнює
одному на лінін x=2+3t, тобто \(u(2+3t, t)\)=1 для всіх t\textgreater0.
Фіг. 28.2 показує характеристики загального випадку та цього конкретного
завдання.

Тепер зрозуміло, що, знаючи характеристики в кожній точці і знаючи, що
розв'язок за характеристиками не змінюється, можна отримати розв'язок
\(u\left(x,\ t\right)\) у будь-який момент часу t. Ми не отримали явного
виразу для \(u\left(x,\ t\right)\) через змінні x і t, але знання
характеристик рівняння дозволяє розв'язати деякі цікаві задачі.

\section{Задача дорожного
движения}\label{ux437ux430ux434ux430ux447ux430-ux434ux43eux440ux43eux436ux43dux43eux433ux43e-ux434ux432ux438ux436ux435ux43dux438ux44f}

Займемся моделью дорожного движения, в которой поток вычисляется по
формуле \(f(u)=u^2\) , а начальная плотность изо- Р И С

\pandocbounded{\includegraphics[keepaspectratio]{images/_page_207_Figure_8.jpeg}}

. 28.3. Начальная плотность автомобилей движущихся слева направо.

бражена на рис. 28.3. Другими словами, нам нужно решить задачу

\[
(\text{YUII}) \quad u_t + 2uu_x = 0, \quad -\infty < x < \infty, \quad 0 < t < \infty,
(\text{HV}) \quad u(x, \ 0) =
\begin{cases} 1, & x \leq 0, \\ 1 - x, & 0 < x < 1, \\ 0, & 1 \leq x. \end{cases}
\]

Начнем с определения характеристик, выходящих из точек вида \((x_0, 0)\)
. Для \(x_0 < 0\) получаем

\[
x = x_0 + g [u(x_0, 0)] t =
\]

= \(x_0 + g [1] t =\) = \(x_0 + 2t\) .

Разрешая относительно t, находим уравнение характеристик

\[
t = \frac{1}{2} (x - x_0).
\]

Эти прямые линии изображены на рис. 28.4. Теперь рассмотрим точки
\(0 < x_0 < 1\) . Здесь характеристики определяются следующими
соотношениями:

\[
x = x_0 + g[u(x_0, 0)] t =
= x_0 + g[(1 - x_0)] t =
= x_0 + 2(1 - x_0) t.
\]

Разремая последнее соотношение относительно t, получаем

\[
t = \frac{x - x_0}{2(1 - x_0)}.
\]

При \(1 ≤ x_0 < ∞\) характеристики

\[
x = x_0 + g[u(x_0, 0)]t =
\]

= \(x_0 + g[0]t =\) = \(x_0\) приводят к вертикальным линиям, выходящим
из точек хо. Все характеристики задачи изображены на рис. 28.4.

\pandocbounded{\includegraphics[keepaspectratio]{images/_page_208_Figure_10.jpeg}}

РИС. 28.4. Характеристики уравнения \(u_t + 2uu_x = 0\) .

Ситуація на дорозі при 0 \textless{} t \textless{} 1/2 досить зрозуміла.
Рис. 28.5 показує розв'язок нашої задачі в різні моменти часу. Зверніть
увагу, що характеристики сходяться в одній точці t=1/2. Отже, щоб
розв'язати задачу на t\textgreater1/2, потрібно використовувати інший
метод. Коли характеристики сходяться в одній точці, говоримо про ударну
хвилю (розрив розв'язку). Тепер потрібно відповісти на питання: на якій
швидкості є передній край

ударной волны будет перемещаться вдоль автострады? Хотя это и не
очевидно, но оказывается, что скорость распространения разрыва находится
по формуле

\[
S = \frac{f(u_R) - f(u_L)}{u_R - u_L},
\]

где \(u_R\) и \(u_L\) ---значения решения справа и слева от волнового
фронта, а \(f(u_R)\) и \(f(u_L)\) ---значения потока при этих
плотностях.

\pandocbounded{\includegraphics[keepaspectratio]{images/_page_209_Figure_5.jpeg}}

РИС. 28.5. Плотность дорожного движения в различные моменты времени.

У нашому прикладі швидкість поширення щілини виявляється рівною одиниці.
Насправді, якщо згадати цей \(f(u) = u^2\) , то одразу отримуємо

\[
S = \frac{0-1}{0-1} = 1
\]

.

Это означает, что при t\textgreater1/2 фронт волны будет перемещаться
слева направо со скоростью единица.

Полное решение задачи

\[
(\text{Y} \, \Pi) \quad u_t + 2uu_x = 0, \quad -\infty < x < \infty,
(\text{HY}) \quad u(x, \ 0) =
\begin{cases} 1, & x \leq 0, \\ 1 - x, & 0 < x < 1, \\ 0, & 1 < x, \end{cases}
\]

изображено на рис. 28.5.

\section{ЗАМЕЧАНИЯ}\label{ux437ux430ux43cux435ux447ux430ux43dux438ux44f-11}

\begin{enumerate}
\def\labelenumi{\arabic{enumi}.}
\item
  Ударна хвиля в нашому прикладі виникає через те, що потік швидко
  зростає зі збільшенням густини u. Якщо потік заданий формулою \(f(u)\)
  = u, то рівняння було б у вигляді \(u_1+u_x=0\) , і розв'язок цього
  рівняння представляє хвилю, що рухається вправо без спотворень. Якщо
  на мить замислитися над значенням формули \(f(u)\) = u, стає очевидно,
  що розв'язок рухатиметься таким чином.
\item
  Прямой подстановкой решения в виде
\end{enumerate}

\[
u = \varphi(x - g(u) t)
\]

можно убедиться в том, что эта формула неявно определяет решение задачи

\[
\begin{array}{lll} u_t + g\left(u\right) u_x = 0, & -\infty < x < \infty, & 0 < t < \infty, \\ u\left(x, \ 0\right) = \varphi\left(x\right), & -\infty < x < \infty. \end{array}
\]

Например, решение задачи Коши

\[
(Y\Pi\Pi) \quad u_t + uu_x = 0,
(HY) \quad u(x, 0) = x
\]

неявно вадается соотношением

\[
u = \varphi [x - g(u) t] =
= x - g(u) t = x - ut.
\]

У нашому конкретному випадку можливо отримати рішення у явній формі (у
багатьох інших випадках це неможливо)

\[
u(x, t) = \frac{x}{1+t}.
\]

Проверьте, что эта функция действительно является решением нашей задачи.

\section{\texorpdfstring{\textbf{ЗАДАЧИ}}{ЗАДАЧИ}}\label{ux437ux430ux434ux430ux447ux438-12}

\begin{enumerate}
\def\labelenumi{\arabic{enumi}.}
\tightlist
\item
  Решите задачу Коши
\end{enumerate}

\[
\begin{array}{lll} (\text{YUII}) & u_t + uu_x = 0, & -\infty < x < \infty, & 0 < t < \infty, \\ (\text{HV}) & u(x, \ 0) = \begin{cases} 0, & x < 0, \\ x, & 0 \leqslant x. \end{cases} \end{array}
\]

Постройте график решения для различных моментов времени. Как вы
интерпретируете это решение? Как выглядит соотношение между потоком и
плотностью в этой задаче?

\begin{enumerate}
\def\labelenumi{\arabic{enumi}.}
\setcounter{enumi}{1}
\tightlist
\item
  Решите задачу
\end{enumerate}

\[
\begin{split} & (\text{УЧП}) \quad u_t + u^2 u_x = 0, \quad -\infty < x < \infty, \quad 0 < t < \infty, \\ & (\text{НУ}) \quad u_t(x, 0) = \begin{cases} 0, & x < 0, \\ x, & 0 \leqslant x. \end{cases} \end{split}
\]

Какова здесь связь между потоком и плотностью? Совнадает ли поведение
полученного решения с тем, что вы ожидали? Сравните решения задачи 1 и
задачи 2.

\begin{enumerate}
\def\labelenumi{\arabic{enumi}.}
\setcounter{enumi}{2}
\tightlist
\item
  Предположим, что невязкая жидкость течет по трубе и при этом
  просачивается через стенки по закону \(F(u)\) = ku (г/см·с). Закон
  сохранения (скорее закон изменения) в этом случае записывается в виде
\end{enumerate}

\[
u_t + f_x = -F(u)
\]

.

Знайдіть розв'язок цього рівняння, якщо \(f(u)\) = u і
\(\varphi(x) = u\) (x, 0) (початковий розподіл задається загальним
чином). Дайте тлумачення отриманого рішення.

\begin{enumerate}
\def\labelenumi{\arabic{enumi}.}
\setcounter{enumi}{3}
\item
  Знайти розв'язок попереднього рівняння, якщо втрати визначаються
  функцією \(F(x, t)\) = 1/x. Перевірте розв'язок. Ви розумієте його
  фізичне значення?
\item
  Проверьте, что функция и, неявно заданная соотношением
\end{enumerate}

\[
u = \varphi[x - g(u) t],
\]

является решением нелинейной задачи

\[
(\text{УЧП}) \quad \begin{array}{ll} u_1 + g(u) u_x = 0, \\ (\text{НУ}) & u(x, 0) = \varphi(x). \end{array}
\]

\# СИСТЕМЫ УРАВНЕНИЙ С ЧАСТНЫМИ ПРОИЗВОДНЫМИ

ЦЕЛЬ ЛЕКЦИИ: Познакомить читателя с тем фактом, что многие физические
системы (особенно в динамике жидкости) нельзя описать одним уравнением.
Для описания таких систем требуются системы связанных уравнений. В эти
системы входят некоторые неизвестные функции, такие, как давление
\(p(x, y, z, t)\), плотность \(\rho(x, y, z, t)\) , температура
\(T(x, y, z, t)\) и их частные производные. Связи между этими функциями
определяются физическими законами и наша задача --- найти все эти
функции одновременно.

В лекции будет показано, как можно решить линейную систему уравнений

\[
u_t + Au_x = 0
\]

если преобразовать ее в новую систему независимых уравнений

\[
v_1 + \Lambda v_x = 0,
\]

а затем независимо решить каждое уравнение этой системы.

У багатьох галузях науки існують системи величин, які взаємопов'язані не
одним, а кількома відношеннями одночасно. Наприклад, у гідродинамікі це
чотири рівняння

\[
u_t + uu_x + vu_y + \frac{1}{\rho} p_x = 0
\]

(сохранение импульса вдоль оси \(x\) ), (29.1)
\(v_t + uv_x + vv_y + \frac{1}{\rho} p_y = 0\) (сохранение импульса
вдоль оси \(y\) ), (сохранение массы), (сохранение массы), (сохранение
массы),
\(\left(\frac{p}{\rho^{\gamma}}\right)_t + u\left(\frac{p}{\rho^{\gamma}}\right)_x + v\left(\frac{p}{\rho^{\gamma}}\right)_y = 0\)
(сохранение энергии), \(p(x, y, z, t)\) --- давление в жидкости,
\(u(x, y, z, t)\) --- \(x\) -компонента скорости, \(v(x, y, z, t)\) ---
\(y\) -компонента скорости, \(\rho(x, y, z, t)\) --- плотность жидкости.

Приведенная нелинейная система уравнений называется уравнениями Эйлера
движения жидкости. Задача состоит в одновременном отыскании всех
неизвестных функций р, и, v и р, которые удовлетворяли бы всем четырем
уравнениям (вместе с начальными и граничными условиями).

Існують й інші причини вивчати системи рівнянь. У теорії звичайних
диференціальних рівнянь (якщо читач пам'ятає) показано, що одне рівняння
другого порядку можна представити як систему з двох рівнянь першого
порядку. Хоча теорія диференціальних рівнянь у похідних не така проста,
дуже часто одне рівняння з похідними у похідних порядків зводиться до
системи диференціальних рівнянь першого порядку. Наприклад, телеграфне
рівняння

\[
u_{tt} = c^2 u_{xx} + a u_t + b u
\]

путем введения трех новых переменных

\[
u_1 = u_1
u_2 = u_2
u_3 = u_4
\]

можно свести к системе уравнений первого порядка

(29.2)

\[
\begin{cases} \frac{\partial u_1}{\partial x} = u_2, \\ \frac{\partial a_1}{\partial t} = u_3, \\ \frac{\partial u_3}{\partial t} = c^2 \frac{\partial u_2}{\partial x} + au_3 + bu_4. \end{cases}
\]

Для того чтобы решать системы типа (29.2), необходимо хорошо знать такие
понятия, как собственные значения и собственные векторы матрицы. По эфой
причине мы предлагаем вниманию читателя краткий обзор линейной алгебры.

\section{Обзор линейной
алгебры}\label{ux43eux431ux437ux43eux440-ux43bux438ux43dux435ux439ux43dux43eux439-ux430ux43bux433ux435ux431ux440ux44b}

Общий случай

Частный случай

Определение. Если А---квадратная матрица порядка n, то собственными
значениями матрицы называются n корней характе-

Пусть

\[
A =
\begin{bmatrix} 1 & 1 \\ 4 & 1 \end{bmatrix}
\]

.

ристического уравнения

\[
\det\left(A-\lambda I\right)=0,
\]

где \det(A-\lambda I) --- определитель матрицы A-\lambda I .

Некоторые собственные значения могут совпадать.

Тогда собственные значения являются корнями уравнения

\[
\det
\begin{bmatrix} 1-\lambda & 1\\ 4 & 1-\lambda \end{bmatrix}
= (1-\lambda)^2 - 4 =
= \lambda^2 - 2\lambda - 3 = 0.
\]

Откуда

\[
\lambda_1 = -1
\]

, \(\lambda_2 = 3\) .

Визначення. Якщо \(\lambda\) --- власне значення матриці A, то
відповідний власний вектор є таким ненульовим вектором, що задовольняє
відношення

\[
Ax = \lambda x
\]

.

Пример. Вектор X =

\[
u_t = -Vu_x
\]

\(является собственным вектором, соответствующим собственному значению\)
\lambda\_2 = 3 , так как

\[
\begin{bmatrix} 1 & 1 \\ 4 & 1 \end{bmatrix}
\begin{bmatrix} 1 \\ 2 \end{bmatrix}
= 3
\begin{bmatrix} 1 \\ 2 \end{bmatrix}
\]

.

Определение. Матрица \(A^{-1}\) называется обратной к матрице A, если

\[
AA^{-1} = A^{-1}A = I_1
\]

Пример, Обратной к

\[
A =
\begin{bmatrix} 2 & -2 \\ 1 & 1 \end{bmatrix}
\]

является матрица

\[
A^{-1} =
\begin{bmatrix} 0.25 & 0.5 \\ -0.25 & 0.5 \end{bmatrix}
\]

,

где 1-единичная матрица.

поскольку

\[
AA^{-1} = A^{-1}A =
\begin{bmatrix} 1 & 0 \\ 0 & 1 \end{bmatrix}
\]

.

\section{Диагонализация
матрицы}\label{ux434ux438ux430ux433ux43eux43dux430ux43bux438ux437ux430ux446ux438ux44f-ux43cux430ux442ux440ux438ux446ux44b}

\begin{verbatim}
Пусть A — квадратная матрина порядка n и пусть все ее собственные значения  \lambda_1  ,  \lambda_2  ...  \lambda_n  различны. Составим матрицу P так, чтобы у нее *столбцами* были координаты собственных векторов (в k-м столбце стоит собственный вектор  X_k  , соответствующий собственному значению  \lambda_k  ), т. е.
\end{verbatim}

\[
P =
\begin{bmatrix} X_1 & X_{\tilde{2}} & \dots & X_n \\ X_{1} & X_{2} & \dots & X_n \end{bmatrix}
=
\begin{bmatrix} x_{11} & x_{1\tilde{2}} & \dots & x_{1n} \\ x_{21} & x_{22} & \dots & x_{2n} \\ \vdots & \vdots & \ddots & \vdots \\ x_{n1} & x_{n\tilde{3}} & \dots & x_{nn} \end{bmatrix}
\]

.

Тогда матрица Л, определяемая соотношением

\[
\Lambda = P^{-1}AP,
\]

где \(P^{-1}\) --- обратная к матрице P, будет диагональной матрицей
вида

\[
\Lambda =
\begin{bmatrix} \lambda_1 & & & \\ & \lambda_2 & 0 & \\ & & \ddots & \\ & & & \lambda_n \end{bmatrix}
\]

.

Например, у матрицы

\[
A =
\begin{bmatrix} 0 & 8 \\ 2 & 0 \end{bmatrix}
\]

собственные числа равны \(\lambda_ 1 = 4\) и \(\lambda_2 = -4\) , а
собственные векторы соответственно

\[
X_1 =
\begin{bmatrix} 2 \\ 1 \end{bmatrix}
in  X_2 =
U(x) = \int_{0}^{\infty} u(x, t) e^{-st} dt
\]

.

Значит, произведение

\[
P^{-1}AP =
\begin{bmatrix} 2 & -2 \\ 1 & 1 \end{bmatrix}
^{-1}
\begin{bmatrix} 0 & 8 \\ 2 & 0 \end{bmatrix}
\begin{bmatrix} 2 & -2 \\ 1 & 1 \end{bmatrix}
= \\ =
\begin{bmatrix} 1/4 & 1/2 \\ -1/4 & 1/2 \end{bmatrix}
\begin{bmatrix} 0 & 8 \\ 2 & 0 \end{bmatrix}
\begin{bmatrix} 2 & -2 \\ 1 & 1 \end{bmatrix}
\]

равно диагональной матрице

\[
\Lambda =
\begin{bmatrix} 4 & 0 \\ 0 & -4 \end{bmatrix}
\]

.

(Читатель может проверить это самостоятельно.)

Теперь можно решить простую систему двух уравнений с частными
производными (с соответствующими начальными условиями).

\section{\texorpdfstring{Решение линейной системы \(u_t\) + Au\_x = 0
.}{Решение линейной системы u\_t + Au\_x = 0 .}}\label{ux440ux435ux448ux435ux43dux438ux435-ux43bux438ux43dux435ux439ux43dux43eux439-ux441ux438ux441ux442ux435ux43cux44b-u_t-au_x-0-.}

Рассмотрим задачу Коши для системы, содержащей два уравнения и два
начальных условия:

(УЧП 1)

\[
\frac{\partial u_1}{\partial t} + 8 \frac{\partial u_2}{\partial x} = 0
\]

,

\begin{verbatim}
(29.3) (УЧП 2)  \frac{\partial u_2}{\partial t} + 2 \frac{\partial u_1}{\partial x} = 0  ,  -\infty < x < \infty  ,  0 < t < \infty  ,
\end{verbatim}

\[
(HУ 1)  u_1(x, 0) = f(x)  ,
(HУ 2)  u_2(x, 0) = g(x)  ,  -\infty < x < \infty  .
\]

Эта задача может соответствовать определению давления \(u_1\)(x, t) и
плотности \(u_2\)(x, t) как функции пространственной координаты x и
времени t по известным распределениям этих величин в начальный момент.

Запишем систему уравнений в матричной форме

\[
\begin{bmatrix} \frac{\partial u_1}{\partial t} \\ \frac{\partial u_2}{\partial t} \end{bmatrix}
\]

\begin{itemize}
\tightlist
\item
\end{itemize}

\[
\begin{bmatrix} 0 & 8 \\ 2 & 0 \end{bmatrix}
\begin{bmatrix} \frac{\partial u_1}{\partial x} \\ \frac{\partial u_2}{\partial x} \end{bmatrix}
=
\begin{bmatrix} 0 \\ 0 \end{bmatrix}
\]

или \((29 \ 4)\)

\[
u_t + Au_s = 0
\]

где

\[
A =
\begin{bmatrix} 0 & 8 \\ 2 & 0 \end{bmatrix}
, \quad u_1 =
\begin{bmatrix} \frac{\partial u_1}{\partial t} \\ \frac{\partial u_2}{\partial t} \end{bmatrix}
, \quad u_x =
\begin{bmatrix} \frac{\partial u_1}{\partial x} \\ \frac{\partial u_2}{\partial x} \end{bmatrix}
, \quad 0 =
\begin{bmatrix} 0 \\ 0 \end{bmatrix}
\]

.

Введем новые неизвестные функции

\[
v =
\begin{bmatrix} v_1 \\ v_2 \end{bmatrix}
\]

с помощью преобразования

\[
u = Pv
\]

.

де P --- матриця, стовпці якої містять власні вектори матриці A (матриця
P, яку вже знають). Виявляється, що після такого перетворення буде
отримана дуже проста система для визначення v (два рівняння відносно
нового невідомого \(v_1\) і \(v_2\) виявляються незалежними). Це
означає, що \(v_1\) і \(v_2\) легко знайти. Після \(v_1\) і \(v_2\)
знайдено, згідно з формулою

\[
u = Pv
\]

находятся искомые функции и, и и2.

Сначала, однако, выясним, как выглядит система для определения v. Чтобы
разобраться в этом, читатель должен обратиться к матричному исчислению.
Продифференцировав обе части соотношения u = Pv, получаем

(29 5)

\[
\frac{\partial u}{\partial t} = P \frac{\partial v}{\partial t},
\frac{\partial u}{\partial x} = P \frac{\partial v}{\partial x}.
\]

(Чтобы лучше понять дальнейшее изложение, читателю рекомендуется
записать эти уравнен я в развернутой форме.)

Теперь подставим соотношения (29.5) в систему

\[
u_t + Au_x = 0
\]

и получим

\[
Pv_t + APv_x = 0.
\]

Если умножить теперь обе части последнего уравнения на \(P^{-1}\) , то
получим

\[
v_t + P^{-1}APv_x = 0,
\]

или

\[
(29.6) v_t + \Lambda v_x = 0.
\]

Раньше мы уже видели, что для нашей матрицы

\[
\Lambda =
\begin{bmatrix} 4 & 0 \\ 0 & -4 \end{bmatrix}
\]

.

Запишем (29.6) в развернутой форме

(29.7)

\begin{verbatim}
\frac{\frac{\partial v_1}{\partial t} + 4 \frac{\partial v_1}{\partial x} = 0, \\ \frac{\partial v_2}{\partial t} - 4 \frac{\partial v_2}{\partial x} = 0.
\end{verbatim}

Результатом є система з двох непов'язаних рівнянь, які можна розв'язати
незалежно. Читач має знати, що розв'язки цих рівнянь --- це хвилі, що
рухаються

\[
v_1(x, t) = \varphi(x-4t),
\]

\(v_2(x, t) = \psi(x+4t),\) где ф и ф-произвольные дифференцируемые
функции.

Таким образом, мы нашли функции \(v_1\) и \(v_2\) . Для получения общего
решения u нужно вычислить

\[
u = Pv =
\begin{bmatrix} 2 & -2 \\ 1 & 1 \end{bmatrix}
\begin{bmatrix} v_1 \\ v_2 \end{bmatrix}
=
=
\begin{bmatrix} 2 & -2 \\ 1 & 1 \end{bmatrix}
\begin{bmatrix} \varphi(x - 4t) \\ \psi(x + 4t) \end{bmatrix}
=
=
\begin{bmatrix} 2\varphi(x - 4t) - 2\psi(x + 4t) \\ \varphi(x - 4t) + \psi(x + 4t) \end{bmatrix}
\]

.

Другими словами,

(29.8)

\[
u_1(x, t) = 2\varphi(x-4t) - 2\psi(x+4t), u_2(x, t) = \varphi(x-4t) + \psi(x+4t).
\]

Следовательно, мы нашли общее решение системы двух уравнений (29.3).

Типичным решением (а их бесконечное число) могут быть, например, функции

\[
\phi(\xi) = \sin \xi, \ \psi(\xi) = \xi^2
\]

(две произвольные функции)

и, следовательно,

\[
u_1(x, t) = 2 \sin(x-4t) - 2(x+4t)^2,
\]

\(u_2(x, t) = \sin(x-4t) + (x+4t)^2.\) Подставим теперь общее решение
(29.8) в начальные условия

\[
u_1(x, 0) = f(x),
\]

\(u_2(x, 0) = g(x),\) в результате получаем

\[
2\varphi(x) - 2\psi(x) = f(x),
\varphi(x) + \psi(x) = g(x),
\]

откуда функци і ф и ф легко находятся:

\[
\varphi(x) = \frac{1}{4} \left[ f(x) + 2g(x) \right],
\psi(x) = \frac{1}{4} \left[ 2g(x) - f(x) \right].
\]

\textbf{Тепе}рь можно написать решение задачи (29.3) в окончательном
виде

\[
u_{1}(x, t) = 2\varphi(x-4t) - 2\psi(x+4t) =
= \frac{1}{2} \left[ f(x-4t) + 2g(x-4t) \right] - \frac{1}{2} \left[ 2g(x+4t) - f(x+4t) \right],
\]

(29.9)

\[
u_{2}(x, t) = \varphi(x-4t) + \psi(x+4t) =
= \frac{1}{4} \left[ f(x-4t) + 2g(x-4t) \right] + \frac{1}{4} \left[ 2g(x+4t) - f(x+4t) \right].
\]

\section{ЗАМЕЧАНИЯ}\label{ux437ux430ux43cux435ux447ux430ux43dux438ux44f-12}

\begin{enumerate}
\def\labelenumi{\arabic{enumi}.}
\item
  Многие численные методы ориентированы на решение систем уравнений и,
  значит, многне программы для ЭВМ написаны так, чтобы решать системы
  уравнений первого порядка. По этой причине, прежде чем воспользоваться
  программои, приходится преобразовать свое уравнение высокого порядка в
  систему уравнений первого порядка.
\item
  Линейную систему
\end{enumerate}

\[
u_t + A(x, t) u_x + B(x, t) u = 0
\]

можна розв'язати тим самим методом, що й система \(u_t + Au_x = 0\) .
Різниця полягає в тому, що в цьому випадку матриця P із власних векторів
матриці \(A(x, t)\) буде функцією від x і t.

\section{ЗАДАЧИ}\label{ux437ux430ux434ux430ux447ux438-13}

\begin{enumerate}
\def\labelenumi{\arabic{enumi}.}
\tightlist
\item
  Запишите систему уравнений (29.2) в матричной форме
\end{enumerate}

\[
Au_t + Bu_x + Cu = 0,
\]

где A, B и C-матрицы третьего порядка.

\begin{enumerate}
\def\labelenumi{\arabic{enumi}.}
\setcounter{enumi}{1}
\tightlist
\item
  Найдите собственные векторы и собственные значения матрицы
\end{enumerate}

\[
A =
\begin{bmatrix} 1 & 1 \\ 4 & 1 \end{bmatrix}
\]

.

\begin{enumerate}
\def\labelenumi{\arabic{enumi}.}
\setcounter{enumi}{2}
\tightlist
\item
  Воспользовавшись результатами задачи 2, решите систему
\end{enumerate}

\[
\begin{cases} \frac{\partial u_1}{\partial t} + \frac{\partial u_1}{\partial x} + \frac{\partial u_2}{\partial x} = 0, \\ \frac{\partial u_2}{\partial t} + 4 \frac{\partial u_1}{\partial x} + \frac{\partial u_2}{\partial x} = 0. \end{cases}
\]

УКАЗАНИЕ. Запишите сначала систему в матричной форме.

\begin{enumerate}
\def\labelenumi{\arabic{enumi}.}
\setcounter{enumi}{3}
\item
  Проверьте формулы (29.5), записав их в скалярной форме.
\item
  Проверьте, что функции \(u_1\) и \(u_2\) из (29.8) удовлетворяют обоим
  уравнениям (29.3).
\end{enumerate}

\section{КОЛЕБАНИЯ МЕМБРАНЫ (ВОЛНОВОЕ УРАВНЕНИЕ В ПОЛЯРНЫХ
КООРДИНАТАХ)}\label{ux43aux43eux43bux435ux431ux430ux43dux438ux44f-ux43cux435ux43cux431ux440ux430ux43dux44b-ux432ux43eux43bux43dux43eux432ux43eux435-ux443ux440ux430ux432ux43dux435ux43dux438ux435-ux432-ux43fux43eux43bux44fux440ux43dux44bux445-ux43aux43eux43eux440ux434ux438ux43dux430ux442ux430ux445}

ЦЕЛЬ ЛЕКЦИИ: Показать, что колебания мембраны можно описать волновым
уравнением

\[
u_{tt} = c^2 \left( u_{rr} + \frac{1}{r} u_r + \frac{1}{r^2} u_{\theta\theta} \right)
\]

.

Если искать решение этого уравнения в виде стоячих волн:

\[
u(r, \theta, t) = P(r) \Theta(\theta) T(t),
\]

то оно распадается на три обыкновенных дифференциальных уравнения:

\[
T'' + \lambda^2 c^2 T = 0
\]

(уравнение гармонических колебаний),
\(r^2 R'' + rR' + (\lambda^2 r^2 - n^2) R = 0\) (уравнение Бесселя),
(уравнение гармонических колебаний).

Все произведения \(R(r) \Theta(\theta) T(t)\) решений этих трех
обыкновенных дифференциальных уравнений описывают основные колебания
мембраны, а функции \(R(r) \Theta(\theta)\) описывают форму мембраны.
Чтобы найти колебания мембраны при произвольных начальных условиях, мы
составим такую комбинацию фундаментальных решений, чтобы начальные
условия удовлетворялись.

Цель этой лекции ясна---найти колебания круглой мембраны с заданными
начальными и граничными условиями. Для простоты будем считать, что
радиус мембраны равен единице, а смещение по границе равно нулю. Пусть
\(u\left(r,\theta,t\right)\) обозначает возвышение точек мембраны под
плоскостью равновесия. Нам необходимо решить задачу (см. рис. 30.1)

\[
\begin{array}{llll} (\text{УЧП}) & u_{tt} = c^2 \left( u_{rr} + \frac{1}{r} u_r + \frac{1}{r^2} u_{\theta\theta} \right), & 0 < r < 1, \ 0 < t < \infty, \\ (\text{ГУ}) & u = 0 & \text{при } r = 1, & 0 < t < \infty, \\ (\text{НУ}) & \begin{cases} u = f(r, \ \theta) \\ u_t = g(r, \ \theta) \end{cases} & \text{при } t = 0. \end{array}
\]

Щоб розв'язати цю задачу, згадаємо проблему коливань скрипкових струн із
лекції 20, яку розв'язали суперпозицією нескінченної кількості простих
коливань. Якщо ви використовуєте-

ся этим же подходом, то решения следует искать в виде

\[
u(r, \theta, t) = U(r, \theta) T(t)
\]

.

Форма этих колебаний определяется функцией \(U(r, \theta)\) , а характер
осцилляций --- множителем \(T(t)\).

Подставив это представление в волновое уравнение (см. задачу 1), получим
два уравнения

\[
\Delta U + \lambda^2 U = 0
\]

(уравнение Гельм-гольца),

\[
T'' + \lambda^2 c^2 T = 0
\]

(уравнение гармонических колебаний),

\pandocbounded{\includegraphics[keepaspectratio]{images/_page_221_Picture_9.jpeg}}

РИС. 30.1. Қолеблющаяся мембрана (гиперболическая смешанная задача).

где

\[
\Delta U = U_{rr} + \frac{1}{r} U_r + \frac{1}{r^2} U_{\theta\theta}.
\]

Зверніть увагу, що тут ми вимагаємо, щоб константа розділення була
додатною (тому позначимо її \(\lambda^{2-1}\) )). Лише в цьому випадку
функції \(T(t)\) будуть періодичними.

На следующем шаге мы решим уравнение Гельмгольца, но сначала разберемся
с граничным условием для него. Чтобы найти его, подставим
\(u(r, \theta, t) = U(r, \theta) T(t)\) в граничное условие для мембраны

\[
u(1, \theta, t) = U(1, \theta) T(t) = 0, \quad 0 < t < \infty,
\]

или

\[
U(1, \theta) = 0.
\]

Следовательно, для того чтобы найти форму фундаментальных колебаний
мембраны, необходимо решить задачу

\[
\Delta U + \lambda^2 U = 0,
U(1, \theta) = 0.
\begin{cases} \Delta U + \mu u = 0 & \text{в круге,} \\ U = 0 & \text{на границе круга} \end{cases}
\]

при \(\mu \le 0\) имеет только тривиальное решение U=0. --- Прим. ред.

1) Поскольку краевая задача

Это хорошо известная эллиптическая задача на собственные значения, и нам
надо найти все \(\lambda\) , при которых она имеет ненуленые решения.
Решения \(U\left(r,\,\theta\right)\) описывают форму фундаментальных
колебаний мембраны, а собственные числа \(\lambda\) являются нулями
некоторых функций Бесселя и пропорциональны частотам этих колебаний.

Итак, нам предстоит теперь решить задачу на собственные значения для
уравнения Гельмгольца (эта задача очень важна и сама по себе). Будем
решать ее тем же методом, что и другие линейные однородные уравнения с
нулевыми граничными условиями, а именно методом разделения переменных.

\subsection{Решение задачи на собственные значения для уравнения
Гельмгольца}\label{ux440ux435ux448ux435ux43dux438ux435-ux437ux430ux434ux430ux447ux438-ux43dux430-ux441ux43eux431ux441ux442ux432ux435ux43dux43dux44bux435-ux437ux43dux430ux447ux435ux43dux438ux44f-ux434ux43bux44f-ux443ux440ux430ux432ux43dux435ux43dux438ux44f-ux433ux435ux43bux44cux43cux433ux43eux43bux44cux446ux430}

Чтобы решить задачу

\[
\begin{array}{ll} (\mathsf{Y}\mathsf{H}\mathsf{\Pi}) & \Delta U + \lambda^2 U = 0, \\ (\mathsf{\Gamma}\mathsf{Y}) & U(1, \theta) = 0, \end{array}
\]

подставим

\[
U(r, \theta) = R(r) \Theta(\theta)
\]

в задачу Гельмгольца. Сделав это, получаем

\[
r^2R'' + rR' + (\lambda^2r^2 - n^2)R = 0
\]

(уравнение Бесселя), \(R(1) = 0\) , \(\Theta'' + n^2\Theta = 0\) (Зроби
все сам). Зверніть увагу, що ми обрали константу розбиття \(n^2\) ,
\(n=0,. \ 1,\ 2,\ \ldots\) тому, що хочемо, щоб функції
\(\Theta\left(\theta\right)\) були періодичними з періодом \(2\pi\) .
Очевидно, що форма мембрани є періодичною функцією \(\theta\) . Отже,
щоб розв'язати задачу Гельмгольца, потрібно розв'язати два звичайних
диференціальні рівняння

\[
r^2R'' + rR' + (\lambda^2r^2 - n^2) R = 0
\]

, \(0 < r < 1\) , \(R(0) < \infty\) (физическое ограничение),
\(R(1) = 0\) , \(\Theta'' + n^2\Theta = 0\) .

\section{Уравнение
Бесселя}\label{ux443ux440ux430ux432ux43dux435ux43dux438ux435-ux431ux435ux441ux441ux435ux43bux44f}

Уравнение

\[
r^2R'' + rR' + (\lambda^2r^2 - n^2)R = 0
\]

жорошо известно в теории обыкновенных дифференциальных уравнений. Это
так называемое уравнение Бесселя. Оно имеет два

линейно-независимых решения \(R_{1}(r) = A J_{n}(\lambda r)\) функция
Бесселя \emph{первого} рода \emph{п}-го порядка,
\(R_{2}(r) = BY_{n}(\lambda r)\) функция Бесселя второго рода n-го
порядка,

и, следовательно, общее решение этого уравнения представимо в виде

\[
R(r) = AJ_n(\lambda r) + BY_n(\lambda r).
\]

Зверніть увагу, що рішення залежать від параметрів n і \(\lambda\) ,
включених у рівняння. Деякі графіки функції Бесселя показані на рис.
30.2. Щоб визначити функції \(J_ n(\lambda r)\) і \(Y_n(\lambda r)\) ,
можливо

\pandocbounded{\includegraphics[keepaspectratio]{images/_page_223_Figure_8.jpeg}}

\pandocbounded{\includegraphics[keepaspectratio]{images/_page_223_Figure_9.jpeg}}

Рис. 30 2 Функции Бесселя. a --- функции Бесселя первого рода;
\(\pmb{\delta}\) --- функции Бесселя второго рода.

використовуйте метод Фробеніуса, який дозволяє знаходити \(R(r)\) у
вигляді силових рядів. Цей метод містить два лінійно незалежні ряди
ступенів \(J_n(\lambda r)\) та \(Y_n(\lambda r)\) . Оскільки функція
\(Y_n(\lambda r)\) не обмежена в r=0, ми обираємо розв'язок у вигляді

\[
R(r) = AJ_n(\lambda r).
\]

Последний шаг в определении \(R\left( r\right)\) связан с использованием
граничного условия \(R\left( \mathbf{l}\right) =0\) для определения всех
\(\lambda\) (в данный

Таблица 30.1

{\def\LTcaptype{none} % do not increment counter
\begin{longtable}[]{@{}lll@{}}
\toprule\noalign{}
Корни & уравнений. & \(J_n(r) = 0\) \\
\midrule\noalign{}
\endhead
\bottomrule\noalign{}
\endlastfoot
------- & ------------ & -------------- \\
\end{longtable}
}

{\def\LTcaptype{none} % do not increment counter
\begin{longtable}[]{@{}lllllllll@{}}
\toprule\noalign{}
Γ & 1 & n & & & & & & \\
\midrule\noalign{}
\endhead
\bottomrule\noalign{}
\endlastfoot
& & 0 & 1 & 2 & 3 & 4 & & \\
& 1 & 2,40 & 3,83 & 5,13 & 6,38 & 7,59 & & \\
ı & 2 & 5,52 & 7,02 & 8,42 & 9,76 & 11,06 & & \\
ľ & 3 & 8,65 & 10,17 & 11,62 & 13,02 & 14,37 & & \\
& 4 & 11,79 & 13,32 & 14,80 & 16,22 & 17,62 & & \\
& 5 & 14,93 & 16,47 & 17,96 & 19,41 & 20,83 & & \\
& . & • & & • & & • & & \\
& & & • & & • & • & & \\
& • & • & • & • & • & • & & \\
\end{longtable}
}

момент нас не интересует коэффициент A). Подставляя \(R (1)\) = 0, в
\(R (r)\) = \(J_n\) ( \(\lambda r\) ), получаем \(J_n\) ( \(\lambda\) )
= 0. Другими словами, для того чтобы функция \(R (r)\) обращалась в нуль
на границе круга, мы должны выбрать константу разделения так, чтобы она
была корнем уравнения \(J_n\) (r) = 0, \(\tau\) . e.

\[
\lambda = k_{nm}
\]

где \(k_{nm}\) --- m-й корень уравнения \(J_n(r) = 0\) .

Таблиці цих коренів добре відомі, але їх можна обчислити на комп'ютері.
Значення перших кількох коренів наведені в Таблиці 1. 30.1. Знаючи ці
корені, ми можемо одразу побудувати розв'язок задачі граничного значення
для рівняння Гельмгольца. Якщо власні значення \(\lambda = k_{nm}\) , то
власні функції \(U_{ nm}(r)\) визначаються формулою

\[
U_{nm}(r, \theta) = J_n(k_{nm}r) [A \sin(n\theta) + B \cos(n\theta)].
\]

На рис. 30.3 изображены некоторые из этих функций для различных значений
n и m. Общая форма функции \(U_{nm}(r,\theta)\) не зависит от значений
констант A и B. Эти константы влияют только на амплитуду колебаний и
расположение линии узлов относительно \(\theta=0\) .

Каждая функция представляет собой фундаментальное колебание круглой
мембраны с частотой

\[
f_{nm} = k_{nm}c/2\pi.
\]

Эти частоты мы нашли из решений

\[
T_{nm}(t) = A \sin(k_{nm}ct) + B \cos(k_{nm}ct)
\]

\section{уравнения}\label{ux443ux440ux430ux432ux43dux435ux43dux438ux44f}

\[
T'' + k_{nm}^2 c^2 T = 0.
\]

Интересно отметить, что отношение частоты колебания
\(U_{nm}(\mathbf{r}, \mathbf{\theta})\) к частоте основного колебания
U\_\{nt\}(\ma

\[
sU(x) = -V \frac{dU}{dx}, \quad 0 < x < \infty,
\begin{cases} u(R_1, \theta) = g_1(\theta), \\ u(R_2, \theta) = g_2(\theta). \end{cases}
U(0) = P/s.
\]

\begin{verbatim}
\begin{cases} 2/3, & n=1, \\0 & \text{для всіх інших } n, \\ d_n=\begin{cases} -2/3, & n=1, \\0 & \text{для всіх інших} n. \end{cases}
\end{verbatim}

\[
H(\xi) =
\begin{cases} 1, & 0 \le \theta < \pi, \\ 0, & \pi \le \theta < 2\pi. \end{cases}
\begin{cases} 0, & \xi < 0, \\ 1, & \xi \geqslant 0. \end{cases}
\]

INL1448X (\Gamma Y) u(1, \(\varphi\)) = \cos(3\(\varphi\)) .

Используя тригонометрию, попытайтесь выразить \(\cos{(3\phi)}\) через
\(\cos{\phi}, \, \cos^2{\phi}, \, \cos^3{\phi}, \, a\) затем
скомбинируйте их так, чтобы по-

лучилось разложение

\[
\cos(3\varphi) = a_0 P_0(\cos\varphi) + a_1 P_1(\cos\varphi) + \dots
\]

\begin{enumerate}
\def\labelenumi{\arabic{enumi}.}
\setcounter{enumi}{4}
\tightlist
\item
  Решите задачу Дирихле в шаре
\end{enumerate}

\[
\Delta u = 0, \quad 0 < r < 1, \quad 0 \le \varphi \le \pi, \quad 0 \le \theta < 2\pi.
u(1, \varphi) =
\begin{cases} 1, & 0 \le \varphi \le \pi/2, \\ -1, & \pi/2 < \varphi \le \pi. \end{cases}
\]

В этой задаче верхняя полусфера нагрета (+1), а нижняя охлаждена (-1).

\begin{enumerate}
\def\labelenumi{\arabic{enumi}.}
\setcounter{enumi}{5}
\tightlist
\item
  Найдите решение задачи Дирихле в шаре
\end{enumerate}

(УЧП)

\[
\Delta u = 0
\]

, \(1 < r < \infty\) , \(0 \le \varphi \le \pi\) ,
\(0 \le \theta < 2\pi\) , \((\Gamma Y)\)
\(u(1, \varphi) = 1 + \cos \varphi\) .

Проверьте ответ.

Лекция 36

\section{НЕОДНОРОДНАЯ ЗАДАЧА ДИРИХЛЕ (ФУНКЦИЯ
ГРИНА)}\label{ux43dux435ux43eux434ux43dux43eux440ux43eux434ux43dux430ux44f-ux437ux430ux434ux430ux447ux430-ux434ux438ux440ux438ux445ux43bux435-ux444ux443ux43dux43aux446ux438ux44f-ux433ux440ux438ux43dux430}

МЕТА ЛЕКЦІЇ: Показати, як гетерогенну задачу Діріхле можна розв'язати за
допомогою функції Гріна (вихідної функції). У цьому важливому методі
права частина рівняння розглядається як певна вхідна дія і розкладається
на неперервний набір дельтаподібних джерел, розподілених по заданій
площі. Потім визначається відповідь системи на кожне таке джерело
(функція Гріна), і всі відповіді підсумовуються (інтегровані). У
результаті ми отримуємо повне розв'язання проблеми.

Достаточно общей задачей прикладной математики является определение
потенциала в некоторой области пространства, если задано распределение
источников \(f(x, y)\) енутри этой области. В электростатике потенциал
(в вольтах) в области D обусловлен распределением зарядов с плотностью
\(f(x, y)\) по этой области. Типичная задача --- определение потенциала
в круге в двумерном случае.

В такой задаче потенциал должен удовлетворять уравнению

Пуассона (например, с нулевыми граничными условиями)

(36.1) (УЧП)

\[
u_{rr} + \frac{1}{r} u_r + \frac{1}{r^2} u_{\theta\theta} = f(r, \theta), \quad 0 \le r < 1, \quad 0 \le \theta < 2\pi,
\]

(ГУ) \(u(1, \theta) = 0, \quad 0 \le \theta \le 2\pi.\) Ми не випадково
обрали нульові умови. Якщо б ми хотіли розв'язати задачу в загальному
випадку, з ненульовими граничними умовами та неоднорідним рівнянням, то
внесок неоднорідних граничних умов можна було б знайти за допомогою
формули інтегралу Пуассона з лекції 33.

Чтобы получить хоть какое-то представление о неоднородных
дифференциальных уравнениях, давайте рассмотрим графически решение
следующей задачи для уравнения Пуассона:

(УЧП)

\[
\Delta u = -q
\]

, \(0 < r < 1\) , \(0 \le \theta \le 2\pi\) ( \(q\) ---положительная
константа),

(HY)

\[
u(1, \theta) = 0, \quad 0 \le \theta \le 2\pi.
\]

На межі потенціал (або температура, якщо ви бажаєте) утримується на
нулі, а лапласіан дорівнює q у всіх точках області. Оскільки
\(\Delta u\left(p\right)\) слугує мірою різниці між функцією
\(u\left(p\right)\) і її середнім, рівняння Пуассона стверджує, що
поверхня \(u\left(r\right)\) завжди опукла вниз. Інакше кажучи, вона
виглядатиме як тонка мембрана з нерухомим краєм, яка здувається зверху
вниз потоком повітря.

Якщо права частина рівняння \(f(x, y)\) змінюється від точки до точки,
то опуклість також буде змінною функцією точки цієї області.

Перейдем теперь к главной части нашей лекции: определению функции Грина
и решению уравнения (36.1).

Сначала, однако, мы должны ввести понятие потенциала точеных источников
и стоков.

\section{Потенциалы точечных источников и
стоков}\label{ux43fux43eux442ux435ux43dux446ux438ux430ux43bux44b-ux442ux43eux447ux435ux447ux43dux44bux445-ux438ux441ux442ux43eux447ux43dux438ux43aux43eux432-ux438-ux441ux442ux43eux43aux43eux432}

Щоб розв'язати неоднорідне лінійне рівняння, достатньо розв'язати це
рівняння за допомогою точкового джерела, оскільки розв'язок задачі в
загальному випадку можна знайти, підсумувавши внески з точкових джерел.
Наше завдання --- визначити потенціал, який створює точкове джерело (або
поглинач) у певній області. Ці джерела можна інтерпретувати по-різному.
У теорії теплопровідності джерело можна розглядати як точку, де походить
тепло, а дренаж --- як точку, де поглинає тепло. В електростатиці
джерело можна розглядати як окремий позитивний заряд (протон), а дренаж
--- як окремий негативний заряд (електрон). У будь-якому разі, незалежно
від інтерпретації, ми знайдемо двовимірний потенціал і,

создаваемый единичным точечным источником (поиск трехмерного потенциала
мы оставляем читателю в качестве упражнения).

Припустимо, що єдине точкове джерело +q розташоване в початку Clear that
heat (або щось інше)

будет вытекать из области вдоль радиальных линий, и, следовательно, для
везничны полного потока, протекающего через окружность радиуса r можно
записать

Полный поток, вытекающий через окруж-

HOCTD

\[
= -\int_{0}^{2\pi} u_{r}(r) r d\theta = -2\pi r u_{r}(r).
\]

Но полный поток должен совпадать с количеством тепла, созданного внутри
круга (закон сохранения), т. е.

\[
-2\pi r u_r(r) = q.
\]

Решая это диф реренциальное уравнение относительно функции \(u(r)\),
получаем Р И С

\pandocbounded{\includegraphics[keepaspectratio]{images/_page_273_Picture_10.jpeg}}

36 1 Радиальный поток тепла от точечного источника; a ---тепло
выделяется в начале координат.

\[
u(r) = \frac{-q}{2\pi} \ln r = \frac{q}{2\pi} \ln \frac{1}{r}.
\]

(Решение изображено на рис. 362)

В электростатике разность потенциа тов \(u(B)\) - \(u(A)\) равна работе,
которую необходимо совершить, чтобы переместить единичный положительный
заряд из точки A в точку B (рис. 36.2).

\pandocbounded{\includegraphics[keepaspectratio]{images/_page_273_Figure_15.jpeg}}

РИС 36 2 Потенциал точечного заряда в двумерном случае. \(\alpha\) ---
работа, необходимая для перемещения единичного заряда из A в B

Стік можна інтерпретувати як від'ємне джерело, так що стік -q створює
потенціал

\[
u\left(r\right)=-\frac{q}{2\pi}\ln\frac{1}{r}.
\]

Этим завершается обсуждение потенциалов точечных зарядов. Теперь мы
готовы к решению неоднородного уравнения методом функции Грина.

\section{Уравнение Пуассона в
круге}\label{ux443ux440ux430ux432ux43dux435ux43dux438ux435-ux43fux443ux430ux441ux441ux43eux43dux430-ux432-ux43aux440ux443ux433ux435}

Займемся решением очень важной задачи

(36.2)

\[
(Y \cup \Pi)
\]

\(u_{rr} + \frac{1}{r} u_r + \frac{1}{r^2} u_{\theta\theta} = f(r, \theta), \quad 0 < r < 1, \quad 0 \le \theta \le 2\pi,\)
\((\Gamma Y)\) \(u(1, \theta) = 0, \quad 0 \le \theta \ le 2\pi.\)
\textbf{М}етод функции Грина (метод функции источника) состоит из двух
шагов.

\begin{enumerate}
\def\labelenumi{\arabic{enumi}.}
\item
  Помещаем в точку \((\rho, \varphi)\) точечный единичный заряд и
  находим потенциал \(G(r, \theta, \rho, \varphi)\) , который создается
  этим зарядом и обращается в нуль на границе.
\item
  Суммируем индивидуальные отклики \(G(r, \theta, \rho, \phi)\) ,
  взвешенные правой частью (плотностью зарядов) \(f(r, \theta)\) , по
  всему кругу. В результате получаем искомое решение
\end{enumerate}

\[
u\left(r,\;\theta\right) = \int_{0}^{2\pi} \int_{0}^{1} G\left(r,\;\theta,\;\rho,\;\varphi\right) f\left(\rho,\;\varphi\right) \rho \; d\rho \; d\varphi.
\]

Найдем теперь функцию источника для нашей задачи.

\section{\texorpdfstring{Определение функции источника
\(G(r, \theta, \rho, \phi)\) Замінимо праву частину \(f(r,\theta)\)
рівняння точковим джерелом величини 1, розташованим у деякій довольно
обраній точці \((\rho, \varphi)\) . Математично точкове джерело
представлене дельта-функцією \(\delta(r-\rho, \theta-\varphi)^1\) ). Ми
вважаємо, що дельта-функція дорівнює нулю всюди, крім точки
\((\rho, \varphi)\) , де розташований одиничний заряд. У мові сил
дельта-функцію можна інтерпретувати як точкова одиниця сили, прикладена
до точки \((\rho, \varphi)\) . Наше завдання --- знайти потенціал
точкового заряду, якщо відомо, що він дорівнює нулю на межі. Функція,
яка задовольняє ці умови, називається функцією Гріна або функцією
джерела. Вона дорівнює \emph{відповіді} системи в точці \((r, \theta)\)
на \emph{джерело}, розташоване в точці \((\rho, \varphi)\) . Складність
розв'язання цієї проблеми полягає в тому, що бажаний потенціал має
зрівнятися з нулем на кордоні. Якби не ця умова, потенціал можна було б
легко знайти, адже ми вже знаємо,
що}{Определение функции источника G(r, \textbackslash theta, \textbackslash rho, \textbackslash phi) Замінимо праву частину f(r,\textbackslash theta) рівняння точковим джерелом величини 1, розташованим у деякій довольно обраній точці (\textbackslash rho, \textbackslash varphi) . Математично точкове джерело представлене дельта-функцією \textbackslash delta(r-\textbackslash rho, \textbackslash theta-\textbackslash varphi)\^{}1 ). Ми вважаємо, що дельта-функція дорівнює нулю всюди, крім точки (\textbackslash rho, \textbackslash varphi) , де розташований одиничний заряд. У мові сил дельта-функцію можна інтерпретувати як точкова одиниця сили, прикладена до точки (\textbackslash rho, \textbackslash varphi) . Наше завдання --- знайти потенціал точкового заряду, якщо відомо, що він дорівнює нулю на межі. Функція, яка задовольняє ці умови, називається функцією Гріна або функцією джерела. Вона дорівнює відповіді системи в точці (r, \textbackslash theta) на джерело, розташоване в точці (\textbackslash rho, \textbackslash varphi) . Складність розв'язання цієї проблеми полягає в тому, що бажаний потенціал має зрівнятися з нулем на кордоні. Якби не ця умова, потенціал можна було б легко знайти, адже ми вже знаємо, що}}\label{ux43eux43fux440ux435ux434ux435ux43bux435ux43dux438ux435-ux444ux443ux43dux43aux446ux438ux438-ux438ux441ux442ux43eux447ux43dux438ux43aux430-gr-theta-rho-phi-ux437ux430ux43cux456ux43dux438ux43cux43e-ux43fux440ux430ux432ux443-ux447ux430ux441ux442ux438ux43dux443-frtheta-ux440ux456ux432ux43dux44fux43dux43dux44f-ux442ux43eux447ux43aux43eux432ux438ux43c-ux434ux436ux435ux440ux435ux43bux43eux43c-ux432ux435ux43bux438ux447ux438ux43dux438-1-ux440ux43eux437ux442ux430ux448ux43eux432ux430ux43dux438ux43c-ux443-ux434ux435ux44fux43aux456ux439-ux434ux43eux432ux43eux43bux44cux43dux43e-ux43eux431ux440ux430ux43dux456ux439-ux442ux43eux447ux446ux456-rho-varphi-.-ux43cux430ux442ux435ux43cux430ux442ux438ux447ux43dux43e-ux442ux43eux447ux43aux43eux432ux435-ux434ux436ux435ux440ux435ux43bux43e-ux43fux440ux435ux434ux441ux442ux430ux432ux43bux435ux43dux435-ux434ux435ux43bux44cux442ux430-ux444ux443ux43dux43aux446ux456ux454ux44e-deltar-rho-theta-varphi1-.-ux43cux438-ux432ux432ux430ux436ux430ux454ux43cux43e-ux449ux43e-ux434ux435ux43bux44cux442ux430-ux444ux443ux43dux43aux446ux456ux44f-ux434ux43eux440ux456ux432ux43dux44eux454-ux43dux443ux43bux44e-ux432ux441ux44eux434ux438-ux43aux440ux456ux43c-ux442ux43eux447ux43aux438-rho-varphi-ux434ux435-ux440ux43eux437ux442ux430ux448ux43eux432ux430ux43dux438ux439-ux43eux434ux438ux43dux438ux447ux43dux438ux439-ux437ux430ux440ux44fux434.-ux443-ux43cux43eux432ux456-ux441ux438ux43b-ux434ux435ux43bux44cux442ux430-ux444ux443ux43dux43aux446ux456ux44e-ux43cux43eux436ux43dux430-ux456ux43dux442ux435ux440ux43fux440ux435ux442ux443ux432ux430ux442ux438-ux44fux43a-ux442ux43eux447ux43aux43eux432ux430-ux43eux434ux438ux43dux438ux446ux44f-ux441ux438ux43bux438-ux43fux440ux438ux43aux43bux430ux434ux435ux43dux430-ux434ux43e-ux442ux43eux447ux43aux438-rho-varphi-.-ux43dux430ux448ux435-ux437ux430ux432ux434ux430ux43dux43dux44f-ux437ux43dux430ux439ux442ux438-ux43fux43eux442ux435ux43dux446ux456ux430ux43b-ux442ux43eux447ux43aux43eux432ux43eux433ux43e-ux437ux430ux440ux44fux434ux443-ux44fux43aux449ux43e-ux432ux456ux434ux43eux43cux43e-ux449ux43e-ux432ux456ux43d-ux434ux43eux440ux456ux432ux43dux44eux454-ux43dux443ux43bux44e-ux43dux430-ux43cux435ux436ux456.-ux444ux443ux43dux43aux446ux456ux44f-ux44fux43aux430-ux437ux430ux434ux43eux432ux43eux43bux44cux43dux44fux454-ux446ux456-ux443ux43cux43eux432ux438-ux43dux430ux437ux438ux432ux430ux454ux442ux44cux441ux44f-ux444ux443ux43dux43aux446ux456ux454ux44e-ux433ux440ux456ux43dux430-ux430ux431ux43e-ux444ux443ux43dux43aux446ux456ux454ux44e-ux434ux436ux435ux440ux435ux43bux430.-ux432ux43eux43dux430-ux434ux43eux440ux456ux432ux43dux44eux454-ux432ux456ux434ux43fux43eux432ux456ux434ux456-ux441ux438ux441ux442ux435ux43cux438-ux432-ux442ux43eux447ux446ux456-r-theta-ux43dux430-ux434ux436ux435ux440ux435ux43bux43e-ux440ux43eux437ux442ux430ux448ux43eux432ux430ux43dux435-ux432-ux442ux43eux447ux446ux456-rho-varphi-.-ux441ux43aux43bux430ux434ux43dux456ux441ux442ux44c-ux440ux43eux437ux432ux44fux437ux430ux43dux43dux44f-ux446ux456ux454ux457-ux43fux440ux43eux431ux43bux435ux43cux438-ux43fux43eux43bux44fux433ux430ux454-ux432-ux442ux43eux43cux443-ux449ux43e-ux431ux430ux436ux430ux43dux438ux439-ux43fux43eux442ux435ux43dux446ux456ux430ux43b-ux43cux430ux454-ux437ux440ux456ux432ux43dux44fux442ux438ux441ux44f-ux437-ux43dux443ux43bux435ux43c-ux43dux430-ux43aux43eux440ux434ux43eux43dux456.-ux44fux43aux431ux438-ux43dux435-ux446ux44f-ux443ux43cux43eux432ux430-ux43fux43eux442ux435ux43dux446ux456ux430ux43b-ux43cux43eux436ux43dux430-ux431ux443ux43bux43e-ux431-ux43bux435ux433ux43aux43e-ux437ux43dux430ux439ux442ux438-ux430ux434ux436ux435-ux43cux438-ux432ux436ux435-ux437ux43dux430ux454ux43cux43e-ux449ux43e}

См.: Владимиров В. С. Обобщенные функции в математической физике --- М.:
Наука, 1976. --- Прим ред.

потенциал точечного заряда равен

\[
\frac{1}{2\pi} \ln \left( \frac{1}{r} \right)
\]

,

где r --- расстояние от заряда до точки наблюдения.

Искомый потенциал может иметь следующий физический смысл:

\begin{enumerate}
\def\labelenumi{\arabic{enumi}.}
\tightlist
\item
  В теории теплопроводности --- это равновесная температура внутри
  круга, если тепловой источник размещается в точке (о. ф),
\end{enumerate}

а на границе поддерживается ну-

левая температура.

\begin{itemize}
\tightlist
\item
  \begin{enumerate}
  \def\labelenumi{\arabic{enumi}.}
  \setcounter{enumi}{1}
  \tightlist
  \item
    Він описує відхилення мембрани від положення рівноваги, якщо
    відхилення на межі дорівнюють нулю, а в точці (p.~f) мембрана
    зігнута до дуже великої висоти.
  \end{enumerate}
\item
  \begin{enumerate}
  \def\labelenumi{\arabic{enumi}.}
  \setcounter{enumi}{2}
  \tightlist
  \item
    В електростатиці це розподіл потенціалу всередині кола, якщо
    точковий позитивний заряд розташований у точці \((\rho, \phi)\) , а
    коло заземлене, а його потенціал підтримується на нулі.
  \end{enumerate}
\end{itemize}

\pandocbounded{\includegraphics[keepaspectratio]{images/_page_275_Picture_11.jpeg}}

Единичная окружность РИС. 36 3. Функция Гряна 6 (г. 0. р, ф) точечного
источника, расположенного в точке (р, ф).

Теперь найдем функцию Грина; она будет похожа на функцию, изображенную
на рис. 36.3.

\section{Построение
решения}\label{ux43fux43eux441ux442ux440ux43eux435ux43dux438ux435-ux440ux435ux448ux435ux43dux438ux44f}

ШАГ 1. Поскольку функция

\[
\frac{1}{2\pi} \ln \frac{1}{R}
\]

является потенциалом поля в точке \(p = (r, \varphi) \theta)\) , которое
создано единичным точечным зарядом, помещенным в точку
\(Q = (\rho, \phi)\) , \((R-paccmonnue\ between\ points\ P\ and\ Q)\)
единственное, что нам осталось сделать---так изменить эту функцию, чтобы
она обратилась в нуль на границе.

ЩАГ 2. Из эксперимента известно, что эквипотенциальные линии точечных
зарядов (положительных или отрицательных) представляют собой окружности
(рис. 36.4).

Основна ідея побудови функції Гріна така: потрібно розмістити інший
(від'ємний) заряд поза колом так, щоб загальний потенціал на колі r = 1
був сталим. Для гему можна відняти цю константу з потенціалу і отримати
нульовий потенціал на межі. Очевидно. Таким чином знайдений потенціал
матиме всі властивості функції \(G(r, \theta, \rho, \phi)\) .
Залишається незрозумілим, у якій точці за межами кола слід розмістити
негативний заряд, щоб потенціал на колі став сталим. Можливо порівняльно
P И С

\pandocbounded{\includegraphics[keepaspectratio]{images/_page_276_Picture_3.jpeg}}

. 36 4. Потенциал поля, образованного двумя противоположно заряженными
частицами; a, b--- линии (окружности) постоянного потенциала.

\pandocbounded{\includegraphics[keepaspectratio]{images/_page_276_Picture_5.jpeg}}

\textbf{РИС.} 36.5. Заряды Q в \(\overline{Q}\) дают постоянный
потенциал на окружности r=1.

Легко побачити, що якщо негативний заряд розміщений у точці
\(\bar{Q} = (\bar{\rho}, \bar{\phi}) = (1/\rho, \phi)\) , то повний
потенціал

\[
u(r, \theta) = \frac{1}{2\pi} \ln 1/R - \frac{1}{2\pi} \ln 1/\overline{R}
\]

буде сталою на колі r=1 Тут відстані від зарядів до точки (g, c)
позначаються R і \(\overline{R}\) (див. рис. 36.5). Можна показати, що
величина потенціалу на колі r = 1 визначається виразом

\[
-\frac{1}{2\pi}\ln\rho
\]

(положительная константа).

С учетом этого результата строим функцию Грина

(36.3)

\[
G(r, \theta, \rho, \phi) = \frac{1}{2\pi} \ln 1/R - \frac{1}{2\pi} \ln 1/\overline{R} + \frac{1}{2\pi} \ln \rho,
\]

Потенциал положентельного заряда, расположенного в тенциал границы
постоянного в тенциал границы постоянного в тенциал границы постоянного
в тенциал границы постоянного в тенциал границы постоянного в тенциал
границы постоянного в тенциал границы постоянного в тенциал границы
постоянного в тенциал границы постоянного в тенциал границы постоянного
в тенциал границы постоянного в тенциал границы постоянного в тенциал
границы постоянного в тенциал границы постоянного в тенциал границы
постоянного в тенциал границы постоянного в тенциал границы постоянного
в тенциал границы постоянного в тенциал границы постоянного в тенциал
границы постоянного в тенциал границы постоянного в тенциал границы
постоянного в тенциал границы постоянного в тенциал границы постоянного
в тенциал границы постоянного в тенциал границы постоянного в тенциал
границы постоянного в тенциал границы постоянного в тенциал границы
постоянного в тенциал границы постоянного в тенциал границы постоянного
в тенциал границы постоянного в тенциал границы постоянного в тенциал
границы постоянного в тенциал границы постоянного в тенциал границы
постоянного в тенциал границы постоянного в тенциал границы постоянного
в тенциал границы постоянного в тенциал границы постоянного в тенциал
границы постоянного в тенциал границы постоянного в тенциал границы
постоянного в тенциал границы постоянного в тенциал границы постоянного
в тенциал границы постоянного в тенциал границы постоянного в тенциал
границы постоянного в тенциал границы постоянного в тенциал границы
постоянного в тенциал границы постоянного в тенциал границы постоянного
в тенциал границы постоянного в тенциал границы постоянного в тенциал
границы постоянного в тенциал границы постоянного в тенциал границы
постоянного в тенциал границы постоянного в тенциал границы постоянного
в тенциал границы постоянного в тенциал границы постоянного в тенциал
границы постоян границы постоян границы постоян границы постоян границы
постоян

где

\[
R = V \frac{r^2 - 2\rho r \cos(\theta - \phi) + \rho^2}{R},
\overline{R} = V \frac{r^2 - 2\frac{r}{\rho}\cos(\theta - \phi) + 1/\rho^2}{r^2 - 2\frac{r}{\rho}\cos(\theta - \phi) + 1/\rho^2}
\]

формулы, выражающие расстояние между двумя точками в полярных
координатах. Чтобы найти решение исходной задачи, осталось провести
наложение всех импульсных откликов системы. Итак, мы переходим к
последнему шагу.

ШАГ 3. Суперпозиция импульсных откликов Этот шаг совсем прост. Из
соотношения

\[
u(r, \theta) = \int_{0}^{2\pi} \int_{0}^{1} G(r, \theta, \rho, \varphi) f(\rho, \varphi) \rho d\rho d\varphi
\]

получаем

(36.4)

\[
u(r, \theta) = \frac{1}{2\pi} \int_{0}^{2\pi} \int_{0}^{1} \ln (\rho \overline{R}/R) f(\rho, \varphi) \rho d\rho d\varphi.
\]

Формула (36.4) виражає розв'язок задачі Діріхле для рівняння Пуассона
всередині кола через функцію Гріна цієї задачі. Якщо щільність заряду
f(r, \(\theta\)) відома, то інтеграл можна обчислити, наприклад,
чисельно.

\section{ЗАМЕЧАНИЯ}\label{ux437ux430ux43cux435ux447ux430ux43dux438ux44f-13}

\begin{enumerate}
\def\labelenumi{\arabic{enumi}.}
\tightlist
\item
  Задачу
\end{enumerate}

\[
\begin{array}{lll} (\text{УЧ\Pi}) & \Delta u = 0, & 0 < r < 1, & 0 \leqslant \theta < 2\pi, \\ (\text{ГУ}) & u(1, \theta) = g(\theta), & 0 \leqslant \theta \leqslant 2\pi, \end{array}
\]

також можна розв'язати за допомогою методу \emph{функції Гріна}. У цьому
випадку рішення фіксується у формі

\[
u(r, \theta) = \int_{0}^{2\pi} \frac{\partial G}{\partial r} (r, \theta, 1, \varphi) g(\varphi) d\varphi.
\]

Останнє відношення можна надати більш обчислювально корисну форму, явно
виставляючи вираз у \partial G/\partial r . Внаслідок цього ми отримуємо
інтегральну формулу Пуассона, знайдену раніше в лекції 33.

(36.5)

\[
u(r, \theta) = \frac{1}{2\pi} \int_{0}^{2\pi} \left[ \frac{1 - r^2}{1 - 2r\cos(\theta - \varphi) + r^2} \right] g(\varphi) d\varphi.
\]

\begin{enumerate}
\def\labelenumi{\arabic{enumi}.}
\setcounter{enumi}{1}
\tightlist
\item
  Решение общей задачи Дирихле
\end{enumerate}

\[
\begin{array}{ll} (\text{УЧП}) & \Delta u = f(r, \theta), & 0 < r < 1, & 0 \le \theta < 2\pi, \\ (\text{ГУ}) & u(1, \theta) = g(\theta), & 0 < \theta \le 2\pi, \end{array}
\]

находится как сумма решений (36.4) и (36.5).

\begin{enumerate}
\def\labelenumi{\arabic{enumi}.}
\setcounter{enumi}{2}
\item
  Метод функцій Гріна дозволяє отримати розв'язки багатьох задач у
  областях різних форм. Однак для кожної області і (точніше, для кожного
  оператора зліва від граничної умови) для кожного рівняння потрібно
  знайти власну функцію Гріна, і це не завжди легко зробити.
\item
  Чтобы реально найти решение по формуле (36.4), в большинстве случаев
  необходимо вычислять интегралы на ЭВМ.
\end{enumerate}

\section{ЗАДАЧИ}\label{ux437ux430ux434ux430ux447ux438-14}

Найти потенциал точечного источника в трехмерном пространстве.

\begin{verbatim}
2. Найти функцию Грина  G(x, y, \xi, \eta)  для уравнения Лапласа в верхней полуплоскости y > 0. Другими словами, найти потенциал в точке (x, y) в верхней полуплоскости, если заряд расположен в точке  (\xi, \eta)  , а потенциал на границе y = 0 равен нулю (см. рис.).
\end{verbatim}

\pandocbounded{\includegraphics[keepaspectratio]{images/_page_278_Figure_11.jpeg}}

\begin{verbatim}
УКАЗАНИЕ. Если отрицательный заряд поместить в точку  $\overline{Q}$ = (\xi,  -\eta)  , очевидно, що загальний потенціал на лінії y = 0 буде нульовим. Отже, функція Гріна збігається з потенціалом цих двох зарядів.

- 3. Використовуючи результати задачі 2, знайдемо розв'язок рівняння Пуассона  $\Delta$ u = -k  у верхній півплощині з умовою нульової межі.
\end{verbatim}

\begin{itemize}
\tightlist
\item
  \begin{enumerate}
  \def\labelenumi{\arabic{enumi}.}
  \setcounter{enumi}{3}
  \tightlist
  \item
    Як би ви побудували функцію Зеленого для першого квадранта
    x\textgreater0,~y\textgreater0 ?
  \end{enumerate}
\item
  \begin{enumerate}
  \def\labelenumi{\arabic{enumi}.}
  \setcounter{enumi}{4}
  \tightlist
  \item
    Інший підхід до розв'язання рівняння Пуассона такий. Припустимо, ви
    хочете розв'язати проблему
  \end{enumerate}
\end{itemize}

(УЧП)

\[
\Delta u = 1
,  0 < r < 1  ,  0 \le \theta < 2\pi  ,  (\Gamma Y)   u(1, \theta) = \sin \theta  ,  0 \le \theta \le 2\pi  .
\]

Сначала попытайтесь найти частное решение уравнения с помощью
подстановки

\[
u_p = (r, \theta) = Ar^2
\]

.

\begin{verbatim}
Затем подставьте это выражение в уравнение Пуассона и определите константу  A_{\cdot}  Після того, як  u_p  знайдено часткове рішення, необхідно шукати повне рішення у вигляді  u=w+u_p  . Яку граничну умову має задовольняти функція w? Знайдіть  w(r, $\theta$)  . Яке рішення початкової проблеми  u(r, $\theta$)  ? Подивись. Ретельно проаналізуйте відповідь. Дайте інтерпретацію кожного члена.
\end{verbatim}

\section{Часть 5. ЧИСЛЕННЫЕ И ПРИБЛИЖЕННЫЕ
МЕТОДЫ}\label{ux447ux430ux441ux442ux44c-5.-ux447ux438ux441ux43bux435ux43dux43dux44bux435-ux438-ux43fux440ux438ux431ux43bux438ux436ux435ux43dux43dux44bux435-ux43cux435ux442ux43eux434ux44b}

Лекция 37

\section{ЧИСЛЕННЫЕ РЕШЕНИЯ (ЭЛЛИПТИЧЕСКИЕ
ЗАДАЧИ)}\label{ux447ux438ux441ux43bux435ux43dux43dux44bux435-ux440ux435ux448ux435ux43dux438ux44f-ux44dux43bux43bux438ux43fux442ux438ux447ux435ux441ux43aux438ux435-ux437ux430ux434ux430ux447ux438}

МЕТА ЛЕКЦІЇ: Показати, як диференціальні рівняння в похідних похідних
можна звести до систем алгебраїчних рівнянь, якщо похідні з частковими
похідними замінити на наближення з скінченними різницями. Розв'язання
цієї системи алгебраїчних рівнянь ітеративним методом дозволяє отримати
наближене розв'язок

Дается информация о пакете прикладных программ ELLPACK, предназначенном
для решения эллиптических

задач общего вида на ЭВМ.

К настоящему моменту мы познакомились с несколькими методами решения
линейных уравнений с частными производными. Однако большинство
уравнений, с которыми мы имели дело, были очень простыми. Так же просты
были и формы областей, в которых решались задачи. Тем не менее многие
задачи нельзя упростить до такой степени, чтобы они свелись к заданному
набору шаблонных задач. Такие задачи следует решать приближенно
численными методами. Прогресс в области высокопроизводительных ЭВМ
привел к созданию новых численных методов. Сегодня уже решены такие
нелинейные задачи динамики жидкости, теории упругости и теории
потенциала, о решении которых десять лет назад никто даже не
задумывался.

Под общим названием «численные методы» объединены несколько различных
подходов к решению задач. Подробное описание всех этих подходов можно
найти в книге {[}4{]} из рекомендованного списка литературы. В этой и
двух последующих лекциях будет показано, как решать эллиптические,
гиперболические и параболические уравнения методом конечных разностей.

Сначала познакомимся с понятием конечных разностей, а затем покажем, как
с их помощью решить задачу Дирихле в квадрате.

\section{Конечно-разностные
аппроксимации}\label{ux43aux43eux43dux435ux447ux43dux43e-ux440ux430ux437ux43dux43eux441ux442ux43dux44bux435-ux430ux43fux43fux440ux43eux43aux441ux438ux43cux430ux446ux438ux438}

Вспомним ряд Тейлора для функции \(f(x)\):

\[
f(x+h) = f(x) + f'(x)h + \frac{f''(x)}{2!}h^2 + \dots
\]

Если оборвать этот ряд на втором члене, то получим

\[
f(x+h) \cong f(x) + f'(x) h
\]

.

Откуда

\[
(37.1) f'(x) \cong \frac{f(x+h)-f(x)}{h}.
\]

Выражение, стоящее в правой части, называется правой разностной
производной. Она аппроксимирует первую производную f'(x) в точке x.

\begin{verbatim}
В разложении Тейлора функции  f(\mathbf{r})  можно заменить h на -h и получить левую разностную производную
\end{verbatim}

(37.2)

\[
f'(x) \cong \frac{f(x) - f(x - h)}{h}
\]

.

\[
Buyutan  f(x - h) \cong f(x) - f'(x) h  \nif  f(x + h) \cong f( x) + f'(x) h  ,
\]

получаем центральную разностную производную

(37.3)

\[
f'(x) \simeq \frac{1}{2h} [f(x+h) - f(x-h)].
\]

Якщо в ряду Тейлора залишиться ще один член, то центральну різницю
похідну для наближення f'\,'(x) можна отримати точно таким самим
способом

(37.4)

\[
f''(x) \cong \frac{1}{h^2} [f(x+h) - 2f(x) + f(x-h)].
\]

Тепер ми можемо розширити поняття наближення скінченних різниць до
\emph{часткових похідних}. Якщо виходити з розкладу Тейлора функції двох
змінних

\[
u(x+h, y) = u(x, y) + u_x(x, y)h + u_{xx}(x, y)\frac{h^2}{2!} + \dots,
u(x-h, y) = u(x, y) - u_x(x, y)h + u_{xx}(x, y)\frac{h^2}{2!} - \dots,
\]

можно получить следующие аппроксимации частных производных:

\[
\begin{split} u_x(x,\ y) & \cong \frac{u(x+h,\ y)-u(x,\ y)}{h}, \\ u_{xx}(x,\ y) & \cong \frac{1}{h^2} \big[ u(x+h,\ y)-2u(x,\ y)+u(x-h,y) \big], \\ u_u(x,\ y) & \cong \frac{u(x,\ y+k)-u(x,\ y)}{k}, \\ u_{uy}(x,\ y) & \cong \frac{1}{h^2} \big[ u(x,\ y+k)-2u(x,\ y)+u(x,\ y-k) \big]. \end{split}
\]

В этих формулах частные производные аппроксимируются правыми,
центральными и левыми разностными производными, но в данной лекции мы
будем пользоваться только центральными аппроксимациями.

\pandocbounded{\includegraphics[keepaspectratio]{images/_page_282_Figure_5.jpeg}}

РИС. 37.1 Сетка для задачи Дирихле внутри квадрата.

Чтобы познакомиться с основными правилами использования
конечно-разностных аппроксимаций, рассмогрим простую задачу Дирихле.

Решение задачи Дирихле методом конечных разностей

Пусть необходимо решить задачу Дирихле

(ЗЧП)

\[
u_{xx} + u_{yy} = 0
\]

, \(0 < x < 1\) , \(0 < y < 1\) , \(u = 0\) , на верхней и боковых
сторонах квадрата, \(u(x, 0) = \sin(x\pi)\) , \(0 \le \le x \le 1\) ,
\(y = 0\) .

Построим в плоскости (x, y) сетку, как показано на рис. 37.1. Удобно
(особенно, если мы собираемся применять ЭВМ)

непользовать следующие обозначения:

\[
\begin{split} u\left(x,\,y\right) &= u_{ij}\,,\\ u\left(x,\,y+k\right) &= u_{i+1,j}\,,\\ u\left(x,\,y-k\right) &= u_{i-1,j}\,,\\ u\left(x-h,\,y\right) &= u_{i,j-1}\,,\\ u\left(x+h,\,y\right) &= u_{i,j+1}\,,\\ u_{x}\left(x,\,y\right) &= \frac{1}{2h}\left(u_{i,j+1}-u_{i,j-1}\right)\,,\\ u_{y}\left(x,\,y\right) &= \frac{1}{2k}\left(u_{i+1,j}-u_{i-1,j}\right)\,,\\ u_{xx}\left(x,\,y\right) &= \frac{1}{h^{2}}\left(u_{i,j+1}-2u_{i,j}+u_{i,j-1}\right)\,,\\ u_{yy}\left(x,\,y\right) &= \frac{1}{h^{2}}\left(u_{i+1,j}-2u_{i,j}+u_{i-1,j}\right)\,. \end{split}
\]

Основная идея решения задачи Дирихле основана на замене частных
производных в уравнении Лапласа

\[
u_{xx} + u_{yy} = 0
\]

их конечно-разностными аппроксимациями. Проделав это и воспользовавшись
компактными обозначениями, получаем следующее разностное уравнение:

\[
\Delta u = \frac{1}{h^2} (u_{i,j+1} - 2u_{i,j} + u_{i,j-1}) + \frac{1}{k^2} (u_{i+1,j} - 2u_{i,j} + u_{i-1,j}) = 0.
\]

В случае когда h=k, уравнение Лапласа приводится к виду (37.6)
\((u_{i+1,j}+u_{i-1,j}+u_{i,j+1}+u_{i,j+1}-4u_{ij})=0.\) Разрешив его
относительно \(u_{II}\) , получим

\[
u_{i,j} = \frac{1}{4} (u_{i+1,j} + u_{i-1,j} + u_{i,j+1} + u_{i,j-1}).
\]

Зверніть увагу, що в цьому співвідношенні всі значення \(u_{tj}\) беруть
для внутрішніх вузлів сітки. Згідно з останнім співвідношенням,
розв'язок \(u_{ij}\) апроксимується середнім значенням розв'язку у
чотирьох сусідніх точках. Відповідно, ми можемо розробити чисельний
метод для розв'язання цієї задачі.

Алгоритм численного решения задачи Дирихле (метод Либманна)

ШАГ 1. Присвоим величинам \(u_{ij}\) во внутренних узлах численное
вначение, равное среднему всех значений граничных условий.

КРОК 2. Ми перерахуємо значення у всіх внутрішніх точках сітки,
замінивши старе значення на середнє значення чотирьох суміжних точок. Не
дуже важливо, як буде організовано підрахунок процесу, але зазвичай він
здійснюється рядками (або колонками). Після кількох операцій процес
збігається до приблизного розв'язання задачі. Швидкість збіжності цього
алгоритму низька, але її можна збільшити кількома способами.
Рекомендуємо зацікавленим читачам звернутися до книги {[}4{]} зі списку
рекомендованої літератури.

\section{ЗАМЕЧАНИЯ}\label{ux437ux430ux43cux435ux447ux430ux43dux438ux44f-14}

\begin{enumerate}
\def\labelenumi{\arabic{enumi}.}
\tightlist
\item
  Если взять сетку с чегырьмя внутренними точками (m=n=4), то система
  (37.6) запишется в виде
\end{enumerate}

\[
-4u_{22} + 0 + \sin(\pi/3) + u_{23} + u_{32} = 0,
-4u_{23} + u_{22} + \sin(2\pi/3) + 0 + u_{33} = 0,
-4u_{32} + 0 + u_{22} + u_{33} + 0 = 0,
-4u_{33} + u_{32} + u_{23} + 0 + 0 = 0.
\]

Це потрібно вирішувати у зв'язку з \(u_{22}\) , \(u_{21}\) , \(u_{32}\)
та \(u_{33}\) . Розв'язок такої системи можна знайти за допомогою
\emph{ітеративного методу}. Метод Лібмана також є одним із ітеративних
методів.

\begin{itemize}
\tightlist
\item
  \begin{enumerate}
  \def\labelenumi{\arabic{enumi}.}
  \setcounter{enumi}{1}
  \tightlist
  \item
    Якщо зменшити кроки сітки на h і k (кількість вузлів сітки
    збільшиться), то отримаємо систему рівнянь типу (37,7), але більших
    розмірів.
  \end{enumerate}
\item
  \begin{enumerate}
  \def\labelenumi{\arabic{enumi}.}
  \setcounter{enumi}{2}
  \tightlist
  \item
    Систему рівнянь (37.7) можна переписати у матричній формі
  \end{enumerate}
\end{itemize}

\[
\begin{bmatrix} -4 & 1 & 1 \\ 1 & -4 & 0 & 1 \\ 1 & 0 & -4 & 1 \\ 0 & 1 & 1 & -4 \end{bmatrix}
\begin{bmatrix} u_{22} \\ u_{23} \\ u_{32} \\ u_{33} \end{bmatrix}
=
\begin{bmatrix} -\sin{(\pi/3)} \\ -\sin{(2\pi/3)} \\ 0 \\ 0 \end{bmatrix}
=
\begin{bmatrix} -0.86 \\ -0.86 \\ 0 \\ 0 \end{bmatrix}
\]

.

С ростом числа уравнений (например, до 1000) матрица коэффициентов
становится разреженной, т. е. содержит много нулей. Системы с
разреженными матрицами решаются специальными методами. Обычно
используются итерационные методы, такие, как метод Якоби, метод
Гаусса---Зейделя или метод последовательной верхней релаксации.

\begin{enumerate}
\def\labelenumi{\arabic{enumi}.}
\setcounter{enumi}{3}
\item
  При решении задачи Неймана те \emph{производные}, которые входят в
  граничное условие, также следует заменить их разностными
  аппроксимациями.
\item
  Подобным образом можно решать задачи
\end{enumerate}

\begin{enumerate}
\def\labelenumi{\alph{enumi})}
\tightlist
\item
\end{enumerate}

\[
u_{xx} + u_{yy} = f(x, y)
\]

\begin{enumerate}
\def\labelenumi{\arabic{enumi})}
\setcounter{enumi}{5}
\tightlist
\item
\end{enumerate}

\[
xu_{xx} + u_{yy} + 2u = \sin(x - y)
\]

\begin{enumerate}
\def\labelenumi{\Alph{enumi})}
\setcounter{enumi}{1}
\tightlist
\item
\end{enumerate}

\[
\sin xu_{xx} + u_{xy} + 3u = 0
\]

(неоднородные уравнения с переменными коэффициентами), (уравнения с
переменными коэффициентами).

\begin{enumerate}
\def\labelenumi{\arabic{enumi}.}
\setcounter{enumi}{5}
\tightlist
\item
  Если область, в которой решается задача, имеет неправильную форму,
  можно покрыть ее сеткой и аппроксимировать решение в ближайших к
  границе точках интерполяцией граничных условий. После того как это
  сделано, задача решается обычным образом (см. рис. 37.2).
\end{enumerate}

\pandocbounded{\includegraphics[keepaspectratio]{images/_page_285_Figure_8.jpeg}}

□ новые ГУ находятся интерполяцией ГУ на кривой

РИС, 37 2.

\begin{itemize}
\tightlist
\item
  \begin{enumerate}
  \def\labelenumi{\arabic{enumi}.}
  \setcounter{enumi}{6}
  \tightlist
  \item
    В некоторых журналах публикуются тексты программ для решения
    уравнений с частными производными на ЭВМ. Ниже приведен краткий
    перечень этих журналов:
  \end{enumerate}
\item
  \begin{enumerate}
  \def\labelenumi{\alph{enumi})}
  \tightlist
  \item
    ACM Transactions on Mathematical Software,
  \end{enumerate}
\item
  \begin{enumerate}
  \def\labelenumi{\arabic{enumi})}
  \setcounter{enumi}{5}
  \tightlist
  \item
    Computer Journal,
  \end{enumerate}
\item
  \begin{enumerate}
  \def\labelenumi{\Alph{enumi})}
  \setcounter{enumi}{1}
  \tightlist
  \item
    Numerische Mathematik,
  \end{enumerate}
\item
  \begin{enumerate}
  \def\labelenumi{\alph{enumi})}
  \setcounter{enumi}{17}
  \tightlist
  \item
    BIT
  \end{enumerate}
\end{itemize}

Кроме того, недавно был разработан пакет прикладных программ ELLPACK,
предназначенный для решения достаточно широкого класса краевых
эллиптических задач. Этот пакет позволит решать разнообразные двух- и
трехмерные задачи в

различных системах координат для произвольных границ и общих граничных
условий. Пользователю представляется большой ассортимент различных
методов решения 1).

\subsubsection{ЗАДАЧИ}\label{ux437ux430ux434ux430ux447ux438-15}

\begin{itemize}
\item
  1 Получите аппроксимацию (39.4) для второй производной f'\,`(x)
  f'\,'(x) = \{\(f(x+h)\) 2f(x) + \(f(x-h)\)\}/h\^{}2 .
\item
  \begin{enumerate}
  \def\labelenumi{\arabic{enumi}.}
  \setcounter{enumi}{1}
  \tightlist
  \item
    Проделайте две итерации для задачи Дирихле (37.5), используя
    итерационный метод Либманна.
  \end{enumerate}
\item
  \begin{enumerate}
  \def\labelenumi{\arabic{enumi}.}
  \setcounter{enumi}{2}
  \tightlist
  \item
    К какой алгебранческой системе сведется задача Дирихле для уравнений
    Пуассона в квадрате
  \end{enumerate}
\end{itemize}

(УЧ\Pi)

\(u_{xx}\) +
\(u_{yy} =\)f(x,y)\(, \quad 0 < x < 1, \quad 0 < y < 1, (\Gamma Y)\)u(x,y)\(=\)g(x,y)\$\$на
границе, если ее решать ме-

тодом конечных разностей? 4. Решите задачу 3, если

(УЧП)

\[
u_{xx} + u_{yy} + 2u = 0
\]

, 0 \textless{} x \textless{} 1 , 0 \textless{} y \textless{} 1 , (ГУ)
\(u(x, y)\) = \(g(x, y)\) на границе.

\begin{enumerate}
\def\labelenumi{\arabic{enumi}.}
\setcounter{enumi}{4}
\tightlist
\item
  Как бы вы решали задачу Неймана внутри квадрата
\end{enumerate}

(УЧП)

\[
u_{xx} + u_{yy} = 0
\]

, 0 \textless{} x \textless{} 1 , 0 \textless{} y \textless{} 1 , на
верхней, нижней и левой стороне квадрата,

\[
\begin{cases} u = 0 & \text{квадрата,} \\ \frac{\partial u}{\partial x}(1, y) = 1, & 0 \leqslant y \leqslant 1, \end{cases}
\]

методом конечных разностей?

\begin{enumerate}
\def\labelenumi{\arabic{enumi}.}
\setcounter{enumi}{5}
\tightlist
\item
  Постройте блок-схему алеоритма решения задачи Дирихле в квадрате
\end{enumerate}

\[
(\text{УЧП})
\]

\(u_{xx}\) + \(u_{yy} =\)f(x,
y)\texttt{,\ \textbackslash{}quad\ 0\ \textless{}\ x\ \textless{}\ 1,\ \textbackslash{}quad\ 0\ \textless{}\ y\ \textless{}\ 1,\ (\textbackslash{}text\{}ГУ\})
\(u(x, y)\) = \(g(x, y)\) на границе,

\begin{itemize}
\item
  с произвольным числом узлов сетки. Если вы знаете какойнибудь язык
  программирования, то напишите программу для выполнения этих
  вычислений.
\item
  \begin{enumerate}
  \def\labelenumi{\arabic{enumi})}
  \tightlist
  \item
    На русском языке работы по численным методам решения уравнений с
    частными производными публикуются в многочисленных журналах и
    сборниках. Например: «Журнал вычислительной математики и
    математической физики», изд-во «Наука», «Инженерно-физический
    журнал», Изд-во АН БССР, Сборники серии «Вычислительные методы и
    программирование», Изд-во МГУ. Пакетам прикладных задач посвящено
    много публикаций, в основном статьи. В качестве примера можно
    указать: Горбунов-Посадов М. М., Карпов В. Я, Корягин и др. Пакет
    САФРА. Программное обеспечение вычислительного эксперимента. В кн.:
    Пакеты прикладных программ. Вычислительный эксперимент. М.: Наука,
    1983. Прим. перев.
  \end{enumerate}
\end{itemize}

\section{ЯВНЫЕ РАЗНОСТНЫЕ
СХЕМЫ}\label{ux44fux432ux43dux44bux435-ux440ux430ux437ux43dux43eux441ux442ux43dux44bux435-ux441ux445ux435ux43cux44b}

ЦЕЛЬ ЛЕКЦИИ: Ввести понятие явной разностной схемы и показать, как с ее
помощью можно решать гипер-болические и параболические задачи. Основная
идея состоит в том, что после замены уравнения типа

\[
u_t = u_{xx}
\]

его конечно-разностной аппроксимацией получаются формулы, явно
выражающие значения решения для одного момента времени через значения
решения в предыдущий момент времени. Таким образом, смешанную задачу для
параболического или гиперболического уравнения можно решать,
последовательно вычисляя решения для всех последующих моментов времени.

Явная схема не свободна от недостатков. Если необходимо увеличить
точность аппроксимации производных, то при изменении сегки растет не
только объем вычислений, но и ошибки округления.

Предыдущая лекция была посвящена решению краевых эллиптических
(стационарных задач). В таких задачах необходимо найти решение уравнения
с частными производными в данной области пространства, если на границе
области решение или его производная заданы, Приближенное решение
эллиптических задач сводится к решению системы алгебраических уравнений
для значений функции во внутренних узлах сетки. Другими словами,
значения решения во всех внутренних узлах определяются одновременно.

\begin{verbatim}
В этой лекции мы займемся разностными методами решения задач, зависящих от времени. Основная идея состоит в том, что если известно решение в начальный момент времени, то можно по схеме бегущего счета найти решение при  t=$\Delta$ t,\ 2$\Delta$ t,\ 3$\Delta$ t,\ldots  Заменяя частные производные по времени и по пространственной переменной конечно-разностными производными, можно получить явные выражения для  u_{tj}  через значения функции u в предыдущие моменты времени. Такой процесс называется явной схемой бегущего счета.
\end{verbatim}

Чтобы показать этот метод в действии, рассмотрим типичную задачу
теплопроводности.

\section{Явная схема для уравнения
теплопроводности}\label{ux44fux432ux43dux430ux44f-ux441ux445ux435ux43cux430-ux434ux43bux44f-ux443ux440ux430ux432ux43dux435ux43dux438ux44f-ux442ux435ux43fux43bux43eux43fux440ux43eux432ux43eux434ux43dux43eux441ux442ux438}

Рассмотрим задачу теплопроводности в стержне, начальная температура
которого равна нулю. Пусть температура левого конца фиксирована, а на
правом конце происходит теплообмен с окружающей средой, так что тепловой
полок пропорционален Р

\pandocbounded{\includegraphics[keepaspectratio]{images/_page_288_Figure_3.jpeg}}

ИС. 38.1. Сетка для уравнения теплопроводности.

\begin{verbatim}
разности температур конца стержия и среды. Пусть температура среды определяется функцией  g\left(t\right) . Другими словами, мы решаем задачу
\end{verbatim}

задачу (УЧП)

\[
u_t = u_{xx}, \quad 0 < x < 1, \quad 0 < t < \infty,
\]

(38.1) (ГУ)

\[
\begin{cases} u(0, t) = 1, \\ u_x(1, t) = -[u(1, t) - g(t)], \\ u(x, 0) = 0, \quad 0 \le x \le 1. \end{cases}
\]

Для решения этой задачи методом конечных разностей построим
прямоугольную сетку, узлы которой определяются формулами (см. рис. 38.1)

\[
x_j = jh
,  j = 0, 1, 2, ..., n  ,  y_i = ik  ,  i = 0, 1, 2, ..., m  .
\]

Отметим, что значения \(u_{ij}\) на левой и нижней сторонах сетки на
рис. 38.1 известны из граничных и начальных условий и наша вадача
состоит в отыскании остальных значений \(u_{ij}\) . Для решения этой
задачи заменим частные производные в уравнении теплопроводности их
конечно-разностными аппроксимациями

\[
u_{t} = \frac{1}{k} [u(x, t+k) - u(x, t)] = \frac{1}{k} [u_{t+1,j} - u_{t,j}],
u_{xx} = \frac{1}{h^{2}} [u(x+h, t) - 2u(x, t) + u(x-h, t)] = \frac{1}{h^{2}} [u_{t,j+1} - 2_{t,j} + u_{t,j-1}].
\]

Подставим эти выражения в уравнение \(u_t = u_{xx}\) и разрешим
получившееся уравнение относительно значений функции на верхнем
временном слое. В результате получаем

(38.2)

\[
u_{t+1,j} = u_{t,j} + \frac{k}{h^2} \left[ u_{t,j+1} - 2u_{t,j} + u_{t,j-1} \right].
\]

Це та формула, яку ми шукаємо, оскільки вона виражає розв'язок у певний
момент часу через розв'язок у попередній момент часу (індекс \emph{i}
означає часову змінну). Фіг. 38.1 підкреслює значення, включені до цієї
формули.

Тепер можна починати розрахунки. Однак спочатку потрібно апроксимувати
похідну в граничній умові праворуч

\[
a_x(1, t) = -[u(1, t) - g(t)].
\]

В результате аппроксимации получаем

(38.3)

\[
\frac{1}{h}[u_{t,n} - u_{t,n-1}] = -[u_{t,n} - g_t].
\]

где значения \(g_i = g(ik)\) известны. Здесь мы заменили \(u_x(1, t)\)
левой разностной производной, поскольку правая разностная производная
потребовала бы значений функции за пределами сетки. Из (38.3) находим

(38.4)

\[
u_{t,n} = \frac{u_{t,n-1} - hg_t}{1 + h}.
\]

Формулы (38.2) и (38.4) позволяют начать вычисления.

Алгоритм вычислений по явной схеме

ШАГ 1. Находим решение на сеточном слое \(t = \Delta t\) , используя
явную формулу (см. рис. 38.2)

\[
u_{2,j} = u_{1,j} + \frac{k}{h^2} [u_{1,j+1} - 2u_{1,j} + u_{1,j-1}], \quad j = 2, 3, \ldots, n-1.
\]

ШАГ 2. Величину \(u_{2,n}\) находим по формуле (38.4)

\[
u_{2,n} = \frac{u_{2,n-1} + hg_2}{1+h}.
\]

Після виконання кроків 1 і 2 ми отримуємо рішення для \(t=\Delta t\) .
Щоб отримати розв'язок у \(t=2\Delta t\) (другий рядок знизу на рис.
38.2), повторіть кроки 1 і 2, рухаючись на одну лінію вгору, тобто
збільшуючи t на 1 і використовуючи \(u_{t,j}\) з попередньої лінії.
Розв'язок обчислюється таким самим способом у наступні
\(t=3\Delta t,\ 4\Delta t\dots\) Для того чтобы помочь читателю в
проведении расчетов по этому методу, на рис. 38.3 приведена достаточно
подробная блоксхема алгоритма. Те, кто не очень хорошо знаком с
блок-схемами.

\pandocbounded{\includegraphics[keepaspectratio]{images/_page_290_Figure_3.jpeg}}

РИС. 38.2. Иллюстрация явной разностной схемы: n--- число узлов сетки
вдоль оси \(x;\ m\) --- число узлов сетки вдоль оси
\(y;\ h=1/(m-1),\ k=T/(m-1)\) могут счигать их промежуточным звеном
между вычислительными алгоритмами и подробными программами для ЭВМ.
Блок-схемы дают строгое описание процесса вычислений.

\section{ЗАМЕЧАНИЯ}\label{ux437ux430ux43cux435ux447ux430ux43dux438ux44f-15}

\begin{itemize}
\tightlist
\item
  Явна схема має серйозний недолік. Якщо часовий крок достатньо великий
  порівняно з кроком x, похибки округлення можуть стати настільки
  великими, що отриманий розв'язок втрачає сенс. Відношення кроків у t і
  x залежить від рівняння та граничних умов, але в загальному випадку
  крок у часі має бути значно меншим за крок у координаті. У книзі
  {[}4{]} доведено, що для застосування явної схеми необхідно виконати
  умову k/h2 ≤ 0,5.
\item
  \begin{enumerate}
  \def\labelenumi{\arabic{enumi}.}
  \setcounter{enumi}{1}
  \tightlist
  \item
    Існує таке правило: якщо зменшити кількість кроків \(\Delta t\) і
    \(\Delta x\) , то похибка у наближенні часткових похідних за
    скінченними різницями також зменшиться, але чим менша сітка, тим
    більше обчислень потрібно виконати, і, відповідно, тим більшими
    будуть похибки округлення. Це явище ілюструється на рис. 38.5.
  \end{enumerate}
\item
  \begin{enumerate}
  \def\labelenumi{\arabic{enumi}.}
  \setcounter{enumi}{2}
  \tightlist
  \item
    Для гіперболічної задачі
  \end{enumerate}
\end{itemize}

(УЧП)

\[
u_{tt} = u_{xx}, \quad 0 < x < 1, \quad 0 < t < \infty,
\]

(ГУ)

\[
\begin{cases} u(0, t) = g_1(t), \\ u(1, t) = g_2(t), \end{cases}
\quad 0 < t < \infty,
\]

(HY)

\[
\begin{cases} u(x, 0) = \varphi(x), \\ u_t(x, 0) = \psi(x), \end{cases}
\quad 0 \leqslant x \leqslant 1,
\]

Також можливо побудувати явну різницю схеми. Для цього ми апроксимуємо
\(u_{tt}\) і \(u_{xy}\) центральних різницьких пропозицій И С

\pandocbounded{\includegraphics[keepaspectratio]{images/_page_291_Figure_4.jpeg}}

. 38.3. Алгоритм решения задачи по явной схеме.

водными

\[
u_{tt} \simeq \frac{1}{h^2} [u(x, t+k) - 2u(x, t) + u(x, t-k)],
u_{xx} \simeq \frac{1}{h^2} [u(x+h, t) - 2u(x, t) + u(x-h, t)].
\]

а начальное условие--- по схеме

\[
u_t(x, 0) \approx \frac{1}{k} [u(x, k) - u(x, 0)] = \frac{1}{k} [u(x, k) - \varphi(x)].
\]

В результате для вычисления величины \(u(x, t+k)\) получаем следующую
явную схему:

(38.5)

\[
u(x, t+k) = 2u(x, t) - u(x, t-k) + + (k/h)^2 [u(x+h, t) - 2u(x, t) + u(x-h, t)].
\]

З (38.5) видно, що для обчислення розв'язку на заданому часовому шарі
необхідно знати розв'язок на двох попередніх шарах. Отже, щоб почати
рахувати, потрібно використовувати P И С

\pandocbounded{\includegraphics[keepaspectratio]{images/_page_292_Figure_6.jpeg}}

. 38.4. Использование массивов при реализации явной схемы.

начальным условнем на скорость

\[
\frac{1}{k}\left[u\left(x,\,k\right)-\varphi\left(x\right)\right]=\psi\left(x\right),
\]

з якої ми виводимо \(u(x, k)\) = \varphi(x) + k\psi(x) , тобто цінність
розв'язку на t = \Delta t . Розв'язок у наступні часи можна знайти за
явною формулою (38.5).

\section{ЗАДАЧИ}\label{ux437ux430ux434ux430ux447ux438-16}

\begin{enumerate}
\def\labelenumi{\arabic{enumi}.}
\tightlist
\item
  Постройте явную конечно-разностную схему для задачи
\end{enumerate}

(УЧП)

\[
u_t = u_{xx}
,  0 < x < 1  ,  0 < t < \infty  ,
\]

(ГУ)

\[
\begin{cases} u(0, t) = 0, \\ u(1, t) = 0, \\ u(x, 0) = \sin(\pi x), \quad 0 \le x \le 1. \end{cases}
(HУ)  u(x, 0) = \sin(\pi x), \quad 0 \le x \le 1.
\]

\begin{verbatim}
Полагая  h = $\Delta$ x = 0.1 , найдите решение при  t_1 = 0.005 ,  t_2 = 0.010 ,  t_3 = 0.015 . Постройте график полученного решения на сетке x = 0; 0.1; 0.2; 0.3; . .; 0.9: 1 при t = 0.015.
\end{verbatim}

\pandocbounded{\includegraphics[keepaspectratio]{images/_page_293_Figure_6.jpeg}}

РИС. 38.5. Зависимость полной ошибки от шага сетки.

\begin{itemize}
\tightlist
\item
  \begin{enumerate}
  \def\labelenumi{\arabic{enumi}.}
  \setcounter{enumi}{1}
  \tightlist
  \item
    Решите задачу 1 аналитически (методом разделения переменных).
    Найдите значение решения в узлах сетки x=0; 0,1; 0,2; 0,3; \ldots;
    0,9; 1 при t=0,015. Сравните эти результаты с результатами
    численного решения задачи 1. (Для табулирования аналитического
    решения можно воспользоваться калькулятором или написать маленькую
    программу для ЭВМ.)
  \end{enumerate}
\item
  \begin{enumerate}
  \def\labelenumi{\arabic{enumi}.}
  \setcounter{enumi}{2}
  \tightlist
  \item
    В соответствии с замечанием 3 постройте блок-схему решения
    гиперболической задачи.
  \end{enumerate}
\item
  \begin{enumerate}
  \def\labelenumi{\arabic{enumi}.}
  \setcounter{enumi}{3}
  \tightlist
  \item
    Решите задачу 1 с новым граничным условием при x=1 \(u_x\)(1, t) =
    -{[}\(u(1, t)\) 1{]}.
  \end{enumerate}
\end{itemize}

\section{НЕЯВНЫЕ РАЗНОСТНЫЕ СХЕМЫ (СХЕМА
КРАНКА---НИКОЛЬСОНА)}\label{ux43dux435ux44fux432ux43dux44bux435-ux440ux430ux437ux43dux43eux441ux442ux43dux44bux435-ux441ux445ux435ux43cux44b-ux441ux445ux435ux43cux430-ux43aux440ux430ux43dux43aux430ux43dux438ux43aux43eux43bux44cux441ux43eux43dux430}

ЦЕЛЬ ЛЕКЦИИ: Показать, как можно решать задачи, содержащие функции,
зависящие от времени, с помощью неявной разностной схемы. В этой схеме,
как и рачьше, частные производные заменяются конечно-разностными
аппроксимациями, но \(u_{l+1,\ J}\) уже не выражаются в явном виде через
значения на предыдущих слоях. Теперь для определения величин
\(u_{l+1,\ J}\) необходимо решить систему уравнений. Другими словами,
для определения решения на каждом временном наге необходимо решать
систему уравнений.

Преимущество неявных схем перед явными в том, что в неявных схемах шаг
сетки можно сделать достаточно большим, не опасаясь, что ошибки
округления ``разрушат'' решение.

Для решения параболической задачи мы воспользуемся широко известной
\emph{схемой Кранка} --- \emph{Никольсона}.

Как уже говорилось в предыдущей лекции, для надежной работы явной схемы
необходимо, чтобы временной шаг сетки был мал. В частности, если простая
задача для уравнения теплопроводности

(39.1)

\begin{verbatim}
\begin{array}{ccccccccccccccccccccccccccccccccccc
\end{verbatim}

решается по явной схеме, то шаги сетки \(\Delta\) t и \(\Delta\) x
должны удовлетворять условию

\[
\frac{\Delta t}{(\Delta x)^2} \leqslant 0.5.
\]

В противном случае метод будет численно неустойчив (и ошибки округления
будут расти неограниченно). (Вопросы численной устойчивости обсуждаются
в работе {[}53{]}.) Другими словами, если шаг сетки по оси x выбрать
\(\Delta\) x=0,1 , то шаг по времени \(\Delta\) t не может быть больше,
чем 0.5\(\Delta\) x\^{}2=0.005 . Это значит, что от t=0 до t=1 можно
дойти, совершив 200 шагов.

Существуют, однако, такие схемы (неявные схемы), которые позволяют вести
вычисления с гораздо большим шагом, при этом,

правда, растет объем вычислений, совершаемых на одном шаге. По таким
схемам можно вести вычисления с достаточно большим шагом, но на каждом
шаге приходится решать систему алгебраимеских уравнений. Для иллюстрации
этого метода решим следующую задачу теплопроводности.

\section{Неявная схема для уравнения
теплопроводности}\label{ux43dux435ux44fux432ux43dux430ux44f-ux441ux445ux435ux43cux430-ux434ux43bux44f-ux443ux440ux430ux432ux43dux435ux43dux438ux44f-ux442ux435ux43fux43bux43eux43fux440ux43eux432ux43eux434ux43dux43eux441ux442ux438}

Рассмотрим задачу

(39 2)

\begin{verbatim}
\begin{array}{ccccccccccccccccccccccccccccccccccc
\end{verbatim}

Воспользуемся следующими конечно-разностными аппроксимациями для частных
производных \(u_t\) и \(u_{xx}\) :

\[
\begin{split} u_t(x,\ t) &= \frac{1}{k} \big[ u(x,\ t+k) - u(x,\ t) \big], \\ u_{xx}(x,\ t) &= \frac{\lambda}{h^2} \big[ u(x+h,\ t+k) - 2u(x,\ t+k) + u(x-h,\ t+k) \big] + \\ &\quad + \frac{(1-\lambda)}{h^2} \big[ u(x+h,\ t) - 2u(x,\ t) + u(x-h,\ t) \big], \end{split}
\]

где \lambda  выбирается из отрезка {[}0,1{]}. Отметим, что \(u_{xx}\)
аппроксимируется взвешенным средним центральных разностных производных в
момент t и t+k. При \lambda=0.5 получается обычное среднее этих двух
центральных производных, а при \lambda=0.75 ---одна из разностных
производных берется с весом 0.75, а вторая---с весом 0.25. При \lambda=0
получается обычная явная схема, о которой шла речь в предыдущей лекции.

После замены частных производных \(u_t\) и \(u_{xx}\) в задаче (39.2) мы
получаем разностную задачу (см. рис. 39.1).

Разностное уравнение

\[
\frac{1}{k} (u_{i+1, j} - u_{i, j}) = \frac{\lambda}{h^2} (u_{i+1, j+1} - 2u_{i+1, j} + u_{i+i, j-1}) + \frac{(1-\lambda)}{h^2} (u_{i, j+1} - 2u_{i, j} + u_{i, j-1}),
\]

(39.3)

\[
(\Gamma Y) \begin{cases} u_{i, 1} = 0, & i = 1, 2, ..., m, \\ u_{i, n} = 0, & i = 1, 2, ..., m, \\ u_{i, j} = 1, & j = 2, ..., n-1. \end{cases}
\]

298

Перенесем все неизвестные значения u с верхнего временного слоя (с
индексом i+1) в левую часть уравнения (39.3) и получим (39.4)

\[
-\lambda r u_{i+1, j+1} + (1+2r\lambda) u_{i+1, j} - \lambda r u_{i+1, j-1} = = r(1-\lambda) u_{i, j+1} + [1-2r(1-\lambda)] u_{i, j} + r(1-\lambda) u_{i, j-1},
\]

где введено обозначение r = k/h\^{}2 . Отметим, что если i фиксирован, а
j изменяется от 2 до n-1, соотношения (39.4) определяют систему n-2
уравнений с n-2 неизвестными \(u_{t+1,2}\), \(u_{t+1,3}\),

\pandocbounded{\includegraphics[keepaspectratio]{images/_page_296_Figure_5.jpeg}}

РИС. 39.1. Сетка для неявной разностной схемы.

\pandocbounded{\includegraphics[keepaspectratio]{images/_page_296_Figure_7.jpeg}}

РИС. 39.2. Шаблон для неявной схемы.

\begin{verbatim}
u_{i+1, 4}, \ldots, u_{i+1, n-1} , которые являются решением задачи во внут-

ренних узлах сетки на временном слое  t=(i+1)$\Delta$ t .
\end{verbatim}

Рис. 39.2 дает наглядное представление о ссруктуре каждого уравнения
системы (39.4). Перейдем теперу к решению системы (39.4).

\section{Алгоритм решения задачи
(39.2)}\label{ux430ux43bux433ux43eux440ux438ux442ux43c-ux440ux435ux448ux435ux43dux438ux44f-ux437ux430ux434ux430ux447ux438-39.2}

\begin{verbatim}
ШАГ 1. Выбираем некоторое значение  \lambda  ( 0 \le \lambda \le 1 ). Если  \lambda = 0 , то уравнения (39.4), переходят в явные формулы из лекции 38.

ШАГ 2. Пусть, например,  h=$\Delta$ x=0,2  и  k=$\Delta$ t=0,08  (при этом  r=k/h^2=2 ). В данном случае сетка содержит 6 узлов вдоль осн x (4 внутренних узла) (см. рис. 39.1). Возьчем весовой параметр  \lambda=0,5  (получающаяся при этом схема называется схемой Кранка—Никольсона). В соответствии с вычислительной схемой (ее принято называть шаблоном), изображенной на рис. 39.2, двигаясь слева направо (j=2,3,4,5) по первым двум слоям (i=1), получаем следующие четыре уравнения:
\end{verbatim}

\[
\begin{array}{l} -u_{21}+3u_{22}-u_{23}=u_{11}-u_{12}+u_{13}=1,\\ -u_{22}+3u_{23}-u_{24}=u_{12}-u_{13}+u_{14}=1,\\ -u_{23}+3u_{24}-u_{25}=u_{13}-u_{14}+u_{15}=1,\\ -u_{24}+3u_{25}-u_{26}=u_{14}-u_{15}+u_{16}=1. \end{array}
\]

Перепишем их в матричной форме

(39.5)

\[
\begin{bmatrix} 3 & -1 & 0 & 0 \\ -1 & 3 & -1 & 0 \\ 0 & -1 & 3 & -1 \\ 0 & 0 & -1 & 3 \end{bmatrix}
\begin{bmatrix} u_{2\bar{z}} \\ u_{\bar{z}a} \\ u_{24} \\ u_{\bar{z}x} \end{bmatrix}
=
\begin{bmatrix} 1 \\ 1 \\ 1 \\ 1 \end{bmatrix}
\]

.

Матрица этой системы называется трехдиагональной. Для того чтобы решить
трехдиагональную систему

\[
\begin{bmatrix} b_1 & c_1 & 0 & 0 & \dots & & & 0 \\ a_1 & b_2 & c_2 & 0 & \dots & & & 0 \\ 0 & a_{\bar{z}} & b_3 & c_3 & \dots & & & 0 \\ & & & & & & & & \\ \vdots & & & & & & & & \\ 0 & 0 & 0 & \dots & & & & a_{n-1} & b_n \end{bmatrix}
\begin{bmatrix} x_1 \\ x_{\bar{z}} \\ x_{\bar{s}} \\ \vdots \\ x_n \end{bmatrix}
=
\begin{bmatrix} d_t \\ d_{\bar{z}} \\ d_{\bar{s}} \\ \vdots \\ \vdots \\ d_n \end{bmatrix}
\]

,

преобразуем ее в эквивалентную систему вида

\[
\begin{bmatrix} 1 & c_1^* & 0 & 0 & \dots & 0 \\ 0 & 1 & c_2^* & 0 & \dots & 0 \\ 0 & 0 & 1 & c_2^* & \dots & 0 \\ \vdots & \vdots & & & \vdots & \vdots \\ 0 & 0 & \dots & & 1 \end{bmatrix}
\begin{bmatrix} x_1 \\ x_2 \\ x_3 \\ \vdots \\ x_n \end{bmatrix}
=
\begin{bmatrix} d_1^* \\ d_2^* \\ d_3^* \\ \vdots \\ d_n^* \end{bmatrix}
\]

.

где

\[
c_1^* = c_1/b_1
,  c_{j+1}^* = \frac{c_{j+1}}{b_{j+1} - a_j c_j^*}  ,  j = 1, 2, \ldots, n-2  , и  d_1^* = d_1/b_1  ,  d_{j+1}^* = \frac{d_{j+1} - a_j d_j^*}{b_{j+1} - a_j c_j^*}  ,  j = 1,  2, \ldots, n-1  .
\]

В этом преобразовании нет ничего таниственного, поскольку вторая система
полностью эквивалентна первой. Матрица новой системы устроена так, что
эта система очень легко решается. Решая уравнения последовательно снизу
вверх, получаем

\[
x_n = d_n^*, \quad x_j = d_j^* - c_j^* x_{j+1}, \quad j = n-1, \ n-2, \ \dots, \ 2, \ 1.
\]

Застосовуючи цей метод до системи чотирьох рівнянь (39,5), отримуємо
розв'язок \(u_{22}=0,60\); \(u_{23}=0,80\); \(u_{24}=0,80\);
\(u_{23}=0,60\).Ми отримали наближені значення розв'язку у внутрішніх
точках сітки на t=\Delta t . Тепер ми можемо зробити наступний крок у
часі, але для цього нам доведеться розв'язати нову систему рівнянь.

В неявной схеме объем вычислений на каждом шаге больше, чем в явной, но
хорошую точность можно получить даже при гораздо большем шаге.

\section{ЗАДАЧИ}\label{ux437ux430ux434ux430ux447ux438-17}

\begin{itemize}
\item
  \begin{enumerate}
  \def\labelenumi{\arabic{enumi}.}
  \tightlist
  \item
    Из дифференциального уравнения (39.3) получите уравнение (39.4).
  \end{enumerate}
\item
  \begin{enumerate}
  \def\labelenumi{\arabic{enumi}.}
  \setcounter{enumi}{1}
  \tightlist
  \item
    Постройте неявную конечно-разностную схему для задачи
  \end{enumerate}

  \begin{array}{ll} (\mathrm{Y}\mathrm{H}\Pi) & u_t = u_{xx}, \quad 0 < x < 1, \quad 0 < t < \infty, \\ (\Gamma\mathrm{Y}) & \left\{ \begin{array}{ll} u_{\,0}(0, \ t) = 1, \\ u_{\,x}(1, \ t) + u_{\,}(1, \ t) = g_{\,}(t), \\ u_{\,x}(1, \ t) = 0, \quad 0 \leqslant \kappa \leqslant 1. \end{array}

    \right. & 0 < t < \infty, \end{array}
\end{itemize}

\begin{enumerate}
\def\labelenumi{\arabic{enumi}.}
\setcounter{enumi}{2}
\tightlist
\item
  Как с помощью неявной схемы решить задачу
\end{enumerate}

(УЧП)

\[
u_t = u_{xx} + u, \quad 0 < x < 1,
\]

(ГУ)

\begin{itemize}
\tightlist
\item
  \begin{enumerate}
  \def\labelenumi{\arabic{enumi}.}
  \setcounter{enumi}{3}
  \tightlist
  \item
    Как выглядит шаблон для уравнения (39.4), когда \lambda = 1 ?
  \end{enumerate}
\item
  \begin{enumerate}
  \def\labelenumi{\arabic{enumi}.}
  \setcounter{enumi}{4}
  \tightlist
  \item
    Постройте блок-схему решения уравнения теплопроводности (39.2). Если
    есть возможность, то напишите программу для ЭВМ Было бы очень
    полезно решить эту задачу численно при различных значениях λ с
    начальным условием и (x, 0) = sin (πx). Сравните полученные при
    различных λ численные результаты с аналитическим решением, которое в
    данном случае имеет вид
  \end{enumerate}
\end{itemize}

\[
u(x, t) = e^{-\pi^2 t} \sin(\pi x).
\]

6 Воспользовавшись формулами, приведенными в тексте, решите систему
алгебраических уравнений (39.5).

Лекция 40

\section{СРАВНЕНИЕ АНАЛИТИЧЕСКИХ РЕШЕНИЙ С
ЧИСЛЕННЫМИ}\label{ux441ux440ux430ux432ux43dux435ux43dux438ux435-ux430ux43dux430ux43bux438ux442ux438ux447ux435ux441ux43aux438ux445-ux440ux435ux448ux435ux43dux438ux439-ux441-ux447ux438ux441ux43bux435ux43dux43dux44bux43cux438}

ЦЕЛЬ ЛЕКЦИИ: Сравнить достоинства и недостатки аналитических и численных
методов решения уравнений с частными производными. Приведена важная
математическая задача об идентификации физических величин
(параметрическая идентификация). В связи с этой задачей рассматривается
один важный пример из биологии.

Вероятно, пришло время обсудить достоинства и недостатки аналитических и
численных решений уравнений с частными производными. Сначала выясним,
что мы подразумеваем под этими двумя типами решений.

\section{Аналитические
решения}\label{ux430ux43dux430ux43bux438ux442ux438ux447ux435ux441ux43aux438ux435-ux440ux435ux448ux435ux43dux438ux44f}

Под аналитическими решениями мы подразумеваем такие решения, в которых
неизвестная функция и выражена через независимые переменные и параметры
системы в виде формул, бесконечных рядов и интегралов.

\section{Численные
решения}\label{ux447ux438ux441ux43bux435ux43dux43dux44bux435-ux440ux435ux448ux435ux43dux438ux44f}

Під чисельним ми зрозуміємо чисельні розв'язки, отримані чисельно після
приблизної заміни початкового рівняння іншим, простішим рівнянням.
Наприклад, у методі скінченної різниці похідні виражаються через
скінченні різниці, а розв'язок диференціального рівняння в похідних у
похідних апроксимується розв'язком різницького рівняння. Результатом
такої процедури зазвичай є таблиця значень розв'язків, а для деяких
значень пояснювальних змінних.

Таблина 40.1

~~~~~~~~~~~~~~~~~~~~~~~~~~\textbar{} × \textbar{} \textbar{} \textbar{}
\textbar{} \textbar{} \textbar{} \textbar{} \textbar{} \textbar{}
\textbar{} \textbar{}

\textbar---------------------------\textbar-------------\textbar------------------\textbar-----------\textbar----------\textbar----------\textbar-----\textbar-----\textbar-----\textbar-----------\textbar------------------\textbar------------------\textbar{}

~~~~~~~~~~~~~~~~~~~~~~~~~~\textbar{} 0 \textbar{} 0,1 \textbar{} 0,2
\textbar{} 0,3 \textbar{} 0,4 \textbar{} 0,5 \textbar{} 0,6 \textbar{}
0,7 \textbar{} 0,8 \textbar{} 0,9 \textbar{} 1 \textbar{}

00,010,020,03 \textbar{} 100 \textbar{} 10,20,15 \textbar{} 10,34
\textbar{} l: \textbar{} : \textbar{} : \textbar{} ! \textbar{} :
\textbar{} 10,34 \textbar{} 10,20,15 \textbar{} 1000 \textbar{}

; \textbar{} \textbar{} : \textbar{} : \textbar{} : \textbar{} :
\textbar{} : \textbar{} : \textbar{} : \textbar{} : \textbar{} :
\textbar{} : \textbar{}

Теперь, когда мы знаем, что подразумевается под решением каждого из этих
типов, давайте сравним их между собой.

\section{\texorpdfstring{\textbf{Сравнительный анализ численного} и
аналитического
решений}{Сравнительный анализ численного и аналитического решений}}\label{ux441ux440ux430ux432ux43dux438ux442ux435ux43bux44cux43dux44bux439-ux430ux43dux430ux43bux438ux437-ux447ux438ux441ux43bux435ux43dux43dux43eux433ux43e-ux438-ux430ux43dux430ux43bux438ux442ux438ux447ux435ux441ux43aux43eux433ux43e-ux440ux435ux448ux435ux43dux438ux439}

Рассмотрим простейшую смещанную параболическую задачу

(40.1)

\begin{verbatim}
\begin{array}{ccccccccccccccccccccccccccccccccccc
\end{verbatim}

Розв'язок якого показано на рис. 40.1. Що краще --- мати аналітичне
рішення цієї проблеми?

(40.2)

\[
u(x, t) = \frac{2}{\pi} \left[ e^{-(\pi\alpha)^2 t} \sin(\pi x) + \frac{1}{3} e^{-(3\pi\alpha)^2 t} \sin(3\pi x) + \dots \right]
\]

или численное, представленное при \(\alpha\) = 1 в табл. 40.1?

Вопрос хороший, а ответ на него зависит от того, что мы хотим делать с
этим решением дальше. Однако каждому из типов решений присущи очевидные
преимущества.

\section{Преимущества аналитического
решения}\label{ux43fux440ux435ux438ux43cux443ux449ux435ux441ux442ux432ux430-ux430ux43dux430ux43bux438ux442ux438ux447ux435ux441ux43aux43eux433ux43e-ux440ux435ux448ux435ux43dux438ux44f}

\begin{enumerate}
\def\labelenumi{\arabic{enumi}.}
\tightlist
\item
  Очевидно, що розв'язок (40.2) є більш інформативним, ніж таблиця
  чисел. Якщо ми хочемо обчислити розв'язок у певній точці (x, t), ми
  можемо зробити це так точно, як хочемо, просто збільшуючи кількість
  членів ряду, які розглядаються. Водночас легко отримати оцінку зверху
  щодо значення допущеної помилки.
\end{enumerate}

\pandocbounded{\includegraphics[keepaspectratio]{images/_page_301_Figure_4.jpeg}}

РИС. 40.1. Решение уравнения теплопроводности (40.1) в различные моменты
времени.

\begin{itemize}
\tightlist
\item
  \begin{enumerate}
  \def\labelenumi{\arabic{enumi}.}
  \setcounter{enumi}{1}
  \tightlist
  \item
    Аналітичне розв'язання завжди дозволяє обчислити значення розв'язку
    в одній точці (x, t) без необхідності обчислювати значення
    розв'язків в інших точках, як це відбувається при розв'язанні задач
    із явними або неявними різницею схем.
  \end{enumerate}
\item
  \begin{enumerate}
  \def\labelenumi{\arabic{enumi}.}
  \setcounter{enumi}{2}
  \tightlist
  \item
    Аналітичне рішення дозволяє визначити розв'язок у будь-якій точці, а
    не лише на вузлах сітки.
  \end{enumerate}
\item
  \begin{enumerate}
  \def\labelenumi{\arabic{enumi}.}
  \setcounter{enumi}{3}
  \tightlist
  \item
    Найважливішою перевагою аналітичного розв'язку для нас є можливість
    простежити вплив фізичних параметрів, початкових і граничних умов на
    природу розв'язку.
  \end{enumerate}
\end{itemize}

Чисельні методи не виявляють цих регулярностей, оскільки дозволяють
знаходити розв'язок лише за заданих параметрів --- початкових і
граничних умовах. Іноді важливо знати взаємозв'язок між параметрами
моделі та розв'язком, особливо коли йдеться про оцінку фізичних
параметрів за типом розв'язку. Наприклад, припустимо, що розв'язок і
визначається експериментально, і ми знаємо аналітичний розв'язок

u = функция от параметров.

Тогда можно поставить вопрос об определении параметров нак функции
входных данных по схеме

Параметры = функция от u = функция исходных данных.

Такая задача называется параметрической идентификацией. Для
параметрической идентификации необходимо уметь решать уравнение с
частными производными. Несколько позже в этой лекции мы приведем пример
параметрической идентификации в биологии, но сначала посмотрим, чем же
хороши численные решения.

\section{Преимущества численных
решений}\label{ux43fux440ux435ux438ux43cux443ux449ux435ux441ux442ux432ux430-ux447ux438ux441ux43bux435ux43dux43dux44bux445-ux440ux435ux448ux435ux43dux438ux439}

Головна перевага чисельних розв'язків полягає в тому, що їх можна
отримати навіть якщо аналітичне рішення неможливо. Майже всі нелінійні
диференціальні рівняння в похідних похідних потрібно розв'язувати
чисельними методами, і більшість реальних моделей фізики, хімії,
біології тощо є нелінійними. Як правило, лінійні моделі наближаються до
нелінійних лише якщо в них відкинути нелінійні члени. Наприклад, ось
кілька нелінійних рівнянь:

\begin{enumerate}
\def\labelenumi{\arabic{enumi})}
\item
  нелинейное волновое уравнение \(u_{tt} = u_{xx} +\)f(u)\$\$;
\item
  уравнение реакции с нелинейной диффузией \(u_t = u_{xx} +\)f(u)\$\$;
\item
  система Ходжкина---Хаксли
\end{enumerate}

\[
u(x, t) =
\]

Аналитические решения для этих уравнений не известны ни при каких
нелинейных функциях f и \(g^{(1)}\) . Поэтому общий подход к решению
нелинейных (а в ряде случаев и линейных) задач обычно базируется на
численных решениях.

Тепер давайте розглянемо приклад того, як можна використати аналітичне
рішення для визначення фізичних параметрів. Детальніший опис цих завдань
можна знайти у рекомендованій літературі.

\section{Параметрическая идентификация (в
биологии)}\label{ux43fux430ux440ux430ux43cux435ux442ux440ux438ux447ux435ux441ux43aux430ux44f-ux438ux434ux435ux43dux442ux438ux444ux438ux43aux430ux446ux438ux44f-ux432-ux431ux438ux43eux43bux43eux433ux438ux438}

Припустимо, що біолог намагається визначити швидкість, з якою (K\^{}+)
іони калію дифундують у екзоплазмовому розчині. Якщо ви знаєте цей
коефіцієнт дифузії, можна багато сказати про те, як нервові імпульси
передаються вздовж аксонів. Проблема в тому, що це співвідношення майже
неможливо знайти прямим вимірюванням. Однак можливо знайти математичний

\begin{verbatim}
<sup>1)</sup> Это не так. В качестве примера можно привести уравнение  u_{ii} = u_{xx} - \sin u_x - \Pi pum  , ped.

связь между концентрацией поташа  u\left(x,\,t\right)  и коэффициентом диффузии  D,\,\,a\,\,  затем по измерениям величины  u\left(x,\,t\right)  определить величину D.
\end{verbatim}

Біологи Ходжкін і Кейз виявили, що після розміщення гігантського аксону
у спеціальному соляному розчині концентрація радіохктивного калію (**2K)
вздовж аксону приблизно описується кривою в початковий момент

\[
u(x, 0) = Ae^{-x^2/a}
\]

.

Отже, параметри A і A слід обирати так, щоб крива відповідала
експериментальним даним (наприклад, у сенсі методу найменших квадратів).
Також було встановлено, що P И С

\pandocbounded{\includegraphics[keepaspectratio]{images/_page_303_Figure_6.jpeg}}

\begin{verbatim}
. 40.2. Начальная концентрация ионов (  ^{42}K  ): a—теоретическая кривая;  $\delta$  —экспериментальные точки.
\end{verbatim}

с течением времени ноны поташа под действием диффузии и конвекции
растекаются вдоль аксона. Ходжкин и Кейс предположили, что концентрация
и является решением следующей задачи с начальными данными для уравнения
диффузии:

(40.3)

\[
\begin{aligned} (\text{УЧ}\Pi) & u_t = Du_{xx} - Vu_x, & -\infty < x < \infty, \\ (\text{HY}) & u(x, 0) = Ae^{-x^2, a}, & -\infty < x < \infty. \end{aligned}
\]

Цю проблему вирішили у лекції 15, переключившись на рухому систему
координат. Якщо читач пам'ятає, загальна схема методу виглядає так.
Спочатку відкидаємо конвективний член і отримуємо розв'язок суто
дифузійної задачі (при V = 0):

\[
u(x, t) = \frac{A}{2\sqrt{\pi Dt}} \int_{-\infty}^{+\infty} e^{-\xi^2/a} e^{-(x-\xi)^2/4Dt} d\xi =
= \frac{A\sqrt{a}}{\sqrt{a+4Dt}} e^{-[x^2/(a+4Dt)]},
\]

а затем в этом решении заменяем x на x-Vt и получаем решение задачи
(40.3):

(40.4)

\[
u(x, t) = \frac{A\sqrt{a}}{\sqrt{a+4Dt}}e^{-[(x-Vt)^{2}/(a+4Dt)]}.
\]

В опыте Ходжкина и Кейса проводилось прямое измерение скорости, а
формула (40.4) устанавливала связь между концентрацией ионов (
\(^{42}K\) ) и коэффициентом диффузии D (величины A, a и V были
известны). Экспериментальная установка позволяла измерять Р И С

\pandocbounded{\includegraphics[keepaspectratio]{images/_page_304_Figure_3.jpeg}}

. 40.3. Кривая \(u\left(x,\ t\right)\) соответствует экспериментальным
точкам: a---кривая, наилучшим образом соответствующая экспериментальным
данным; b---экспериментально найденная концентрация в точке.

концентрацию поташа \(u(x, t)\) в одной точке \(x_0\) аксона, но в
различные моменты времени. Оставалось найти коэффициент D таким образом,
чтобы теоретическая кривая (40.4) наилучшим образом соответствовала
экспериментальным данным (рис. 40.3).

Іншими словами, для різних коефіцієнтів D ми отримаємо різні криві
\(u\left(x_{0},\ t\right)\) , і тому значення D потрібно вибрати так,
щоб розв'язок рівняння відповідав досвіду. У своєму експерименті Ходжкін
і Кейсі виявили, що коефіцієнт дифузії \((K^{+})\) іонів в екзоплазмі
становить \(D=1.5\times10^{-5}\) см²/с.

\section{ЗАДАЧИ}\label{ux437ux430ux434ux430ux447ux438-18}

\begin{itemize}
\tightlist
\item
  \begin{enumerate}
  \def\labelenumi{\arabic{enumi}.}
  \tightlist
  \item
    Який експеримент ви б порадили для оцінки параметра α у задачі
    (40.1), якщо у вас є аналітичне рішення, наприклад (40.2)?
  \end{enumerate}
\item
  \begin{enumerate}
  \def\labelenumi{\arabic{enumi}.}
  \setcounter{enumi}{1}
  \tightlist
  \item
    У методі найменших квадратів сума квадратів відхилень кривої від
    експериментальних точок мінімізується. Наприклад, згідно з методом
    найменших квадратів, пряма \(y(x)\) = a + bx мінімізує (див. рисунок
    40.4) функцію
  \end{enumerate}
\end{itemize}

\[
SS = \sum_{i=1}^{n} [y_i - (a + bx_i)]^2.
\]

Выразите константы a и b через начальные данные из условия минимума
функции SS.

\begin{enumerate}
\def\labelenumi{\arabic{enumi}.}
\setcounter{enumi}{2}
\tightlist
\item
  Одной из важных задач в биохимии является определение молекулярного
  веса макромолекул многлобина. В одном из способов кровь помещают в
  ультрацентрифугу и вращают в течение заданного промежутка времени.
  Концентрация жидкости в центрифуге описывается уравнением Ламма
\end{enumerate}

\[
u_t = \frac{1}{r} \frac{\partial}{\partial r} \left( Dr u_r - s \omega^2 r^2 u \right), \quad 0 < r < 1,
\]

где

r-расстояние от центра центрифуги,

D---коэффициент диффузии (зависит от молек лярного веса миоглобина),

s -- коэффициент седиментации (определяется экспериментально),

\pandocbounded{\includegraphics[keepaspectratio]{images/_page_305_Figure_8.jpeg}}

РИС. 40.4. Аппрокенмация экспериментальных данных по методу наименьших
квадратов. \(\omega\) --- угловая скорость вращения центрифуги
(известна), \(u(r, t)\) --- концентрация вещества в центрифуге.
Молекулярный вес миоглобина определяется с помощью стационарного (
\(u_t \equiv 0\) ) решения \(u(r, \infty)\) уравнения Ламма. Для
определения \(u(r, \infty)\) получается обыкновенное дифференциальное
уравнение. Вопрос ставится следующим образом: как по экспериментально
найденному стационарному состоянию оценить молекулярный вес многлобина?

\section{\texorpdfstring{КЛАССИФИКАЦИЯ УРАВНЕНИЙ (ПАРАБОЛИЧЕСКИЕ И
ЭЛЛИПТИЧЕСКИЕ
\textbf{УРАВНЕНИЯ})}{КЛАССИФИКАЦИЯ УРАВНЕНИЙ (ПАРАБОЛИЧЕСКИЕ И ЭЛЛИПТИЧЕСКИЕ УРАВНЕНИЯ)}}\label{ux43aux43bux430ux441ux441ux438ux444ux438ux43aux430ux446ux438ux44f-ux443ux440ux430ux432ux43dux435ux43dux438ux439-ux43fux430ux440ux430ux431ux43eux43bux438ux447ux435ux441ux43aux438ux435-ux438-ux44dux43bux43bux438ux43fux442ux438ux447ux435ux441ux43aux438ux435-ux443ux440ux430ux432ux43dux435ux43dux438ux44f}

ЦЕЛЬ ЛЕКЦИИ: Для линейного дифференциального уравнения второго порядка

\[
Au_{xx} + Bu_{xy} + Cu_{yy} + Du_x + Eu_y + Fu = G
\]

найти координатную систему (~ , η), у якій рівняння зводиться до однієї
з двох канонічних форм

\[
\begin{cases} 0, & t < x/V, \\ P, & t \ge x/V. \end{cases}
\]

(параболічний канонічний

форма),
\(u_{\xi\xi} + u_{\eta\eta} = q \; (\xi, \, \eta, \, u, \, u_{\xi}, \, u_{\eta})\)
(еліптична канонічна форма)

залежно від значення \(B^{2} - 4AC = 0\) є параболічним випадком,
\(B^2 - 4AC < 0\) є еліптичним випадком.

У лекції 23 ми розділили лінійні рівняння другого порядку двома змінними
вигляду

(41.1)

Конечно, все это очень просто! Не сложнее, чем сбрасывать что-нибудь на
ленту транспортера, а затем наблюдать за его движением. Однако все
становится очень интересным, если примесь диффундирует в среде. Чтобы
разобраться в том, что про-исходит с движущейся волной при наличии
диффузии, решим

следующую задачу

(15.2)

на три основні класи та знайшли дві канонічні форми рівнянь
гіперболічного типу.

У цій лекції ми покажемо, як канонічно зводити параболічні та еліптичні
рівняння.

\section{Зведення параболічних рівнянь до канонічної
форми}\label{ux437ux432ux435ux434ux435ux43dux43dux44f-ux43fux430ux440ux430ux431ux43eux43bux456ux447ux43dux438ux445-ux440ux456ux432ux43dux44fux43dux44c-ux434ux43e-ux43aux430ux43dux43eux43dux456ux447ux43dux43eux457-ux444ux43eux440ux43cux438}

Розглянемо рівняння виду (41.1) з \(B^2 - 4AC = 0\) і введемо нові
координати (E, p) таким чином, щоб звести рівняння до вигляду

З формальною заміною змінних \xi = \xi(x, y) та \eta = \eta(x, y) (41.1)
зводиться до тієї ж форми

(41.2)

, рей

где, как обычно, \(H(x)\) --- функция Хевисайда Начальное распределение
концентрации изображено на рис 15 4.

\pandocbounded{\includegraphics[keepaspectratio]{images/_page_114_Figure_4.jpeg}}

\begin{verbatim}
РИС. 15.3 Конвективная волна a—передиції фронт примеси,  $\delta$ —на одну единицу вперед волнового фронта

Обратим внимание на то, что в задаче (15.2) мы отодвинули границу в —  \infty  (теперь решается задача Коши), так что она не оказывает теперь влияние на оценку решение. Чтобы решить задачу (15.2), можно воспользоваться преобразованием Лапласа Р
\end{verbatim}

\pandocbounded{\includegraphics[keepaspectratio]{images/_page_114_Figure_7.jpeg}}

ИС. 15 4 Начальные условия в задаче конвективной диффузии.

\begin{verbatim}
ло переменной t или преобразованием Фурье по переменной x. Однако в этом случае гораздо интереснее поступить совсем по-другому Введем новую систему координат, которая движется вдоль старой со скоростью V. Другими словами, вместо системы координат, привязанной к берегу реки, мы рассмотрим систему координат, которая движется со скоростью фронта примеси (конечно, при налични диффузии фронт будет «смазан»). С точки зрения математики это означает, что мы заменяем пространственную координату x на новую  \xi = x - Vt .
\end{verbatim}

Теперь ясно, что

\begin{verbatim}
когда  \xi = 0 , мы находимся па фронте распространяющейся примеси.

когда  \xi = 1 , мы находимся на одну единицу длины впереди фронта.

когда  \xi = -1 , мы находимся на одну единицу длины позади фронта.
\end{verbatim}

Наша задача --- преобразовать исходную задачу с НУ

\[
\overline{B} = 2A\xi_{x}\eta_{x} + B(\xi_{x}\eta_{y} + \xi_{y}\eta_{x}) + 2C\xi_{y}\eta_{y},
\overline{D} = A\xi_{xx} + B\xi_{xy} + C\xi_{yy} + D\xi_{x} + E\xi_{y},
\overline{F} = F, \quad \overline{G} = G.
[\xi_x/\xi_y] = -B/2A
\xi = x - Vt,
\xi(x, y) = c
\tau = t.
\frac{dy}{dx} = -\left[\xi_x/\xi_y\right] = B_t 2A.
u_{t} = u_{\xi} \, \xi_{t} + u_{\tau} \, \tau_{t} = -V u_{\xi} + u_{\tau},
\overline{B} = 2A\xi_x\eta_x + B(\xi_x\eta_y + \xi_y\eta_x) + 2C\xi_y\eta_y,
u_{x} = u_{\xi} \, \xi_{x} = u_{\xi},
\]

1) Завдяки рівності

І з B\^{}2 - 4AC = 0 році , то

\[
u_{xx} = (u_{\xi})_{x} = u_{\xi\xi} \xi_{x} = u_{\xi\xi}.
\]

Оскільки \xi\_x/\xi\_y = -B ~2A = -2 \sqrt{A\tilde{C}}/2A = -\sqrt{C/A}
останній вираз для B буде записаний як

\[
-Vu_{\xi}+u_{\tau}=Du_{\xi\xi}-Vu_{\xi},
\]

І оскільки \(\overline{A}\) дорівнює нулю, то \(\overline{B}\) , звісно,
також дорівнює нулю.

\begin{verbatim}
Оскільки при нашому виборі  \xi  обидва коефіцієнти  $\overline{A}$  і  $\overline{B}$  повертаються до нуля, ми можемо взяти будь-яку функцію як  \eta  (якщо вона не пропорційна координаті  \xi  ). У прикладі можна вибрати  \eta = y  .

Залишається знайти канонічну форму рівняння в нових координатах. Для цього потрібно просто підставити формули  \xi  і  \eta  (41.3) і знайти всі шанси  $\overline{A}$  ,  $\overline{B}$  ,  $\overline{C}$  ,  $\overline{D}$  ,  $\overline{E}$  ,  $\overline{F}$  і  $\overline{G}$  . На цьому вивчення параболічного випадку завершується, але перш ніж перейти до іншої теми, розглянемо простий приклад.
\end{verbatim}

\section{\texorpdfstring{Перетворення параболічного рівняння \(u_{xx}\)
+ 2u\_\{xy\} + \(u_{yy} = 0\)у канонічну
форму}{Перетворення параболічного рівняння u\_\{xx\} + 2u\_\{xy\} + u\_\{yy\} = 0у канонічну форму}}\label{ux43fux435ux440ux435ux442ux432ux43eux440ux435ux43dux43dux44f-ux43fux430ux440ux430ux431ux43eux43bux456ux447ux43dux43eux433ux43e-ux440ux456ux432ux43dux44fux43dux43dux44f-u_xx-2u_xy-u_yy-0ux443-ux43aux430ux43dux43eux43dux456ux447ux43dux443-ux444ux43eux440ux43cux443}

Розглянемо рівняння

\[
u_{\tau} = Du_{\xi\xi}
\]

Його коефіцієнти мають такі значення: A=1, B=2, C=1, D=E=F=G=0. Отже,
B\^{}2-4AC=0 для всіх значень x і y. Щоб знайти нові координати \xi  і
\eta  та канонічну форму рівняння, ми рухаємося наступним чином.

КРОК 1. Запишемо \emph{характеристичне рівняння} (тепер лише одне, а не
два)

\[
u_{\tau} = Du_{\xi\xi}, \quad -\infty < \xi < \infty,
\]

Розв'язуючи це рівняння відносно y, отримуємо

\[
u(\xi, \tau) = \frac{1}{2\sqrt{D\pi\tau}} \int_{-\infty}^{\infty} \varphi(\beta) e^{-(\xi-\beta)^2/4D\tau} d\beta,
\]

,

\begin{verbatim}
Де — характеристика  \xi = y - x  визначена. Використання цієї характеристики забезпечує досягнення рівності  $\overline{A}$ = 0  . Координату  \eta  можна вибрати довільно (якщо вона не залежить від  \xi  ). Ми оберемо його наступним чином:
\end{verbatim}

\[
u(\xi, \tau) = \frac{1}{2\sqrt{D\pi\tau}} \int_{-\infty}^{0} e^{-(\xi-\beta)^2/4D\tau} d\beta.
\]

.

Нові координати

\[
\overline{\beta} = \frac{\xi - \beta}{2\sqrt{D\tau}}, \quad d\overline{\beta} = \frac{-1}{2\sqrt{D\tau}}d\beta
\]

изображены на рис. 41 1.

\pandocbounded{\includegraphics[keepaspectratio]{images/_page_309_Figure_3.jpeg}}

\begin{verbatim}
РИС. 41.1. Новая система координат   \xi = y - x  ,   \eta = y  .

ШАГ 2. Осталось найти каноническую форму уравнения. Подставляя   \xi   и   \eta   в формулы для коэффициентов   $\overline{A}$  ,   $\overline{B}$  ,   $\overline{C}$  ,   $\overline{D}$  ,   $\overline{E}$  ,   $\overline{F}$   и   $\overline{G}$  , получаем

$\overline{A}$ = 0   (так и должно быть; мы находили   \xi   из условия равенства нулю этого коэффициента),

$\overline{B}$=0   (мы уже раньше показали, что этог коэффициент также обращается в нуль),
\end{verbatim}

(41.5)

\[
u(\xi, \tau) = \frac{1}{2} \left[ \frac{2}{\sqrt{\pi}} \int_{\frac{\xi}{\sqrt{D\tau}}}^{\infty} e^{-\overline{p}^{2}} d\overline{\beta} \right] =
\overline{D} = A \xi_{xx} + B \xi_{xy} + C \eta_{yy} + D \xi_x + E \xi_y = 0,
\overline{F} = F = 0,
\]

Отже, нове рівняння

перетворюється на \(u_{nn} = 0\) .

На цьому наш приклад завершено, але перш ніж рухатися далі, звернемо
увагу на одну обставину у позначення канонічної форми. Це рівняння
настільки просте, що можна знайти загальне розв'язання (тобто всі
розв'язки). Інтегруючи \(u_{\eta\eta} = 0\) до \(\eta\) , ми отримуємо

\[
\xi = x - Vt
\]

де \(f(\xi)\) --- довільна функція. Реінтеграція веде до спільного
рішення

\[
u(x, t) =
\]

де \(g(\xi)\) --- інша довільна функція.

Читач легко може переконати себе, що будь-яка функція такого роду
задовольняє рівняння \(u_{nn} = 0\) .

Залишається зробити останній крок і замінити початкові координати x і y
у загальний розв'язок.

Внаслідок підстановки отримуємо, що

є загальним розв'язком рівняння

Наприклад, якщо вибрати \(f(x) = \sin x\) і \(g(x) = x^2\) , то

Перед нами решение задачи конвективной диффузии (15.2). Его очень легко
интерпретировать, если представить, что мы движемся относительно
графика, изображенного на рис. 15.6.

\pandocbounded{\includegraphics[keepaspectratio]{images/_page_117_Figure_3.jpeg}}

РИС. 15.6. Диффузня из области высокой концентрации в область низкой
концентрации. Чем больше коэффициент диффузии, тем быстрее установится
стационарное значение.

Другими словами, в зависимости от относительной величины D (коэффициента
диффузии) и V (скорости потока) решение движется слева направо со
скоростью V и в то же самое время передний фронт расплывается со
скоростью, определяемой величиной D.

\pandocbounded{\includegraphics[keepaspectratio]{images/_page_117_Figure_6.jpeg}}

РИС. 15.7. Решение задачи конвективнои диффузии. Вещество одновременно
движется и диффундирует.

(На рис. 15.7 показано, как расплывается передний фронт концентрации.)

\subsubsection{ЗАМЕЧАНИЯ}\label{ux437ux430ux43cux435ux447ux430ux43dux438ux44f-16}

Преобразование координат является важным методом решения уравнений с
частными производными. Выбрав подходящую систему координат, можно
существенно упростить уравнение.

\subsubsection{ЗАДАЧИ}\label{ux437ux430ux434ux430ux447ux438-19}

\begin{enumerate}
\def\labelenumi{\arabic{enumi}.}
\tightlist
\item
  Решите задачу Коши:
\end{enumerate}

буде одним із нескінченної кількості розв'язків (звісно, розв'язок, що
моделює певну фізичну задачу, залежить від початкових і граничних умов).

\section{Редукція еліптичних рівнянь до канонічної
форми}\label{ux440ux435ux434ux443ux43aux446ux456ux44f-ux435ux43bux456ux43fux442ux438ux447ux43dux438ux445-ux440ux456ux432ux43dux44fux43dux44c-ux434ux43e-ux43aux430ux43dux43eux43dux456ux447ux43dux43eux457-ux444ux43eux440ux43cux438}

Розглянемо загальне рівняння

Але тепер у B\^{}2-4AC\textless0 . Переходячи до нових пояснювальних
змінних, ми хочемо перетворити її у нову форму

Щоб знайти ці нові змінні, ми зробимо ті ж обчислення, що й у двох
попередніх випадках: отримаємо трансформоване рівняння

\begin{enumerate}
\def\labelenumi{\arabic{enumi}.}
\setcounter{enumi}{1}
\tightlist
\item
  Найдите решение задачи Коши для уравнения конвективной диффузии:
\end{enumerate}

і вимагати, щоб \(\xi\) і \(\eta\) задовольняли рівняння
\(\overline{A}=\overline{C}\) і \(\overline{B}=0\) . На жаль, ці
рівняння не дозволяють одразу знайти \(\xi\) і \(\eta\) , як це було у
попередніх випадках, тому ми продовжимо трохи пізніше.

Ми будемо шукати необхідне перетворення як композицію двох перетворень.

\section{Конверсія
1}\label{ux43aux43eux43dux432ux435ux440ux441ux456ux44f-1}

За першим перетворенням у нові координати \textbf{ξ і ŋ ми} надаємо
нашому рівнянню таку форму

Це можливо, але лише шляхом введення складних координат. Щоб знайти ці
комплексні координати \xi  і \eta  , розв'яжемо характеристичні
рівняння, так само як у гіперболічному випадку:

, (нагадуємо, що \(B^2 - 4AC < 0\) ),

і ми отримаємо

Воспользуйтесь преобразованием из лекции 8.

\begin{enumerate}
\def\labelenumi{\arabic{enumi}.}
\setcounter{enumi}{2}
\tightlist
\item
  Найдите решение задачи переноса:
\end{enumerate}

Наприклад, характеристичні рівняння

\[
\frac{dy}{dx} = \sqrt{-4x^2} = 2ix
\xi = y + ix^2, \quad \eta = y - ix^2.
\\ -\infty < x < \infty, \quad 0 < t < \infty,
\alpha = \frac{\xi + \eta}{2}
u_t = Du_{xx} - Vu_x, \quad -\infty < x < \infty, \quad 0 < t < \infty,
u_{\xi\eta} = \psi(\xi, \eta, u, u_{\xi}, u_{\eta})
u(x, 0) = e^{-x^2}, \quad -\infty < x < \infty.
u_{\alpha\alpha} + u_{\beta\beta} = \varphi(\alpha, \beta, u, u_{\alpha}, u_{\beta}),
y^2 u_{xx} + x^2 u_{yy} = 0,
\begin{array}{ll} (\text{УЧП}) & u_{\tau} = Du_{\xi\xi}, & -\infty < \xi < \infty, & 0 < \tau < \infty, \\ (\text{HV}) & u(\xi, 0) = e^{-\frac{\epsilon}{2}}, & -\infty < \xi < \infty, \end{array}
\frac{dy}{dx} = \frac{B - \sqrt{B^2 - 4AC}}{2A} = -\frac{\sqrt{-4x^2y^2}}{2y^2} = -i\frac{x}{y},
\]

Крім того, S. S. Розв'язуючи їх, отримуємо перші інтеграли

\[
u(\xi, \tau) = \frac{1}{2\sqrt{D\pi\tau}} \int_{-\infty}^{+\infty} e^{-\beta^2} e^{-(\xi-\beta)^2/4D\tau} d\beta.
y^2 - ix^2 = \text{const},
u_{tt} = \alpha^2 u_{xx}
\xi(x, y) = y^2 + ix^2,
u_{\xi\eta} = \psi(\xi, \eta, u, u_{\xi}, u_{\eta}).
\alpha = \frac{\xi + \eta}{2} = y^2
\xi(x, y) = y^2, \quad \eta(x, y) = x^2.
\overline{A}u_{28} + \overline{B}u_{2n} + \overline{C}u_{nn} + \overline{D}u + \overline{E}u_{n} + \overline{F}u = \overline{G}
T \sin \theta_2 - T \sin \theta_1 =
u_{s\xi} + u_{\eta\eta} = \frac{-\epsilon u_{\eta} - \eta u_{\xi}}{2\xi\eta}.
\Delta x \rho u_{tt} = T \left[ u_x \left( x + \Delta x, \ t \right) - u_x \left( x, \ t \right) \right] + \Delta x F \left( x, \ t \right) - \frac{1}{\Delta x \rho u_t} \left( x, \ t \right) - \Delta x \gamma \dot{u} \left( x, \ t \right)
u_t = u_{xx} + u_{yy},
u_{tt} = \alpha^{3} u_{xx} - \beta u_{t} - \gamma u + F(x, t).
u_{tt} = u_{xx} + u_{yy}
(16.2) u_{tt} = \alpha^2 u_{xx},
I = \int_0^1 x^2 dx
u_{tt} = ku_{xx}
I = \int_{a}^{b} f(x) \, dx
u_{tt} = \frac{\partial}{\partial x} [\alpha^2(x) u_x],
\]

\widehat{I} =
\Big[ \begin{array}{ll} \mbox{относительное число дротнков} \mbox{ дротнков} \mbox{ } \mbox{]} \times (Площадь области $R$) \mbox{ } \end{array}

$$
u(x, 0) = f(x)
x_i = a + (b - a) r_i.
i_x + Cv_t + Gv = 0, \\ v_x + Li_t + Ri = 0,
Mr_0 + K = 571
i_{xx} + Gv_x - CLi_{tt} - CRi_t = 0.
Mr_2 + K = 1607
v_{x} = -Li_{t} - Ri,
r_2 = 43
i_{t,s} = CLi_{tt} + (CR + GL)i_t + GRi.
$$

Оскільки P = 100, усі випадкові числа лежать між нулем і 100 (так званий клас дедукції за модулем 100). Розміщуючи десяткову крапку перед кожним із цих чисел, ми отримуємо послідовність випадкових чисел, що лежать між нулем і одним $r_0=0.15;\ r_1=0.71;\ r_2=0.43;\ r_3=0.07.$ (Ми також позначали ці числа $r_L$ )

# ПРИМІТКИ

1. Оскільки наш генератор випадкових чисел має P=100, дедукції можуть приймати лише значення $0,1,2,\ldots,99$ Тому випадкові числа в нашому процесі дуже скоро почнуть повторюватися. Насправді, якщо продовжити наведений вище приклад, отримаємо

15, 71, 43, 7, 43, 7, 43, 7... (період із двох чисел).

Це означає, що наш метод не дуже хороший. Ідеально, якщо алгоритм випадково видає весь клас відрахувань $\{0, 1, 2, \ldots, 99\}$ до початку повторних спроб. Математично можна довести, що якщо числа M, K і P обираються за певними правилами, то незалежно від вибору першого випадкового числа $r_o$ алгоритм генеруватиме весь клас дедукцій. Тож якщо ви оберете P дуже великим (наприклад, $2^{40}$ ), можна бути впевненим, що (з практичної точки зору) цей процес ніколи не повториться.

2. Звичайні чисельні методи (наприклад, метод Сімпсона), загалом вищі за якістю за метод Монте-Карло. Але якщо ми говоримо про обчислення багатовимірних інтегралів типу

$$
v_{xx} = CLv_{tt} + (CR + GL)v_t + GRv.
$$

тоді краще використовувати метод Монте-Карло.

3. Можливо отримати випадкові числа з різними законами розподілу (і не лише рівномірними на сегменті [0, 1]).
Комп'ютерні програми можуть генерувати випадкові числа з біноміальними,
нормальними, гамма-розподілами та кратними

дистрибуції. Ці розподіли зустрічаються в різних галузях природничих
наук і стають необхідними для моделювання реального світу.

\section{ЗАВДАННЯ}\label{ux437ux430ux432ux434ux430ux43dux43dux44f-11}

\begin{enumerate}
\def\labelenumi{\arabic{enumi}.}
\tightlist
\item
  Напишіть програму для оцінки інтегралу
\end{enumerate}

\[
v_{tt} = \alpha^2 v_{xx}, \\ i_{tt} = \alpha^2 i_{xx}, \qquad \alpha^2 = 1/CL,
\]

Яка точність вашої відповіді? Побудуйте графік точності та кількості
кидків. Чи збігається цей процес?

\begin{enumerate}
\def\labelenumi{\arabic{enumi}.}
\setcounter{enumi}{1}
\tightlist
\item
  Отримати послідовність випадкових чисел за допомогою алгоритму
\end{enumerate}

\[
u_{tt} = u_{xx} - u_t
\]

\begin{enumerate}
\def\labelenumi{\arabic{enumi}.}
\setcounter{enumi}{2}
\tightlist
\item
  Напишіть програму (або принаймні блок-схему) для оцінки інтегралу
\end{enumerate}

Що ви можете сказати про точність?

\begin{enumerate}
\def\labelenumi{\arabic{enumi}.}
\setcounter{enumi}{3}
\tightlist
\item
  Як отримати послідовність випадкових точок, що лежать всередині
  трикутника T?
\end{enumerate}

\pandocbounded{\includegraphics[keepaspectratio]{images/_page_319_Figure_13.jpeg}}

\pandocbounded{\includegraphics[keepaspectratio]{images/_page_319_Figure_14.jpeg}}

\begin{enumerate}
\def\labelenumi{\arabic{enumi}.}
\setcounter{enumi}{4}
\tightlist
\item
  Як отримати послідовність випадкових чисел, розподіл яких показаний на
  рисунку?
\end{enumerate}

\pandocbounded{\includegraphics[keepaspectratio]{images/_page_319_Figure_16.jpeg}}

Іншими словами, як отримати послідовність чисел \{0, 1, 2\}, якщо
ймовірність виникнення 0 і 2 дорівнює 0,25, а ймовірність появи 1 ---
0,5?

Лекція 43

\section{РОЗВ'ЯЗАННЯ ДИФЕРЕНЦІАЛЬНИХ РІВНЯНЬ У ПОХІДНИХ ПОХІДНИХ
МОНТЕ-КАРЛО}\label{ux440ux43eux437ux432ux44fux437ux430ux43dux43dux44f-ux434ux438ux444ux435ux440ux435ux43dux446ux456ux430ux43bux44cux43dux438ux445-ux440ux456ux432ux43dux44fux43dux44c-ux443-ux43fux43eux445ux456ux434ux43dux438ux445-ux43fux43eux445ux456ux434ux43dux438ux445-ux43cux43eux43dux442ux435-ux43aux430ux440ux43bux43e}

МЕТА ЛЕКЦІЇ: Показати, як побудувати таку гру, результатом якої є
приблизне розв'язання диференціального рівняння. Одна з цих ігор (модель
випадкової ходьби) веде до скінченної різниці апроксимації задачі
Діріхле в квадраті. Цю модель можна узагальнити для отримання розв'язків
інших задач.

У попередній лекції було сказано, що можливо винайти такі ігри на удачу,
результатом яких є розв'язки (можливо, приблизні) задач для рівнянь з
частковими коефіцієнтами проникнення. У цій лекції ми покажемо, як
побудувати таку гру для наближеного розв'язання задачі Діріхле (PPP)
\(u_{xx}\) + \(u_{yy} = 0 , 0 < x < 1, 0 < y < 1\),

(GI)

\[
\begin{cases} 1, & \text{на верхней стороне квадрата,} \\ 0, & \text{на боковых и пижней сторонах квадрата.} \end{cases}
R(A) = g_1 P_A(p_1) + g_2 P_A(p_2) + \dots + g_{12} P_A(p_{12}).
\]

(HУ)

\[
R(A) = 1 \cdot (0.04) + 1 \cdot (0.15) + 1 \cdot (0.03) + 0 \cdot (0.06) + \dot + 0 \cdot (0.04).
R(A) = \frac{1}{4} [R(B) + R(C) + R(D) + R(E)].
R(A) = \frac{1}{4} [R(B) + R(C) + R(D) + R(E)]
\]

Как выглядит в различные моменты времени решение следующей задачи

\[
u_{i,j} = \frac{1}{4} \left[ u_{i-1,j} + u_{i+1,j} + u_{i,j-1} + u_{i,j+1} \right],
\]

\begin{verbatim}
\begin{aligned} & (\text{УЧ}\Pi) & u_{tt} = u_{xx}, & 0 < x < 1, \\ & (\text{ГУ}) & \begin{cases} u(0, t) = 0, & 0 < t < \infty, \\ u(1, t) = \sin t, & 0 < t < \infty, \end{cases} \\ & (\text{НУ}) & \begin{cases} u(x, 0) = 0, & 0 \le x \le 1, \\ u_t(x, 0) = 0 \end{cases}
\end{verbatim}

(PPP)

\[
u(x, t) = e^{ax+bt}
u_{xx}=u_{l,\ j+1}-2u_{l,\ j}+u_{l,\ j-1},
(\text{УЧП}) \qquad u_{tt} = c^2 u_{xx}, \quad -\infty < x < \infty, \quad 0 < t < \infty,
u_{i, j} = \frac{u_{i, j+1} + u_{i, j-1} + \sin x_j (u_{i+1, j} + u_{i-1, j})}{2 (1 + \sin x_j)}.
(\text{HY}) \quad \begin{cases} u(x, 0) = \int (x), \\ u_t(x, 0) = g(x), \end{cases} \quad -\infty < x < \infty.
\]

(i, j+1)

\[
u(x, t) = \frac{1}{2} [f(x - ct) + f(x + ct)] + \frac{1}{2\varepsilon} \int_{x-ct}^{x+ct} g(\xi) d\xi
u(A) = g_1 P_A(p_1) + g_2 P_A(p_2) + \dots + g_N P_A(p_N),
u_{tt} = c^{2}u_{xx}, \qquad -\infty < x < \infty, \quad 0 < t < \infty,
u_{xx} + u_{yy} = 0
(17.1) \quad (HY) \quad \begin{cases} u(x, 0) = f(x), \\ u_{t}(x, 0) = g(x), \end{cases} \quad -\infty < x < \infty,
u_{xx} + x^2 u_{yy} = 0
\xi = x + ct,
(\text{UBP}) \qquad
\eta = x - ct,
\]

\begin{enumerate}
\def\labelenumi{\arabic{enumi}.}
\setcounter{enumi}{4}
\tightlist
\item
  Разработайте схему метода Монте-Карло для решения смешанной
  параболической задачи
\end{enumerate}

\[
u_{ti} = c^2 u_{xx}
\]

\begin{verbatim}
\right. \end{array}
\end{verbatim}

\[
u_{\xi n} = 0
J[y] = \int_{b}^{a} F(x, y, y') dx,
u_{x} = u_{\xi} + u_{\eta}, u_{t} = c (u_{\xi} - u_{\eta}), u_{xx} = u_{\xi\xi} + 2u_{\xi\eta} + u_{\eta\eta}, u_{tt} = c^{2} (u_{\xi\xi} - 2u_{\xi\eta} + u_{\eta\eta}).
J[y] = \int_{0}^{1} [y^{2}(x) + y'^{2}(x)] dx.
u_{\xi\eta} = 0
\frac{df(x)}{dx} = 0
u_{\eta}(\xi, \eta) = \varphi_{1}(\eta)
T = \int_{0}^{T} dt = \int_{a}^{L} \frac{dt}{ds} ds = \int_{0}^{L} \frac{ds}{v} = \frac{1}{\sqrt{2mg}} \int_{0}^{L} \frac{ds}{\sqrt{y}} = \frac{1}{\sqrt{2mg}} \int_{a}^{b} \sqrt{\frac{1+y'^{2}}{y}} dy
u(\xi, \eta) = \varphi(\eta) + \psi(\xi),
J[y] = \int_{a}^{b} F(x, y, y') dx.
u_{\xi\eta} = 0
\frac{df(x)}{dx} = 0
u(\xi, \eta) = \varphi(\eta) + \psi(\xi),
J[y] = \int_{a}^{y} F(x, y, y') dx
u(\xi, \eta) = \sin \eta + \xi^2,
J[y] = \int_{a}^{b} F(x, y, y') dx
\xi = x - ct,
y(a) = A
n = x + ct
J[\overline{y}] \leqslant J[\overline{y} + \varepsilon \eta]
u(\xi, \eta) = \varphi(\eta) + \psi(\xi).
\varphi(\varepsilon) = J[\bar{y} + \varepsilon \eta]
u(x, t) = \varphi(x-ct) + \psi(x+ct).
\frac{d\varphi\left(\varepsilon\right)}{d\varepsilon} = \frac{d}{d\varepsilon} J\left[\overline{y} + \varepsilon\eta\right]|_{\varepsilon=0} =
u(x, t) = \sin(x-ct)
\]

(Читач має робити це самостійно.) Інтегруючи на частини, ми отримуємо

\[
u(x, t) = (x + ct)^2
\]

Оскільки цей інтеграл є нульовим у будь-якій функції \eta(x) , що
задовольняє граничні умови \eta(a) = \eta(b) = 0 , ми Р и с

\pandocbounded{\includegraphics[keepaspectratio]{images/_page_330_Figure_6.jpeg}}

\begin{verbatim}
. 44 3. Граф функції  J[y+] \in \eta  ] в околиці  \varepsilon = 0  .
\end{verbatim}

Ми приходимо до висновку, що друга частина підоб'єктивного виразу має
дорівнювати нулю, тобто

(44 2)

\[
u(x, t) = \sin(x - ct) + (x + ct)^2
\]

(Рівняння Ейлера--Лагранжа).

\begin{verbatim}
Рівняння (44 2) називається рівнянням Ейлера-Лагранжа, і хоча загалом воно здається комплексним, коли в нього підставляється конкретна функція F(x, y, y'), воно стає звичайним диференціальним рівнянням другого порядку відносно невідомої функції  $$\overline{y}$(x)$  . Отже, щоб визначити мінімізуючу функцію  $\overline{y}$  , необхідно розв'язати рівняння Ейлера-Лагранжа.
\end{verbatim}

Отже, ми показали, що

\begin{verbatim}
якщо функція y(x) мінімізує функціональний  $J[y] = \int_a^b F(x, y, y') dx  (у класі гладких функцій з граничними умовами y(a) = A; y(b) = B), то вона повинна задовольняти рівняння$
\end{verbatim}

\[
u(x, t) = \varphi(x - ct) + \psi(x + ct)
\]

(з граничною умовою).

\begin{verbatim}
(Щоб спростити позначення, ми прибрали лінію над  $\bar{y}$  .) Щоб проілюструвати теорію, наведемо приклад.

Пошук мінімальної функціональності  $J[y] = \int_{0}^{1} [y^2 + y'^2] dx$.
\end{verbatim}

Спробуємо знайти функцію \(y(x)\), яка проходить через точки (0,0) і
(1,1) і мінімізує функціональність

\[
u(x, 0) = f(x),
\]

З контексту зрозуміло, що шукана функція має бути диференційовною,
оскільки інтегративний вираз залежить від y'. Щоб визначити \bar\{y\}(x)
, запишемо рівняння Ейлера--Лагранжа (разом із граничною умовою
\(y(0)\)=0 та \(y(1)\)=1).

\[
\varphi(x) + \psi(x) = f(x), \\ -c\varphi'(x) + c\psi'(x) = g(x).
\]

\(y(0)\) = 0, \(y(1)\) = 1. Але F(x, y, y') = y\^{}2 + y'\^{}2 і, отже,

\[
-c\varphi(x) + c\psi(x) = \int_{x_0}^{x} g(\xi) d\xi + K.
\]

(Продифференцировали F по y и y'. )

Уравнение Эйлера --- Лагранжа принимает вид

\[
\psi(x) = \frac{1}{2} f(x) + \frac{1}{2c} \int_{x_0}^{x} g(\xi) d\xi + K/2c,
\]

откуда

\[
\varphi(x) = \frac{1}{2} f(x) - \frac{1}{2c} \int_{x_0}^{x} g(\xi) d\xi - K/2c.
\]

Решая это простое дифференциальное уравнение с граничными условиями
\(y(0)\)=0 и \(y(1)\)=1, получаем

\[
u(x, t) = \frac{1}{2} \left[ f(x-ct) + f(x+ct) \right] + \frac{1}{2c} \int_{x-ct}^{x+ct} \mathbf{g}(\xi) d\xi.
\]

.

График этой функции изображен на рис. 44.4. Ясно, что любая другая
гладкая функция, удовлетворяющая тем же граничным условиям, даст большее
значение функционала \(J(y)\).

\pandocbounded{\includegraphics[keepaspectratio]{images/_page_332_Figure_5.jpeg}}

\begin{verbatim}
РИС. 44 4. Все допустимые гладкие кривые удовлетворяют граничным условиям y(0) = 0 и y(1) = 1: a—допустимая кривая;   $\delta$  —минимизирующая функция   y(x) = 0.42e^x - 0.42e^{-x}  .
\end{verbatim}

\section{ЗАМЕЧАНИЯ}\label{ux437ux430ux43cux435ux447ux430ux43dux438ux44f-17}

\begin{enumerate}
\def\labelenumi{\arabic{enumi}.}
\item
  Уравнение Эйлера---Лагранжа аналогично равенству нулю производной в
  дифференциальном исчислении. Если читатель помнит, то на самом деле
  равенство нулю производной не является достаточным условием
  экстремума. Например, производная функции \(f(x)\) = x\^{}3 при x = 0
  равна нулю, но эта точка не является нп точкой минимума, ни точкой
  максимума. Так же и с уравнением Эйлера---Лагранжа: оно является
  только необходимым условием экстремума, но не достаточным. Значит,
  экстремум может достигаться на других функциях. Очень часто, однако,
  решение уравнения Эйлера---Лагранжа по существу задачи дает локальный
  или даже глобальный минимум.
\item
  Если у (х) минимизирует функционал
\end{enumerate}

\[
u(x, 0) = \sin x,
\]

то площадь поверхности, образованной вращением кривой \(y(x)\) вокруг
оси x, будет минимальной (рис. 44.5). Решением уравнения Эйлера ---
Лагранжа в этом случае будет

цепная линия

\[
u(x, t) = \frac{1}{2} [\sin(x-ct) + \sin(x+ct)].
\]

где константы \alpha   и \beta   определяются из условий \(y(a)\) = A и
\(y(b)\) = B (т. е. цепная линия должна проходить через заданные точки).
Решить уравнение Эйлера---Лагранжа в данном случае достаточно трудно,
поэтому мы рекомендуем читателю только проверить, что цепная линия ему
удовлетворяет.

\begin{enumerate}
\def\labelenumi{\arabic{enumi}.}
\setcounter{enumi}{2}
\tightlist
\item
  В этой лекции мы нашли, что функция
\end{enumerate}

\[
\frac{\sin x}{2}
\]

доставляет минимум функционалу

\[
u(x, 0) = \begin{cases} 1, & -1 < x < 1, \\ 0, & \end{cases}
\]

При этом оказывается, что J{[}\(\overline{y}\){]} = 0.46 . Если
подставить в функционал любую другую гладкую функцию, график которой
проходит через точки (0, 0) и (1, 1), то значение функционала J{[}y{]}
станет больше.

\pandocbounded{\includegraphics[keepaspectratio]{images/_page_333_Picture_8.jpeg}}

РИС. 44 5. Минимальная поверхность вращения.

\begin{enumerate}
\def\labelenumi{\arabic{enumi}.}
\setcounter{enumi}{3}
\tightlist
\item
  Основные законы физики чаще всего формулируются на языке вариационных
  принципов, а не дифференциальных уравнений. В качестве примеров можно
  привести принцип Ферма (свет при распространении из одной точки в
  другую выбирает путь, которому соответствует наименьшее время
  распространения) или принцип Гамильтона (в консервативном поле частица
  движется так, что интеграл действия
\end{enumerate}

\[
u(x, 0) = 0,
\]

(кинетическая энергия--- потенциальная энергия) dt

будет минимальным). Таким образом, физические явления развиваются только
так, что эти функционалы принимают минимальные значения.

Основные идеи вариационного исчисления можно распространить на
функционалы, зависящие от функции нескольких независимых переменных,
вида

\[
u_t(x, 0) = \sin x.
\]

Соответствующее уравнение Эйлера --- Лагранжа для такого функционала
будет уже уравнением с частными производными

\[
u(x, t) = \frac{1}{2c} \int_{x-ct}^{x+2t} \sin \xi \, d\xi = \frac{1}{2c} \left[ \cos (x+ct) - \cos (x-ct) \right].
\]

Такого рода функционалы мы рассмотрим в следующей лекции. Однако общая
философия метода будет иной. Некоторым новым способом мы найдем функцию
\(u(x, y)\), которая минимизирует функционал, а поскольку она является
решением уравнения Эйлера --- Лагранжа, то тем самым мы найдем решение
уравнения с частными производными. Другими словами, мы будем искать
решение уравнения с частными производными методом минимизации
функционала. В этой лекции все было наоборот: задачу минимизации
функционала, зависящего от функции одной переменной, мы сводили к
решению уравнения Эйлера --- Лагранжа. Методы решения дифференциальных
уравнений путем минимизации соответствующих функционалов принято
называть прямыми методами вариационного исчисления. Методы минимизации
функционалов путем решения соответствующих уравнений Эйлера --- Лагранжа
называются непрямыми методами вариационного исчисления. Следующая лекция
посвящена хорошо известному прямому методу Ритца.

\section{ЗАДАЧИ}\label{ux437ux430ux434ux430ux447ux438-20}

\begin{enumerate}
\def\labelenumi{\arabic{enumi}.}
\item
  Среди всех кривых, удовлетворяющих условиям \(y(0)\) = 0, и \(y(1)\) =
  1, найдите ту, которая минимизирует функционал

  (\text{УЧП}) \qquad u\_\{tt\} = u\_\{x+\}, \quad -\infty \textless{} x
  \textless{} \infty, \quad 0 \textless{} t \textless{} \infty,
  \textbackslash{} (\text{HV}) \qquad \left\{

  \begin{array}{ll} u(x, \ 0) = e^{-x^2} \\ u_t(x, \ 0) = 0. \end{array}

  \right. \quad -\infty \textless{} x \textless{} \infty,
  \textbackslash{} \textbackslash end\{array\}
\end{enumerate}

Как вы интерпретируете получелный результат? Чему равна величина
J{[}\(\overline{y}\){]} ? Каков смысл величины J{[}\(\overline{y}\){]} ?

\begin{verbatim}
2. Кинетическая энергия колеблющейся материальной точки определяется выражением   KE = \frac{1}{2} m\dot{y}^2  , где y = dy/dt, а потенциальная равна   PE = \frac{1}{2} ky^2  .
\end{verbatim}

\pandocbounded{\includegraphics[keepaspectratio]{images/_page_335_Picture_2.jpeg}}

Принцип Гамильтона утверждает, что материальная точка будет двигаться
так, что интеграл

\[
u_{tt} = u_{xx}, -\infty < x < \infty, 0 < t < \infty,
\]

будет минимален. Если этот закоп справедлив, то получите
дифференциальное уравнение движения материальной точки.

\begin{verbatim}
3. Покажите, что в классе гладких функций, удовлетворяющих граничным условиям y(0)=0 и   y($\pi$/2)=1  , минимум функционала
\end{verbatim}

\[
\begin{cases} u(x, 0) = 0, \\ u_t(x, 0) = xe^{-\lambda^2}, -\infty < x < \infty, \end{cases}
\]

достигается на функции \$\$\(\overline{y}\)(x)\$ = \sin x\$ . Вычислите
J(\sin x) . 4. Исходя из функционала

\[
u_{tt} = c^2 u_{xx}, \quad 0 < x < \infty, \quad 0 < \ell < \infty,
\]

получите уравнение Эйлера --- Лагранжа

\[
\begin{cases} u(x, 0) = f(x), \\ u_t(x, 0) = g(x). \end{cases}
\]

\section{ВАРИАЦИОННЫЕ МЕТОДЫ РЕШЕНИЯ УРАВНЕНИЙ С ЧАСТНЫМИ
ПРОИЗВОДНЫМИ}\label{ux432ux430ux440ux438ux430ux446ux438ux43eux43dux43dux44bux435-ux43cux435ux442ux43eux434ux44b-ux440ux435ux448ux435ux43dux438ux44f-ux443ux440ux430ux432ux43dux435ux43dux438ux439-ux441-ux447ux430ux441ux442ux43dux44bux43cux438-ux43fux440ux43eux438ux437ux432ux43eux434ux43dux44bux43cux438}

ЦЕЛЬ ЛЕКЦИИ: Показать, как можно решать дифференциальное уравнение,
рассматривая его как уравнение Эйлера---Лагранжа для некоторого
функционала и находя некоторым новым методом функцию, минимизирующую
этот функционал. Тогда минимизирующая функция будет решением уравнения с
частными производными. Задача, конечно, состоит в том, чтобы найти такой
функционал, для которого исходное уравнение является уравнением
Эйлера---Лагранжа. Хорошо известен результат (теорема о минимуме
энергии), который гласит, что нахождение решения и некоторой
эллиптической краевой задачи

\[
u(x, t) = \frac{1}{2} [f(x-ct) + f(x+ct)] + \frac{1}{2c} \int_{x-ct}^{x+ct} g(\xi) d\xi
\]

, в области D , u = 0 , на границе D ,

эквивалентно нахождению такой функции u (тоже обращающейся в нуль на
границе области D), которая минимизирует функционал потенинальной
энергии

\[
\begin{aligned} & (\text{Y-I}\Pi) & u_{tt} = c^2 u_{xx}, & -\infty < x < \infty, & 0 < t < \infty \\ & (\text{HY}) & \begin{cases} x(x, 0) = f(x), \\ u_t(x, 0) = g(x), \end{cases} & -\infty < x < \infty, \end{aligned}
\]

То есть, \(\Delta\) u = f ---уравнение Эйлера---Лагранжа для J{[}u{]}.
Функция, приближенно минимизирующая функционал J{[}u{]}, ищется методом
Ритца; тем самым мы получаем решение (приближениое) уравнения с частными
производными. В этой лекции мы познакомимся с методом Ритца и посмотрим,
как с помощью этого метода можно и и и и и и и уравнения.

Очень хорошо решать граничные задачи (как, например, задачу о
закрепленной мембране) путем отыскания гладкой поверхности, которая
минимизирует потенциальную энергию мембраны. То есть, если мы
рассматриваем исходное уравнение как уравнение Эйлера --- Лагранжа для
искоторого функционала J{[}u{]}, то решение дифференциального уравнения
можно получить путем минимизации функционала (поскольку минимизирующая
функция функционала является решением соответствующего уравнения

Эйлера --- Лагранжа). В лекции 44 мы рассматривали такие функционалы, у
которых уравнение Эйлера --- Лагранжа было обыкновенным дифференциальным
уравнением. В этой лекции рассматриваются функционалы, для которых
уравнение Эйлера --- Лагранжа относится к классу уравнений с частными
производными. Например для функционала

дается формулой Даламбера

уравнение Эйлера---Лагранжа (подобно случаю обыкновенного
дифференциального уравнения в последней лекции) имеет вид

Дадим интерпретацию этой формулы для двух частных случаев.

СЛУЧАЙ 1 (Начальное смещение отлично от нуля, начальная скорость равна
нулю).

Пусть начальные условия имеют вид

и, следовательно, для решения задачи Дирихле в единичном квадрате

(УЧП)

\(u_1(x, 0) = 0.\) Тогда по формуле Даламбера получаем

, 0 \textless{} x \textless{} 1 , 0 \textless{} y \textless{} 1 , (ГУ) u
= g , на границе квадрата,

можно, наоборот, найти такую функцию \(u(x, y)\), которая минимизирует
J{[}u{]} и равна g на границе. Вероятно, не стоит удивляться тому, что
функционал

Обратимся к рис. 181. Решение u в точке \((x_0, t_0)\) можно
интерпретировать как среднее значение смещений \(f(x)\) в точках Р

\pandocbounded{\includegraphics[keepaspectratio]{images/_page_132_Figure_10.jpeg}}

ИС. 18.1. Интерпретация решения
\(u(x, t) = \frac{1}{2} [f(x-ct)+f(x+ct)]\) в xt-плоскости.

\((x_{\mathrm{ }e}-ct_{\mathrm{ }o},\ 0)\) и
\((x_{\mathrm{ }e}+ct_{\mathrm{ }o},\ 0)\) . Эти точки являются точками
пересечения прямых

представляет собой потенциальную эпергию мембраны и фактически мы
находим поверхность с минимальной потенциальной энерешей. Вопрос,
конечно, состоит в следующем: если дано дифференциальное уравнение, как
мы находим функционал, выражающий потенциальную энергию решения? Ответ
на этот вопрос дает хорошо известная теорема (теорема о минимуме
энергии), которая утверждает:

Решение и задачи Дирихле

, \(x+ct=x_0+ct_0\)

е осью х Указанные прямые называются характеристиками волнового
уравнения.

Используя такую интерпретацию формулы Даламбера, решим задачу

(18.3)

\^{}\{\Lambda\}u = f в области \^{}\{D\} (ГУ) \^{}\{U\}=0 на границе
\^{}\{D\}

является той же функциен, которая минимизирует (среди функций,
удовлетворяющих граничному условию u=0) энергегический функционал

Подобные теоремы справедливы и для других типов граничных условий. Чтобы
помочь читателю понять теорему, рассмотрим пример.

Замена уравнения Пуассона функционалом потенциальной энергии

Рассмотрим задачу Дирихле

(45.1)

\[
\begin{cases} u_{tx}, \quad 0 = \begin{cases} 1, \quad -1 < x < 1, \\ 0 \quad \text{в остальных точках,} \end{cases} \end{cases}
\]

, \(0 < x < 1\) , \(0 < y < 1\) , \(0 < y < 1\) , \(0 < y < 1\) ,
\(0 < y < 1\) , \(0 < y < 1\) , \(0 < y < 1\) , \(0 < y < 1\) ,
\(0 < y < 1\) , \(0 < y < 1\) , \(0 < y < 1\) , \(0 < y < 1\) ,
\(0 < y < 1\) , \(0 < y < 1\) , \(0 < y < 1\) , \(0 < y < 1\) ,
\(0 < y < 1\) , \(0 < y < 1\) , \(0 < y < 1\) , \(0 < y < 1\) ,
\(0 < y < 1\) , \(0 < y < 1\) , \(0 < y < 1\) , \(0 < y < 1\) ,
\(0 < y < 1\) , \(0 < y < 1\) , \(0 < y < 1\) , \(0 < y < 1\) ,
\(0 < y < 1\) , \(0 < y < 1\) , \(0 < y < 1\) , \(0 < y < 1\) ,
\(0 < y < 1\) , \(0 < y < 1\) , \(0 < y < 1\) , \(0 < y < 1\) ,
\(0 < y < 1\) , \(0 < y < 1\) , \(0 < y < 1\) , \(0 < y < 1\) ,
\(0 < y < 1\) , \(0 < y < 1\) , \(0 < y < 1\) , \(0 < y < 1\) ,
\(0 < y < 1\) , \(0 < y < 1\) , \(0 < y < 1\) , \(0 < y < 1\) ,
\(0 < y < 1\) , \(0 < y < 1\) , \(0 < y < 1\) , \(0 < y < 1\) ,
\(0 < y < 1\) , \(0 < y < 1\) , \(0 < y < 1\) , \(0 < y < 1\) ,
\(0 < y < 1\) , \(0 < y < 1\) , \(0 < y < 1\) , \(0 < y < 1\) ,
\(0 < y < 1\) , \(0 < y < 1\) , \(0 < y < 1\) , \(0 < y < 1\) ,
\(0 < y < 1\) , \(0 < y < 1\) , \(0 < y < 1\) , \(0 < y < 1\) ,
\(0 < y < 1\) , \(0 < y < 1\) , \(0 < y < 1\) , \(0 < y < 1\) ,
\(0 < y < 1\) , \(0 < y < 1\) , \(0 < y < 1\) , \(0 < y < 1\) ,
\(0 < y < 1\) , \(0 < y < 1\) , \(0 < y < 1\) , \(0 < y < 1\) ,
\(0 < y < 1\) , \(0 < y < 1\) , \(0 < y < 1\) , \(0 < y < 1\) ,
\(0 < y < 1\) , \(0 < y < 1\) , \(0 < y < 1\) , \(0 < y < 1\) ,
\(0 < y < 1\) , \(0 < y < 1\) , \(0 < y < 1\) , \(0 < y < 1\) ,
\(0 < y < 1\) , \(0 < y < 1\) , \(0 < y < 1\) , \(0 < y < 1\) ,
\(0 < y < 1\) , \(0 < y < 1\) , \(0 < y < 1\) , \(0 < y < 1\) ,
\(0 < y < 1\) , \(0 < y < 1\) , \(0 < y < 1\) , \(0 < y < 1\) ,
\(0 < y < 1\) , \(0 < y < 1\) , \(0 < y < 1\) , \(0 < y < 1\) ,
\(0 < y < 1\) , \(0 < y < 1\) , \(0 < y < 1\) , \(0 < y < 1\) ,
\(0 < y < 1\) , \(0 < y < 1\) , \(0 < y < 1\) , \(0 < y < 1\) ,
\(0 < y < 1\) , \(0 < y < 1\) , \(0 < y < 1\) , \(0 < y < 1\) ,
\(0 < y < 1\) , \(0 < y < 1\) , \(0 < y < 1\) , \(0 < y < 1\) ,
\(0 < y < 1\) , \(0 < y < 1\) , \(0 < y < 1\) , \(0 < y < 1\) ,
\(0 < y < 1\) , \(0 < y < 1\) , \(0 < y < 1\) , \(0 < y < 1\) ,
\(0 < y < 1\) , \(0 < y < 1\) , \(0 < y < 1\) , \(0 < y < 1\) ,
\(0 < y < 1\) , \(0 < y < 1\) , \(0 < y < 1\) , \(0 < y < 1\) ,
\(0 < y < 1\) , \(0 < y < 1\) , \(0 < y < 1\) , \(0 < y < 1\) ,
\(0 < y < 1\) , \(0 < y < 1\) , \(0 < y < 1\) , \(0 < y < 1\) ,
\(0 < y < 1\) , \(0 < y < 1\) , \(0 < y < 1\) , \(0 < y < 1\) ,
\(0 < y < 1\) , \(0 < y < 1\) , \(0 < y < 1\) ,

(см. рис. 45.1).

\pandocbounded{\includegraphics[keepaspectratio]{images/_page_338_Figure_5.jpeg}}

РИС. 45.1. Уравнение Пуассона внутри квадрата.

Для этой задачи функционал энергии J{[}u{]} записывается в виде

(45.2)

и, следовательно, для решения задачи (45.1) достаточно найти функцию
\(\overline{u}\) (среди функций, обращающихся в нуль на граннце),
минимизирующую J\{u\} . Этим замечанием завершается первая часть лекции.
В оставшейся части мы покажем, как находить минимизирующую функцию
\(\overline{u}\) для J{[}u{]}, используя метод Ритца.

\section{Метод Ритца минимизации
функционалов}\label{ux43cux435ux442ux43eux434-ux440ux438ux442ux446ux430-ux43cux438ux43dux438ux43cux438ux437ux430ux446ux438ux438-ux444ux443ux43dux43aux446ux438ux43eux43dux430ux43bux43eux432}

Этот метод --- один из многих вариационных методов, описанных в
литературе по вариационному исчислению. Идея, предложенная математиком
В. Ритцем, очень проста. Метод состоит из следующих шагов.

ШАГ 1. Выбираем n и функцию u, минимизирующую функционал

\[
u_{t}(x, \quad 0) = 0.
\]

ищем в виде

\[
u(x, 0) = 0,
\]

где функции \phi\_1, \phi\_2, \phi\_3, \ldots, \phi\_n принадлежат к
классу достаточно хороших функций и все обращаются в нуль на границе,
так что из них можно построить разумное приближение для решения задачи.
Эти функции принято называть пробными функциями. Типичный набор пробных
функций для задачи Дирихле в единичном квадрате приведен ниже:

\[
u(x, t) = \frac{1}{2c} \int_{x-ct}^{x+ct} g(\xi) d\xi.
\]

--- обращаются в нуль на границе, \varphi\_2(x, y) = x\varphi\_1(x, y),
\varphi\_3(x, y) = y\varphi\_1(x, y), \varphi\_4(x, y) =
x\^{}2\varphi\_1(x, y), \varphi\_5(x, y) = xy\varphi\_1(x, y),
\varphi\_6(x, y) = y\^{}2\varphi\_1(x, y), \vdots     \vdots

Иными словами, первые четыре приближения имеют вид

\[
u_{2}(x, y) = xy(1-x)(1-y)[a_{1} + a_{2}x],
\begin{aligned} & u_{ti} = c^{2}u_{xx}, & -\infty < x < +\infty, & 0 < t < \infty, \\ & (\text{НУ}) & \begin{cases} u_{t}(x, 0) = 0, \\ u_{t}(x, 0) = \begin{cases} 1, & -1 < x < 1, \\ 0, & \text{в остальных случаях.} \end{cases} \end{aligned} \end{cases}
u_{4}(x, y) = xy(1-x)(1-y)[a_{1} + a_{2}x + a_{3}y + a_{4}x^{2}]
\]

в координатах (x, t).

\pandocbounded{\includegraphics[keepaspectratio]{images/_page_134_Figure_7.jpeg}}

РИС. 18.4. Решение задачи (18.4) в xt-плоскости.

Задача (18.4) описывает поведение струны, которой сообщается рачальная
единичная скорость на отрезке -1 \textless{} x \textless{} 1. Смещение .
.

описывается формулой Даламбера

ШАГ 2. Функционал

является теперь функцией коэффициентов \(a_1a_2\)\dots \(a_n\) .
Следовательно, для того чтобы найти минимум функционала J, приравняем
нулю частные производные

\[
\frac{\partial J\left[u_{n}\right]}{\partial a_{n}} = 2 \int_{0}^{1} \int_{0}^{1} \left\{ \sum_{j=1}^{n} \left[ \frac{\partial \varphi_{j}}{\partial x} \frac{\partial \varphi_{n}}{\partial x} + \frac{\partial \varphi_{j}}{\partial y} \frac{\partial \varphi_{n}}{\partial y} \right] a_{j} + f \varphi_{n} \right\} dx dy = 0.
Aa = b
A_{ij} = \int_{0}^{1} \int_{0}^{1} \left[ \frac{\partial \varphi_{i}}{\partial x} \frac{\partial \varphi_{j}}{\partial x} + \frac{\partial \varphi_{i}}{\partial y} \frac{\partial \varphi_{j}}{\partial y} \right] dx dy,
b_i = -\int_0^1 \int_0^1 f(x, y) \varphi_i(x, y) dx dy,
u_n(x, y) = a_1 \varphi_1(x, y) + a_2 \varphi_2(x, y) + \dots + a_n \varphi_n(x, y),
\]

Это решение в различные моменты времени изображено на рис. 18.5.

\pandocbounded{\includegraphics[keepaspectratio]{images/_page_135_Figure_5.jpeg}}

РИС. 18.5. Решение задачи (18.4) в различные моменты времени.

\begin{verbatim}
Этим завершается наша интерпретация формулы Даламбера в плоскости переменных  \boldsymbol{x}  и t. В оставшейся части лекции мы займемся решением смешанной задачи для полубесконечной струны:
\end{verbatim}

(18.6)

\[
J[y] = \int_{0}^{1} [y^{2} + y'^{4}] dx,
\]

Единственное отличие в том, что аппроксимирующие функции \phi\_1,
\phi\_2, \phi\_3, \ldots, \phi\_n должны удовлетворять граничным
условиям:

\(\varphi_i(1) = 1,\) \(i = 1, 2, ..., n.\)

\begin{enumerate}
\def\labelenumi{\arabic{enumi}.}
\setcounter{enumi}{1}
\tightlist
\item
  В этой лекции не было докачано, что уравнением Эйлера --- Лагранжа для
  функционала, зависящего от функции двух переменных
\end{enumerate}

является иравнение Лапласа, но это доказательство проводится аналогично
тому, как это делалось в лекции 44 для функционала, зависящего от
функции одной переменной.

\begin{enumerate}
\def\labelenumi{\arabic{enumi}.}
\setcounter{enumi}{2}
\item
  В книгах по вариационному исчислению показано, как строить
  энергетический функционал не только для уравнения Пуассона с граничным
  условием Дирихле u=0, но и для многих других дифференциальных
  уравнений и многих других типов граничных условий. Следовательно,
  большое число краевых задач можно решить, минимизируя соответствующие
  энергетические функционалы.

  \begin{enumerate}
  \def\labelenumii{\arabic{enumii}.}
  \setcounter{enumii}{3}
  \tightlist
  \item
    Чем большее значение п мы выбираем, тем меньшее значение
    J{[}\(\bar{u}_n]\) получается. Значит, с увеличением n растет
    точность л приближенного решения дифференциального уравнения. Один
    из методов определения и, основан на последовательном вычислении
    функционала J{[}u\_n{]} для возрастающих значений n.~Вычисления
    прекращаются, если с ростом п функционал перестает практически
    уменьшаться.
  \end{enumerate}
\item
  Если и велико, то проведение вычислений по методу Ритца требует
  привлечения ЭВМ. На рис. 45.2 (с. 345) приведена блок-
\end{enumerate}

схема решения краевой задачи методом Ритца.

Пользователь должен позаботиться о подпрограмме, которая должна
вычислять значения пробных функций. Такая подпрограмма может иметь вид:

SUBROUTINE BC (X.Y.PHI)

\begin{itemize}
\tightlist
\item
  SUBROUTINE PROVIDED BY THE USER TO C - EVALUATE THE - FUNCTIONS
  PHI(1), PHI(2), \ldots, PHI(N) DIMENSION PHI(20)
\end{itemize}

\[
PHI(1) = X\*Y\*(1-X)\*(1-Y)
\]

PHI(2) = X\^{}*PHI(1) (используемые функции)

\[
PHI(3) = Y\*PHI(1)
\]

PHI(N) = (что получится)RETURN END

\section{\texorpdfstring{\textbf{ЗА}ДАЧИ}{ЗАДАЧИ}}\label{ux437ux430ux434ux430ux447ux438-21}

Найдите энергетический функционал для задачи

(УЧП)

\[
u(x, t) = \varphi(x - ct) + \psi(x + ct)
\]

, 0 \textless{} x \textless{} 1 , 0 \textless{} y \textless{} 1 , (ГУ) u
= 0 на границе.

\begin{enumerate}
\def\labelenumi{\arabic{enumi}.}
\setcounter{enumi}{1}
\tightlist
\item
  Как минимизировать функционал
\end{enumerate}

\[
\varphi(x-ct) = \frac{1}{2} f(x-ct) - \frac{1}{2c} \int_{x-ct}^{x-ct} g(\xi) d\xi - K/2c,
\]

методом Ритца?

УКАЗАНИЕ. Если ввести новую функцию \(z(x)\):

\[
\psi(x+ct) = \frac{1}{2} f(x+ct) + \frac{1}{2c} \int_{x-ct}^{x+ct} g(\xi) d\xi + K/2c.
\]

то можно заметить, что она удовлетворяет граничным условиям \(z(0)\) = 0
и \(z(1)\) = 0.1

\begin{itemize}
\tightlist
\item
  \begin{enumerate}
  \def\labelenumi{\arabic{enumi}.}
  \setcounter{enumi}{2}
  \tightlist
  \item
    Напишите программу для проведения вычислений по блоксхеме,
    изображенной на рис. 45 2.
  \end{enumerate}
\item
  \begin{enumerate}
  \def\labelenumi{\arabic{enumi}.}
  \setcounter{enumi}{3}
  \tightlist
  \item
    Покажите, что уравнение
  \end{enumerate}
\end{itemize}

\[
\phi(x-ct)
\]

является уравнением Эйлера --- Лагранжа для функционала

\[
u(x, t) = \varphi(x - ct) + \psi(x + ct)
\]

\begin{enumerate}
\def\labelenumi{\arabic{enumi}.}
\setcounter{enumi}{4}
\tightlist
\item
  Задачу Дирихле
\end{enumerate}

, \(0 < x < 1\) , \(0 < y < 1\) , \(u = 0\) , на границе квадрата,

можно решить методом конечного синус-преобразования (по переменной х). В
результате получим

где A=0.06 и B=0.04. Как найти потенциальную энергию этого решения? Мы
рекомендуем читателю вспомнить синуспреобразование и самостоятельно
получить решение \(u\left(x,y\right)\) .

*) При такой замене возникает особенность в т. x=1. Лучше ввести функцию
\(z(x)\) по формуле \(z(x)\)=\(y(x)\)=x. Прим. ред.

\pandocbounded{\includegraphics[keepaspectratio]{images/_page_343_Figure_2.jpeg}}

\begin{verbatim}
Далсе. Приближенное решение   u_n(x,y) = a_1 \varphi_1 + a_2 \varphi_2 + \dots + a_n \varphi_n  .
\end{verbatim}

В это время можно численно найти эначение функционала J{[}\(u_n\){]} .
Если необходимо найти значение решения \(u_n\)(x,y) в некоторой точке,
можно написать небольшую программу эля быполнения отой операции. При
необходимости можно повторить вычисления по этой программе с другими
значениями п

РИС 452. Блок-схема метода Ритца.

\section{РЕШЕНИЕ УРАВНЕНИЙ С ЧАСТНЫМИ ПРОИЗВОДНЫМИ МЕТОДАМИ ТЕОРИИ
ВОЗМУЩЕНИЙ}\label{ux440ux435ux448ux435ux43dux438ux435-ux443ux440ux430ux432ux43dux435ux43dux438ux439-ux441-ux447ux430ux441ux442ux43dux44bux43cux438-ux43fux440ux43eux438ux437ux432ux43eux434ux43dux44bux43cux438-ux43cux435ux442ux43eux434ux430ux43cux438-ux442ux435ux43eux440ux438ux438-ux432ux43eux437ux43cux443ux449ux435ux43dux438ux439}

ЦЕЛЬ ЛЕКЦИИ: Показать, как различные сложные задачи (нелинейные
уравнения с переменными коэффициентами, области неправильной формы и т.
д.) можно решить методом возмущения более простых задач. Иными словами,
мы хотим показать, как следует изменить решение простой задачи, чтобы
оно давало приближенное решение трудной задачи.

Например, мы покажем, как нелинейную задачу

(УЧП)

, \(0 < r < 1\) , \(0 \le \theta < 2\pi\) , \((\Gamma Y)\)
\(u(1, \theta) = \cos \theta\) , \(0 \le \theta \le 2\pi\) ,

можно приближенно решить методом возмущения решения линейной задачи
Дирихле 1)

(УЧП)

\[
,   0 < r < 1  ,
\]

\begin{verbatim}
(ГУ)   u(1, $\theta$) = \cos $\theta$  ,   0 \le $\theta$ \le 2$\pi$  .
\end{verbatim}

Очень часто задачу, решение которой не известно, можно путем
непрерывного изменения каких-то параметров свести к близкой, но легко
решаемой задаче. Тогда и решение исходной задачи будет близко решению
модифицирующей задачи.

Например, уравнение Лапласа

\[
\varphi(x-ct) = -\frac{1}{2}f(ct-x) - \frac{1}{2c}\int_{x_0}^{ct-x}g(\xi)\,d\xi - K/2c.
\]

в помощью семейства уравнений

\[
u(x, t) = \frac{1}{2} [f(x+ct) - f(ct-x)] + \frac{1}{2c} \int_{ct-x}^{x+ct} g(\xi) d\xi, \quad 0 < x < ct.
,   0 \le \varepsilon \le 1  ,
(18.8) \quad u(x, t) =
,   0 < r < 1  ,   0 \le \theta < 2\pi  ,   u(1, \theta) = g(\theta)  ,   0 \le \theta < 2\pi  ,
\]

\begin{verbatim}
при одних граничных условиях (например,   g($\theta$) = \cos $\theta$  ) имеет два различных вещественных решения, а при других (например,   g($\theta$) = A \cos $\theta$   с A > 20,65) эта задача не имеет ни одного вещественного решения. — Прим. ред.
\end{verbatim}

4) Краевая задача

\begin{verbatim}
можно непрерывно модифицировать в нелинейное уравнение   $\Delta$ u = u^2 = 0   1).
\end{verbatim}

На рис. 46.1 изображена схема решения нелинейной задачи путем сведения
ее к решению уравнения Лапласа2)

( \(\varepsilon = 0\) ) (0 \textless{} \(\varepsilon\) \textless{} 1) (
\(\varepsilon = 1\) )

Решение Решение Решение

РИС. 46.1. Схема метода возмущении.

Чтобы найти решение возмущенного уравнения Лапласа

мы добавим возмущение к решению \(u_{\mathrm{ }e}\) невозмущенного
уравнения Лапласа. Представляется естественным, что решение возмущенного
уравнения

\[
u(0, t) = f(t),
\]

следует искать в виде

\[
u_{x}(0, t) = 0.
\]

Заметим, что степенной ряд по \(\epsilon\) (46.1) переходит в решение
уравнения Лапласа \(u_{\bullet}\) при \(\epsilon=0\) и в решение
нелинейного уравнения

\[
x + ct = \text{const},
\]

при \(\varepsilon = 1\) . Другими словами, разложение (46.1) действует
как ракета, которая переносит нас от уравнения Лапласа к нелинейному
уравнению 3). Задача состоит в том, чтобы найти функ-

\begin{enumerate}
\def\labelenumi{\arabic{enumi})}
\tightlist
\item
  В целом это не так, хотя бы потому, что задача
\end{enumerate}

\[
0 < r < 1,     0 \le \theta < 2\pi,     u(1, 0) = 0,     0 \le \theta < 2\pi,
\]

нмеет два решения (одно --- нулсьое, другое --- нетривиальное), в то
время как задача

\begin{verbatim}
\begin{array}{ll} (\text{УЧП}) & u_{tt} = u_{xx}, & 0 < x < \infty, & 0 < t < \infty, \\ (\Gamma \text{У}) & u(0, t) = 0, \\ (\text{HY}) & \begin{cases} u(x, 0) = xe^{-x^2}, \\ u_t(x, 0) = 0. \end{cases} & 0 < x < \infty, \end{cases}
\end{verbatim}

\[
u(1, \theta) = 0, \quad 0 \le \theta < 2\pi,
u(x, 0) = xe^{-x^2}, -\infty < x < 0,
\]

\begin{verbatim}
\begin{array}{ccccccccccccccccccccccccccccccccccc
\end{verbatim}

\[
\begin{array}{lll} (\mathrm{Y} \mathrm{H} \Pi) & u_{tt} = c^2 u_{xx}, & 0 < x < \infty, & 0 < t < \infty, \\ (\mathrm{\Gamma} \mathrm{Y}) & u_x (0, \, t) = 0, & 0 < t < \infty, \\ (\mathrm{H} \mathrm{Y}) & \begin{cases} u \, (x, \, 0) = f \, (x), \\ u_t \, (x, \, 0) = 0. \end{cases} & 0 \leqslant x < \infty, \\ \end{array}
\begin{array}{ccc} (\text{УЧ}\Pi) & \Delta u = 0, & 0 < r < 1, \\ (\text{ГУ}) & u(1, \theta) = \cos \theta, & 0 \leqslant \theta \leqslant 2\pi, \end{array}
u(r, \theta) = u_0(r, \theta) + \varepsilon u_1(r, \theta) + \varepsilon^2 u_2(r, \theta) + \dots
u_{tt} = c^2 u_{xx}, \quad 0 < x < L, \quad 0 < t < \infty,
\Delta u + \varepsilon u^2 = 0
u(0, t) = g_1(t),
\Delta (u_0 + \varepsilon u_1 + \varepsilon^2 u_2 + \ldots) + \varepsilon (u_0 + \varepsilon u_1 + \varepsilon^2 u_2 + \ldots)^2 = 0,
u_x(0, t) = g_1(t),
\]

Виконуючи прості алгебраїчні перетворення та прирівнюючи коефіцієнти за
однаковими ступенями \(\varepsilon\) , ми отримуємо послідовність задач
\(P_0,\ P_1,\ P_2,\ \dots\) розв'язання, \(u_0,\ u_1,\ \dots\) невідомих
функцій (усі ці задачі лінійні).

Ця серія, очевидно, сходиться на \textbar{} є \textbar{} \textless{} 1
443.--- Ред. Ед

.

гетерогенні):

\[
u_x(0, t) - \gamma_1 u(0, t) = g_1(t),
\begin{cases} \Delta u_{0} = 0; & 0 < r < 1, \\ u_{0}(1, \theta) = \cos \theta, \end{cases}
u(0, t) = g_1(t),
P_{1}
u_x(0, t) = g_1(t),
P_{2}
u_x(0, t) - \gamma_1 u(0, t) = g_1(t),
\]

Почнемо з визначення функції \(u_1\left(r,\;\theta\right)\) з завдання
\(P_1\) (ця функція вже відома \(u_n\) ):

\[
\begin{cases} \Delta u_{1} = -r^{2} \cos^{2} \theta = -\frac{r^{2}}{2} \left[1 + \cos(2\theta)\right] = -\frac{r^{2}}{2} - \frac{r^{2}}{2} \cos(2\theta), \\ u_{1}(1, \theta) = 0. \end{cases}
u_h(r, \theta) = \sum_{n=0}^{\infty} r^n [a_n \cos(n\theta) + b_n \sin(n\theta)].
\Delta u = -\frac{r^2}{2} - \frac{r^3}{2} \cos{(2\theta)}
\]

в которых концевые точки совершают заданные движения (рис. 19.1).

\pandocbounded{\includegraphics[keepaspectratio]{images/_page_141_Figure_14.jpeg}}

РИС. 19.1. Концы колеблющенся струны совершают заданные движения.

Характерная задача этого типа---крутильные колебания стержня, левый
конец которого закреплен, а правый поворачивается на некоторый угол
(рис. 19.2)

\pandocbounded{\includegraphics[keepaspectratio]{images/_page_142_Figure_2.jpeg}}

РИС. 19,2. Крутильные колебания стержия.

Одна из важных задач математической теории управления состоит в
определении такой функции \(g_2(t)\) , чтобы за минимальное время
погасить колебан я в струне.

\section{2. Заданные силы в граничных
точках}\label{ux437ux430ux434ux430ux43dux43dux44bux435-ux441ux438ux43bux44b-ux432-ux433ux440ux430ux43dux438ux447ux43dux44bux445-ux442ux43eux447ux43aux430ux445}

Вертикальные силы, действующие на левый и правый конец струны,
определяются выражениями \(Tu_{\mathbf{x}}(0,\,t)\) и
\(Tu_{\mathbf{x}}(L,\,t)\) . Если Р

\pandocbounded{\includegraphics[keepaspectratio]{images/_page_142_Figure_7.jpeg}}

ИС. 19.3. Струна со свободными концами.

конец перемещается вдоль вертикальных направляющих без трения (см. рис.
19.3), то граничные условия принимают вид

\[
u_p(r, \theta) = Ar^4 + Br^4 \cos(2\theta).
\]

Граничные условия этого типа возникают в следующих двух примерах:

\begin{itemize}
\tightlist
\item
  а) Продольные колебания пружины со свободным концом. Рассмотрите
  колебания пружины, изображенной на рис. 19.4.
\item
  б) Колебания пружины под действием силы, приложенной к одному концу.
\end{itemize}

Если сила величины \(v(t)\) дин приложена к концу пружины x=1
(направленная вниз сила счигается положительной), то граничное условие
записывается в виде

\[
u_p(r, \theta) = -\frac{r^4}{32} - \frac{r^4}{24} \cos(2\theta).
\]

В случае граничных условий 2-го рода концы струны (или пружины) не
обязаны находиться в заданном положении. Под действием сил концевые
точки участвуют в движении. Подобные задачи возникают в физике, напри- Р

\pandocbounded{\includegraphics[keepaspectratio]{images/_page_143_Picture_5.jpeg}}

ИС, 19 4. Пружина со свободным концом.

мер при изучении колебаний электронов в электрическом поле.

\section{3. Упругое закрепление
концов.}\label{ux443ux43fux440ux443ux433ux43eux435-ux437ux430ux43aux440ux435ux43fux43bux435ux43dux438ux435-ux43aux43eux43dux446ux43eux432.}

Рассмотрим, наконец, скрипичную струну, концы которой упруго закреплены.
Схема упругого закрепления изображена на Р

\pandocbounded{\includegraphics[keepaspectratio]{images/_page_143_Picture_10.jpeg}}

ИС, 19.5. Схема упругого закрепления: h---коэффициент упругости пружниы
(растянутая пружниа).

рис. 19.5. Пружинки, закрепленные к концам струны, создают вертикальные
силы, пропорциональные смещениям концов:

Смещение левого конца = \(u(0, t)\), Смещение правого конца =
\(u(L, t)\).

Усилия, которые создаются пружинками на концах струны:

Вертикальное усилие на левом конце = \(Tu_x(0, t)\) , Вертикальное
усилие на правом конце = \(-Tu_x(L, t)\) , (T- натяжение струны)

с точностью до множителя h(h-коэффициент упругости пружинки) равны
перемещениям. Отсюда получаем искомые гранич-

ные условия

(19.2)

\[
\sum_{n=0}^{\infty} [a_n \cos(n\theta) + b_n \sin(n\theta)] - \frac{1}{32} - \frac{1}{24} \cos(2\theta) = 0.
u_1(r, \theta) = \frac{1}{32} + \frac{1}{24}r^2\cos(2\theta) - \frac{r^4}{24}\cos(2\theta) - \frac{r^4}{32} = \frac{(r^4 - 1)}{32} - \frac{(r^4 - r^2)}{24}\cos(2\theta).
\]

Отметим, что если величина \(u(0, t)\) положительна, то величина
\(u_x\)(0, t) также положительна, если \(u(L, t)\) положительна, то
\(u_x\)(L, t)

\pandocbounded{\includegraphics[keepaspectratio]{images/_page_144_Picture_4.jpeg}}

РИС. 19.6. Возникновение неоднородных граничных условий упругого
закрепления: h---коэффициент упругости пружины.

отрицательна. Однородные граничные условия (19.2) можно записать в виде

(19.3)

\[
u = u_0 + u_1 = r \cos \theta - \frac{(r^4 - 1)}{32} - \frac{(r^4 - r^2)}{24} \cos (2\theta).
P_{2}
\]

Если обе прикрепленные пружинки совершают заданные движения
\(\theta_1(t)\) и \(\theta_2(t)\) (см. рис. 19.6), то приходим к
неоднородным граничным условиям

\[
\Delta u = 0, \quad 0 < r < 1 + \frac{1}{4} \sin \theta,
\]

Цю задачу можна розглядати як задачу визначення форми мильної плівки,
натягнутої на деформоване коло, висота якого вище цієї межею дорівнює
\(\cos\theta\) (див. рис. 46.2).

\pandocbounded{\includegraphics[keepaspectratio]{images/_page_349_Figure_5.jpeg}}

Рис. 46.2. Рівняння Лапласа всередині деформованого кола.

Оскільки загальна ідея теорії збурень полягає в тому, щоб звести складну
задачу до простої, припустимо, що наша межа умовна від деформованого
кола

досягається в кінцевій точці сімейства граничних умов виду (див. рисунок
46 3)

Ця ідея приводить нас до серії про Тейлора. Використання серії Тейлора

\[
Tu_x(0, t) = -\beta u_t(0, t).
\]

Розкладемо складну граничну умову на ряд відповідно до простої межі Р и
с

\pandocbounded{\includegraphics[keepaspectratio]{images/_page_350_Figure_2.jpeg}}

. 46 3. Схематичне представлення методу збурення.

умови, тобто

\[
Tu_{x}(0, t) = \varphi[u(0, t)],
\]

Підставляючи цю декомпозицію в початкову задачу, отримуємо нову
еквівалентну задачу

(46.7)

\[
Tu_{x}(0, t) = -hu^{3}(0, t)
u(1, \theta) + u_r(1, \theta) (\varepsilon \sin \theta) + u_{rr}(1, \theta) \frac{(\varepsilon \sin \theta)^2}{2!} + \dots = \cos \theta.
mu_{tt}(L, t) = -ku_x(L, t) + mg.
(46.8) u = u_0 + \varepsilon u_1 + \varepsilon^2 u_2 + \dots.
P_{0}
P_{1}
\]

Отже, розв'язуючи кожну з задач Дірікла всередині кола для функцій
\(u_0\) , \(u_1\) , \(u_2\) , отримуємо розв'язок

(46.9)

\begin{verbatim}
\right. & 0 < t < \infty, \\ (\mathrm{H} \, \mathrm{Y}) & \left\{
\end{verbatim}

Задачі Діріхле в деформованій області (зверніть увагу, що тут

\begin{verbatim}
Ми вже взяли  \epsilon = 1/4  :)
\end{verbatim}

Як вправу, читача просять знайти пертурбацію першого порядку \(u_1\) і
перевірити, наскільки добре є приблизне розв'язання

задовольняє рівняння (46 b)

\section{ЗАУВАЖЕННЯ}\label{ux437ux430ux443ux432ux430ux436ux435ux43dux43dux44f-7}

ŗ

Методи теорії збурень можна використовувати не лише для розв'язання
нелінійних рівнянь із частковими похідними. Наприклад, можливо
розв'язати задачу Коші.

(46 10)

якщо розглянути рівняння з параметром

(46.11)

в различные моменты времени.

\begin{enumerate}
\def\labelenumi{\arabic{enumi}.}
\setcounter{enumi}{1}
\tightlist
\item
  Изобразить в общих чертах решения задачи
\end{enumerate}

і побачити його розв'язок у вигляді

\[
\begin{array}{ll} (\text{YHI}) & u_{tt} = u_{xx}, \quad 0 < x < 1, \quad 0 < t < \infty, \\ (\text{ГУ}) & \begin{cases} u(0, t) = 0, \\ u_x(1, t) = 0, \\ 0 < t < \infty, \end{cases} \\ (\text{HY}) & \begin{cases} u(x, 0) = \sin(\pi x/2), \\ u_t(x, 0) = 0, \end{cases} \quad 0 \le x \le 1, \end{array}
\]

Підставляючи цю декомпозицію на (46.11), отримуємо таку послідовність
завдань

\begin{verbatim}
\begin{cases} \frac{\partial u_{0}}{\partial t} = \frac{\partial^{2} u_{1}}{\partial x^{2}}, \\ u_{0}(x, 0) = $\varphi$(x), \\ P_{1} \begin{cases} \frac{\partial u_{1}}{\partial t} - \frac{\partial^{2} u_{1}}{\partial x^{2}} = x \frac{\partial^{2} u_{0}}{\partial x^{2}}, \\ u_{1}(x, 0) = 0 \\ \vdots & \vdots \end{cases}
\end{verbatim}

\[
u_{x}(0, t) = \frac{h}{T} [u(0, t) - \theta_{1}(t)]
\begin{array}{ll} \Delta u + u^2 = 0, & 0 \leqslant r < 1, \quad 0 \leqslant \theta < 2\pi, \\ u(1, \theta) = \cos \theta, & 0 \leqslant \theta < 2\pi, \end{array}
u_{tt} = u_{xx}, \quad 0 < x < 1, \quad 0 < t < \infty,
u(r, \theta) = u_0(r, \theta) + \frac{1}{4}u_1(r, \theta)
\Delta u = 0,
q_{xx} + \varphi_{uu} = 0
q_{uu} + q_{vv} = 0
\]

(HУ)

\[
w = f(z)
u + iv = (x + iy)^2 = x^2 - y^2 + 2ixy
\]

\begin{verbatim}
\left\{
\end{verbatim}

Лекция 20

\section{КОЛЕБАНИЯ ОГРАНИЧЕННОЙ СТРУНЫ (СТОЯЧИЕ
ВОЛНЫ)}\label{ux43aux43eux43bux435ux431ux430ux43dux438ux44f-ux43eux433ux440ux430ux43dux438ux447ux435ux43dux43dux43eux439-ux441ux442ux440ux443ux43dux44b-ux441ux442ux43eux44fux447ux438ux435-ux432ux43eux43bux43dux44b}

ЦЕЛЬ ЛЕКЦИИ: Показать, что решение смешанной задачи для поперечных
колебаний струны

\[
\Rightarrow.
\]

\begin{verbatim}
\begin{array}{ll} (\text{УЧ}\Pi) & u_{tt} = $\alpha$^2 u_{xx}, & 0 < x < L, & 0 < t < \infty, \\ (\text{ГУ}) & \begin{cases} u \ (0, \ t) = 0, & 0 < t < \infty, \\ u \ (L, \ t) = 0, & 0 < t < \infty, \end{cases} \\ (\text{HY}) & \begin{cases} u \ (x, \ 0) = f \ (x), & 0 \leqslant x \leqslant L, \\ u_1 \ (x, \ 0) = g \ (x), & 0 \leqslant x \leqslant L, \end{cases}
\end{verbatim}

\[
\begin{aligned} (Y & \Pi) & \varphi_{xx} + \varphi_{yy} = 0, & -\infty < x < \infty, & 0 < y < \infty, \\ (47.1) & \varphi(x, 0) = \begin{cases} 0, & |x| > 1, \\ 1, & |x| \le 1. \end{cases} \end{aligned}
u(x, t) = \sum_{n=1}^{\infty} X_n(x) T_n(t),
w = \ln\left\{\frac{z-1}{z+1}\right\}
\begin{array}{ll} (\mathrm{Y}\mathrm{H}\Pi) & u_{tt} = \alpha^2 u_{xx}, & 0 < x < L, & 0 < t < \infty, \\ .(\mathrm{\Gamma}\mathrm{Y}) & \begin{cases} u \ (0, \ t) = 0, & 0 < t < \infty \\ u \ (L, \ t) = 0, & 0 < t < \infty \end{cases} \\ (\mathrm{H}\mathrm{Y}) & \begin{cases} u \ (x, \ 0) = f \ (x), & 0 \leqslant x \leqslant L, \\ u_t \ (x, \ 0) = g \ (x). & 0 \leqslant x \leqslant L, \end{cases} \end{array}
\begin{cases} (\text{УЧП}) & \varphi_{uu} + \varphi_{vv} = 0, \quad -\infty < u < \infty, \quad 0 < v < \pi, \\ (\text{47.2}) & \begin{cases} \varphi(u, 0) = 0, \\ \varphi(u, \pi) = 1. \end{cases} \end{cases}
w = u + iv = \ln\left[\frac{z-1}{z+1}\right] =
u(x, t) = \sum_{n=1}^{\infty} c_n X_n(x) T_n(t),
\]

і, відповідно,

\[
u_{tt} = \alpha^2 u_{xx}, \quad 0 < x < L, \quad 0 < t < \infty,
= \arg\left[\frac{x^2+y^2-1+i2y}{(x+1)^2+y^2}\right] =
(\Gamma Y)
\]

(згідно з загальним правилом \(\arg{(x+iy)}=\arg{(y^{\dagger}x)}\) ).
Отже, розв'язок задачі (47.1) має вигляд

Щоб візуально зрозуміти, як працює ця функція, ми радимо читачеві
побудувати її за кількома лініями y=c.

У другому прикладі ми перетворюємо область між двома неконцентричними
колами у кільце.

Задача Діріхле в області між двома некопоцентричними колами

Припустимо, ви хочете знайти потенціал між двома колами

\[
(x-1)^{2} + y^{2} = 9,
\quad 0 < t < \infty,
\begin{aligned} & (\text{УЧП}) & \phi_{xx} + \phi_{yy} = 0, & \text{внутри } D, \\ & (\text{ГУ}) & \begin{cases} \varphi(x, y) = 1 & \text{на } x^2 + y^2 = 1, \\ \varphi(x, y) = 2 & \text{на } (x - 1)^2 + y^2 = 9. \end{cases} \end{aligned}
\]

(HY)

\[
w = 2t \left[ \frac{z-s}{z-t} \right],
u = \gamma [(x-s)(x-t)-\gamma^2],
\]

де

\[
\quad 0 \le x \le L,
\]

Під впливом цього конформного відображення область між неконцентричними
колами переходить у кільце, межі якого визначаються рівняннями
\(u^2 + v^2 = 1\) і \(u^2 + v^2 = 6,86\) (читач може перевірити це сам)

Якщо застосувати це конформне відображення до задачі (47 3), отримаємо
задачу Діріхле всередині кільця

(47.6)

\[
u(x, t) = X(x) T(t).
\]

Розв'язати цю задачу досить просто, оскільки розв'язок має бути
радіально симетричним, а радіально симетричний розв'язок рівняння
Лапласа, як відомо, має вигляд

\[
T'' - \alpha^2 \lambda T = 0,
\]

Підставляючи це рівняння у граничну умову (47,6), отримуємо

\[
X'' - \lambda X = 0,
\]

Якщо повернутися до початкових координат x і y відповідно до
відображення (47.5), отримаємо розв'язок задачі (47.3)

\[
u(x, t) = [C \sin(\beta x) + D \cos(\beta x)][A \sin(\alpha \beta t) + B \cos(\alpha \beta t)]
\]

де w і v визначаються за рівняннями (47 5).

Бачиш. Див. також Коппенфельс В., Штальман Ф. Практика конформних
відображень. --- Москва, Іллінойс, 1963. --- Примітка редактора.

\section{ПРИМІТКИ}\label{ux43fux440ux438ux43cux456ux442ux43aux438-3}

\begin{enumerate}
\def\labelenumi{\arabic{enumi}.}
\item
  Для побудови конформного відображення області з достатньо загальним
  виглядом у площині z на верхній площині w використовується
  відображення Крістоффеля-Шварца (коли область відображена у верхню
  півплощину, тоді цю напівплощину можна відобразити в іншу область).
  Інтеграл Крістоффеля--Шварца можна знайти у книзі {[}28{]} з
  рекомендованої літератури.
\item
  Метод конформних відображень має обмежені можливості, оскільки op
  застосовний лише до двовимірного рівняння Лапласа (і, з незначними
  змінами, до рівняння
\end{enumerate}

Пуассон).

\section{ЗАВДАННЯ}\label{ux437ux430ux432ux434ux430ux43dux43dux44f-12}

\begin{enumerate}
\def\labelenumi{\arabic{enumi}.}
\tightlist
\item
  Визначити область відповідності відображення
\end{enumerate}

\[
\beta_n = \frac{n\pi}{L}, \quad n = 0, 1, 2, \ldots
\]

Переконайтеся, що це відображення переносить верхню половину площини
змінної z у зону --- \(\infty < u < \infty\) , \(0 < v < \pi\) ,
w-площину.

\begin{enumerate}
\def\labelenumi{\arabic{enumi}.}
\setcounter{enumi}{1}
\tightlist
\item
  Яке зображення першого квадранта при відображенні \(w=z^{8}\) ? Щоб
  розв'язати задачу, корисно перейти до експоненціальної форми
  комплексного числа
\end{enumerate}

\[
u_n(x, t) = X_n(x) T_n(t) = \sin(n\pi x/L) [a_n \sin(n\pi at/L) + b_n \cos(n\pi at/L)],
\]

.

\begin{enumerate}
\def\labelenumi{\arabic{enumi}.}
\setcounter{enumi}{2}
\tightlist
\item
  За допомогою конформного відображення \(w=z^2\) розв'язати задачу
  Діріхле у першому квадранті
\end{enumerate}

\[
u_n(x, t) = R_n \sin(n\pi x/L) \cos(n\pi\alpha (t - \delta_n)/L],
\right.
u(x, t) = \sum_{n=1}^{\infty} \sin(n\pi x/L) \left[ a_n \sin(n\pi \alpha t/L) + b_n \cos(n\pi \alpha t/L) \right]
H(x-a) =
u(x, \theta) = f(x),
\]

Функція Хевісайду,

\[
\sum_{n=1}^{\infty} b_n \sin(n\pi x/L) = f(x),
\begin{cases} 0, & x < a, \\ 1, & x \geqslant a, \end{cases}
\sum_{n=1}^{\infty} a_n (n\pi \alpha/L) \sin(n\pi x/L) = g(x).
f(x) = \mathcal{F}^{-1}[F] = \frac{1}{\sqrt{2\pi}} \int_{-\infty}^{+\infty} F(\omega) e^{-i\omega x} d\omega \qquad F(\omega) = \mathcal{F}[f] = \frac{1}{\sqrt{2\pi}} \int_{-\infty}^{+\infty} f(x) e^{-i\omega x} dx
\int_{0}^{L} \sin(m\pi x/L) \sin(n\pi x/L) dx =
\begin{cases} 1, |x| < a \\ 0, |x| > a \end{cases}
\begin{cases} 1, |x| < 1 \\ 0, |x| > 1 \end{cases}
\delta(x-a)
\]

\(f(x)\) * \(g(x)\)

\[
a_{n} = \frac{2}{n\pi\alpha} \int_{0}^{L} g(x) \sin(n\pi x/L) dx, \\ b_{n} = \frac{2}{L} \int_{0}^{L} f(x) \sin(n\pi x/L) dx.
(1+x^2)^{-1}
u(x, t) = \sum_{n=1}^{\infty} \sin(n\pi x/L) \left[ a_n \sin(n\pi \alpha t/L) + b_n \cos(n\pi \alpha t/L) \right],
xe^{-a+x^4}, a > 0
u(x, t) = \sum_{n=1}^{\infty} b_n \sin(n\pi x/L) \cos(n\pi\alpha t/L)
\]

\(H(x+a)\) - \(H(x-a)\)

\[
u(x, \theta) = f(x)
\frac{a}{x^2+a^2}
f(x) = \sum_{n=1}^{\infty} b_n \sin(n\pi x/L).
\frac{2ax}{(x^2+a^2)^2}
u_n(x, t) = b_n \sin(n\pi x/L) \cos(n\pi\alpha t/L).
\begin{cases} \cos(ax), |x| > \pi/2a \\ 0, |x| > \pi/2a \end{cases}
u(x, t) = \sin(\pi x/L)\cos(\pi \alpha t/L) + 0.5\sin(3\pi x/L)\cos(3\pi \alpha t/L) + 0.25\sin(5\pi x/L)\cos(5\pi \alpha t/L).
\begin{cases} 1-|x|, |x| < 1 \\ 0, |x| > 1 \end{cases}
R_n \sin(n\pi x/L) \cos[n\pi\alpha(t-\delta_n)/L],
\cos(ax)
\omega_n = \frac{n\pi\alpha}{L} = \frac{n\pi}{L} \sqrt{\frac{T}{\rho}},
\frac{2}{\pi} \frac{\sin a\omega}{\omega}
u(x, 0) = \sin(\pi x/L) + 0.5 \sin(3\pi x/L),
\frac{2}{\pi} \frac{\sin a\omega}{(\omega^2+a^2)^2}
u(x, 0) = 0,
\frac{\pi}{2} e^{-a+\omega}
u(x, 0) = \sin(3\pi x/L),
\frac{\pi}{2} e^{-a+\omega}
u(x, 0) =
\frac{\pi}{2} e^{-a+\omega}
\begin{cases} 2hx, & 0 \le x \le 0.5, \\ 2h(1-x), & 0.5 \le x \le 1. \end{cases}
\frac{\pi}{2} e^{-a+\omega}
\]

\pandocbounded{\includegraphics[keepaspectratio]{images/_page_153_Picture_8.jpeg}}

Какие движения будет совершать струна, если ее отпустить? 6. Решите
задачу о затухающих колебаниях струны

\[
\frac{\pi}{2} e^{-a+\omega}
\frac{2ax}{(x^2+a^2)^2}
\frac{\cos(ax), |x| > \pi/2a}{0, |x| > \pi/2a}
\]

Представляется ли вам разумным полученное решение? Удовлетворяет ли оно
уравнению с частными производными, граничным и начальным условиям?

\begin{enumerate}
\def\labelenumi{\arabic{enumi}.}
\setcounter{enumi}{6}
\tightlist
\item
  Как вы будете решать неоднородное уравнение с заданными начальными и
  граничными условиям?
\end{enumerate}

\[
\frac{1-|x|, |x| < 1}{0, |x| > 1}
\cos(ax)
\cos(ax)
\frac{\pi}{2} [\delta(\omega+a) + \delta(\omega-a)]
u_{tt} = -\alpha^2 u_{xxxx},
f(x) = \int_{0}^{\infty} F(\omega) \sin(\omega x) d\omega
u(0, t) = 0,
0 < x < \infty
u_{xx}(0, t) = 0,
\]

\begin{enumerate}
\def\labelenumi{\arabic{enumi}.}
\tightlist
\item
\end{enumerate}

\[
-\omega^2 F(\omega) + \frac{2}{\pi} \omega f(0)
\frac{1}{\alpha} F\left(\frac{\omega}{\alpha}\right)
\frac{2\omega}{\pi (\alpha^2 + \omega^2)}
x^{-1/2}
u(x, t) = X(x) [A \sin(\omega t) + B \cos(\omega t)].
\]

\begin{enumerate}
\def\labelenumi{\arabic{enumi}.}
\setcounter{enumi}{4}
\tightlist
\item
  \(H(a-x)\)
\end{enumerate}

\[
X^{(1V)} - \omega^2 X = 0,
\]

\begin{enumerate}
\def\labelenumi{\arabic{enumi}.}
\setcounter{enumi}{5}
\tightlist
\item
  \(x^{-1}\)
\end{enumerate}

\[
X(x) = C\cos\sqrt{\omega}x + D\sin\sqrt{\omega}x + E \cosh\sqrt{\omega}x + F \sinh\sqrt{\omega}x.
\]

\begin{enumerate}
\def\labelenumi{\arabic{enumi}.}
\setcounter{enumi}{6}
\item
  \(\frac{x}{x^2+a^2}\)

  \begin{enumerate}
  \def\labelenumii{\arabic{enumii}.}
  \setcounter{enumii}{7}
  \item
    \frac{x}{x^4+4}
  \end{enumerate}
\item
  \(\arctan \frac{a}{x}\)

  \begin{enumerate}
  \def\labelenumii{\arabic{enumii}.}
  \setcounter{enumii}{9}
  \tightlist
  \item
    -x\^{}2/(x)
  \end{enumerate}
\end{enumerate}

Откуда получаем два уравнения

\section{ТАБЛИЦЯ C. Перетворення
Фур'є}\label{ux442ux430ux431ux43bux438ux446ux44f-c.-ux43fux435ux440ux435ux442ux432ux43eux440ux435ux43dux43dux44f-ux444ux443ux440ux454}

\[
0 < x < \infty \qquad \qquad 0 < \omega < \infty
\]

Решая их, находим

j'\,'(x)

,

\[
\sin \sqrt{\omega} = 0 \Rightarrow \omega = (n\pi)^2 ,  n = 1, 2, \dots
\]

Другими словами, \emph{собственные частоты} свободно опирающейся балки
равны

\[
S_n = \frac{2}{\pi} \int_0^{\pi} f(x) \sin(nx) dx.
\]

а фундаментальные решения и, (удовлетворяют уравнению и граничным
условиям (21.2)) записываются в виде

\[
y = \frac{\pi}{b - a} (x - a).
\]

А теперь, поскольку уравнение и граничные условия линейны и однородны,
можно утверждать, что сумма

(21.4)

\[
f(x) = \sum_{n=1}^{\infty} S_n \sin(nx).
\]

также удовлетворяет уравнению и граничным условиям. Единственное, что
нам осталось---это выбрать коэффициенты \(a_n\) и \(b_n\) так, чтобы
удовлетворить начальным условиям. Подставляя (21.4) в начальные условия,
получаем

(21.5)

\[
S_n = \frac{2}{\pi} \int_{0}^{\pi} f(x) \sin(nx) dx
n = 1, 2, ...
\]

Воспользовавшись ортогональностью семейства функции (sin (nnx)) на
отреже {[}0, 1{]}, находим формулы для коэффициентов

(21.6)

f'\,'(x)

\[
\frac{\pi}{2} e^{ax}
\]

Таким образом, решение задачи имеет вид (21.4), а коэффициенты \(a_n\) и
\(b_n\) вычисляются по формулам (21.6)

Чтобы помочь читателю разобраться в этой задаче, приведем простой
пример.

\section{Модель колеблющейся
балки}\label{ux43cux43eux434ux435ux43bux44c-ux43aux43eux43bux435ux431ux43bux44eux449ux435ux439ux441ux44f-ux431ux430ux43bux43aux438}

Рассмотрим свободно опирающуюся балку, изображенную на рис. 21.2. Пусть
начальные условия таковы:

\[
C_n = \frac{2}{\pi} \int_0^{\pi} f(x) \cos(nx) dx.
\]

\(u_1\)(x, 0) = 0 Мы могли бы найти решение, подставив функции \(f(x)\)
и \(g(x)\) в формулы (21.6), однако легко видеть, что

\[
y = \frac{\pi}{b-a}(x-a).
\]

для всех n=1,,2,,\ldots,
\(b_1=1,\,b_2=0,\,b_3=0,5 b_n=0,\,n=4,\,5,\,\ldots\)

Следовательно, решение имеет вид

\[
f(x) = \frac{C_0}{2} + \sum_{n=1}^{\infty} C_n \cos(nx).
\]

.

Интересно сравнить это решение с решением для струны при тех же
начальных условиях. Если заглянуть в лекцию 20, там можно найти решение
соответствующей задачи для струны:

\[
0 \le x \le \pi \qquad n = 0, 1, 2, \dots
\]

.

\pandocbounded{\includegraphics[keepaspectratio]{images/_page_158_Figure_1.jpeg}}

РИС 21.2. Простые колебания свободно опирающейся балки.

Другими словами, балка колеблется на более высоких частотах, чем струна.
Отметим, что в том и другом случаях высшие частоты кратны основным
частотам колебаний.

\section{ЗАМЕЧАНИЯ}\label{ux437ux430ux43cux435ux447ux430ux43dux438ux44f-18}

\begin{itemize}
\item
  \begin{enumerate}
  \def\labelenumi{\arabic{enumi}.}
  \tightlist
  \item
    Вообще говоря, балку можно закрепить тремя способами:
  \end{enumerate}
\item
  а) оставить свободной (незакрепленной),
\item
  б) свободно опереть.
\item
  в) жестко закрепить.
\end{itemize}

Схематически эти способы закрепления изображены на рис. 21.3.

\begin{enumerate}
\def\labelenumi{\arabic{enumi}.}
\setcounter{enumi}{1}
\tightlist
\item
  Еще одна важная задача --- колебания балки с заделанным концом (см.
  рис. 213). Благодаря нестандартным граничным условиям
\end{enumerate}

f'\,'(x)

\(u_x\)(0, t) = 0, \(u_{xx}\)(1, t) = 0, \(u_{xx}\)(1, t) = 0, решение
уже не будет выражаться в виде ряда из произведений синусов и косинусов.
Оно будет иметь более сложную структуру

\begin{enumerate}
\def\labelenumi{\arabic{enumi}.}
\setcounter{enumi}{1}
\tightlist
\item
  \(\frac{a_{0}}{2} + \sum_{n=1}^{\infty} a_{n} \cos(nx)\) \(a_{n}\) 3.
  \(f(n-x)\) \((-1)^{n} \frac{2}{\pi} C_{n}\) 4. 1
\end{enumerate}

где собственные функции \(X_n\)(x) являются линейными комбинациями
синусов, косинусов, гиперболических синусов и гиперболических косинусов.
Решение этой задачи можно найта в руководствах по теории упругости.

\pandocbounded{\includegraphics[keepaspectratio]{images/_page_159_Figure_2.jpeg}}

РИС. 21.3. Типичные задачи для балки: a---защемленная балка; b---балка,
жестко заделанная с обоих концов; a---левый конец жестко заделан, правый
свободно опирается.

\subsubsection{ЗАДАЧИ}\label{ux437ux430ux434ux430ux447ux438-22}

\begin{enumerate}
\def\labelenumi{\arabic{enumi}.}
\item
  Решите задачу для балки с закрепленным концом.
\item
  \(\cos(mx), m=1, 2, \dots\)
\item
  \textbackslash begin\{cases\}
\end{enumerate}

ТАБЛИЦЯ F. Перетворення Лапласа

{\def\LTcaptype{none} % do not increment counter
\begin{longtable}[]{@{}l@{}}
\toprule\noalign{}
F(s) = \mathscr{L}{[}f(t){]} \\
\midrule\noalign{}
\endhead
\bottomrule\noalign{}
\endlastfoot
1/с, с \textgreater{} 0 \\
\frac{1}{s-a} , s\textgreater a \\
\frac{a}{s^2 + a^2}, s \textgreater{} 0 \\
\frac{s}{s^2 + a^2}, s \textgreater{} 0 \\
\frac{a}{s^2 - a^2}, s \textgreater{} a \\
\frac{s}{s^2-a^2}, s\textgreater{} a \\
\frac{b}{(s-a)^2+b^2}, s\textgreater a \\
\frac{s-a}{(s-a)^2+b^2}, s\textgreater a \\
\frac{n!}{8^{n+1}} \textless\textgreater0 \\
\frac{n!}{(s-a)^{n+1}}, s \textgreater{} a \\
\frac{e^{-as}}{s} , s\textgreater0 \\
e\^{}\{-as\}F(s) . \\
F(s-a) \\
F(s) G(s) \\
\end{longtable}
}

\_ \textbar{}

{\def\LTcaptype{none} % do not increment counter
\begin{longtable}[]{@{}
  >{\raggedright\arraybackslash}p{(\linewidth - 2\tabcolsep) * \real{0.5902}}
  >{\raggedright\arraybackslash}p{(\linewidth - 2\tabcolsep) * \real{0.4098}}@{}}
\toprule\noalign{}
\begin{minipage}[b]{\linewidth}\raggedright
15. f\^{}\{(n)\}(t) , (n-похідна)
\end{minipage} & \begin{minipage}[b]{\linewidth}\raggedright
s\^{}\{n\}F(s) - s\^{}\{n-1\}f(0) - \dots - f\^{}\{(n-1)\}(0)
\end{minipage} \\
\midrule\noalign{}
\endhead
\bottomrule\noalign{}
\endlastfoot
16. ж. (at) & \frac{1}{a}F\left(\frac{s}{a}\right), a\textgreater0 \\
17. \int\_\{0\}\^{}\{t\} f(\tau) d\tau & \frac{1}{s} F(s) \\
18. ERF (T/2A) & \frac{1}{s} e\textsuperscript{\{a}2s\^{}2\} erfe
(AS) \\
19. Ерік (A/2 V T ) & \frac{1}{s}e\^{}\{-a\} \\
20. J\_\{\(\theta\)\}(at) & (s\textsuperscript{2+a}2)\^{}\{1/2\} \\
21. \(\delta\)(t-a) & e\^{}\{-sa\} \\
22. \frac{1}{\sqrt{$\pi$ t}} \exp\left(-\frac{a^2}{4t}\right) &
\frac{e^{-a/s}}{\sqrt{s}}, a \ge 0 \\
23. \frac{1}{\sqrt{nt}} - ae\textsuperscript{\{a}2t\}
\operatorname{eric} (a \sqrt{t}) & \frac{1}{\sqrt{s}+a} \\
\end{longtable}
}

\section{Додаток 2}\label{ux434ux43eux434ux430ux442ux43eux43a-2}

\section{ПРЕДСТАВЛЕННЯ ЛАПЛАСІАНА В РІЗНИХ СИСТЕМАХ
КООРДИНАТ}\label{ux43fux440ux435ux434ux441ux442ux430ux432ux43bux435ux43dux43dux44f-ux43bux430ux43fux43bux430ux441ux456ux430ux43dux430-ux432-ux440ux456ux437ux43dux438ux445-ux441ux438ux441ux442ux435ux43cux430ux445-ux43aux43eux43eux440ux434ux438ux43dux430ux442}

\[
\Delta u = u_{xx} + u_{yy} \Delta u = u_{rr} + \frac{1}{r}u_r + \frac{1}{r^2}u_{\theta\theta} \Delta u = u_{xx} + u_{yy} + u_{zz} \Delta u = u_{rr} + \frac{1}{r}u_r + \frac{1}{r^2}u_{\theta\theta} + u_{zz} \Delta u = u_{rr} + \frac{2}{r}u_r + \frac{1}{r^2}u_{\theta\theta} + \frac{\cot\theta}{r^2} \Delta u = u_{rr} + \frac{2}{r}u_r + \frac{1}{r^2}u_{\theta\theta} + \frac{\cot\theta}{r^2}u_0 + \frac{1}{r^2\sin\theta}u_{\phi\phi}
\]

У тривимірній сферичній системі координат

\pandocbounded{\includegraphics[keepaspectratio]{images/_page_370_Figure_1.jpeg}}

\pandocbounded{\includegraphics[keepaspectratio]{images/_page_370_Picture_2.jpeg}}

\pandocbounded{\includegraphics[keepaspectratio]{images/_page_370_Figure_3.jpeg}}

\section{Декартова система
координат}\label{ux434ux435ux43aux430ux440ux442ux43eux432ux430-ux441ux438ux441ux442ux435ux43cux430-ux43aux43eux43eux440ux434ux438ux43dux430ux442}

\[
-\infty < x < +\infty -\infty < y < +\infty -\infty < z < +\infty
\]

\section{Циліндрична система
координат}\label{ux446ux438ux43bux456ux43dux434ux440ux438ux447ux43dux430-ux441ux438ux441ux442ux435ux43cux430-ux43aux43eux43eux440ux434ux438ux43dux430ux442}

\[
\begin{array}{lll}
 -\infty & < x < +\infty & r \geqslant 0, \ 0 \leqslant \theta \leqslant 2\pi, \\
 -\infty & < y < +\infty & -\alpha \leqslant z \leqslant \alpha, \\
 -\infty & < z < +\infty & x = r \cos \theta, \\
 y = r \sin \theta, \\
 z = z.
 \end{array}
\]

\subsection{Сферична система
координат}\label{ux441ux444ux435ux440ux438ux447ux43dux430-ux441ux438ux441ux442ux435ux43cux430-ux43aux43eux43eux440ux434ux438ux43dux430ux442}

\[
r \geqslant 0, 0 \leqslant \theta \leqslant 2\pi , 0 \leqslant \varphi \leqslant \pi , x = r \sin \varphi \cos \theta , y = r \sin \varphi \sin \theta , z = r \cos \varphi .
\]

\section{Еліптичні типові
рівняння}\label{ux435ux43bux456ux43fux442ux438ux447ux43dux456-ux442ux438ux43fux43eux432ux456-ux440ux456ux432ux43dux44fux43dux43dux44f}

\[
\Delta u = 0 \Delta u + \lambda^2 u = 0 \Delta u = k \Delta u + k(E - V)u = 0
\]

Рівняння Лапласа Рівняння Гельмгольца Рівняння Пуассона-Шредінгера

\section{Рівняння гіперболічного
типу}\label{ux440ux456ux432ux43dux44fux43dux43dux44f-ux433ux456ux43fux435ux440ux431ux43eux43bux456ux447ux43dux43eux433ux43e-ux442ux438ux43fux443}

\[
\begin{array}{l} u_{tt} = c^2 u_{xx} \\ u_{tt} = c^2 u_{xx} - h u_t \\ u_{tt} = c^2 u_{xx} - h u_t - k u \\ u_{tt} = c^2 u_{xx} + f(x, t) \end{array}
u_{tt} = c^2 \Delta u u_{tt} = c^2 \Delta u - hu_t
\]

Рівняння коливання струни Рівняння коливання з тертям

Телеграфне рівняння

Рівняння коливань під впливом зовнішньої сили

Багатовимірне хвильове рівняння з тертям

\section{Параболічні типові
рівняння}\label{ux43fux430ux440ux430ux431ux43eux43bux456ux447ux43dux456-ux442ux438ux43fux43eux432ux456-ux440ux456ux432ux43dux44fux43dux43dux44f}

\(u_t = \alpha^2 u_{xx} u_t = \alpha^2 u_{xx} - hu_x u_t = \alpha^2 u_{xx} - ku u_t = \a\)lpha\^{}2
\(u_{xx}\) + \(f(x, t)\) Уравнение одномерной диффузии Уравнение
конвективной диффузии

Уравнение теплопроводности с поглощением Уравнение теплопроводности с
источником

\section{Джерела ()}\label{ux434ux436ux435ux440ux435ux43bux430}

\section{Головне}\label{ux433ux43eux43bux43eux432ux43dux435}

\begin{itemize}
\tightlist
\item
  Бідадзе А. В. Рівняння математичної фізики.--- Москва: Наука, 1976.
\item
  Володимиров В. С. Рівняння математичної фізики.--- Москва: Наука,
  1971. Кошляков Н. С., Глінер Е. Б., Смірнов М. М. Рівняння в зокрема
  математичній фізиці.--- Москва: Вища школа, 1970.
\item
  \begin{enumerate}
  \def\labelenumi{\arabic{enumi}.}
  \setcounter{enumi}{3}
  \tightlist
  \item
    Самарський, А. А. Вступ до чисельних методів. Москва, Наукове
    видавництво, 1982.
  \end{enumerate}
\item
  \begin{enumerate}
  \def\labelenumi{\arabic{enumi}.}
  \setcounter{enumi}{4}
  \tightlist
  \item
    Смирнов В. І. Курс, виший математіки, т. 2.--- М.: Наука, 1967.
  \end{enumerate}
\item
  \begin{enumerate}
  \def\labelenumi{\arabic{enumi}.}
  \setcounter{enumi}{5}
  \tightlist
  \item
    Тихонов А. Н., Самарський А. А. Рівняння математичної фізики.
    Москва, Наукове видання, 1972.
  \end{enumerate}
\end{itemize}

\section{Додатково}\label{ux434ux43eux434ux430ux442ux43aux43eux432ux43e}

\begin{itemize}
\tightlist
\item
  \begin{enumerate}
  \def\labelenumi{\arabic{enumi}.}
  \setcounter{enumi}{6}
  \tightlist
  \item
    Задача Адамара Дж. Коші для лінійних диференціальних рівнянь у
    похідних похідних гіперболічного типу. Москва, Наукове видання,
    1978.
  \end{enumerate}
\item
  Арсенін В. Я. Методи математичної фізики та спеціальні функції.
  Москва; Наука, 1974.
\item
  \begin{enumerate}
  \def\labelenumi{\arabic{enumi}.}
  \setcounter{enumi}{8}
  \tightlist
  \item
    Арфкен Г. Математичні методи у фізиці.--- Москва: Атомиздат, 1970.
  \end{enumerate}
\item
  \begin{enumerate}
  \def\labelenumi{\arabic{enumi}.}
  \setcounter{enumi}{9}
  \tightlist
  \item
    Берс Л., Джон Ф., Шехтер Б. Диференціальні рівняння з частинними
    похідними. Москва, видавництво Мир, 1966.
  \end{enumerate}
\item
  \begin{enumerate}
  \def\labelenumi{\arabic{enumi}.}
  \setcounter{enumi}{10}
  \tightlist
  \item
    Бінадзе, А. В. Крайові задачі для еліптичних рівнянь другого
    порядку.--- Москва: Наука, 1966.
  \end{enumerate}
\item
  \begin{enumerate}
  \def\labelenumi{\arabic{enumi}.}
  \setcounter{enumi}{11}
  \tightlist
  \item
    Бідадзе А. В. Деякі класи диференціальних рівнянь у похідних
    похідних. Москва, Наукове видання, 1981.
  \end{enumerate}
\item
  \begin{enumerate}
  \def\labelenumi{\arabic{enumi}.}
  \setcounter{enumi}{12}
  \tightlist
  \item
    Біцадзе А. В., Калініченко Д. К. Збірка задач з рівнянь математичної
    фізики. Москва, Наука, 1977.
  \end{enumerate}
\item
  \begin{enumerate}
  \def\labelenumi{\arabic{enumi}.}
  \setcounter{enumi}{13}
  \tightlist
  \item
    Будак Б. М., Самарський А. А., Тихонов А. Н. Сборник задач по
    математичній фізиці {[}Збірка задач математичної фізики{]}. Москва,
    Наукове видання, 1972.
  \end{enumerate}
\item
  Вазов В., Форсайт Дж. Методи різниці для розв'язання диференціальних
  рівнянь у похідних похідних.--- Москва: Іл, 1963
\item
  Нові методи розв'язання еліптичних рівнянь {[}Нові методи розв'язання
  еліптичних рівнянь{]}. Москва, Гостехіддат, 1948. Векуа, І. Н.
  Узагальнені аналітичні функції. Москва, Фізматгиз, 1959.
\item
  Володимиров В. С. Узагальнені функції в математичній фізиці, Москва,
  Наука, 1976.
\item
  \begin{enumerate}
  \def\labelenumi{\arabic{enumi}.}
  \setcounter{enumi}{18}
  \tightlist
  \item
    Володимиров В. С. та ін. Збірка задач з рівнянь математичної фізики.
    --- Москва: Наука, 1974.
  \end{enumerate}
\item
  Годунов С. К. Рівняння математичної фізики.--- Москва: Наука, 1971.
  Годунов С. К., Рябенський В. С. Схеми різниці.--- Москва: Наука, 1977.
\item
  Гордін: задача Л. Коші для гіперболічних рівнянь.--- Москва: IL, 1961.
  Джеффріс Г., Спірлс Б. Методи математичної фізики. т. 1-3, --- Москва:
  Мир,
\item
  \begin{enumerate}
  \def\labelenumi{\arabic{enumi}.}
  \setcounter{enumi}{1969}
  \tightlist
  \item
  \end{enumerate}
\item
  \begin{enumerate}
  \def\labelenumi{\arabic{enumi}.}
  \setcounter{enumi}{23}
  \tightlist
  \item
    Зоммерфельд А. Рівняння в часткових похідних математичної фізики.---
    Москва: Іл, 1950. 25. Кализкін, Н. Н. Численільна методика.--- М.:
    Наука, 1978.
  \end{enumerate}
\item
  25 Курант Р., Гілберт Д. Методи математичної фізики. т. 1-2, - Москва:
  Гостехиздат, 1951.
\item
  \begin{enumerate}
  \def\labelenumi{\arabic{enumi}.}
  \setcounter{enumi}{26}
  \tightlist
  \item
    Курант Р. Рівняння з частковими виробництвами. Москва, Видавництво
    Мир, 1964.
  \end{enumerate}
\item
  \begin{enumerate}
  \def\labelenumi{\arabic{enumi}.}
  \setcounter{enumi}{27}
  \tightlist
  \item
    Лаврентьєв М. А., Шабат Б. В. Методи теорії функцій комплексної
    змінної. - Москва: Наука, 1965.
  \end{enumerate}
\item
  \begin{enumerate}
  \def\labelenumi{\arabic{enumi}.}
  \setcounter{enumi}{28}
  \tightlist
  \item
    Ладиженська О. А. Прикордонні задачі математичної фізики. Москва:
    Наука,
  \end{enumerate}
\item
  Ладиженська О. А. Змішана задача для гіперболічних рівнянь.--- Москва:
  Фізматгиз, 1953.
\end{itemize}

Список був складений редактором перекладу.

\begin{itemize}
\tightlist
\item
  Ладиженська О. А., Солоняков В. А., Уралцева Н. Н. Лінійні та
  кварілінні рівняння параболічного типу.--- Москва: Наука, 1967.
\item
  \begin{enumerate}
  \def\labelenumi{\arabic{enumi}.}
  \setcounter{enumi}{31}
  \tightlist
  \item
    Ладиженська О. А., Уралцева Н. Н. Лінійні та квазілінійні рівняння
    еліптичного типу. - Москва: Наука, 1964, 1973.
  \end{enumerate}
\item
  Лебедєв Н. Н., Скальська М. П., Уфляйнд Ю. С. Сборник задах по
  математичній фізиці.--- Москва: Гостехиддат, 1965.
\item
  \begin{enumerate}
  \def\labelenumi{\arabic{enumi}.}
  \setcounter{enumi}{33}
  \tightlist
  \item
    Лайонс, Дж.--- Л. Деякі методи розв'язання задач з нелінійними
    краями.--- М.З. Мір. 1972.
  \end{enumerate}
\item
  \begin{enumerate}
  \def\labelenumi{\arabic{enumi}.}
  \setcounter{enumi}{34}
  \tightlist
  \item
    Лайонс Дж.---Л., Maggenes E. Neobreed Grancinum Problems та їх
    застосування. Москва, Видавництво Мир, 1971. 36. Марчук Г. І. Методи
    обчислювальної математики. Москва, Наукове видання, 1980.
  \end{enumerate}
\item
  Мізохата С. Теорія рівнянь з приватними виробниками. --- Москва: Мир,
  1977.
\item
  \begin{enumerate}
  \def\labelenumi{\arabic{enumi}.}
  \setcounter{enumi}{37}
  \tightlist
  \item
    Міранда К. Рівняння з частковими похідними еліптичного типу.---
    Москва: Іл, 1957.
  \end{enumerate}
\item
  Михайлов, В. П. Диференціальні рівняння в похідних з похідними. ---
  Москва: Наука, 1976.
\item
  Міклін С. Г. Курс математичної фізики.--- Москва: Наука, 1968. Міхлін,
  С. Г. Лінійні диференціальні рівняння. --- Москва; Вишня школа, 1977.
\item
  Морс П. М., Фецибах X, Методи теоретичної фізики. т. 1-2.--- Москва:
  IL, 1958.
\item
  \begin{enumerate}
  \def\labelenumi{\arabic{enumi}.}
  \setcounter{enumi}{42}
  \tightlist
  \item
    Меттьюз Дж., Вокер Р. Математичні методи фізики.--- Москва:
    Атомиздат, 1972.
  \end{enumerate}
\item
  Лекції Нагумо М. з сучасної теорії рівнянь у приватних виробництвах
  Москва, видавництво Мвр, 1967. 45. Обен, Дж.-П. Приблизне розв'язання
  еліптичних граничних задач.--- Москва:
\item
  Мір, 1977. 46. Петровський І. Г. Лекції з рівнянь із частковими
  похідними. --- Москва: Фізматгиз, 1961.
\item
  \begin{enumerate}
  \def\labelenumi{\arabic{enumi}.}
  \setcounter{enumi}{46}
  \tightlist
  \item
    Положий Г. Н. Еквалізація математичної фізики.--- Москва: Вища
    школа.
  \end{enumerate}
\item
  \begin{enumerate}
  \def\labelenumi{\arabic{enumi}.}
  \setcounter{enumi}{47}
  \tightlist
  \item
    Рід М., Саймон Б. Методи сучасної математичної фізики.--- Москва:
    Мир,
  \end{enumerate}
\item
  т. 1-1977, т. 2-1978, т. 3, 4-1982. Ріхамаєр Р. Принципи сучасної
  математичної фізики.--- Москва: Мир,
\item
  \begin{enumerate}
  \def\labelenumi{\arabic{enumi}.}
  \setcounter{enumi}{49}
  \tightlist
  \item
    Ріхтмаєр Р., Мортон К. Методи різниці для розв'язання задач з
    граничними значеннями.--- Москва:
  \end{enumerate}
\item
  Мір, 1972,
\item
  Рождественський Б. Л., Япенко, Н. Н. Системи квазіполіцейських рівнянь
  та їх застосування в газовій динаміці.--- Москва: Наука, 1978.
\item
  52, Самарський А. А. Теорія різницьких схем. Москва, Наука, 1977.
\item
  \begin{enumerate}
  \def\labelenumi{\arabic{enumi}.}
  \setcounter{enumi}{52}
  \tightlist
  \item
    Самарський А. А., Андрєєв В. Б. Різниці для еліптичних рівнянь---
    Москва, Наука, 1976.
  \end{enumerate}
\item
  \begin{enumerate}
  \def\labelenumi{\arabic{enumi}.}
  \setcounter{enumi}{53}
  \tightlist
  \item
    Смірнов В. І. Курс вищої математики.--- Москва: Наука, 1974.
  \end{enumerate}
\item
  \begin{enumerate}
  \def\labelenumi{\arabic{enumi}.}
  \setcounter{enumi}{54}
  \tightlist
  \item
    Смірнов М. М. Дегенеративні еліптичні та гіперболічні рівняння.--
    Москва: Наука, 1966.
  \end{enumerate}
\item
  56, Смірнов М. М. Диференціальні рівняння у похідних другого порядку,
  -- Мінськ: Видавництво БДУ, 1974.
\item
  \begin{enumerate}
  \def\labelenumi{\arabic{enumi}.}
  \setcounter{enumi}{56}
  \tightlist
  \item
    Смірнов М. М. Проблеми рівнянь математичної фізики. Москва, Наукове
    видавництво, 1965.
  \end{enumerate}
\item
  \begin{enumerate}
  \def\labelenumi{\arabic{enumi}.}
  \setcounter{enumi}{57}
  \tightlist
  \item
    Соболєв С. Л. Пекаки primeneniya funktsional'nogo analiza v
    matematicheskoi fizike {[}Застосування функціонального аналізу в
    математичній фізиці{]}. Ленінград: Видавництво ЛГУ, 1950,
    Новосибірськ: СБ АН РСРСР, 1962.
  \end{enumerate}
\item
  Соболєв, С. Л. Рівняння математичної фізики.--- М · Наука, 1966.
\item
  Траптер К. Дж. Інтегральні перетворення у математичній фізиці.---
  Москва: Гостехиддат, 1957.
\item
  \begin{enumerate}
  \def\labelenumi{\arabic{enumi}.}
  \setcounter{enumi}{60}
  \tightlist
  \item
    Трев Ф. Лекції з теорії диференціальних рівнянь у похідних
    похідних.--- Москва: Мир, 1965.
  \end{enumerate}
\item
  Трікомі Ф. Лекції з диференціальних рівнянь у похідних похідних.
  Москва, IL Publ., 1957. Фрідман А. Рівняння з частковими похідними
  параболічного типу. --- Москва: Мир, 1968.
\item
  Германдер Л. Лінійні диференціальні оператори з частковими
  похідними.--- Москва: Мир. 1965.
\item
  \begin{enumerate}
  \def\labelenumi{\arabic{enumi}.}
  \setcounter{enumi}{64}
  \tightlist
  \item
    Шварц Л. Математичні методи фізичних наук.--- Москва: Мир, 1965.
  \end{enumerate}
\item
  \begin{enumerate}
  \def\labelenumi{\arabic{enumi}.}
  \setcounter{enumi}{65}
  \tightlist
  \item
    Шилов Г. Е. Математичний аналіз. Спеціальний курс. Москва,
    Фізмат-гіз видавництво, 1960.
  \end{enumerate}
\item
  \begin{enumerate}
  \def\labelenumi{\arabic{enumi}.}
  \setcounter{enumi}{66}
  \tightlist
  \item
    Щилов, Г. Е. Математичний аналіз. Другий спеціальний курс.---
    Москва: Наука, 1965.
  \end{enumerate}
\item
  \begin{enumerate}
  \def\labelenumi{\arabic{enumi}.}
  \setcounter{enumi}{67}
  \tightlist
  \item
    Яненко Н. Н. Метод дробових щагів для розв'язання багатовимірних
    задач математичної фізики.--- Новосибірськ: Наука, 1967.
  \end{enumerate}
\end{itemize}

\section{ІНДЕКС
ІМЕНІ}\label{ux456ux43dux434ux435ux43aux441-ux456ux43cux435ux43dux456}

Бернуллі 329, Бессель 106, Буффон 322, Ніколсон 296, Ньютон 296

Гамільтон 335, 336 Гаусс 286 Гельмгольц 223 Грін 10 Гук 147 Гюйгенс 177,
181 Пуассон 180, 235

Д'Аламбер 126--129, Діріхле 85, 240, Дюамель 107, 108, 110--112, Рітц
338, 340, Робін 240, Родюг 268

** Seidel 286

Сімсон 320 Стокс 8, 180

Кіз 305, 306 Крхгоф 124

Кірхгоф 124 Кранк 296 Крістоффель 363 Ферма 335 Фур'є 24, 25, 36, 41,
74, 84, 91

Гакслі 304, Ханкел 105, 106, Хевісайд 82, 115, 116, Ходжкін 304-306

Лагранж 329 Ламм 307 Лаплас 8, 50, 98-107 Лежандр 235 Лельбніц 180
Лібман 285 Ліувіль 51, 58, 66, 148

Шварц 363 Шредінгер 8 Штурм 51, 58, 66, 148

Ейлер 84, 215, 329

\textbf{М}аксвелл 8 Меллін 105, 106

10sr (Молодий) 124

Navier 8 Neumann 240, 286

Jacobi 286

\section{ІНДЕКС
ПРЕДМЕТІВ}\label{ux456ux43dux434ux435ux43aux441-ux43fux440ux435ux434ux43cux435ux442ux456ux432}

Матрична алгебра 215, 216 Аналітичне розв'язання 301 Наближення сплайну
10

Позачасовий час 165

Калькулюс варіацій 329 Методи варіацій 10 Зовнішня похідна 79 Хвилі, що
рухаються, 149

-- квартира 177 -- стоячи 222

сферичний 178

--- ударні 210 циліндричні 177, 182

Гармонік 84, 154 Сферичний 265 градієнів 22 Граничні умови 17, 22, 23,
26

--- лінійний 26

--- однорідні 36, 46 --- гетерогенні 26, 49, 50

Дельта-функція 276 Діагоналізація матриці 216 Задача дифузії 113, 163
Дифузійний потік 114 --- член 113 Дифузія 113

Задача ** 147 -- Діріхле 242, 243, 246

\begin{itemize}
\item
  Діріхле - Нойман 364 - Кірхгоф 124
\item
  Коуу 69, 95, 116, 126, 207
\end{itemize}

--- мінімізація 10 --- Нойман 243---245

Ньютон 52, 121

--- на брахістохроні 329 Штурму --- Ліувіль 51, 58, 59, 66,

Закон збереження енергії 31, 114, 205

--- Теплопередача Ньютона 8 246

--- Фур'є 24, 25 Звукові хвилі 124 Момент згинання 157 Інтеграл Фур'є 87
Інтегральні перетворення 9, 73.76. Рівняння 10

Класифікація UCP 308 Крутильні осциляції 123 Поздовжні 123 Конвективні
хвилі 116 Конвективний транспорт 113 --- Член 113 Конвекція 113
Скінченні різниці 282 Канонічні координати 127, 132 --- Полярні 233,
235-238 Сферичні 233, 235, 236, 239, 266 Циліндричні 233, 235, 236, 239
--- еліптичний 239 Неперервний набір 96 Косинусне перетворення 73 ---
скінченне 185, 186 --- Обернене 75, 185, 186 Коефіцієнт дифузії 119 ---
теплопередача 23 Задача граничного значення 240, 241

Лапласіан 233 Лінійність 11

--- --- Еліптичний тренажер 241

Метод дедукції 319 -- Лібман 285, 286 -- Монте-Карло 316 -- Рітц 338,
340 -- Сімпсон 320 спуск 182 -- Фур'є 36 Режим коливання 154 Модуль Юнга
124, 145

Початковий зсув 231 Початкові умови 18

Оператор Лапласа 233 Ортогональність 41, 84, 90

\begin{itemize}
\tightlist
\item
  Ходжкіна - Гакслі 304
\end{itemize}

Власні вектори 67, 215, 217

Випадкові номери 319

Змішане завдання 19

--- enachenia 54, 215, 216 Відгук 96 функція 10, 65 Конформне
відображення 253, 356 Крістоффель-Шварц 363 --- частота 162 Спектральна
функція 76, 77 --- неперервна 77 --- частота 85 Stock 274 Параметрична
ідентифікація 304 Змінна залежна 8 Стоячі хвилі 148, 149 Суперпозиція
97, 193, 223 --- незалежна 8 Схема підрахунку тікання 289 Щільність
матеріалу 30 --- явна 289 Многочлен Лежандра 265 Кранк--Ніколсон 296,
299, 328 Порядок рівняння 10 Потенціал Pewtop 266 з круговою симетрією
240. сферично-симетрична 240 циліндрично симетрична 270 Температурна
характеристика 69 Теорема Діріхле 85 Правило Лейбніца 180--Стокс 180
Перетворення залежних змінних Теорія збурень 10 Тепловий імпульс 96
інтеграл 9, 73, 76 Теплоємність 30 --- координати 9 --- Лаплас 50, 75,
98---107, 115, 116, теплопровідність 25, 30, 33 126 Точкове джерело 274
--- обернене 75 --- Меллін 105, 106 --- Ханкель 76, 105, 106 ---
обернене 76. Питома теплоємність 34 --- Фур'є 74, 75, 83, 85, 89, 91
Вирівнювання Бесселя 106, 224, 239 --- хвиля 8, 142, 169, 234 --- ---
скінченне 185 --- нелінійне 304 --- пряме 88, 116, 126 --- обернене 88,
74, 75 гармонічні коливання 223 --- Гельмгольц 223, 224, 235 Принцип
Гамільтона 335, 336 --- Гюйгенс 177, 181 --- Дюамель 107, 108, 110, 112
Рух Ньютона 8 дифузія 18, 197 - Суперлосоцитності 191, 196, 250
конвектива 19, 63, 113, 119, 120 ** - Ферма 335** - Ламм 307 - Лаплас 8,
170, 235, 258, 285, 363 - Лежандр 239, 268 Змінне розділення 9, 191
лінійне 11 - гіперболічне 11, 12, 121, 169 Розклад власних функцій 10,
66 - параболічне 11, 12, 14, 169 - еліптичне 11, 12, 169 - Максвелл 8 -
Нав'є - Стокс 8 Згортка 90, 93 неліцейних {[}{]} гетерогенних 11, 19,
65. --- скінченна 102 Гравітаційна сила 122 гомогенна P Пуассон 235. --
--- ранкового перетворення 122 синусоїдального перетворення 75 постпост
285 скінченне 185, 186, 193 зворотне 75, 185, 186 реакції з нелінійною
дифузією Система телеграфних рівнянь 124 --- з прямокутними
коефіцієнтами 60.

170

з частковими похідними 7

Телеграф 8, 123, 125, 215

Зберігання тепла 16

{\def\LTcaptype{none} % do not increment counter
\begin{longtable}[]{@{}llllllll@{}}
\toprule\noalign{}
\_ & tep & Lonre & Доступність води & 8, & 16, & 30, & 169. \\
\midrule\noalign{}
\endhead
\bottomrule\noalign{}
\endlastfoot
& & 234. & & • & • & • & , \\
\end{longtable}
}

\begin{itemize}
\tightlist
\item
  характеристики 202, 208
\end{itemize}

--- Шредінгер, 8 років

--- Ейлер 215, 256, 268

--- Ейлер-Лагранж 329, 330, \textbf{33}2

Стан Діріхле 240

\begin{verbatim}
- Neumann 240 - Robin 240
\end{verbatim}

Сингулярний диференціал 268

Гріна 10, 96, 196, 273, 276

д'Аламбер 126-132

\begin{itemize}
\tightlist
\item
  додатково до ймовірнісного інтеграла 80
\end{itemize}

Неточника 96, 276

-- неперіодичні 77 -- періодичні 77 -- Хевісайд 82, 115, 116, 365.
Фур'є-образ 88, 89, 91, 93

Формула д'Аламбера 180, 183

--- Лейбніц 184

Джерела 99, 100

\begin{itemize}
\tightlist
\item
  Пуассон 180
\end{itemize}

--- --- інтеграл 248, 249, 251

--- Родріга 268

--- Ейлер 84, 85 Фундаментальне рішення 37 Функціональність 329

Функція Бесселя 224, 225

Чисельні методи 9 Кількість персемних 11

Рівнопотенціальні лінії 277

Ядро перетворення 73 -- перетворення Пуассона 252

\section{ЗМІСТ}\label{ux437ux43cux456ux441ux442}

{\def\LTcaptype{none} % do not increment counter
\begin{longtable}[]{@{}
  >{\raggedright\arraybackslash}p{(\linewidth - 4\tabcolsep) * \real{0.1856}}
  >{\raggedright\arraybackslash}p{(\linewidth - 4\tabcolsep) * \real{0.1237}}
  >{\raggedright\arraybackslash}p{(\linewidth - 4\tabcolsep) * \real{0.6907}}@{}}
\toprule\noalign{}
\begin{minipage}[b]{\linewidth}\raggedright
З редагування
\end{minipage} & \begin{minipage}[b]{\linewidth}\raggedright
Гора
\end{minipage} & \begin{minipage}[b]{\linewidth}\raggedright
Переклад
\end{minipage} \\
\midrule\noalign{}
\endhead
\bottomrule\noalign{}
\endlastfoot
Передмова & viye & ****** \\
& & \\
Частина 1. & Введено & \\
Лекція & 1. & Вступ до теорії диференціальних рівнянь з похідними
похідними \\
Частина 2. & Діф & Завдання з синтезу \\
Лекція & 2. & Задачі типу дифузії (параболічні рівняння) \\
Лекція & 3. & Граничні умови в задачах дифузії \\
Лекція & 4. & Виведення рівняння теплопровідності , \\
Лекція & 5. & Розщеплення змінних \\
Лекція & 6. & Перетворення неоднорідних граничних умов у одно- \\
& & Родичі \\
Лекція & 7. & Розв'язання складніших задач методом розділення змін- \\
Тторене \textasciitilde{} & o & Перетворити комплексні рівняння у просту
форму \\
Лекція-лекція & o. & Перетворення комплексних рівнянь у простий
вигляд \\
Мексика & σ, & Власні функції , \\
Pection & 10 & Інтегральні перетворення (перетворення синусів і
козіпус- \\
Боулз & 10. & ) \\
Лекція & 11. & Перетворення рядів і четвірок \\
Лекіня & 12 & Перетворення Фур'є та його застосування до розв'язку \\
* 10HM 11 /1 & , & Рівняння з похідними похідними \\
Лекція & 13. & Перетворення Лапласа \\
Лекція & 14. & Принцип Дюамеля \\
Лекчич & 15. & Принцип Дюамеля \\
& & \\
Частина 3. & Guipe & Проблеми \\
Лекпея & 16. & Одномірне хвильове рівняння (гіперболічні еквалізації) \\
& & \\
Лекція & 17. & Формула Д'Аламбера \\
Лекція & 18. & Формула Д'Аламбера (продовження) \\
Лекція & 19. & Рівняння води та граничні умови \\
Лекція & 20. & Вібрації обмеженої струни (стоячі хвилі) 14 \\
Лекція & 21. & Осциляції пучка (рівняння часткового виробника \\
& & Четвертий порядок) \\
& & Перехід на безрозмірні змінні \\
Лекція & 23. & Класифікація диференціальних рівнянь у похідних \\
& & (Канонічна форма гіперболічного рівняння) 10 \\
Лекція & 24. & Хвильове рівняння у вільному просторі (двовимірне \\
\_ & \textasciitilde{} - & та тривимірні задачі) \\
Лекція & 25. & Скінченні перетворення Фур'є (синус і косинус до - \\
& 00 & Освіта) \\
& & Принцип суперпозиції --- основи теорії динамічних систем 19 \\
& & Рівняння першого порядку (характеристичний метод) 19 \\
Лекція & 28. & Нединамічні рівняння першого порядку (закони
збереження \\
t. & 00 & Ньяя) , \\
Лекція & 29. & Системи диференціальних рівнянь у похідних 2 \\
Розділ & & Мембранні осциляції (покращення хвиль у Polar 2000
Godinates) \\
& & БУУУУММаалака аа аа аа аа аа \\
\end{longtable}
}

{\def\LTcaptype{none} % do not increment counter
\begin{longtable}[]{@{}
  >{\raggedright\arraybackslash}p{(\linewidth - 2\tabcolsep) * \real{0.7658}}
  >{\raggedright\arraybackslash}p{(\linewidth - 2\tabcolsep) * \real{0.2342}}@{}}
\toprule\noalign{}
\begin{minipage}[b]{\linewidth}\raggedright
Частина 4. Еліптичні задачі
\end{minipage} & \begin{minipage}[b]{\linewidth}\raggedright
233
\end{minipage} \\
\midrule\noalign{}
\endhead
\bottomrule\noalign{}
\endlastfoot
Лекція 31. Лапласнан (інтутивна олісанація) & 233240248255 \\
гармоніки) & 264273 \\
Частина 5. Чисельні та наближені методи & 282 \\
Лекція 37. Чисельні розв'язки (еліптичні задачі) & 282289 \\
Lsson) & 296301 \\
Рівняння) & 308316 \\
Монте-Карло & 322329 \\
Лекція 45. Варіаційні методи розв'язання диференціальних рівнянь у
похідних похідних & 338 \\
Лекція 46. Розв'язок рівнянь часткових похідних теоретиків збурень &
346 \\
Лекція 47. Розв'язання диференціальних рівнянь у похідних похідних за
методом конформних відображень & 355 \\
Додаток 1 & 365 \\
Додаток 2 & 371 \\
Spinning 3 & 373 \\
Джерела & 375 \\
Індекс назв & 378 \\
Індекс предметів & 379 \\
\end{longtable}
}

\section{Дорогий
читачу!}\label{ux434ux43eux440ux43eux433ux438ux439-ux447ux438ux442ux430ux447ux443}

Будь ласка, надсилайте свої коментарі щодо змісту приладів, його
конструкції, якості перекладу та іншого за адресою: 129820, Москва,
I-110, GSP. 1-й Рижський провулок, 2, видавництво «Мир».

\section{Фарлоу
Стейпле}\label{ux444ux430ux440ux43bux43eux443-ux441ux442ux435ux439ux43fux43bux435}

\section{ДИФЕРЕНЦІАЛЬНІ РІВНЯННЯ В ПОХІДНИХ ПОХІДНИХ ДЛЯ НАУКОВЦІВ І
ІНЖЕНЕРІВ}\label{ux434ux438ux444ux435ux440ux435ux43dux446ux456ux430ux43bux44cux43dux456-ux440ux456ux432ux43dux44fux43dux43dux44f-ux432-ux43fux43eux445ux456ux434ux43dux438ux445-ux43fux43eux445ux456ux434ux43dux438ux445-ux434ux43bux44f-ux43dux430ux443ux43aux43eux432ux446ux456ux432-ux456-ux456ux43dux436ux435ux43dux435ux440ux456ux432}

Науковий редактор М. В. Хатунцева-молодший науковий редактор Е. І.
Качулина Художник В. Е. Карпов Художній редактор В. І. Шаповалов
Технічний редактор Т. А. Максимова Коректор Т. І. Стіфсева

\section{UB No 5289}\label{ub-no-5289}

Набор 27.12.84. Подано до публікації 01.07.85. Формат \(60\times90^{1}\)
є. Папір книги. журналу. Печатка висока. Літературний Гарінтура. Том L2
--- бум. n, Usl. Печ. с. 24. Стан. Cr.-Ott. 24. Уч.-изд. л. 19.69, Ізд.
No 1/3807. Тираж 5 000 примірників. Замовлення L6 241. Ціна 1 ранд. 70
тис.

MIR 129820, GSP, Москва, I-110, 1-й Рижський провулок, 2

Надруковано в Ленінградській друкарні No 6 Ордена Червоного Прапора
Праці Ленінградської асоціації «Технічний король» імені Євга. Соколова
Союзполіграф при Державному комітеті СРСР з видавничої справи,
друкарства та книжкової торгівлі, 193144, Ленінград, вулиця Моісеєнко,
10, з матрицями Ордена Жовтневого Регулюта та Ордена Червоного Прапора
Праці МПО «Перша зразкова друкарня» імені А. А. Жданога
Союзполіграфірому при Державному комітеті СССГ з видавничої, друкарської
та книжкової торгівлі, 113054, Москам, Валовія, 28

\end{document}
